%Version 3 December 2023
% See section 11 of the User Manual for version history
%
%%%%%%%%%%%%%%%%%%%%%%%%%%%%%%%%%%%%%%%%%%%%%%%%%%%%%%%%%%%%%%%%%%%%%%
%%                                                                 %%
%% Please do not use \input{...} to include other tex files.       %%
%% Submit your LaTeX manuscript as one .tex document.              %%
%%                                                                 %%
%% All additional figures and files should be attached             %%
%% separately and not embedded in the \TeX\ document itself.       %%
%%                                                                 %%
%%%%%%%%%%%%%%%%%%%%%%%%%%%%%%%%%%%%%%%%%%%%%%%%%%%%%%%%%%%%%%%%%%%%%

%%\documentclass[referee,sn-basic]{sn-jnl}% referee option is meant for double line spacing

%%=======================================================%%
%% to print line numbers in the margin use lineno option %%
%%=======================================================%%

%%\documentclass[lineno,sn-basic]{sn-jnl}% Basic Springer Nature Reference Style/Chemistry Reference Style

%%======================================================%%
%% to compile with pdflatex/xelatex use pdflatex option %%
%%======================================================%%

%%\documentclass[pdflatex,sn-basic]{sn-jnl}% Basic Springer Nature Reference Style/Chemistry Reference Style


%%Note: the following reference styles support Namedate and Numbered referencing. By default the style follows the most common style. To switch between the options you can add or remove “Numbered” in the optional parenthesis. 
%%The option is available for: sn-basic.bst, sn-vancouver.bst, sn-chicago.bst%  
 
%%\documentclass[pdflatex,sn-nature]{sn-jnl}% Style for submissions to Nature Portfolio journals
%%\documentclass[pdflatex,sn-basic]{sn-jnl}% Basic Springer Nature Reference Style/Chemistry Reference Style
\documentclass[pdflatex,sn-mathphys-num]{sn-jnl}% Math and Physical Sciences Numbered Reference Style 
%%\documentclass[pdflatex,sn-mathphys-ay]{sn-jnl}% Math and Physical Sciences Author Year Reference Style
%%\documentclass[pdflatex,sn-aps]{sn-jnl}% American Physical Society (APS) Reference Style
%%\documentclass[pdflatex,sn-vancouver,Numbered]{sn-jnl}% Vancouver Reference Style
%%\documentclass[pdflatex,sn-apa]{sn-jnl}% APA Reference Style 
%%\documentclass[pdflatex,sn-chicago]{sn-jnl}% Chicago-based Humanities Reference Style

%%%% Standard Packages
%%<additional latex packages if required can be included here>

\usepackage{graphicx}%
\usepackage{multirow}%
\usepackage{amsmath,amssymb,amsfonts}%
\usepackage{amsthm}%
\usepackage{mathrsfs}%
\usepackage[title]{appendix}%
\usepackage{xcolor}%
\usepackage{textcomp}%
\usepackage{manyfoot}%
\usepackage{booktabs}%
\usepackage{algorithm}%
\usepackage{algorithmicx}%
\usepackage{algpseudocode}%
\usepackage{listings}%
\usepackage{placeins}%
\usepackage{float}%
\usepackage{tabularx}%
\usepackage{array}%
\usepackage{placeins}
\usepackage[utf8]{inputenc}
\usepackage[T5]{fontenc}
\usepackage[vietnamese]{babel}
\usepackage{booktabs}


\theoremstyle{thmstyleone}%
\newtheorem{theorem}{Theorem}%
\newtheorem{proposition}[theorem]{Proposition}%

\theoremstyle{thmstyletwo}%
\newtheorem{example}{Example}%
\newtheorem{remark}{Remark}%

\theoremstyle{thmstylethree}%
\newtheorem{definition}{Definition}%

\raggedbottom
%%\unnumbered% uncomment this for unnumbered level heads

\begin{document}

\title[Semi-supervised Detection on Chest X-rays]{Semi-supervised Object Detection For Anomaly Detection On Chest X-rays
}
\author[1]{Nguyen Thai Quach}\email{qtnguyen.ist@gmail.com}
\author*[2]{Hai Thanh  Nguyen}\email{nthai.cit@ctu.edu.vn}

\affil*[1]{\orgname{FPT University}, \orgaddress{\city{Can Tho}, \country{Vietnam}}}
\affil[2]{\orgname{Can Tho University}, \orgaddress{\city{Can Tho}, \country{Vietnam}}}

\abstract{
Semi-supervised object detection for anomaly detection on chest X-rays
}
\keywords{
Semi-supervised Learning,
Object Detection,
Chest X-ray,
Anomaly Detection,
Pseudo-labeling,
Teacher--Student Learning,
VinDr-CXR,
Medical Image Analysis
}

\maketitle

\section{Giới Thiệu}\label{sec:intro}
\subsection{Bối cảnh phát hiện bất thường trên ảnh X-quang ngực}

Chụp X-quang ngực (Chest X-ray -- CXR) là một trong những phương thức chẩn đoán hình ảnh được sử dụng phổ biến trong thực hành lâm sàng nhờ tính sẵn có, chi phí tương đối thấp, thời gian thực hiện nhanh và liều bức xạ thấp. Ảnh CXR thường được sử dụng trong sàng lọc và hỗ trợ phát hiện nhiều bất thường thuộc phổi, màng phổi, tim--trung thất và hệ cơ xương lồng ngực, chẳng hạn như viêm phổi, tràn dịch màng phổi, khối mờ, tổn thương xương và các biến đổi bất thường khác \cite{wang2017chestx,irvin2019chexpert,nguyen2022vindr}.

Tuy nhiên, phân tích ảnh CXR vẫn là một nhiệm vụ đầy thách thức. Các cấu trúc giải phẫu trong lồng ngực bị chiếu chồng lên nhau trên ảnh hai chiều, trong khi nhiều bất thường cũng có thể xuất hiện với kích thước, vị trí và hình thái khác nhau, từ các tổn thương nhỏ, khu trú đến những vùng bất thường lan rộng. Trong cùng một ảnh, nhiều loại bất thường có thể xuất hiện đồng thời và bị che lấp bởi mô lành xung quanh \cite{nguyen2022vindr,liu2020chestx}.

Trong bối cảnh đó, các mô hình học sâu có tiềm năng hỗ trợ quá trình phân tích ảnh CXR bằng cách cung cấp dự đoán định lượng và gợi ý các vùng nghi ngờ cần được xem xét. Các hệ thống này không nhằm thay thế quyết định chuyên môn của bác sĩ mà đóng vai trò như công cụ hỗ trợ sàng lọc và đọc ảnh. Đối với nhiều ứng dụng, một dự đoán ở mức toàn ảnh chỉ cho biết ảnh có hay không có một loại bất thường, nhưng không cung cấp vị trí cụ thể của tổn thương. Vì vậy, object detection là một lựa chọn phù hợp khi mục tiêu nghiên cứu yêu cầu mô hình đồng thời xác định loại bất thường và định vị vùng liên quan bằng bounding box.

Các nghiên cứu supervised object detection gần đây đã cải thiện khả năng định vị bất thường thông qua nhiều hướng phát triển về kiến trúc và biểu diễn đặc trưng. Tuy nhiên, các cải tiến này phần lớn vẫn dựa trên giả định rằng một lượng đủ lớn ảnh huấn luyện đã được gán nhãn bounding box đầy đủ và chính xác. Trong ảnh y tế, giả định này khó được đáp ứng vì mỗi vùng bất thường cần được xác định vị trí và gán lớp bởi người có chuyên môn. Do đó, bên cạnh độ khó của bản thân bài toán detection, sự phụ thuộc vào nhãn bounding box chất lượng cao trở thành một nút thắt quan trọng đối với việc mở rộng các hệ thống phát hiện bất thường trên ảnh CXR \cite{liu2020chestx,nguyen2022vindr,ji2022benchmark}.


\subsection{Phân biệt classification, segmentation và object detection trong ảnh X-quang ngực}

Trong phân tích ảnh y tế bằng học sâu, image classification, object detection và image segmentation là ba dạng bài toán có mục tiêu dự đoán, mức độ định vị và yêu cầu gán nhãn khác nhau. Việc phân biệt rõ các dạng bài toán này là cần thiết vì kết quả của chúng không thể được đánh giá bằng cùng một hệ metric và cũng không đáp ứng cùng một nhu cầu phân tích hình ảnh.

Với \textbf{image classification}, mô hình nhận toàn bộ ảnh làm đầu vào và dự đoán một hoặc nhiều nhãn ở mức ảnh. Trong bối cảnh ảnh X-quang ngực, classification có thể cho biết một loại bất thường có khả năng xuất hiện trong ảnh hay không, nhưng không bắt buộc mô hình chỉ ra vị trí cụ thể của vùng bất thường. Vì vậy, classification phù hợp với các nhiệm vụ sàng lọc hoặc nhận diện ở mức toàn ảnh, nhưng còn hạn chế khi mục tiêu nghiên cứu yêu cầu xác định vị trí tổn thương.

Với \textbf{image segmentation}, mô hình dự đoán nhãn tại từng điểm ảnh để tạo ra mặt nạ biểu diễn cơ quan hoặc vùng bất thường. So với bounding box, segmentation cung cấp thông tin hình học chi tiết hơn về hình dạng và ranh giới tổn thương. Tuy nhiên, việc tạo nhãn pixel-level thường đòi hỏi chuyên gia xác định tương đối chính xác đường biên của từng vùng bất thường, làm tăng đáng kể thời gian và chi phí gán nhãn.

Với \textbf{object detection}, mô hình đồng thời thực hiện hai nhiệm vụ: định vị một hoặc nhiều vùng bất thường bằng bounding box và phân loại từng vùng vào lớp tương ứng. So với classification, object detection cung cấp thêm thông tin về vị trí của bất thường; so với segmentation, nó không yêu cầu chú thích chi tiết đến từng điểm ảnh. Do đó, object detection tạo ra một mức biểu diễn trung gian giữa dự đoán ở cấp độ toàn ảnh và phân vùng pixel-level, phù hợp với các nghiên cứu cần đồng thời nhận diện và định vị bất thường nhưng chỉ có nhãn dạng bounding box.

Trong nghiên cứu này, bài toán được xác định là \textbf{semi-supervised multi-class object detection} trên ảnh X-quang ngực. Mô hình cần phát hiện các vùng bất thường bằng bounding box và gán mỗi bounding box vào một trong 14 lớp bất thường được xác định trong bộ dữ liệu. Nhãn \textit{No Finding} không được xem là một lớp detection, mà được sử dụng để biểu diễn ảnh âm tính không có bounding box bất thường. Vì vậy, hiệu năng mô hình được đánh giá bằng các metric dành cho object detection, bao gồm bbox mAP@[0.50:0.95], AP50, AP75 và AP theo từng lớp. Các chỉ số như Accuracy hoặc AUC ở mức ảnh, cũng như Dice hoặc mask IoU trong segmentation, không được dùng làm metric chính, vì chúng không đánh giá được đồng thời chất lượng phân loại và độ chính xác định vị bounding box.


\subsection{Vấn đề thiếu nhãn bounding box trong ảnh y tế}

Hiệu quả của các mô hình học sâu trong phân tích ảnh y tế phụ thuộc đáng kể vào quy mô, chất lượng và tính nhất quán của dữ liệu gán nhãn. Tuy nhiên, việc tạo nhãn chi tiết trong lĩnh vực này thường tốn nhiều thời gian, chi phí và đòi hỏi sự tham gia của người có chuyên môn. Đối với image classification, mỗi ảnh thường chỉ cần một hoặc một tập nhãn ở mức toàn ảnh. Ngược lại, trong object detection, người gán nhãn phải xác định riêng từng vùng bất thường, vẽ bounding box và gán lớp tương ứng cho mỗi tổn thương.

Do đó, chi phí gán nhãn trong object detection không chỉ tăng theo số lượng ảnh mà còn phụ thuộc vào số lượng tổn thương xuất hiện trong từng ảnh. Một ảnh chứa nhiều bất thường có thể yêu cầu nhiều bounding box, trong khi mỗi bounding box cần được xác định về vị trí, kích thước và nhãn lớp. Với các ảnh có tổn thương nhỏ, ranh giới không rõ hoặc nhiều bất thường chồng lấp, quá trình này còn đòi hỏi nhiều thời gian kiểm tra hơn. Đây là điểm khác biệt quan trọng so với classification, nơi chi phí gán nhãn thường gắn chủ yếu với số lượng ảnh thay vì số lượng đối tượng trong ảnh.

Trong ảnh y tế, nhãn bounding box còn cần được tạo ra hoặc kiểm tra bởi người có chuyên môn về hình ảnh y khoa. Ngoài chi phí nhân lực, sự khác biệt trong cách diễn giải ảnh giữa các bác sĩ cũng có thể dẫn đến bất nhất về vị trí, kích thước hoặc lớp của bounding box. Vì vậy, việc mở rộng một bộ dữ liệu detection không chỉ là thu thập thêm ảnh, mà còn phải bảo đảm chất lượng và độ tin cậy của từng chú thích. Điều này làm cho bounding-box annotation trở thành một nút thắt quan trọng trong quá trình xây dựng các hệ thống phát hiện bất thường trên ảnh X-quang ngực \cite{nguyen2022vindr,le2023learning}.

Trong nhiều kịch bản nghiên cứu, số lượng ảnh y tế có sẵn có thể lớn hơn đáng kể so với số ảnh được gán nhãn đầy đủ ở mức bounding box. Nếu chỉ sử dụng supervised learning, detector chỉ có thể khai thác phần dữ liệu đã được chú thích. Khi tỷ lệ ảnh có nhãn thấp, mô hình có thể không học đủ biến thiên về hình thái, kích thước và vị trí của tổn thương, đồng thời dễ quá khớp với tập nhãn nhỏ hoặc thiên lệch về các mẫu phổ biến.

Vấn đề này trở nên nghiêm trọng hơn trong bài toán detection đa lớp có phân phối long-tailed. Các bất thường phổ biến thường chiếm phần lớn số bounding box, trong khi các lớp hiếm chỉ xuất hiện trong một số ít ảnh. Khi ngân sách nhãn bị giới hạn, xác suất các lớp hiếm được đại diện đầy đủ trong tập huấn luyện càng thấp. Hệ quả là mô hình không chỉ suy giảm hiệu năng tổng thể mà còn có nguy cơ bỏ sót các lớp ít gặp, ngay cả khi mAP trung bình vẫn có thể bị chi phối bởi các lớp phổ biến \cite{nguyen2022vindr,liu2023anomaly}.

Vì vậy, thách thức trọng tâm không chỉ là cải thiện kiến trúc detector, mà còn là nâng cao hiệu quả sử dụng nhãn, tức đạt được hiệu năng phát hiện tốt hơn với một lượng bounding-box annotation hạn chế. Bài toán này đặt nghiên cứu vào bối cảnh \textbf{label-efficient learning}, đặt ra câu hỏi: trong nhiều hướng giảm phụ thuộc nhãn hiện có, hướng nào phù hợp nhất để cải thiện detection mà không làm gia tăng đáng kể chi phí chú thích. Mục tiếp theo phân tích và lựa chọn hướng tiếp cận này.


\subsection{Label-efficient learning và lý do lựa chọn semi-supervised learning}

Như đã trình bày ở mục trước, khó khăn của object detection trong ảnh y tế không chỉ nằm ở việc xây dựng một mô hình có khả năng phát hiện chính xác các bất thường, mà còn ở chi phí tạo dữ liệu gán nhãn. Khi số lượng ảnh tăng lên, việc yêu cầu chuyên gia xác định và gán nhãn từng bounding box cho mọi tổn thương trở thành một quá trình tốn nhiều thời gian, chi phí và nguồn lực. Do đó, cùng với sự phát triển của các kiến trúc detector, một hướng nghiên cứu quan trọng khác là \textit{label-efficient learning}, tức các phương pháp hướng tới mục tiêu đạt hiệu năng học cao trong điều kiện dữ liệu gán nhãn bị hạn chế hoặc chi phí tạo nhãn lớn \cite{ji2022benchmark,sheng2024barlowtwins}.

Trong các bài toán image classification, mỗi ảnh thường chỉ cần được gán một hoặc một số ít nhãn ở mức toàn ảnh. Vì vậy, chi phí annotation nhìn chung tăng gần tỷ lệ thuận với số lượng ảnh cần gán nhãn. Ngược lại, trong object detection, mỗi đối tượng xuất hiện trong ảnh đều cần được chuyên gia xác định vị trí, hiệu chỉnh kích thước bounding box và gán đúng lớp bất thường. Do đó, chi phí annotation không chỉ phụ thuộc vào số lượng ảnh mà còn phụ thuộc vào số lượng tổn thương xuất hiện trong mỗi ảnh. Một ảnh X-quang ngực có thể chứa nhiều bất thường với kích thước, vị trí và hình thái khác nhau; mỗi bất thường đều yêu cầu một bounding box riêng. Điều này làm cho chi phí annotation trong object detection tăng theo số lượng tổn thương thay vì chỉ theo số lượng ảnh, khiến việc xây dựng các bộ dữ liệu detection trong ảnh y tế tốn kém và khó mở rộng hơn đáng kể so với các bài toán image classification.

Ngoài chi phí về thời gian, annotation trong ảnh y tế còn đòi hỏi chuyên môn của bác sĩ chẩn đoán hình ảnh hoặc các annotator được đào tạo. Việc xác định chính xác vị trí và ranh giới của từng tổn thương thường cần được kiểm tra hoặc hiệu chỉnh bởi chuyên gia, làm cho chi phí annotation tăng không chỉ về số lượng công việc mà còn về nguồn lực chuyên môn cần huy động. Do đó, trong object detection y tế, nút thắt không chỉ nằm ở việc phát triển detector có độ chính xác cao, mà còn ở khả năng sử dụng hiệu quả ngân sách nhãn vốn rất hạn chế. Đây chính là động lực thúc đẩy sự phát triển của các phương pháp thuộc nhóm label-efficient learning.

Để giảm sự phụ thuộc vào dữ liệu gán nhãn đầy đủ, nhiều hướng tiếp cận thuộc nhóm label-efficient learning đã được đề xuất. Các hướng và chiến lược thường được cân nhắc để giảm phụ thuộc vào nhãn bao gồm Active Learning, Weakly Supervised Learning, Self-supervised Learning, Few-shot Learning, Foundation Models, Synthetic Data, Data Augmentation và Semi-supervised Learning. Mặc dù đều có thể góp phần giảm chi phí hoặc tăng hiệu quả sử dụng dữ liệu, các hướng này dựa trên những giả định, nguồn tín hiệu giám sát và mục tiêu huấn luyện khác nhau.

\textbf{Active Learning} hướng tới việc tối ưu hóa quá trình gán nhãn bằng cách lựa chọn những mẫu được đánh giá là mang nhiều thông tin nhất để chuyên gia tiếp tục annotation. Thay vì gán nhãn toàn bộ dữ liệu, mô hình và người gán nhãn tương tác theo nhiều vòng lặp; ở mỗi vòng, mô hình lựa chọn một tập ảnh cần được chú thích bổ sung nhằm cải thiện hiệu quả học với chi phí thấp hơn. Trong object detection, các mẫu được lựa chọn vẫn cần chuyên gia xác định và kiểm tra từng bounding box. Vì vậy, mặc dù Active Learning có thể giảm tổng số ảnh cần annotation, phương pháp này vẫn phụ thuộc vào khả năng tổ chức các vòng gán nhãn bổ sung và sự tham gia liên tục của chuyên gia trong suốt quá trình huấn luyện.

Trong phạm vi nghiên cứu này, mục tiêu không phải là xây dựng một quy trình annotation lặp nhiều vòng hay đánh giá chiến lược lựa chọn mẫu để gán nhãn. Tập dữ liệu có nhãn được cố định ngay từ đầu và không phát sinh thêm chi phí annotation trong quá trình thực nghiệm. Do đó, Active Learning không phản ánh trực tiếp thiết kế nghiên cứu được lựa chọn trong luận văn.

\textbf{Weakly Supervised Learning} là một hướng tiếp cận khác nhằm giảm chi phí annotation bằng cách thay thế bounding box đầy đủ bằng các dạng chú thích yếu hơn, chẳng hạn như nhãn ở mức ảnh (image-level labels), điểm đánh dấu (point annotations), scribbles hoặc các tín hiệu giám sát không đầy đủ khác. Thay vì yêu cầu chuyên gia xác định chính xác vị trí và kích thước của từng tổn thương, Weakly Supervised Learning khai thác các nhãn đơn giản hơn để suy diễn vị trí bất thường trong quá trình huấn luyện. Cách tiếp cận này có thể giảm đáng kể chi phí tạo nhãn so với object detection truyền thống \cite{ji2022benchmark}.

Tuy nhiên, Weakly Supervised Learning vẫn giả định rằng dữ liệu huấn luyện có sẵn một dạng \textit{weak annotation}. Mặc dù bộ dữ liệu VinDr-CXR cung cấp các nhãn ở mức ảnh bên cạnh annotation dạng bounding box, nghiên cứu này chủ động không sử dụng các nhãn global trong quá trình huấn luyện. Thay vào đó, phần dữ liệu ngoài tập có nhãn được xem như dữ liệu chưa gán nhãn nhằm mô phỏng kịch bản chỉ có một lượng nhỏ ảnh được gán bounding box đầy đủ và phần lớn ảnh còn lại không được sử dụng bất kỳ thông tin nhãn nào. Thiết kế này giúp cô lập và đánh giá trực tiếp giá trị của dữ liệu chưa gán nhãn, thay vì khai thác thêm tín hiệu từ các dạng weak annotation. Vì vậy, Weakly Supervised Learning không phải là hướng tiếp cận chính phù hợp với mục tiêu nghiên cứu của luận văn.

\textbf{Multiple Instance Learning (MIL)} gán nhãn ở mức "túi" (bag) thay vì từng đối tượng: mỗi ảnh được xem như một túi chứa nhiều vùng (instance), và túi được coi là dương tính nếu chứa ít nhất một tổn thương. Mô hình suy diễn instance nào chịu trách nhiệm cho nhãn túi, qua đó định vị bất thường từ nhãn mức ảnh. Cách tiếp cận này giảm chi phí annotation vì không cần bounding box cho từng tổn thương. Tuy nhiên, MIL vẫn là một dạng giám sát: đòi hỏi nhãn mức túi (image-level) cho phần dữ liệu được khai thác và thường chỉ cho định vị thô, khó tạo bounding box chính xác cho nhiều lớp. Vì luận văn giả định phần dữ liệu ngoài tập có nhãn là hoàn toàn chưa gán nhãn (kể cả nhãn mức ảnh), MIL không phù hợp với thiết lập unlabeled-only được lựa chọn.

\textbf{Self-supervised Learning} khai thác dữ liệu chưa gán nhãn để học biểu diễn đặc trưng thông qua các nhiệm vụ tiền huấn luyện không cần nhãn thủ công (pretext tasks hoặc pretraining objectives). Sau giai đoạn này, mô hình được tinh chỉnh (fine-tuning) trên tập dữ liệu có nhãn để thực hiện bài toán đích, chẳng hạn như image classification hoặc object detection. Trong những năm gần đây, self-supervised learning đã chứng minh khả năng cải thiện chất lượng biểu diễn đặc trưng và giảm sự phụ thuộc vào các mô hình được tiền huấn luyện trên ImageNet, đặc biệt trong các miền dữ liệu có khoảng cách lớn với ảnh tự nhiên như ảnh y tế \cite{sheng2024barlowtwins}.

Đối với object detection, self-supervised learning chủ yếu mang lại lợi ích ở giai đoạn khởi tạo đặc trưng (representation learning). Sau khi hoàn thành pretraining, detector vẫn cần được huấn luyện bằng dữ liệu có bounding box để học nhiệm vụ định vị và phân loại đối tượng. Nói cách khác, self-supervised learning giúp cải thiện chất lượng biểu diễn của backbone nhưng không trực tiếp khai thác dữ liệu chưa gán nhãn trong quá trình tối ưu detector thông qua pseudo-label hoặc consistency learning. Vì vậy, mặc dù self-supervised learning là một hướng bổ trợ quan trọng và có thể kết hợp với semi-supervised learning trong các nghiên cứu mở rộng, hướng tiếp cận này không trả lời trực tiếp câu hỏi nghiên cứu của luận văn về khả năng khai thác dữ liệu chưa gán nhãn trong quá trình huấn luyện detector.

\textbf{Transfer Learning (TL)} tận dụng một mô hình đã được tiền huấn luyện trên miền nguồn quy mô lớn (ví dụ ImageNet) rồi tinh chỉnh trên miền đích với lượng nhãn hạn chế, nhờ đó giảm nhu cầu dữ liệu có nhãn ở miền đích. Tuy nhiên, TL chủ yếu giải quyết vấn đề khởi tạo và chuyển giao biểu diễn, chứ không khai thác trực tiếp các ảnh chưa gán nhãn trong quá trình tối ưu detector; giai đoạn tinh chỉnh vẫn cần bounding box. Quan trọng hơn, trong luận văn, backbone tiền huấn luyện được sử dụng \emph{như nhau} cho cả supervised baseline lẫn mô hình semi-supervised, nên TL là một thành phần được giữ cố định giữa các thiết lập, không phải biến số đang khảo sát. Vì vậy TL bổ trợ cho pipeline nhưng không phải hướng tiếp cận trung tâm để trả lời câu hỏi nghiên cứu.

\textbf{Few-shot Learning} hướng tới khả năng học hiệu quả từ số lượng rất nhỏ mẫu có nhãn. Trong object detection, few-shot detection thường giả định chỉ có một số ít bounding box cho mỗi lớp hoặc yêu cầu detector nhanh chóng thích nghi với các lớp đối tượng mới (novel classes) sau khi đã được huấn luyện trên các lớp cơ sở (base classes). Do đó, trọng tâm của few-shot learning là khả năng khái quát hóa từ rất ít ví dụ và chuyển tri thức giữa các lớp đối tượng.

Thiết lập này khác với phạm vi nghiên cứu của luận văn. Trong nghiên cứu này, tập lớp bất thường được giữ cố định gồm 14 lớp và không tồn tại khái niệm base classes hoặc novel classes. Thay vì đánh giá khả năng thích nghi với lớp mới, nghiên cứu tập trung phân tích ảnh hưởng của việc giảm ngân sách nhãn bounding box trên cùng một không gian lớp. Các mức dữ liệu có nhãn được thiết kế lần lượt là 1\%, 5\%, 10\% và 20\% nhằm mô phỏng các mức độ thiếu nhãn khác nhau trong thực tế. Vì vậy, few-shot learning không phản ánh trực tiếp bài toán và mục tiêu thực nghiệm của luận văn.

\textbf{Foundation Models} là các mô hình quy mô lớn được tiền huấn luyện trên lượng dữ liệu đa dạng và có khả năng chuyển giao sang nhiều nhiệm vụ hạ nguồn thông qua fine-tuning, prompt-based adaptation hoặc các cơ chế thích nghi khác. Trong ảnh y tế, foundation models có tiềm năng giảm nhu cầu huấn luyện từ đầu và cung cấp biểu diễn tổng quát hơn cho các bài toán classification, segmentation hoặc detection. Tuy nhiên, khả năng chuyển giao của các mô hình này phụ thuộc đáng kể vào mức độ tương đồng giữa dữ liệu tiền huấn luyện và ảnh X-quang ngực, cũng như khả năng thích nghi với các lớp bất thường nhỏ, có độ tương phản thấp và phân phối long-tailed. Đối với object detection, foundation model vẫn thường cần dữ liệu bounding box hoặc một dạng tín hiệu định vị để thích nghi với nhiệm vụ đích. Ngoài ra, hướng này chủ yếu đặt câu hỏi về khả năng transfer và adaptation của một mô hình tiền huấn luyện quy mô lớn, thay vì đánh giá trực tiếp giá trị bổ sung của dữ liệu chưa gán nhãn trong cùng một giao thức teacher--student. Vì vậy, foundation model có thể được xem là một lựa chọn backbone hoặc hướng mở rộng, nhưng không trùng với câu hỏi nghiên cứu trung tâm của luận văn.

\textbf{Synthetic Data} hướng tới việc mở rộng dữ liệu huấn luyện bằng cách tạo ra ảnh hoặc vùng tổn thương nhân tạo thông qua mô phỏng, image synthesis, generative models hoặc các kỹ thuật chèn tổn thương. Hướng tiếp cận này có thể làm tăng số lượng mẫu huấn luyện và hỗ trợ các lớp hiếm mà không yêu cầu gán nhãn thủ công cho từng ảnh thực. Tuy nhiên, giá trị của synthetic data phụ thuộc mạnh vào mức độ chân thực về hình thái, vị trí và bối cảnh giải phẫu của tổn thương được tạo ra. Nếu dữ liệu tổng hợp không phản ánh đúng phân phối của ảnh X-quang thực, mô hình có thể học các dấu hiệu giả tạo hoặc gặp khoảng cách miền giữa ảnh tổng hợp và ảnh lâm sàng. Bên cạnh đó, việc đánh giá chất lượng dữ liệu tổng hợp và xác nhận tính hợp lý y khoa thường cần thêm quy trình kiểm định chuyên môn. Do đó, synthetic data là một hướng độc lập nhằm mở rộng phân phối huấn luyện, nhưng không trực tiếp trả lời câu hỏi liệu các ảnh X-quang thực chưa gán nhãn có cải thiện detector trong điều kiện ngân sách nhãn bounding box bị giới hạn hay không.

\textbf{Data Augmentation} sử dụng các phép biến đổi trên dữ liệu đã có, chẳng hạn như thay đổi độ sáng, tương phản, tỷ lệ, vị trí hoặc các biến đổi hình học, nhằm tăng tính đa dạng của mẫu huấn luyện và cải thiện khả năng khái quát hóa. Đây là thành phần quan trọng trong cả supervised và semi-supervised object detection. Tuy nhiên, data augmentation không tạo ra thông tin lâm sàng độc lập mới mà chủ yếu biến đổi các mẫu đã quan sát. Khi tập dữ liệu có nhãn quá nhỏ, augmentation có thể làm tăng biến thiên hình thức nhưng không bảo đảm bổ sung đầy đủ các kiểu tổn thương, vị trí hoặc lớp hiếm còn thiếu trong tập huấn luyện. Vì vậy, data augmentation nên được xem là một kỹ thuật hỗ trợ trong quy trình huấn luyện, thay vì một chiến lược label-efficient learning thay thế hoàn toàn cho việc khai thác dữ liệu chưa gán nhãn.



\textbf{Semi-supervised Learning} kết hợp một tập nhỏ dữ liệu được gán nhãn
đầy đủ với một tập lớn dữ liệu chưa gán nhãn trong quá trình huấn luyện.
Trong semi-supervised object detection dựa trên teacher--student
pseudo-labeling, Teacher sinh các pseudo-label gồm bounding box, nhãn lớp và
độ tin cậy cho dữ liệu chưa gán nhãn, trong khi Student được tối ưu từ cả dữ
liệu có nhãn thật và dữ liệu chưa gán nhãn thông qua các pseudo-label đã được
lọc \cite{sohn2020simple,liu2021unbiased,xu2021end}. Tùy từng phương pháp,
Teacher có thể được cập nhật từ Student theo các cơ chế khác nhau. Trong
nghiên cứu này, Teacher và Student được khởi tạo đồng thời và Teacher được cập
nhật bằng exponential moving average trong suốt quá trình huấn luyện SSL;
không sử dụng một giai đoạn supervised burn-in riêng.

So với Active Learning, Semi-supervised Learning không yêu cầu phát sinh thêm các vòng annotation mới trong quá trình thực nghiệm. So với Weakly Supervised Learning, phương pháp này không phụ thuộc vào image-level labels, point annotations hoặc các dạng weak annotation khác. So với Self-supervised Learning, dữ liệu chưa gán nhãn được khai thác trực tiếp trong quá trình tối ưu detector thay vì chỉ phục vụ giai đoạn pretraining. Khác với Few-shot Learning, mục tiêu không phải thích nghi với lớp mới từ rất ít ví dụ. So với Foundation Models, nghiên cứu không tập trung đánh giá khả năng transfer hoặc prompt-based adaptation của một mô hình quy mô lớn. So với Synthetic Data, phương pháp sử dụng trực tiếp ảnh X-quang thực chưa gán nhãn thay vì tạo thêm dữ liệu nhân tạo có nguy cơ sai khác phân phối. Data Augmentation vẫn được sử dụng như một thành phần hỗ trợ huấn luyện, nhưng không được xem là chiến lược chính vì nó không bổ sung tín hiệu từ các ảnh thực ngoài tập có nhãn.

Tóm lại, trong số các hướng được cân nhắc để giảm phụ thuộc vào nhãn, Semi-supervised Learning phù hợp trực tiếp nhất với cấu trúc dữ liệu và câu hỏi nghiên cứu của luận văn: phương pháp này khai thác các ảnh X-quang thực chưa gán nhãn ngay trong quá trình huấn luyện detector mà không yêu cầu thêm vòng annotation hoặc tín hiệu nhãn phụ. Vì vậy nghiên cứu tập trung đánh giá liệu dữ liệu chưa gán nhãn có thể cải thiện hiệu năng phát hiện khi ngân sách dữ liệu có nhãn bị giới hạn.

Việc lựa chọn Semi-supervised Learning không dựa trên giả định rằng phương pháp này luôn tốt hơn mọi hướng thay thế, mà dựa trên mức độ phù hợp giữa phương pháp, cấu trúc dữ liệu và thiết kế thực nghiệm. Cụ thể, Semi-supervised Learning cho phép cô lập và đánh giá ảnh hưởng của dữ liệu chưa gán nhãn đối với hiệu quả phát hiện bất thường trong điều kiện lượng bounding-box supervision được kiểm soát.

Theo cách tiếp cận này, luận văn được định vị là một nghiên cứu về \textbf{label-efficient learning} trong object detection trên ảnh X-quang ngực, thay vì một nghiên cứu đề xuất kiến trúc detector hoặc thuật toán semi-supervised object detection mới. Đóng góp chính của nghiên cứu không nằm ở việc xây dựng một mô hình mới, mà ở việc thiết kế một giao thức thực nghiệm có kiểm soát nhằm đánh giá tác động của dữ liệu chưa gán nhãn dưới các mức ngân sách nhãn khác nhau.

\begin{figure}[!ht]
\centering
\includegraphics[width=\linewidth]{z_label_efficiency_diagram.png}
\caption{
Nghiên cứu được đặt trong hướng học tiết kiệm nhãn (label-efficient learning) và đánh giá có kiểm soát hiệu quả của semi-supervised object detection dưới các mức ngân sách dữ liệu có nhãn khác nhau.
}
\label{research_positioning}
\end{figure}

Cụ thể, nghiên cứu xây dựng các tập dữ liệu có nhãn theo bốn mức ngân sách dữ liệu có nhãn ở cấp độ ảnh (\textit{image-level labeled-data budget}) là 1\%, 5\%, 10\% và 20\% của tập huấn luyện cố định. Một ảnh được xem là có nhãn khi toàn bộ các chú giải \textit{bounding box} thuộc ảnh đó được cung cấp cho quá trình huấn luyện. Ở mỗi mức, mô hình supervised learning và semi-supervised learning được huấn luyện trên cùng tập dữ liệu có nhãn, cùng detector, cùng quy trình tiền xử lý, cùng chiến lược chia dữ liệu và cùng giao thức đánh giá. Khác biệt có chủ đích chính giữa hai thiết lập là mô hình semi-supervised được phép khai thác thêm tập dữ liệu chưa gán nhãn thông qua cơ chế teacher--student và pseudo-labeling; các yếu tố nền tảng như detector, dữ liệu có nhãn, phân chia dữ liệu và giao thức đánh giá được giữ nhất quán.

Tuy nhiên, hiệu quả của Semi-supervised Learning không chỉ phụ thuộc vào số lượng dữ liệu chưa gán nhãn mà còn phụ thuộc vào các giả định học nền tảng của phương pháp. Nếu các giả định này không được thỏa mãn, dữ liệu chưa gán nhãn có thể không cải thiện, thậm chí làm suy giảm chất lượng của detector do hiện tượng pseudo-label noise và confirmation bias. Vì vậy, bên cạnh việc đánh giá hiệu năng phát hiện, luận văn cũng phân tích các điều kiện mà dưới đó Semi-supervised Learning phát huy hiệu quả hoặc gặp hạn chế trong bài toán phát hiện bất thường trên ảnh X-quang ngực. Các giả định và rủi ro của Semi-supervised Learning sẽ được trình bày trong mục tiếp theo.

Trong Semi-supervised Learning đã có nhiều cơ chế khai thác dữ liệu chưa gán
nhãn, trong đó teacher--student pseudo-labeling, consistency regularization và
các phương pháp dựa trên mô hình sinh là những hướng tiêu biểu. Luận văn lựa
chọn \textbf{teacher--student pseudo-labeling} vì ba lý do. Thứ nhất, đây là
một trong những khung được sử dụng rộng rãi trong semi-supervised object
detection \cite{sohn2020simple,liu2021unbiased,xu2021end}, phù hợp với định
hướng đánh giá một phương pháp đại diện thay vì đề xuất thuật toán mới. Thứ hai, pseudo-label trong detection có cấu trúc trùng khớp với nhãn thật (bounding box, lớp, confidence), nên khai thác trực tiếp được dữ liệu chưa gán nhãn mà không thay đổi kiến trúc detector, giúp giữ nguyên điều kiện so sánh giữa supervised và semi-supervised. Thứ ba, và quan trọng nhất với mục tiêu nghiên cứu, cơ chế này sinh ra pseudo-label tường minh và có thể phân tích được — cho phép khảo sát trực tiếp chất lượng pseudo-label, hành vi lớp hiếm, confirmation bias và false positive trên ảnh \textit{No Finding}, vốn là trọng tâm của các câu hỏi nghiên cứu. So với các phương pháp chủ yếu dựa trên consistency regularization, teacher--student pseudo-labeling tạo ra tập pseudo-label tường minh (bounding box, nhãn lớp và confidence score), thuận lợi hơn cho việc phân tích chất lượng pseudo-label và các hiện tượng như confirmation bias hoặc rare-class bias.

\subsection{Giả định và rủi ro của semi-supervised learning}

Mặc dù Semi-supervised Learning cho phép khai thác dữ liệu chưa gán nhãn nhằm giảm phụ thuộc vào bounding-box annotation, hiệu quả của phương pháp này không được bảo đảm trong mọi trường hợp. Khác với supervised learning chỉ học từ dữ liệu có nhãn, Semi-supervised Learning dựa trên một số giả định về cấu trúc dữ liệu và quá trình học. Nếu các giả định này không được thỏa mãn, việc bổ sung dữ liệu chưa gán nhãn có thể không cải thiện, thậm chí làm suy giảm hiệu năng của mô hình.

Một trong những giả định cơ bản là \textit{smoothness assumption}, cho rằng các mẫu dữ liệu nằm gần nhau trong không gian đặc trưng có xu hướng cùng thuộc một lớp. Theo giả định này, nếu một mẫu có nhãn được dự đoán chính xác thì các mẫu chưa gán nhãn lân cận cũng nhiều khả năng mang cùng nhãn. Điều này tạo cơ sở cho việc lan truyền thông tin từ dữ liệu có nhãn sang dữ liệu chưa gán nhãn thông qua pseudo-label hoặc consistency regularization.

Liên quan chặt chẽ với giả định trên là \textit{cluster assumption}, cho rằng dữ liệu tự nhiên có xu hướng hình thành các cụm trong không gian đặc trưng và ranh giới quyết định của mô hình nên đi qua vùng có mật độ dữ liệu thấp thay vì cắt xuyên qua các cụm dữ liệu. Giả định này dẫn đến \textit{low-density separation assumption}, theo đó một bộ phân loại tốt sẽ đặt decision boundary tại những vùng ít mẫu, giúp các mẫu trong cùng một cụm được gán cùng nhãn. Trong Semi-supervised Learning, pseudo-labeling và consistency learning đều khai thác các giả định này để mở rộng tín hiệu giám sát từ tập dữ liệu có nhãn sang tập dữ liệu chưa gán nhãn.

Bên cạnh các giả định về cấu trúc dữ liệu, Semi-supervised Learning còn dựa trên \textit{same-distribution assumption}, tức dữ liệu có nhãn và dữ liệu chưa gán nhãn được giả định đến từ cùng một phân phối xác suất hoặc ít nhất có cùng không gian đặc trưng và cùng tập lớp cần nhận diện. Nếu tập dữ liệu chưa gán nhãn có phân phối khác đáng kể so với tập dữ liệu có nhãn, các pseudo-label do teacher sinh ra có thể không còn đáng tin cậy và việc khai thác dữ liệu chưa gán nhãn có thể làm giảm hiệu quả của detector.

Trong nghiên cứu này, dữ liệu có nhãn và dữ liệu chưa gán nhãn đều được lấy từ cùng bộ dữ liệu VinDr-CXR, cùng modality ảnh X-quang ngực và cùng không gian lớp mục tiêu. Vì vậy, nghiên cứu xem giả định đồng phân phối là hợp lý ở mức thiết kế thực nghiệm. Tuy nhiên, điều này không đồng nghĩa với việc phân phối dữ liệu hoàn toàn đồng nhất. Những khác biệt về mức độ bệnh, chất lượng ảnh, thiết bị chụp hoặc tần suất xuất hiện của từng bất thường vẫn có thể tạo ra sai khác giữa các tập dữ liệu. Do đó, nghiên cứu không giả định rằng same-distribution assumption luôn được thỏa mãn tuyệt đối mà chỉ xem đây là giả định hợp lý để xây dựng giao thức thực nghiệm.

Một rủi ro quan trọng của Semi-supervised Learning là \textit{distribution shift}, xảy ra khi dữ liệu chưa gán nhãn khác đáng kể so với dữ liệu có nhãn. Khi đó, teacher detector có thể tạo ra pseudo-label không phản ánh đúng phân phối thực tế của dữ liệu, làm cho student tiếp tục học từ các dự đoán sai. Hiện tượng này đặc biệt đáng lưu ý trong ảnh y tế khi dữ liệu được thu thập từ các bệnh viện, thiết bị hoặc quần thể bệnh nhân khác nhau.

Ngay cả khi dữ liệu được thu thập từ cùng một nguồn, Semi-supervised Learning vẫn có thể gặp hiện tượng \textit{confirmation bias}. Trong cơ chế teacher--student, student học trực tiếp từ pseudo-label do teacher tạo ra. Nếu teacher dự đoán sai ở giai đoạn đầu, các lỗi này có thể tiếp tục được củng cố qua nhiều vòng huấn luyện, khiến mô hình ngày càng tự tin vào các dự đoán không chính xác \cite{sohn2020simple,liu2021unbiased,xu2021end}. Confirmation bias được xem là một trong những thách thức lớn nhất của các phương pháp pseudo-labeling.

Liên quan trực tiếp đến confirmation bias là \textit{pseudo-label noise}. Không phải mọi pseudo-label đều phản ánh chính xác vị trí và lớp của bất thường. Một pseudo-label có thể sai về vị trí bounding box, sai lớp hoặc bỏ sót tổn thương nhỏ. Nếu các pseudo-label này không được lọc phù hợp thông qua confidence threshold hoặc các chiến lược kiểm soát chất lượng khác, student detector có thể học từ tín hiệu giám sát bị nhiễu, dẫn đến suy giảm hiệu năng.

Một rủi ro khác là \textit{rare-class bias}. Các lớp bất thường hiếm thường xuất hiện với số lượng bounding box rất ít trong tập dữ liệu có nhãn. Teacher detector vì vậy có xu hướng tạo pseudo-label đáng tin cậy hơn cho các lớp phổ biến, trong khi các lớp hiếm dễ bị bỏ sót hoặc dự đoán sai. Điều này có thể làm gia tăng mất cân bằng lớp trong quá trình huấn luyện Semi-supervised Learning, mặc dù lượng dữ liệu chưa gán nhãn được bổ sung nhiều hơn.

Trong bộ dữ liệu VinDr-CXR, ảnh \textit{No Finding} được xem là ảnh âm tính và không có bounding box bất thường. Do đó, một rủi ro cần được xem xét là \textit{false positive trên ảnh No Finding}, tức detector dự đoán các bounding box không tồn tại trên những ảnh không có tổn thương. Nếu teacher sinh ra pseudo-label sai trên các ảnh này, student có thể tiếp tục học các tín hiệu không đúng, làm tăng số lượng false positive và ảnh hưởng đến độ tin cậy của hệ thống trong thực tế lâm sàng.

Những giả định và rủi ro trên cho thấy Semi-supervised Learning không phải luôn mang lại cải thiện so với supervised learning. Hiệu quả của phương pháp phụ thuộc vào mức độ các giả định được thỏa mãn cũng như chất lượng của pseudo-label trong từng điều kiện dữ liệu cụ thể. Vì vậy, thay vì giả định rằng dữ liệu chưa gán nhãn luôn giúp cải thiện detector, luận văn xây dựng một giao thức thực nghiệm có kiểm soát để đánh giá khi nào Semi-supervised Learning phát huy hiệu quả, khi nào xuất hiện các hiện tượng như confirmation bias, pseudo-label noise hoặc rare-class bias, và chúng ảnh hưởng như thế nào đến hiệu năng phát hiện bất thường.


\subsection{Bài toán nghiên cứu trên VinDr-CXR}

Nghiên cứu sử dụng bộ dữ liệu \textbf{VinDr-CXR}, được công bố và sử dụng trong bài toán VinBigData Chest X-ray Abnormalities Detection, làm nguồn dữ liệu chính \cite{nguyen2022vindr}. Bộ dữ liệu này phù hợp với mục tiêu nghiên cứu vì cung cấp các chú thích bounding box cho nhiều loại bất thường trên ảnh X-quang ngực, qua đó cho phép xây dựng và đánh giá mô hình object detection đa lớp.

Theo kết quả phân tích dữ liệu ban đầu, tập huấn luyện gốc gồm 15,000 ảnh với 67,914 dòng chú thích. Trong đó, 4,394 ảnh chứa ít nhất một bất thường, 10,606 ảnh được gán nhãn \textit{No Finding}, và có tổng cộng 36,096 bounding box thuộc 14 lớp bất thường cần phát hiện. Các bước kiểm tra dữ liệu cho thấy không có bounding box vi phạm điều kiện hình học cơ bản, không có dòng \textit{No Finding} chứa tọa độ bounding box, và không có ảnh đồng thời mang nhãn \textit{No Finding} và nhãn bất thường.

Trong phạm vi thực nghiệm có kiểm soát, nghiên cứu sử dụng một working scope gồm 4,894 ảnh, bao gồm toàn bộ 4,394 ảnh bất thường và 500 ảnh \textit{No Finding} được chọn làm ảnh âm tính. Việc giới hạn số lượng ảnh \textit{No Finding} nhằm bảo đảm tính khả thi về tài nguyên tính toán và tránh để số lượng ảnh âm tính lớn trong phân bố gốc chi phối quá mức working scope. Vì vậy, working scope này được xem là một thiết lập thực nghiệm có kiểm soát, không phải là sự tái tạo đầy đủ phân bố của toàn bộ dataset gốc.

Trong nghiên cứu này, \textit{No Finding} không được xem là một lớp cần phát hiện và không được đưa vào danh sách các detection categories. Ảnh \textit{No Finding} được giữ lại dưới dạng negative images, tức là ảnh không chứa bounding box bất thường. Chính sách này cho phép đánh giá riêng xu hướng detector sinh false positive trên các ảnh không có tổn thương, thay vì chỉ kết luận dựa trên mAP tổng thể.

Thiết kế thực nghiệm tập trung so sánh supervised object detection với semi-supervised object detection theo cơ chế teacher--student pseudo-labeling. Ngân sách dữ liệu có nhãn được định nghĩa ở cấp độ ảnh (\textit{image level}), với bốn mức 1\%, 5\%, 10\% và 20\% của tập huấn luyện cố định. Các tập ảnh có nhãn được xây dựng theo nguyên tắc nested subset $1\% \subset 5\% \subset 10\% \subset 20\%$ trên cùng một phân chia train, validation và test cố định. Ở mỗi mức ngân sách nhãn, supervised baseline và mô hình semi-supervised sử dụng cùng tập ảnh có nhãn, cùng detector, cùng quy trình tiền xử lý và cùng giao thức đánh giá. Khác biệt chính là mô hình semi-supervised được phép khai thác thêm phần dữ liệu chưa gán nhãn trong quá trình huấn luyện.

Hiệu năng mô hình được đánh giá chủ yếu bằng các metric dành cho object detection, bao gồm bbox mAP@[0.50:0.95], AP50, AP75 và AP theo từng lớp. Ngoài các metric tổng hợp, nghiên cứu còn phân tích phân bố pseudo-label theo lớp, hành vi của các lớp hiếm và false positive trên ảnh \textit{No Finding}. Accuracy và AUC ở mức toàn ảnh không được sử dụng làm metric chính vì chúng không phản ánh đồng thời chất lượng phân loại và định vị bounding box.

Bên cạnh hiệu năng mô hình, nghiên cứu cũng xem xét chất lượng annotation như một đặc điểm của dữ liệu. Phân tích Fleiss' Kappa được sử dụng để mô tả mức độ đồng thuận nhãn giữa nhiều bác sĩ đọc ảnh, trong khi tính nhất quán của bounding box được đánh giá ở mức mô tả. Các phân tích này đóng vai trò làm rõ giới hạn và độ nhiễu của dữ liệu, không được sử dụng như metric đánh giá mô hình detection.


\subsection{Câu hỏi nghiên cứu và đóng góp dự kiến}
\label{subsec:research_questions_contributions}

Từ bài toán chi phí cao của việc gán nhãn bounding box trên ảnh X-quang ngực,
nghiên cứu tập trung đánh giá khả năng khai thác dữ liệu chưa gán nhãn bằng
semi-supervised object detection trong một thiết kế thực nghiệm có kiểm soát.
Thay vì chỉ xem xét hiệu năng cuối cùng của một mô hình đơn lẻ, nghiên cứu được
xây dựng nhằm phân tích đồng thời ảnh hưởng của ngân sách nhãn, phương thức học
supervised hoặc semi-supervised, chất lượng pseudo-label, hành vi trên các lớp
hiếm, false positive trên ảnh âm tính và sự phụ thuộc của lợi ích semi-supervised
learning vào kiến trúc mô hình.

Trên cơ sở đó, nghiên cứu đặt ra bảy câu hỏi nghiên cứu sau:

\begin{itemize}

    \item \textbf{RQ1.}

    Hiệu năng của supervised multi-class object detection thay đổi như thế nào khi ngân sách dữ liệu có nhãn được giới hạn ở các mức 1\%, 5\%, 10\% và 20\% của tập huấn luyện (training set) được cố định?

    \item \textbf{RQ2.}
    Trong các điều kiện được kiểm soát về tập dữ liệu có nhãn, kiến trúc mô hình,
    phân chia dữ liệu, training seed và giao thức đánh giá, teacher--student
    semi-supervised object detection dựa trên pseudo-labeling có cải thiện hiệu
    năng phát hiện so với supervised baseline hay không?

    \item \textbf{RQ3.}
    Mức cải thiện của semi-supervised learning so với supervised baseline
    có thay đổi giữa các ngân sách nhãn 1\%, 5\%, 10\% và 20\% hay không?

    \item \textbf{RQ4.}
    Chất lượng pseudo-label có mối liên hệ như thế nào với mức cải thiện hiệu năng
    của semi-supervised learning trong các điều kiện kiến trúc và ngân sách nhãn
    đã xác định trước?

    \item \textbf{RQ5.}
    Các lớp bất thường có tần suất xuất hiện thấp trong tập huấn luyện nhận được
    lợi ích như thế nào từ semi-supervised learning so với supervised learning?

    \item \textbf{RQ6.}
    Semi-supervised learning có làm thay đổi mức false positive trên các ảnh
    \textit{No Finding}, được xem là các ảnh âm tính không chứa bounding box bất
    thường, so với supervised baseline hay không?

    \item \textbf{RQ7.}
    Mức lợi ích của semi-supervised learning so với supervised baseline có khác
    nhau giữa kiến trúc CNN truyền thống dựa trên ResNet-50 và kiến trúc Vision Transformer phân cấp dựa trên Swin-T hay không?

\end{itemize}

Các câu hỏi trên được tổ chức theo vai trò khác nhau trong thiết kế nghiên cứu.
RQ1 có tính mô tả, nhằm xác định sự thay đổi của supervised baseline khi ngân
sách nhãn giảm và không gắn với một giả thuyết thống kê khẳng định riêng.
RQ2 là câu hỏi nghiên cứu chính, trực tiếp kiểm tra liệu việc khai thác dữ liệu
chưa gán nhãn bằng semi-supervised learning có tạo ra lợi ích so với supervised
learning hay không. RQ3, RQ4, RQ6 và RQ7 là các câu hỏi nghiên cứu thứ cấp có phân tích
chính thức, lần lượt làm rõ sự phụ thuộc của mức lợi ích semi-supervised learning vào ngân sách nhãn, mối liên hệ với chất lượng pseudo-label, hành vi false positive trên ảnh âm tính và sự khác biệt theo kiến trúc mô hình. RQ5 được xem là phân tích khám phá do số lượng mẫu của các lớp
hiếm thấp hơn đáng kể so với các lớp phổ biến. Các giả thuyết nghiên cứu và phương pháp kiểm định tương ứng được trình bày chi tiết trong phần giao thức thực nghiệm.

Nghiên cứu được định vị là một \textbf{controlled experimental study về label efficiency} trong bài toán
semi-supervised multi-class object detection với điều kiện hạn chế nhãn
bounding box trên ảnh X-quang ngực. Nghiên cứu không nhằm đề xuất một detector
mới hoặc một thuật toán semi-supervised object detection mới. Thay vào đó, mục
tiêu là xây dựng một giao thức thực nghiệm được xác định trước để so sánh
supervised và semi-supervised learning trong các điều kiện có thể đối chiếu,
đồng thời phân tích những yếu tố có thể giải thích sự khác biệt về hiệu năng.

Các đóng góp dự kiến của nghiên cứu gồm năm nhóm chính.

Thứ nhất, nghiên cứu xây dựng một giao thức thực nghiệm có kiểm soát cho
semi-supervised multi-class object detection trên VinDr-CXR. Giao thức sử
dụng một phân chia train/validation/test cố định, các \textit{labeled subset}
lồng nhau theo quan hệ
$1\% \subset 5\% \subset 10\% \subset 20\%$, cùng chính sách training seed
được xác định trước và cùng giao thức đánh giá. Thiết kế này tạo cơ sở để
phân biệt ảnh hưởng của chiến lược học, ngân sách nhãn và kiến trúc khỏi
các nguồn biến thiên thực nghiệm không được kiểm soát.

Thứ hai, nghiên cứu thực hiện so sánh ghép cặp giữa supervised baseline và
semi-supervised teacher--student pseudo-labeling tại từng ngân sách nhãn và
trên từng kiến trúc. Trong mỗi cặp so sánh, hai phương pháp sử dụng cùng
tập dữ liệu có nhãn, cùng kiến trúc, cùng training seed và cùng ngân sách
cập nhật tham số; khác biệt có chủ đích giữa hai phương thức học là việc mô hình semi-supervised được phép khai thác thêm tập dữ liệu chưa gán nhãn. Qua đó, nghiên cứu đánh giá trực tiếp phần lợi ích gia tăng gắn với semi-supervised
learning thay vì chỉ so sánh các mô hình được huấn luyện trong những điều
kiện khác nhau.

Thứ ba, nghiên cứu mở rộng việc đánh giá vượt ra ngoài các chỉ số detection
tổng hợp bằng cách phân tích chất lượng pseudo-label và mối liên hệ giữa chất lượng pseudo-label và mức cải thiện hiệu năng phát hiện của mô hình semi-supervised so với supervised baseline. Chất lượng pseudo-label
được xác định trước bằng
$PL$-$mAP@[0.50:0.95]$ trên tập validation cố định, sử dụng EMA Teacher tại
checkpoint LAST. Việc đánh giá này được thực hiện trên validation thay vì sử
dụng ground truth ẩn của các tập dữ liệu chưa gán nhãn, qua đó bảo toàn nguyên
tắc không khai thác nhãn ẩn của dữ liệu unlabeled trong quá trình phát triển
phương pháp.

Thứ tư, nghiên cứu mở rộng đánh giá ra ngoài hiệu năng tổng thể bằng cách
phân tích riêng các lớp bất thường có mức hỗ trợ thấp và các ảnh
\textit{No Finding}. Các lớp hiếm được đánh giá ở mức từng lớp, trong khi
\textit{No Finding} được xem là nhóm ảnh âm tính có zero ground-truth
bounding box chứ không phải một detection class. Việc đo false positive
trên nhóm ảnh này bổ sung một khía cạnh độ tin cậy mà các chỉ số mAP tổng
thể không phản ánh đầy đủ.

Thứ năm, nghiên cứu đưa kiến trúc mô hình vào thiết kế thực nghiệm như một yếu
tố được xác định trước. Hai kiến trúc được khảo sát gồm Faster R-CNN với ResNet-50 và FPN, đại diện cho kiến trúc CNN truyền thống, và Faster R-CNN với Swin-T và FPN, đại diện cho Vision Transformer phân cấp sử dụng
shifted-window self-attention. Việc sử dụng cùng họ detector và FPN giúp hạn chế
confounding do thay đổi detector paradigm, trong khi cho phép khảo sát liệu mức
lợi ích của semi-supervised learning có phụ thuộc vào kiến trúc hay không. Mục
tiêu của phân tích này không phải chứng minh Swin-T vượt trội hơn ResNet-50 một
cách tổng quát, mà đánh giá \textit{architecture-dependent SSL effect} dưới các
giao thức huấn luyện đặc thù theo kiến trúc đã được xác định trước.

\section{Related Work}\label{sec:related}
 
% ---------------------------------------------------------------------
\subsection{Bộ dữ liệu CXR và mức độ chi tiết của chú thích}
\label{sec:cxr_datasets_annotation_granularity}
 
Các bộ dữ liệu ảnh X-quang ngực có thể được phân nhóm theo mức độ chi tiết của
chú thích (\textit{annotation granularity}), và cách phân nhóm này quyết định
trực tiếp loại bài toán có thể thực hiện, dạng tín hiệu giám sát và hệ metric
đánh giá. Ba nhóm chính gồm: nhãn ở mức toàn ảnh phục vụ classification; chú
thích định vị yếu hoặc chỉ trên một phần dữ liệu; và bounding box cho từng vùng
bất thường phục vụ object detection.
 
Ở nhóm nhãn mức ảnh, các bộ dữ liệu quy mô lớn tiêu biểu gồm ChestX-ray14 (nhãn
bệnh mức ảnh kèm một tập con định vị) \cite{wang2017chestx}, CheXpert (nhãn mức
ảnh có trạng thái không chắc chắn) \cite{irvin2019chexpert} và MIMIC-CXR (ảnh
kèm báo cáo văn bản đã khử định danh) \cite{johnson2019mimic}; nhiều nghiên cứu
classification đa nhãn, long-tailed cũng khai thác các nguồn này cùng các chiến
lược augmentation, ensemble và xử lý mất cân bằng lớp \cite{nguyen2023advanced}.
Điểm chung của nhóm này là chỉ cho biết một bất thường có xuất hiện hay không mà
không cung cấp vị trí; do đó các metric như Accuracy hay AUC không phản ánh chất
lượng định vị, và một mô hình phân loại tốt vẫn có thể bỏ sót tổn thương nhỏ khi
được yêu cầu trả về bounding box.
 
Ở nhóm có tín hiệu định vị, mức độ chi tiết tăng dần. Tập con bounding box của
ChestX-ray14 thường được dùng để đánh giá weak localization, trong đó mô hình
học từ nhãn mức ảnh nhưng chỉ được kiểm tra vị trí trên một phần dữ liệu
\cite{wang2017chestx}; tín hiệu định vị vì vậy vẫn không đầy đủ. Ở mức chú thích
đầy đủ, ChestX-Det10 (xây từ NIH ChestX-ray14; 3.543 ảnh, 10 nhóm bất thường;
Faster R-CNN + FPN đạt AP50 $=$ 40.2 và recall@0.1 FP/image $=$ 0.436 trên 542
ảnh) \cite{liu2020chestx} đại diện cho sự chuyển dịch sang đánh giá localization
trực tiếp, nhưng quy mô và số lớp còn nhỏ. VinDr-CXR (18.000 ảnh: 15.000 train /
3.000 test; nhãn global mức ảnh và bounding box; train đọc độc lập bởi ba bác sĩ,
test đồng thuận năm bác sĩ) \cite{nguyen2022vindr} là bộ dữ liệu bounding box đa
lớp phù hợp trực tiếp nhất với nghiên cứu này, đồng thời bộc lộ các vấn đề đặc
thù như bất nhất giữa annotator, nhiễu nhãn, phân phối long-tailed và cách xử lý
ảnh \textit{No Finding}.
 
\begin{table*}[htbp]
\caption{So sánh các bộ dữ liệu CXR theo mức độ chi tiết của chú thích}
\label{tab:cxr_datasets_annotation_granularity}
\centering
\scriptsize
\setlength{\tabcolsep}{3pt}
\renewcommand{\arraystretch}{1.25}
% Tổng trọng số các cột X phải bằng số cột X (ở đây = 5)
\newcolumntype{s}[1]{>{\hsize=#1\hsize\raggedright\arraybackslash}X}
\begin{tabularx}{\textwidth}{
  s{0.80} % Bộ dữ liệu
  s{1.15} % Loại chú thích
  c       % Bounding box
  s{1.00} % Bài toán chính
  s{1.05} % Metric điển hình
  s{1.00} % Vai trò
  c       % Tài liệu
}
\toprule
\textbf{Bộ dữ liệu} & \textbf{Loại chú thích} & \textbf{BBox} &
\textbf{Bài toán chính} & \textbf{Metric điển hình} &
\textbf{Vai trò với nghiên cứu} & \textbf{Ref.} \\
\midrule
ChestX-ray14 / NIH & Nhãn mức ảnh; tập con có bounding box & Hạn chế &
Classification, weak localization & AUC, AP classification, localization &
Bối cảnh nhãn mức ảnh, định vị yếu & \cite{wang2017chestx} \\
\addlinespace[2pt]
CheXpert & Nhãn mức ảnh, có trạng thái không chắc chắn & Không &
Multi-label classification & AUC, metric classification &
Bối cảnh classification, uncertainty labels & \cite{irvin2019chexpert} \\
\addlinespace[2pt]
MIMIC-CXR & Ảnh CXR kèm báo cáo văn bản tự do & Không &
Classification, representation learning & Phụ thuộc bài toán hạ nguồn &
Nguồn lớn cho học biểu diễn, khai thác báo cáo & \cite{johnson2019mimic} \\
\addlinespace[2pt]
ChestX-Det10 & Bounding box cho 10 nhóm bất thường & Có &
Thoracic abnormality detection & AP50, recall@FP/image &
Benchmark detection CXR liên quan & \cite{liu2020chestx} \\
\addlinespace[2pt]
VinDr-CXR & Bounding box + nhãn global mức ảnh & Có &
Multi-class abnormality detection & bbox mAP@[0.50:0.95], AP50/75, AP theo lớp &
Bộ dữ liệu trung tâm & \cite{nguyen2022vindr} \\
\bottomrule
\end{tabularx}
\end{table*}
 
Tổng hợp lại, annotation granularity tạo ra một sự đánh đổi rõ ràng giữa khả
năng định vị và chi phí tạo nhãn: dữ liệu mức ảnh rẻ nhưng không đánh giá được
localization, còn dữ liệu bounding box như ChestX-Det10 và VinDr-CXR cho phép
detection trực tiếp nhưng đòi hỏi định vị và gán lớp cho từng tổn thương. VinDr-CXR
vừa là điều kiện phù hợp cho detection đa lớp, vừa phản ánh rõ nút thắt phụ thuộc
bounding box đầy đủ — nút thắt mà phần lớn nghiên cứu detection sau đây vẫn mặc
định chấp nhận.
 
% ---------------------------------------------------------------------
\subsection{Supervised object detection trên ảnh X-quang ngực}
\label{sec:supervised_cxr_detection}
 
Khi các bộ dữ liệu bounding box trở nên phổ biến, phần lớn nghiên cứu phát hiện
bất thường trên CXR đi theo hướng supervised và cải thiện hiệu năng qua ba xu
hướng kiến trúc kế tiếp nhau: (i) detector và feature pyramid, (ii) attention và
multi-scale feature fusion, và (iii) detector hiệu quả với mô hình hóa ngữ cảnh
toàn cục. Điểm chung xuyên suốt là các cải tiến đều nhằm nâng độ chính xác định
vị trong khi giả định dữ liệu huấn luyện đã được gán đầy đủ bounding box.
 
\subsubsection*{Detector cơ bản và Feature Pyramid Network}
Giai đoạn đầu chủ yếu điều chỉnh các detector tổng quát cho đặc thù CXR nhằm xử
lý tổn thương đa kích thước qua backbone và feature pyramid. Hai nghiên cứu tiêu
biểu gồm CXR-RefineDet, mở rộng RefineDet bằng RRNet và Information Reuse để tăng
khai thác đặc trưng đa mức (mAP $\approx$ 0.1618; 6.8 FPS) \cite{lin2022cxr}, và
baseline RetinaNet + ResNet101 + FPN kết hợp Weighted Boxes Fusion (mAP@0.5
$\approx$ 0.22 cho toàn bộ 14 lớp và $\approx$ 0.55 khi chỉ xét năm lớp phổ biến)
\cite{ngo2023application}. Nhóm này khẳng định vai trò nền tảng của feature pyramid
trong xử lý biến thiên kích thước tổn thương, nhưng cải thiện chỉ giới hạn ở khả
năng định vị khi đã có đầy đủ nhãn.
 
\subsubsection*{Attention mechanism và multi-scale feature fusion}
Trước hạn chế về độ tương phản thấp và kích thước nhỏ của tổn thương, một nhóm
nghiên cứu tập trung tăng cường biểu diễn đặc trưng thay vì thay đổi toàn bộ kiến
trúc detector. Theo hướng này các công trình tiêu biểu gồm SAR-CNN, kết hợp multi-scale feature
fusion và two-stage attention (tỷ lệ phát hiện tăng từ 12.83\% lên 15.75\% tại
IoU $>$ 0.4) \cite{lin2023lesion}; mô hình dual-attention kết hợp Augmented-FPN và
Seesaw Loss để xử lý đồng thời mất cân bằng lớp và biến thiên kích thước (mAP
$=$ 0.362) \cite{liu2023anomaly}; DualAttNet kết hợp attention mức toàn ảnh và
attention tinh chỉnh theo bệnh, tuy cải thiện hiệu năng nhưng vẫn khó với tổn
thương rất nhỏ \cite{xu2024dualattnet}; và MLLA bổ sung location-aware learning
vào multi-scale fusion (mAP $=$ 0.471; mR $=$ 0.711) \cite{yuan2024multi}. Dù khác
nhau về cơ chế attention hay fusion cụ thể, điểm chung của nhóm này là đều khắc
phục hạn chế \emph{biểu diễn} của detector, không nhằm giảm nhu cầu dữ liệu có nhãn.
 
\subsubsection*{Các detector hiệu quả và mô hình hóa ngữ cảnh toàn cục}
Xu hướng gần đây hướng tới cân bằng giữa độ chính xác định vị và hiệu quả tính
toán, đồng thời mở rộng khả năng mô hình hóa ngữ cảnh toàn cục. Các nghiên cứu
tiêu biểu gồm YOLO-SG, khai thác salience-guided learning cho đối tượng nhỏ
\cite{han2022yolo}; YOLO-CXR, bổ sung xử lý đặc trưng đa mức cho phát hiện đa tổn
thương (mAP@0.5 $=$ 0.338) \cite{hao2024yolo}; Mamba-YOLOvX, kết hợp state-space
modeling để nắm quan hệ không gian dài hạn (AP50 $=$ 0.366) \cite{khalili2025localization};
và NSEC-YOLO, kết hợp adaptive noise suppression với global perception aggregation
(mAP@0.5 $=$ 0.416; $\approx$ 163 FPS) \cite{zhang2025nsec}. Nhóm này phản ánh sự
chuyển dịch từ chỉ cải thiện đặc trưng cục bộ sang khai thác đồng thời đa tỉ lệ
và ngữ cảnh toàn cục.
 
\begin{table*}[htbp]
\caption{Tổng hợp các xu hướng supervised object detection trên ảnh X-quang ngực
theo đóng góp và hạn chế chung của từng nhóm.}
\label{tab:supervised_detection_trends}
\centering
\scriptsize
\setlength{\tabcolsep}{3pt}
\renewcommand{\arraystretch}{1.25}
% Tổng trọng số các cột X = số cột X (= 4)
\newcolumntype{s}[1]{>{\hsize=#1\hsize\raggedright\arraybackslash}X}
\begin{tabularx}{\textwidth}{
  >{\raggedright\arraybackslash}m{1.4cm} % Giai đoạn
  s{1.00} % Xu hướng nghiên cứu
  s{1.15} % Nghiên cứu đại diện
  s{1.00} % Đóng góp chính
  s{0.85} % Hạn chế chung
}
\toprule
\textbf{Giai đoạn} & \textbf{Xu hướng nghiên cứu} & \textbf{Nghiên cứu đại diện}
& \textbf{Đóng góp chính} & \textbf{Hạn chế chung} \\
\midrule
2022--2023 & Detector cơ bản và Feature Pyramid Network &
CXR-RefineDet \cite{lin2022cxr}; RetinaNet + FPN + WBF \cite{ngo2023application} &
Điều chỉnh detector tổng quát cho CXR, cải thiện biểu diễn đa tỉ lệ qua backbone,
FPN và hậu xử lý. &
Phụ thuộc hoàn toàn vào bounding box đầy đủ; chưa giảm chi phí annotation. \\
\addlinespace[2pt]
2023--2024 & Attention và Multi-scale Feature Fusion &
SAR-CNN \cite{lin2023lesion}; Dual Attention + Aug-FPN \cite{liu2023anomaly};
DualAttNet \cite{xu2024dualattnet}; MLLA \cite{yuan2024multi} &
Tăng khả năng tập trung vào vùng bất thường, xử lý tổn thương nhỏ và độ tương phản thấp. &
Tập trung vào khả năng biểu diễn; vẫn giả định có đầy đủ bounding box. \\
\addlinespace[2pt]
2022--2025 & Detector hiệu quả và ngữ cảnh toàn cục &
YOLO-SG \cite{han2022yolo}; YOLO-CXR \cite{hao2024yolo};
Mamba-YOLOvX \cite{khalili2025localization}; NSEC-YOLO \cite{zhang2025nsec} &
Cân bằng tốc độ--độ chính xác bằng kiến trúc YOLO và mô hình hóa ngữ cảnh toàn cục. &
Vẫn phụ thuộc dữ liệu có nhãn đầy đủ; chưa đánh giá khi ngân sách nhãn bounding box bị giới hạn. \\
\bottomrule
\end{tabularx}
\end{table*}
 
\paragraph{Tổng hợp và nhận xét}
Ba xu hướng trên cho thấy hiệu năng detection được cải thiện dần từ feature
pyramid, qua attention/fusion, đến YOLO/Mamba và ngữ cảnh toàn cục. Tuy nhiên,
các con số mAP/AP giữa những công trình này \emph{không so sánh trực tiếp được}
do khác nhau về tập dữ liệu, số lớp, độ phân giải và tiêu chuẩn metric (ví dụ
Detection Rate, AP50 hay bbox mAP@[0.50:0.95]); vì vậy chúng chỉ nên được hiểu như bằng
chứng về xu hướng, không phải bảng xếp hạng. Quan trọng hơn, dù khác nhau về kiến
trúc, tất cả đều chia sẻ một giả định: tập huấn luyện đã được gán đầy đủ bounding
box. Nói cách khác, các nghiên cứu này trả lời câu hỏi \textit{``làm thế nào để
xây dựng detector mạnh hơn''} chứ chưa trả lời \textit{``làm thế nào để giảm phụ
thuộc vào dữ liệu có nhãn''}. Hạn chế chung này dẫn tới hướng semi-supervised
object detection được tổng hợp trong Mục~\ref{sec:ssod}.
 
% ---------------------------------------------------------------------
\subsection{Semi-supervised object detection sử dụng dữ liệu chưa gán nhãn}
\label{sec:ssod}
 
Semi-supervised object detection (SSOD) nhằm giảm phụ thuộc vào dữ liệu có nhãn
bằng cách khai thác đồng thời một tập nhỏ dữ liệu có nhãn và một tập lớn dữ liệu
chưa gán nhãn, thay vì thay đổi kiến trúc detector. Trong phạm vi luận văn, SSOD
được hiểu theo thiết lập \textbf{unlabeled-only}: phần dữ liệu ngoài tập có nhãn
không sử dụng bất kỳ tín hiệu chú thích bổ sung nào (image-level, point hay weak
annotation) mà chỉ được khai thác qua cơ chế teacher--student và pseudo-labeling.
 
\subsubsection*{Teacher--student và pseudo-labeling trên ảnh tự nhiên}
Khung teacher--student sinh pseudo-label (bounding box, lớp, confidence) trên dữ
liệu chưa gán nhãn, lọc theo ngưỡng rồi dùng làm giám sát bổ sung cho student.
STAC đặt nền tảng cho hướng này bằng pseudo-label độ tin cậy cao kết hợp strong
augmentation, cải thiện đáng kể so với supervised baseline ở tỷ lệ nhãn thấp
(MS-COCO 5\%: 18.47 $\rightarrow$ 24.38; 10\%: 23.86 $\rightarrow$ 28.64)
\cite{sohn2020simple}. Các nghiên cứu sau đó không xây pipeline mới mà chủ yếu
ổn định teacher và nâng chất lượng pseudo-label; tiêu biểu gồm Unbiased Teacher
(EMA teacher + class-balance loss) \cite{liu2021unbiased}, Soft Teacher (trọng số
mềm theo confidence, huấn luyện end-to-end) \cite{xu2021end}, ISMT (interactive
self-training với mean teacher) \cite{yang2021interactive}, Instant-Teaching
(giảm độ trễ giữa sinh pseudo-label và cập nhật student) \cite{feng2021instant},
và DTG-SSOD (mở rộng sang dense teacher guidance thay vì chỉ pseudo-box đã lọc)
\cite{li2022dtg}. Điểm chung của nhóm này là cùng hướng tới giảm confirmation
bias và khai thác dữ liệu chưa gán nhãn hiệu quả hơn, nhưng gần như toàn bộ bằng
chứng đến từ ảnh tự nhiên (MS-COCO, PASCAL VOC).
 
\subsubsection*{Chuyển giao SSOD sang ảnh y tế}
Một số nghiên cứu đã mở rộng teacher--student sang ảnh y tế, nhưng chủ yếu ở các
modality và nhiệm vụ khác CXR 2D đa lớp. Tiêu biểu gồm FocalMix cho phát hiện nốt
phổi trên CT 3D, khai thác unlabeled scans và cải thiện đáng kể chỉ số CPM (metric
FROC/CPM, một loại tổn thương) \cite{wang2020focalmix}, và Robust Mean Teacher cho
phát hiện tế bào trên ảnh mô bệnh học, tập trung kiểm soát pseudo-label noise
\cite{wen2024robust}. Hai nghiên cứu này xác nhận dữ liệu chưa gán nhãn có thể hỗ trợ
detection y tế, nhưng khác biệt về modality, số lớp, phân bố không gian và metric
khiến kết quả chưa thể suy rộng sang bài toán đa bất thường trên CXR.
 
\begin{table*}[!ht]
\caption{Các xu hướng chính của semi-supervised object detection và mức độ liên
quan đối với bài toán ảnh X-quang ngực.}
\label{tab:ssod_research_trends}
\centering
\scriptsize
\begin{tabularx}{\textwidth}{
>{\raggedright\arraybackslash}m{2.0cm} >{\raggedright\arraybackslash}m{3.5cm}
>{\raggedright\arraybackslash}X >{\raggedright\arraybackslash}X
>{\raggedright\arraybackslash}X}
\toprule
\textbf{Giai đoạn} & \textbf{Xu hướng nghiên cứu} & \textbf{Nghiên cứu đại diện}
& \textbf{Đóng góp chính} & \textbf{Hạn chế đối với nghiên cứu này} \\
\midrule
2020 & Pseudo-labeling + strong augmentation & STAC \cite{sohn2020simple} &
Baseline SSOD với pseudo-label độ tin cậy cao trên dữ liệu chưa gán nhãn. &
Đánh giá trên MS-COCO/PASCAL VOC; chưa phản ánh đặc điểm CXR (tổn thương nhỏ,
long-tailed, ảnh \textit{No Finding}). \\
2021 & Ổn định teacher và cải thiện pseudo-label &
Unbiased Teacher \cite{liu2021unbiased}; Soft Teacher \cite{xu2021end};
ISMT \cite{yang2021interactive}; Instant-Teaching \cite{feng2021instant} &
EMA teacher, class-balance, soft weighting, self-training end-to-end để giảm bias
và nhiễu pseudo-label. &
Bằng chứng chủ yếu từ ảnh tự nhiên; chưa kiểm chứng trên CXR đa lớp. \\
2022 & Dense supervision từ teacher & DTG-SSOD \cite{li2022dtg} &
Mở rộng supervision từ pseudo-box sang dense teacher guidance. &
Nhạy với chất lượng teacher; chưa đánh giá trực tiếp trên VinDr-CXR. \\
2020--2024 & Chuyển giao SSOD sang ảnh y tế &
FocalMix \cite{wang2020focalmix}; Robust Mean Teacher \cite{wen2024robust} &
Xác nhận teacher--student và dữ liệu chưa gán nhãn hỗ trợ detection y tế, chú trọng
pseudo-label noise. &
Tập trung CT 3D / histopathology, khác modality, không gian lớp và metric so với
CXR 2D đa lớp. \\
\bottomrule
\end{tabularx}
\end{table*}
 
\paragraph{Tổng hợp và nhận xét}
Nhìn chung, SSOD phát triển từ việc chứng minh pseudo-labeling khai thác được dữ
liệu chưa gán nhãn, đến cải thiện độ ổn định teacher và chất lượng pseudo-label,
rồi bước đầu chuyển giao sang ảnh y tế. Tuy vậy, ba hạn chế chung vẫn tồn tại:
bằng chứng chủ yếu trên ảnh tự nhiên khác biệt lớn với CXR; các nghiên cứu y tế
tập trung ở CT/histopathology chứ không phải CXR đa lớp; và hiệu quả phụ thuộc
mạnh vào chất lượng pseudo-label, dễ phát sinh confirmation bias cùng pseudo-label
noise khi dữ liệu mất cân bằng. Do đó, hiệu quả của teacher--student pseudo-labeling
trên VinDr-CXR trong điều kiện ngân sách nhãn hạn chế vẫn chưa được kiểm chứng.
Ngoài SSOD unlabeled-only, ngay trên CXR còn có các hướng label-efficient khác,
được tổng hợp trong Mục~\ref{sec:cxr_label_efficiency}.
 
% ---------------------------------------------------------------------
\subsection{Các hướng label-efficient khác trên ảnh X-quang ngực}
\label{sec:cxr_label_efficiency}
 
Bên cạnh SSOD unlabeled-only, nhiều nghiên cứu trên CXR giảm chi phí annotation
bằng các nguồn tín hiệu bổ sung. Ba hướng chính là weak supervision, self-supervised
pretraining và học từ nhiều chuyên gia gán nhãn; điểm chung là đều dựa trên giả
định dữ liệu khác với thiết lập unlabeled-only của luận văn.
 
\subsubsection*{Weak supervision}
Hướng này thay bounding box đầy đủ bằng chú thích rẻ hơn để suy diễn vị trí tổn
thương. Nghiên cứu tiêu biểu là benchmark \textit{weakly semi-supervised
abnormality localization} của Ji và cộng sự trên VinDr-CXR, kết hợp một phần
bounding box với point annotations cho phần còn lại (mAP đạt 23.3 / 26.7 / 31.9 /
39.5 tương ứng 5\% / 10\% / 20\% / 100\% full labels) \cite{ji2022benchmark}. Tuy
cải thiện hiệu năng khi bounding box hạn chế, thiết lập này \emph{không phải
unlabeled-only} (vẫn dùng point-level annotations) và gộp một số lớp hiếm thành
\textit{Others}, khác với bài toán đầy đủ 14 lớp của luận văn.
 
\subsubsection*{Self-supervised pretraining}
Hướng này khai thác dữ liệu chưa gán nhãn ở giai đoạn tiền huấn luyện để tạo
backbone biểu diễn tốt hơn, thay vì trong quá trình tối ưu detector. BarlowTwins-CXR
là ví dụ tiêu biểu: self-supervised pretraining trên dữ liệu CXR lớn rồi fine-tune
Faster R-CNN + FPN trên VinDr-CXR, cải thiện tổng quát hóa so với khởi tạo thông
thường \cite{sheng2024barlowtwins}. Tuy nhiên, giai đoạn fine-tune vẫn cần bounding
box đầy đủ; dữ liệu chưa gán nhãn chỉ phục vụ học biểu diễn, không trực tiếp tham
gia huấn luyện detector như trong pseudo-labeling.
 
\subsubsection*{Học từ nhiều chuyên gia gán nhãn}
Thay vì giảm lượng nhãn, hướng này nâng chất lượng nhãn hiện có bằng cách xử lý
annotation disagreement. Le và cộng sự khai thác nhiều annotation cho cùng một ảnh,
kết hợp Weighted Boxes Fusion và re-weighting để giảm ảnh hưởng bất đồng giữa các
bác sĩ, giúp detector học ổn định hơn khi độ không chắc chắn cao \cite{le2023learning}.
Hướng này giả định đã tồn tại nhiều annotation và không khai thác dữ liệu chưa gán
nhãn, nên giải quyết một bài toán khác với luận văn.
 
\begin{table*}[!ht]
\caption{Các hướng label-efficient trên ảnh X-quang ngực ngoài SSOD unlabeled-only.}
\label{tab:cxr_label_efficient_methods}
\centering
\scriptsize
\begin{tabularx}{\textwidth}{
>{\raggedright\arraybackslash}m{2.3cm} >{\raggedright\arraybackslash}m{3.2cm}
>{\raggedright\arraybackslash}X >{\raggedright\arraybackslash}X
>{\raggedright\arraybackslash}X}
\toprule
\textbf{Hướng nghiên cứu} & \textbf{Nghiên cứu đại diện} &
\textbf{Nguồn tín hiệu bổ sung} & \textbf{Đóng góp chính} &
\textbf{Khác biệt so với luận văn} \\
\midrule
Weak supervision & Ji \textit{et al.} \cite{ji2022benchmark} &
Point annotations + một phần bounding box &
Giảm nhu cầu bounding box qua chú thích mức điểm. &
Không unlabeled-only; dữ liệu ngoài tập có nhãn vẫn có point labels, một số lớp bị gộp. \\
Self-supervised pretraining & BarlowTwins-CXR \cite{sheng2024barlowtwins} &
Dữ liệu chưa gán nhãn cho pretraining &
Học biểu diễn trước khi fine-tune detector. &
Ảnh chưa gán nhãn chỉ dùng ở pretraining; detector vẫn cần bounding box đầy đủ. \\
Multi-annotator learning & Le \textit{et al.} \cite{le2023learning} &
Nhiều annotation cho cùng một ảnh &
Xử lý disagreement để nâng chất lượng supervision. &
Không giảm lượng nhãn và không khai thác dữ liệu chưa gán nhãn. \\
\bottomrule
\end{tabularx}
\end{table*}
 
\paragraph{Tổng hợp và nhận xét}
Ba hướng trên cùng nhắm giảm phụ thuộc vào bounding box hoặc nâng chất lượng nhãn,
nhưng mỗi hướng dựa trên một giả định dữ liệu riêng: weak supervision cần chú thích
yếu sẵn có; self-supervised chỉ dùng dữ liệu chưa gán nhãn cho pretraining; còn
multi-annotator learning cần nhiều nhãn cho cùng một ảnh. Không hướng nào khai thác
trực tiếp ảnh X-quang thực \emph{hoàn toàn chưa gán nhãn} ngay trong quá trình huấn
luyện detector. Vì vậy, trong phạm vi các tài liệu được khảo sát, chưa có nghiên cứu
nào đồng thời (i) dùng dữ liệu unlabeled-only, (ii) đánh giá trên phát hiện đa bất
thường của VinDr-CXR, và (iii) so sánh có kiểm soát supervised với semi-supervised
dưới nhiều mức ngân sách nhãn bounding box. Khoảng trống này được hệ thống hóa và định vị trong Mục~\ref{sec:research_gap}.
 
% ---------------------------------------------------------------------
\subsection{Khoảng trống nghiên cứu và định vị luận văn}
\label{sec:research_gap}
 
Qua các mục trên, bài toán phát hiện bất thường trên CXR đã tiến triển ở cả ba
khía cạnh dữ liệu, kiến trúc detector và học tiết kiệm nhãn, nhưng vẫn để lại ba
khoảng trống. Thứ nhất, các nghiên cứu supervised trên VinDr-CXR
\cite{lin2022cxr,ngo2023application,lin2023lesion,liu2023anomaly,xu2024dualattnet,yuan2024multi,hao2024yolo,khalili2025localization,zhang2025nsec}
liên tục nâng độ chính xác định vị nhưng đều giả định có đầy đủ bounding box, nên
chưa chạm tới bài toán giảm chi phí annotation. Thứ hai, các nghiên cứu SSOD
\cite{sohn2020simple,liu2021unbiased,xu2021end,yang2021interactive,feng2021instant,li2022dtg}
chứng minh hiệu quả của teacher--student pseudo-labeling nhưng chủ yếu trên ảnh tự
nhiên, còn phần y tế tập trung ở CT/histopathology \cite{wang2020focalmix,wen2024robust},
khác đáng kể so với CXR đa lớp. Thứ ba, các hướng label-efficient trên chính CXR
\cite{ji2022benchmark,sheng2024barlowtwins,le2023learning} đều dựa trên giả định
dữ liệu khác với thiết lập unlabeled-only. Trong phạm vi các tài liệu được khảo sát, chưa tìm thấy nghiên cứu nào đánh giá
teacher--student SSOD trên VinDr-CXR với thiết kế có kiểm soát dưới nhiều mức ngân
sách bounding box.
 
Luận văn được đặt vào khoảng trống này như một \textit{controlled experimental
study} thuộc hướng \textit{label-efficient learning}, không đề xuất kiến trúc
detector hay thuật toán SSOD mới. Thay vào đó, nghiên cứu giữ cố định detector,
tập train/validation/test, quy trình tiền xử lý và giao thức đánh giá; xây dựng các
tập nhãn theo nguyên tắc nested subsets $1\% \subset 5\% \subset 10\% \subset 20\%$,
phần còn lại xem là hoàn toàn chưa gán nhãn và chỉ khai thác qua teacher--student
learning. Thiết kế này cô lập ảnh hưởng của dữ liệu chưa gán nhãn, giảm nhiễu do
thay đổi kiến trúc hay thành phần dữ liệu giữa các thí nghiệm.
 
Bên cạnh so sánh supervised với semi-supervised qua bbox mAP@[0.50:0.95] và AP theo lớp,
luận văn phân tích các yếu tố chi phối hiệu quả pseudo-labeling — chất lượng
pseudo-label, hành vi lớp hiếm, ảnh hưởng của confidence threshold và xu hướng
false positive trên ảnh \textit{No Finding} — nhằm trả lời  \emph
liệu dữ liệu chưa gán nhãn giúp ích gì, mà còn trả lời câu hỏi \emph{khi nào} và \emph{vì sao}. Chất lượng pseudo-label được đánh giá thông qua các chỉ số như độ chính xác của pseudo-label so với ground truth trên tập có nhãn, phân bố pseudo-label theo lớp và ảnh hưởng của confidence threshold trong quá trình sinh pseudo-label. Như vậy,
đóng góp chính là một giao thức thực nghiệm có kiểm soát cho semi-supervised
multi-class object detection trên CXR, cung cấp bằng chứng về vai trò của dữ liệu
chưa gán nhãn khi ngân sách nhãn bị giới hạn và tạo cơ sở cho các nghiên cứu
tiếp theo. Những khoảng trống và định vị này là cơ sở cho phương pháp nghiên cứu
trình bày ở chương tiếp theo.


\section{Materials and Methods}\label{sec:methods}

\subsection{Khung nghiên cứu tổng thể}
\label{subsec:overall_research_framework}

Nghiên cứu được thiết kế như một \textbf{controlled experimental study} nhằm
đánh giá hiệu quả của semi-supervised learning trong bài toán multi-class
object detection trên ảnh X-quang ngực khi lượng dữ liệu có chú giải bounding
box bị giới hạn. Nghiên cứu không đề xuất một detector mới hoặc một thuật toán
semi-supervised object detection mới, mà tập trung xây dựng một giao thức thực
nghiệm được xác định trước để so sánh supervised learning và semi-supervised
learning trong các điều kiện có thể đối chiếu.

Thiết kế thực nghiệm được tổ chức theo ba yếu tố thay đổi có chủ đích:
\textit{training paradigm}, \textit{labeled-data budget} và
\textit{architecture}. Cụ thể, training paradigm gồm supervised learning
(SUP) và semi-supervised learning (SSL); ngân sách dữ liệu có nhãn được xác định
theo tỷ lệ số ảnh của tập huấn luyện cố định, gồm các mức
1\%, 5\%, 10\% và 20\%; kiến trúc gồm Faster
R-CNN với ResNet-50 và FPN, đại diện cho kiến trúc CNN truyền thống, và Faster
R-CNN với Swin-T và FPN, đại diện cho Vision Transformer phân cấp dựa trên cơ chế shifted-window self-attention \cite{liu2021swin}. Do đó,
ma trận thực nghiệm chính gồm

\begin{equation}
2\ \text{training paradigms}
\times
4\ \text{labeled-data budgets}
\times
2\ \text{architectures}
=
16\ \text{experimental cells}.
\label{eq:experimental_matrix_16_cells}
\end{equation}

Mỗi experimental cell được thực hiện với cùng danh sách 10
\textit{training seeds} đã xác định trước. Như vậy, phần thực nghiệm chính ở
các chế độ thiếu nhãn gồm tổng cộng
$16 \times 10 = 160$ lần huấn luyện. Ngoài ra, supervised learning với 100\% tập huấn luyện có
nhãn được sử dụng như điều kiện tham chiếu cho từng kiến trúc, gồm
$2\times10=20$ lần huấn luyện bổ sung. Các điều kiện 100\%-SUP chỉ đóng vai
trò tham chiếu và không được đưa vào các phép so sánh ghép cặp chính giữa SUP
và SSL tại các ngân sách nhãn thấp.

Để bảo đảm tính kiểm soát của thiết kế, các phép so sánh SUP--SSL được thực
hiện theo nguyên tắc matched comparison. Trong mỗi tổ hợp kiến trúc, ngân sách
nhãn và training seed, hai phương thức học sử dụng cùng phân chia dữ liệu,
cùng \textit{labeled subset}, cùng biểu diễn đầu vào, cùng họ detector, cùng
giao thức đánh giá, cùng ngân sách cập nhật tham số và cùng mức tiếp xúc với
dữ liệu có nhãn. Khác biệt có chủ đích giữa hai phương thức học là SSL được
phép khai thác thêm phần dữ liệu chưa gán nhãn thông qua cơ chế
teacher--student pseudo-labeling. Các tham số tối ưu hóa được xác định riêng
cho từng kiến trúc và được giữ nhất quán giữa SUP và SSL trong cùng kiến trúc;
vì vậy, so sánh giữa ResNet-50 và Swin-T được diễn giải như
\textit{architecture-dependent SSL effect under prespecified
architecture-specific training recipes}, không phải như một phép kiểm định
thuần túy về ưu thế nhân quả của backbone.

MMDetection được lựa chọn làm framework thực nghiệm chính vì phù hợp trực tiếp
với thiết kế của nghiên cứu. Đây là một toolbox phát hiện đối tượng mã nguồn
mở dựa trên PyTorch, có kiến trúc mô-đun và sử dụng tệp cấu hình để tổ chức
quá trình nạp dữ liệu, huấn luyện và đánh giá \cite{mmdetection}.
MMDetection hỗ trợ trực tiếp dữ liệu theo định dạng COCO thông qua
\texttt{CocoDataset} và đánh giá theo giao thức COCO thông qua
\texttt{CocoMetric}. Framework này cũng cung cấp các thành phần cần thiết để
tổ chức dữ liệu \textit{labeled}/\textit{unlabeled}, pipeline nhiều nhánh và
mô hình teacher--student. Nhờ đó, supervised baseline và phương pháp phát hiện
đối tượng bán giám sát có thể được triển khai trong cùng một hệ sinh thái phần
mềm, với chung nguyên tắc nạp dữ liệu và đánh giá. Lựa chọn này giúp giảm số
thành phần tùy chỉnh phải tự xây dựng, qua đó hạn chế khác biệt triển khai
không cần thiết và hỗ trợ khả năng tái lập giữa các cấu hình thực nghiệm.

COCO Detection JSON được lựa chọn làm định dạng chuẩn hóa chú giải vì cấu trúc
của định dạng này phù hợp với bài toán phát hiện đối tượng đa lớp và tương thích
trực tiếp với framework đã chọn \cite{lin2014microsoft}. Định dạng COCO tổ chức
dữ liệu thành ba thành phần tường minh gồm \texttt{images},
\texttt{annotations} và \texttt{categories}, đồng thời biểu diễn bounding box
bằng tọa độ tuyệt đối theo dạng $[x,y,w,h]$. Cấu trúc này cho phép một ảnh
được khai báo trong \texttt{images} mà không có chú giải đối tượng tương ứng.
Đặc điểm đó đặc biệt quan trọng đối với nghiên cứu vì 500 ảnh
\textit{No Finding} cần được giữ lại dưới dạng ảnh không chứa đối tượng đích,
thay vì được biểu diễn như một lớp phát hiện riêng. COCO cũng phù hợp với giao
thức đánh giá của nghiên cứu, trong đó chỉ số detection chính là
mAP@[0.50:0.95], cùng với các chỉ số bổ sung như AP50, AP75 và AP theo từng
lớp.

Việc lựa chọn COCO không chỉ dựa trên mức độ phổ biến của định dạng. Dữ liệu
nguồn và lược đồ chú giải chuẩn của nghiên cứu đều sử dụng tọa độ bounding box
tuyệt đối; vì vậy, dạng $[x,y,w,h]$ của COCO tránh được bước chuẩn hóa tọa độ
theo kích thước ảnh như trong định dạng YOLO. So với cách quản lý một tệp XML
cho từng ảnh của Pascal VOC, một tệp COCO JSON tập trung thuận lợi hơn cho việc
kiểm tra tính duy nhất của định danh, tính toàn vẹn tham chiếu giữa
ảnh--chú giải--lớp và khả năng truy vết toàn bộ tập dữ liệu. Đồng thời, việc
tương thích trực tiếp với \texttt{CocoDataset} và \texttt{CocoMetric} giúp
tránh phải xây dựng riêng bộ nạp dữ liệu và bộ đánh giá cho một định dạng tùy
chỉnh. Do đó, COCO được chọn dựa trên bốn yêu cầu của nghiên cứu: duy trì tọa
độ tuyệt đối, biểu diễn hợp lệ ảnh âm tính, hỗ trợ kiểm tra tính toàn vẹn của
chú giải và tương thích với framework cũng như giao thức đánh giá đã chọn.

Khả năng áp dụng của hai lựa chọn này trên dữ liệu nghiên cứu đã được kiểm tra
ở các bước chuẩn hóa tiếp theo. Tệp COCO master sau chuyển đổi chứa 4.894 ảnh,
36.096 chú giải và 14 lớp bất thường; toàn bộ 500 ảnh
\textit{No Finding} được giữ lại với số chú giải bằng không. Không ghi nhận
chú giải không hợp lệ, sai khác số lượng bounding box hoặc lỗi toàn vẹn tham
chiếu giữa ảnh, chú giải và lớp. Việc nạp dữ liệu bằng
\texttt{CocoDataset} của MMDetection cũng xác nhận rằng cấu hình
\texttt{filter\_empty\_gt=false} giữ đủ 4.894 ảnh, trong khi
\texttt{filter\_empty\_gt=true} loại đúng 500 ảnh không chứa bounding box.
Vì nghiên cứu cần duy trì các ảnh âm tính,
\texttt{filter\_empty\_gt=false} được sử dụng trong các cấu hình dữ liệu tiếp
theo.

Khung nghiên cứu tổng thể gồm năm giai đoạn chính.

\textbf{Giai đoạn thứ nhất} xác định phạm vi thực nghiệm và chuẩn hóa dữ liệu
đầu vào. Dữ liệu DICOM và chú giải bounding box được kiểm tra, chuẩn hóa và
liên kết thông qua định danh ảnh. Ảnh DICOM được chuyển thành ảnh grayscale
8-bit và mã hóa JPEG bằng
\textit{DICOM metadata-aware, standard-aligned reference representation
pipeline}, trong khi chú giải được chuyển đổi sang định dạng COCO Detection
JSON. Quy trình không thực hiện crop hoặc resize ở bước tạo biểu diễn tham
chiếu, nhằm bảo toàn hình học ảnh và hệ tọa độ bounding box.

\textbf{Giai đoạn thứ hai} thực hiện phân chia cố định dữ liệu thành các tập
huấn luyện, xác thực và kiểm tra. Việc phân chia được thực hiện ở cấp độ ảnh
trước khi xây dựng các tập dữ liệu theo ngân sách nhãn nhằm ngăn ngừa sự giao
thoa dữ liệu giữa các tập. Tập validation và tập test được giữ cố định trong
toàn bộ thí nghiệm. Tập test được cô lập khỏi quá trình lựa chọn mô hình,
điều chỉnh ngưỡng, lựa chọn checkpoint và các quyết định phát triển khác, và
chỉ được sử dụng ở giai đoạn đánh giá cuối cùng theo giao thức đã xác định
trước.

\textbf{Giai đoạn thứ ba} xây dựng các kịch bản dữ liệu có nhãn và chưa gán
nhãn trong tập huấn luyện. Các ngân sách nhãn gồm 1\%, 5\%, 10\% và 20\%.
Các \textit{labeled subsets} được xây dựng theo quan hệ lồng nhau

\begin{equation}
L_{1\%}
\subset
L_{5\%}
\subset
L_{10\%}
\subset
L_{20\%}.
\label{eq:overview_nested_labeled_subsets}
\end{equation}

trong khi \textit{unlabeled subset} ở mỗi ngân sách là phần bù của
\textit{labeled subset} tương ứng trong tập huấn luyện cố định. Chú giải
bounding box ẩn của các ảnh thuộc \textit{unlabeled subset} không được sử dụng
cho huấn luyện, lọc pseudo-label, điều chỉnh ngưỡng, lựa chọn mô hình hoặc đánh
giá chất lượng pseudo-label.

\textbf{Giai đoạn thứ tư} huấn luyện các điều kiện supervised và
semi-supervised trong ma trận thực nghiệm đã xác định. Đối với mỗi kiến trúc,
supervised baseline chỉ sử dụng \textit{labeled subset}; thiết lập
semi-supervised sử dụng cùng \textit{labeled subset} kết hợp với
\textit{unlabeled subset} theo cơ chế teacher--student pseudo-labeling dựa
trên khung Soft Teacher \cite{xu2021end}. Teacher được cập nhật bằng
exponential moving average và sinh pseudo-label từ dữ liệu chưa gán nhãn;
các pseudo-label đáp ứng các tiêu chí lọc được sử dụng làm tín hiệu huấn luyện
bổ sung cho Student. Mỗi điều kiện được lặp lại với 10
training seeds đã xác định trước để đánh giá biến thiên do tính ngẫu nhiên của
quá trình huấn luyện và cho phép xây dựng các phép so sánh ghép cặp giữa SUP và
SSL.

Bên cạnh ma trận thực nghiệm chính, nghiên cứu thực hiện một
\textit{ablation study} thứ cấp nhằm phân tích vai trò của ba thành phần trong
cấu hình semi-supervised: ngưỡng confidence của pseudo-label, strong
photometric augmentation và chiến lược lọc pseudo-label cho nhánh hồi quy
bounding box. Ablation được thực hiện theo nguyên tắc
\textit{one-factor-at-a-time} trên kiến trúc ResNet-50 tại ngân sách nhãn
10\%, sử dụng cùng 10 training seeds. Phân tích này có mục tiêu giải thích
hành vi của các thành phần đã xác định trước, không được sử dụng như một quy
trình tìm kiếm siêu tham số để thay đổi hậu nghiệm cấu hình semi-supervised
chính.

\textbf{Giai đoạn thứ năm} thực hiện đánh giá và phân tích thống kê theo giao
thức đã xác định trước. Chỉ số detection chính là bbox
mAP@[0.50:0.95]; các chỉ số bổ sung gồm AP50, AP75, recall và AP theo từng lớp.
Đối với các lớp có support thấp, nghiên cứu phân tích riêng AP của các lớp hiếm
đã xác định từ tập huấn luyện. Đối với các ảnh \textit{No Finding}, nghiên cứu
đánh giá riêng false-positive burden thay vì xem \textit{No Finding} như một
detection class. Chất lượng pseudo-label được định nghĩa bằng
$PL$-$mAP@[0.50:0.95]$ trên tập validation cố định, sử dụng EMA Teacher tại
checkpoint LAST, và được phân tích về mối liên hệ với mức cải thiện các thí nghiệm huấn luyện và đánh giá mô hình
của semi-supervised learning.

Đơn vị suy luận thống kê của các phép so sánh chính là
\textit{training seed}, không phải từng ảnh, bounding box hay detection riêng
lẻ. Các phép so sánh SUP--SSL được xây dựng trên hiệu ứng ghép cặp theo cùng
kiến trúc, ngân sách nhãn và training seed. Chi tiết về giả thuyết nghiên cứu,
biến nghiên cứu, metric, mô hình thống kê, kiểm soát đa kiểm định và các phân
tích độ nhạy được trình bày trong các tiểu mục tương ứng ở phần sau.

Khung thực nghiệm này cho phép phân biệt có hệ thống ba nhóm ảnh hưởng chính:
phương thức học, ngân sách nhãn và kiến trúc, đồng thời kiểm soát các nguồn
biến thiên liên quan đến phân chia dữ liệu, labeled subset, training seed, biểu diễn đầu vào, họ detector, ngân sách cập nhật tham số và giao thức đánh giá.
Tuy nhiên, thiết kế chỉ hỗ trợ kết luận trong phạm vi bộ dữ liệu, hai kiến trúc,
phương pháp semi-supervised và giao thức thực nghiệm được lựa chọn; kết quả
không mặc nhiên đại diện cho mọi bộ dữ liệu X-quang ngực, mọi kiến trúc object
detection hoặc mọi phương pháp semi-supervised learning.
\subsection{Dataset and Experimental Scope}
\label{subsec:dataset_scope}

Nghiên cứu sử dụng tập huấn luyện của bộ dữ liệu VinDr-CXR, gồm các ảnh
X-quang ngực ở định dạng DICOM được chú giải bởi các bác sĩ chuyên khoa
chẩn đoán hình ảnh \cite{nguyen2022vindr}. Bộ dữ liệu VinDr-CXR đầy đủ
gồm 18.000 ảnh, trong đó 15.000 ảnh thuộc tập huấn luyện và 3.000 ảnh
thuộc tập kiểm tra. Do nhãn tham chiếu của tập kiểm tra không được công
bố, nghiên cứu này chỉ sử dụng 15.000 ảnh thuộc tập huấn luyện làm nguồn
xây dựng dữ liệu thực nghiệm.

Kết quả kiểm kê bảng chú giải nguồn xác định 67.914 bản ghi chú giải,
tương ứng với 15.000 định danh ảnh duy nhất. Trong số này, 4.394 ảnh có
ít nhất một bất thường được định vị bằng bounding box và 10.606 ảnh mang
nhãn \textit{No Finding}. Sau khi phân loại các bản ghi ở cấp độ ảnh và
đối tượng, dữ liệu ghi nhận 36.096 bản ghi chú giải bounding box thuộc
14 lớp bất thường được sử dụng trong bài toán phát hiện đối tượng.

Trong hệ thống chú giải của VinDr-CXR, \textit{No Finding} biểu thị
trường hợp không ghi nhận bất thường trên ảnh và không phải là một lớp
tổn thương cần định vị \cite{nguyen2022vindr}. Do đó, nghiên cứu xác định
14 loại bất thường là các lớp đối tượng cần phát hiện, trong khi ảnh
\textit{No Finding} được biểu diễn dưới dạng ảnh không chứa đối tượng
đích, tức không có ground-truth bounding box. Cách mã hóa này cho phép
giữ lại các ảnh âm tính trong dữ liệu mà không mô hình hóa
\textit{No Finding} như một lớp đối tượng riêng.

Theo quy trình xây dựng VinDr-CXR, mỗi ảnh trong tập huấn luyện được đọc
độc lập bởi ba bác sĩ, được phân công từ tổng số 17 bác sĩ chuyên khoa
chẩn đoán hình ảnh tham gia chú giải \cite{nguyen2022vindr}. Vì vậy,
nhiều bounding box trong bảng chú giải có thể phản ánh các nhận định độc
lập của nhiều bác sĩ đối với cùng một biểu hiện hình ảnh. Theo đó, số
bản ghi bounding box được báo cáo trong nghiên cứu không nhất thiết
tương ứng với số tổn thương lâm sàng độc lập.

\subsubsection{Phạm vi dữ liệu thực nghiệm}
\label{subsubsec:experimental_scope}

Từ tập huấn luyện ban đầu gồm 15.000 ảnh, nghiên cứu xác lập một phạm vi
thực nghiệm cố định gồm 4.894 ảnh. Cụ thể, nghiên cứu giữ lại toàn bộ
4.394 ảnh có ít nhất một bất thường được định vị bằng bounding box. Đối
với các trường hợp âm tính, 500 ảnh \textit{No Finding} được lựa chọn từ
10.606 ảnh \textit{No Finding} của tập huấn luyện ban đầu.

Trước khi lấy mẫu, trạng thái \textit{No Finding} được xác định ở cấp độ
ảnh dựa trên bảng annotation ban đầu. Đối với mỗi ảnh, nghiên cứu kiểm tra
toàn bộ các bản ghi nhãn (annotation records) có cùng \textit{image\_id}.
Một ảnh chỉ được xác định là \textit{No Finding} khi tất cả các bản ghi
nhãn tương ứng đều mang nhãn \textit{No Finding} và không có bản ghi nào
thuộc một trong 14 lớp bất thường cần phát hiện. Từ 10.606 ảnh đáp ứng
điều kiện này, nghiên cứu lựa chọn ngẫu nhiên 500 \textit{image\_id} khác
nhau, trong đó mỗi \textit{image\_id} tương ứng với một ảnh và chỉ được
lựa chọn một lần. Quá trình lấy mẫu sử dụng \textit{partition seed} cố định là 42 nhằm hỗ trợ khả năng tái lập của việc lựa chọn tập ảnh âm tính.

Sau khi hoàn tất quá trình lựa chọn, phạm vi 4.894 ảnh được cố định trong
tệp kê khai dữ liệu và được kiểm tra đối chiếu với bảng annotation nguồn
cũng như các tệp DICOM đã thu thập. Kết quả kiểm tra xác nhận 4.894
\textit{image\_id} là duy nhất, bao gồm 4.394 ảnh có bất thường và 500 ảnh
\textit{No Finding}. Không ghi nhận \textit{image\_id} trùng lặp, ảnh
không tồn tại trong bảng annotation nguồn hoặc trường hợp thiếu tệp DICOM
tương ứng. Đồng thời, toàn bộ 4.394 ảnh có bất thường đều được giữ lại
trong phạm vi thực nghiệm, và trạng thái có bất thường hoặc
\textit{No Finding} của từng ảnh nhất quán với thông tin nhãn trong bảng
annotation nguồn.

Việc đưa các ảnh \textit{No Finding} vào phạm vi thực nghiệm tạo ra các
trường hợp không chứa đối tượng đích, qua đó cho phép đánh giá hành vi của
mô hình trên ảnh âm tính, đặc biệt đối với các dự đoán false positive.
Tuy nhiên, nghiên cứu chỉ lựa chọn 500 trong tổng số 10.606 ảnh
\textit{No Finding}, trong khi giữ lại toàn bộ 4.394 ảnh có bất thường.
Do đó, thành phần dữ liệu trong phạm vi thực nghiệm khác với tập huấn
luyện VinDr-CXR ban đầu.

Theo thiết kế này, ảnh có bất thường chiếm 89,78\% và ảnh
\textit{No Finding} chiếm 10,22\% trong phạm vi thực nghiệm. Các tỷ lệ
này là kết quả của cách nghiên cứu lựa chọn dữ liệu, cụ thể là giữ lại
toàn bộ 4.394 ảnh có bất thường nhưng chỉ lựa chọn 500 trong tổng số
10.606 ảnh \textit{No Finding}. Vì vậy, tỷ lệ 89,78\% không phản ánh
tỷ lệ người có bất thường trong quần thể lâm sàng và cũng không đại diện
cho phân bố ban đầu của tập huấn luyện VinDr-CXR.

\subsubsection{Thống kê mô tả phạm vi thực nghiệm}
\label{subsubsec:scope_statistics}

Bảng~\ref{tab:dataset_scope_statistics} trình bày các thống kê mô tả của
phạm vi dữ liệu được sử dụng trong nghiên cứu. Trong tổng số 4.894 ảnh,
có 4.394 ảnh chứa ít nhất một bất thường được định vị bằng bounding box chiếm 89,78\%, và 500 ảnh \textit{No Finding} chiếm 10,22\%. Trong phạm
vi này, nghiên cứu sử dụng 14 lớp bất thường cần phát hiện, với tổng cộng
36.096 bản ghi bounding box.

Khi tính trên toàn bộ 4.894 ảnh, số bản ghi bounding box trung bình là
7,38 trên mỗi ảnh. Tuy nhiên, giá trị này bao gồm 500 ảnh
\textit{No Finding}, là những ảnh không có bounding box của các lớp bất
thường cần phát hiện. Khi chỉ tính trên 4.394 ảnh có ít nhất một bất
thường, số bản ghi bounding box trung bình tăng lên 8,21 trên mỗi ảnh.
Việc báo cáo riêng hai giá trị này cho phép mô tả rõ mật độ bounding box
trên toàn bộ phạm vi thực nghiệm và trên riêng nhóm ảnh có bất thường.

Cần lưu ý rằng các giá trị trên phản ánh số bản ghi bounding box trong
dữ liệu annotation, không trực tiếp phản ánh số tổn thương lâm sàng độc
lập trên mỗi ảnh. Một ảnh có thể chứa nhiều bản ghi bounding box, và các
bản ghi này không nhất thiết tương ứng một-một với các tổn thương lâm
sàng độc lập.

\begin{table}[htbp]
    \centering
    \caption{Thống kê mô tả của phạm vi dữ liệu thực nghiệm}
    \label{tab:dataset_scope_statistics}
    \begin{tabular}{lr}
        \toprule
        \textbf{Thống kê} & \textbf{Giá trị} \\
        \midrule
        Số ảnh trong tập huấn luyện ban đầu
            & 15.000 \\
        Số annotation records trong bảng annotation ban đầu
            & 67.914 \\
        Số ảnh trong phạm vi thực nghiệm
            & 4.894 \\
        Ảnh có ít nhất một bất thường
            & 4.394 (89,78\%) \\
        Ảnh \textit{No Finding}
            & 500 (10,22\%) \\
        Số lớp bất thường cần phát hiện
            & 14 \\
        Tổng số bản ghi bounding box
            & 36.096 \\
        Số bounding box trung bình trên toàn bộ ảnh
            & 7,38 \\
        Số bounding box trung bình trên ảnh có bất thường
            & 8,21 \\
        \bottomrule
    \end{tabular}
\end{table}

\subsubsection{Phân bố lớp và mức độ mất cân bằng}
\label{subsubsec:class_distribution}

Phân bố dữ liệu giữa 14 lớp bất thường được trình bày trong
Bảng~\ref{tab:class_distribution} và minh họa trong
Hình~\ref{fig:class_distribution}. Phân bố được thống kê theo hai đơn vị:
số ảnh chứa từng lớp và số bản ghi bounding box thuộc từng lớp. Đối với
mỗi lớp, số ảnh được xác định bằng số \textit{image\_id} duy nhất có ít
nhất một bounding box mang nhãn của lớp đó. Do một ảnh có thể chứa đồng
thời nhiều lớp bất thường, cùng một ảnh có thể được tính vào nhiều lớp.
Vì vậy, tổng số ảnh khi cộng theo từng lớp có thể lớn hơn 4.894 và tổng
tỷ lệ ảnh theo lớp có thể vượt quá 100\%.

Ở cấp độ bounding box, số lượng của mỗi lớp được tính từ toàn bộ các bản
ghi bounding box mang nhãn của lớp tương ứng trong bảng annotation.
Thống kê theo số ảnh và theo số bounding box được báo cáo riêng vì hai
đại lượng phản ánh hai đặc điểm khác nhau của dữ liệu: số ảnh cho biết
mức độ xuất hiện của từng lớp trong phạm vi thực nghiệm, trong khi số
bounding box phản ánh số lượng annotation records của lớp đó.

\begin{table*}[htbp]
    \centering
    \caption{Phân bố ảnh và bounding box theo 14 lớp bất thường trong
    phạm vi thực nghiệm}
    \label{tab:class_distribution}

    \small
    \setlength{\tabcolsep}{4pt}
    \renewcommand{\arraystretch}{1.05}

    \begin{tabular}{lrrrrr}
        \toprule
        \textbf{Lớp bất thường} &
        \shortstack{\textbf{Số}\\\textbf{ảnh}} &
        \shortstack{\textbf{Tỷ lệ}\\\textbf{ảnh (\%)}} &
        \shortstack{\textbf{Số}\\\textbf{bounding box}} &
        \shortstack{\textbf{Tỷ lệ}\\\textbf{bounding box (\%)}} &
        \shortstack{\textbf{Bounding box/}\\\textbf{ảnh có lớp}} \\
        \midrule
        Aortic enlargement
            & 3.067 & 62,67 & 7.162 & 19,84 & 2,34 \\
        Atelectasis
            & 186 & 3,80 & 279 & 0,77 & 1,50 \\
        Calcification
            & 452 & 9,24 & 960 & 2,66 & 2,12 \\
        Cardiomegaly
            & 2.300 & 47,00 & 5.427 & 15,03 & 2,36 \\
        Consolidation
            & 353 & 7,21 & 556 & 1,54 & 1,58 \\
        ILD
            & 386 & 7,89 & 1.000 & 2,77 & 2,59 \\
        Infiltration
            & 613 & 12,53 & 1.247 & 3,45 & 2,03 \\
        Lung Opacity
            & 1.322 & 27,01 & 2.483 & 6,88 & 1,88 \\
        Nodule/Mass
            & 826 & 16,88 & 2.580 & 7,15 & 3,12 \\
        Other lesion
            & 1.134 & 23,17 & 2.203 & 6,10 & 1,94 \\
        Pleural effusion
            & 1.032 & 21,09 & 2.476 & 6,86 & 2,40 \\
        Pleural thickening
            & 1.981 & 40,48 & 4.842 & 13,41 & 2,44 \\
        Pneumothorax
            & 96 & 1,96 & 226 & 0,63 & 2,35 \\
        Pulmonary fibrosis
            & 1.617 & 33,04 & 4.655 & 12,90 & 2,88 \\
        \bottomrule
    \end{tabular}

    \vspace{2pt}

    \begin{minipage}{0.94\textwidth}
        \footnotesize
        \textit{Ghi chú:} Tỷ lệ ảnh được tính trên tổng số 4.894 ảnh trong
        phạm vi thực nghiệm. Do một ảnh có thể chứa nhiều lớp bất thường,
        tổng tỷ lệ ảnh theo lớp có thể vượt quá 100\%. Tỷ lệ bounding box
        được tính trên tổng số 36.096 bản ghi bounding box.
    \end{minipage}
\end{table*}

Kết quả cho thấy phân bố giữa 14 lớp bất thường có mức độ mất cân bằng
đáng kể ở cả cấp độ ảnh và cấp độ bounding box. \textit{Aortic
enlargement} là lớp xuất hiện nhiều nhất, với 3.067 ảnh, tương ứng
62,67\% của 4.894 ảnh trong phạm vi thực nghiệm. Lớp này đồng thời có
số lượng bounding box lớn nhất, với 7.162 bản ghi, chiếm 19,84\% tổng
số 36.096 bounding box. Ngược lại, \textit{Pneumothorax} là lớp có tần
suất thấp nhất theo cả hai đơn vị thống kê, chỉ xuất hiện trong 96 ảnh
(1,96\%) và có 226 bounding box (0,63\%).

Xét theo hai lớp ở hai cực của phân bố, số ảnh của \textit{Aortic
enlargement} cao gấp khoảng 31,95 lần \textit{Pneumothorax}; tỷ số
tương ứng theo số bounding box là khoảng 31,69 lần. Bên cạnh sự khác
biệt về tần suất lớp, số bounding box trung bình trên mỗi ảnh có chứa
lớp cũng không đồng đều. Giá trị này dao động từ 1,50 đối với
\textit{Atelectasis} đến 3,12 đối với \textit{Nodule/Mass}. Như vậy,
mức độ mất cân bằng của dữ liệu không chỉ thể hiện ở số ảnh chứa từng
lớp mà còn ở số lượng bounding box được ghi nhận trên các ảnh của từng
lớp.

Đặc điểm phân bố này là một yếu tố cần được xem xét trong các giai đoạn
thực nghiệm tiếp theo. Đặc biệt, khi xây dựng các \textit{labeled subset}
và \textit{unlabeled subset}, sự chênh lệch về tần suất lớp có thể làm cho các
lớp hiếm chỉ xuất hiện với số lượng hạn chế trong \textit{labeled subset}
ở các mức ngân sách nhãn thấp. Đồng thời, khi đánh giá mô hình, kết quả
AP theo từng lớp cần được xem xét cùng với phân bố lớp để tránh việc
chỉ dựa vào các chỉ số tổng hợp khi diễn giải hiệu quả đối với những
lớp có số lượng mẫu khác nhau.

\begin{figure*}[htbp]
    \centering
    \includegraphics[width=\textwidth]
    {z_class_distribution.png}
    \caption{Phân bố dữ liệu theo 14 lớp bất thường trong phạm vi thực nghiệm.
    Panel (a) biểu diễn số \textit{image\_id} duy nhất có ít nhất một
    bounding box thuộc từng lớp; panel (b) biểu diễn tổng số bản ghi bounding
    box của từng lớp. Do một ảnh có thể chứa nhiều lớp bất thường, cùng một
    ảnh có thể xuất hiện trong thống kê của nhiều lớp; do đó, tổng số ảnh khi
    cộng theo lớp không bằng 4.894 ảnh của phạm vi thực nghiệm.}
    \label{fig:class_distribution}
\end{figure*}

\subsubsection{Mức độ đồng thuận và giới hạn của nhãn}
\label{subsubsec:annotation_reliability}

Nghiên cứu sử dụng Fleiss' kappa để định lượng mức độ đồng thuận giữa
các bác sĩ đối với sự hiện diện của từng lớp bất thường ở cấp độ ảnh
\cite{fleiss1971measuring}. Đối với mỗi ảnh và mỗi lượt đọc hợp lệ của
một bác sĩ, trạng thái của một lớp được mã hóa là hiện diện nếu bác sĩ
ghi nhận ít nhất một annotation thuộc lớp đó. Ngược lại, nếu bác sĩ có
lượt đọc hợp lệ nhưng không ghi nhận annotation nào thuộc lớp tương ứng,
trạng thái của lớp được mã hóa là không hiện diện. Cách mã hóa này chuyển
các annotation của từng bác sĩ thành trạng thái nhị phân
(present/absent) cho từng lớp ở cấp độ ảnh, phục vụ việc tính Fleiss'
kappa.

Fleiss' kappa được tính độc lập cho từng lớp trong số 14 lớp bất thường.
Trung bình số học không trọng số của 14 giá trị kappa theo lớp đạt
\(0{,}4879\). Giá trị trung bình này được báo cáo như một thống kê mô tả
tổng hợp về mức độ đồng thuận giữa các bác sĩ trong nguồn annotation,
không được sử dụng để thay thế các giá trị kappa của từng lớp. Việc xem
xét kết quả theo từng lớp vẫn cần thiết, đặc biệt đối với các lớp hiếm,
do số trường hợp dương tính thấp có thể ảnh hưởng đến giá trị kappa và
cách diễn giải mức độ đồng thuận.

Cần lưu ý rằng phân tích Fleiss' kappa trong nghiên cứu chỉ đánh giá sự
đồng thuận về trạng thái present/absent của từng lớp ở cấp độ ảnh. Phân
tích này không đánh giá mức độ đồng thuận về vị trí, kích thước hoặc ranh
giới của bounding box. Đồng thời, phép đo cũng không thực hiện matching
giữa các bounding box do các bác sĩ khác nhau cung cấp; vì vậy, không thể
từ kết quả kappa suy ra rằng các bounding box của nhiều bác sĩ mô tả cùng
một tổn thương lâm sàng.

Do đó, Fleiss' kappa được sử dụng như một chỉ báo bổ trợ nhằm đặc trưng
mức độ đồng thuận và inter-reader variability của nguồn annotation, thay
vì được xem là thước đo trực tiếp cho độ chính xác của annotation.
Chỉ số này không được sử dụng làm evaluation metric của mô hình, tiêu chí
dataset split hoặc cơ sở để tự động hiệu chỉnh annotation. Vì vậy, các
kết quả thực nghiệm được diễn giải với lưu ý rằng reference annotations
có thể chứa inter-reader variability, thay vì giả định chúng là một
ground truth hoàn toàn không có bất định \cite{le2023learning}.


\subsection{Chuẩn hóa dữ liệu và chuyển đổi sang định dạng COCO}
\label{subsec:data_standardization_coco_conversion}

Dữ liệu được chuẩn hóa nhằm xây dựng một biểu diễn đầu vào thống nhất, có khả năng truy vết và tương thích với khung phát hiện đối tượng được sử dụng trong nghiên cứu. Như minh họa trong Hình~\ref{fig:data_standardization_coco_jpeg}, quy trình gồm hai nhánh xử lý song song. Nhánh thứ nhất chuẩn hóa chú giải phát hiện và xây dựng tệp COCO master; nhánh thứ hai chuyển đổi dữ liệu ảnh DICOM thành ảnh xám 8-bit và mã hóa JPEG. Hai nhánh được liên kết bằng định danh ảnh để tạo bộ dữ liệu COCO sử dụng ảnh JPEG, đồng thời giữ nguyên kích thước ma trận điểm ảnh, ánh xạ lớp và tọa độ bounding box.

\begin{figure}[!htbp]
    \centering
    \includegraphics[width=\linewidth]
    {z_quy_trinh_chuan_hoa_du_lieu_.png}
    \caption{Quy trình chuẩn hóa dữ liệu và xây dựng bộ dữ liệu COCO sử dụng ảnh JPEG. Dữ liệu DICOM và chú giải được xử lý theo hai nhánh song song, sau đó liên kết bằng định danh ảnh trong khi vẫn bảo toàn kích thước ảnh, ánh xạ lớp và tọa độ bounding box.}
    \label{fig:data_standardization_coco_jpeg}
\end{figure}

\subsubsection{Chuẩn hóa annotation schema và kiểm tra biên ảnh}
\label{subsubsec:canonicalization_boundary_validation}

Nghiên cứu sử dụng bộ dữ liệu VinDr-CXR, một bộ dữ liệu X-quang ngực được các bác sĩ chẩn đoán hình ảnh thực hiện annotation \cite{nguyen2022vindr}. Bộ dữ liệu gốc gồm 15.000 ảnh DICOM. Trong đó, 4.394 ảnh có ít nhất một bất thường thuộc 14 lớp được xác định bằng bounding box, trong khi 10.606 ảnh còn lại được gán nhãn \textit{No Finding} và không có ground-truth bounding box.

Trong nghiên cứu này, phạm vi dữ liệu thực nghiệm được xác định gồm toàn bộ 4.394 ảnh bất thường và 500 ảnh \textit{No Finding}, tương ứng tổng cộng 4.894 ảnh. Việc giữ lại các ảnh \textit{No Finding} cho phép tập dữ liệu object detection đồng thời chứa hai trường hợp: ảnh có một hoặc nhiều ground-truth bounding box và ảnh không có ground-truth bounding box (zero-GT). Trên phạm vi 4.894 ảnh này, toàn bộ annotation được chuẩn hóa về một annotation schema thống nhất trước khi thực hiện chuyển đổi sang định dạng COCO.

Các ảnh \textit{No Finding} được biểu diễn dưới dạng ảnh không chứa đối tượng đích, tức là không có ground-truth bounding box (zero-GT), thay vì được định nghĩa như một lớp phát hiện riêng. Cách biểu diễn này cho phép duy trì các ảnh âm tính trong tập dữ liệu thực nghiệm, đồng thời giữ nguyên không gian phát hiện gồm 14 lớp bất thường.

Trước khi chuyển đổi annotation sang định dạng COCO, nghiên cứu thực hiện kiểm tra tính khả dụng của các tệp DICOM và tính hợp lệ của toàn bộ bounding box. Kích thước ảnh được xác định trực tiếp từ ma trận điểm ảnh DICOM nhằm bảo đảm rằng việc kiểm tra tọa độ bounding box được thực hiện theo geometry của ảnh nguồn. Với ảnh có chiều rộng \(W\), chiều cao \(H\), và một bounding box được biểu diễn bởi tọa độ
\((x_{\min}, y_{\min}, x_{\max}, y_{\max})\), bounding box được xác định là hợp lệ khi đồng thời thỏa mãn các điều kiện:

\begin{equation}
0 \leq x_{\min} < x_{\max} \leq W,
\qquad
0 \leq y_{\min} < y_{\max} \leq H.
\end{equation}

Các điều kiện trên bảo đảm bounding box có chiều rộng và chiều cao dương, đồng thời toàn bộ phạm vi của bounding box nằm trong biên của ma trận điểm ảnh DICOM tương ứng. Bounding box không thỏa mãn các điều kiện này được xác định là không hợp lệ và được ghi nhận trong bước kiểm tra chất lượng dữ liệu trước khi chuyển đổi sang định dạng COCO.

Kết quả kiểm tra xác nhận toàn bộ 4.894 ảnh thuộc phạm vi thực nghiệm đều có
tệp DICOM khả dụng. Trên các ảnh này, tổng cộng 36.096 ground-truth bounding
box thuộc 14 lớp bất thường được kiểm tra theo kích thước ma trận điểm ảnh
DICOM tương ứng. Không phát hiện bounding box có chiều rộng hoặc chiều cao
không dương, cũng như bounding box có tọa độ vượt ra ngoài biên ảnh. Do đó,
toàn bộ 36.096 bounding box được xác định là hợp lệ về mặt hình học và không
có annotation nào phải loại bỏ do vi phạm điều kiện về kích thước hoặc biên ảnh.

Bên cạnh kiểm tra tính hợp lệ của bounding box, annotation nguồn được sàng lọc
nhằm phát hiện các trường hợp có khả năng trùng lặp trước khi xây dựng
canonical annotation schema. Việc kiểm tra được thực hiện riêng trong từng
nhóm bounding box thuộc cùng một ảnh và cùng một lớp bất thường. Hai bounding
box có tọa độ hoàn toàn giống nhau được xác định là \textit{exact duplicate}.
Đối với các bounding box không trùng khớp hoàn toàn nhưng có mức độ chồng lấp
rất cao, một cặp bounding box được đánh dấu là \textit{near-duplicate candidate}
khi Intersection over Union (IoU) thỏa mãn
\(\mathrm{IoU} \geq 0{,}95\). Kết quả sàng lọc không phát hiện
\textit{exact duplicate}, trong khi 147 bản ghi bounding box được gắn cờ là
\textit{near-duplicate candidates}.

Các \textit{near-duplicate candidates} không được tự động hợp nhất hoặc loại
bỏ. VinDr-CXR được xây dựng từ annotation của nhiều bác sĩ chẩn đoán hình ảnh;
do đó, mức IoU cao giữa các bounding box thuộc cùng một ảnh và cùng một lớp
không đủ để kết luận rằng các annotation này là bản ghi trùng do lỗi dữ liệu.
Vì vậy, trong nghiên cứu này, duplicate screening được sử dụng như một bước
quality control nhằm phát hiện, đánh dấu và duy trì khả năng truy vết đối với
các trường hợp cần lưu ý, thay vì sử dụng một ngưỡng IoU đơn lẻ làm tiêu chí
để tự động hợp nhất hoặc loại bỏ annotation.

Sau các bước kiểm tra trên, canonical annotation schema được xây dựng để tạo
một biểu diễn thống nhất cho dữ liệu object detection trước khi chuyển đổi
sang định dạng COCO. Schema này gồm ba thành phần chính: class mapping,
canonical image table và canonical bounding box table. Class mapping xác lập
ánh xạ cố định cho 14 lớp bất thường. Trong canonical image table, mỗi ảnh
được gán một image identifier duy nhất và được liên kết với DICOM path,
chiều rộng, chiều cao cùng các thuộc tính kỹ thuật cần thiết cho việc kiểm tra
và truy vết dữ liệu. Trong canonical bounding box table, mỗi bounding box
được liên kết với đúng một image identifier và một lớp trong class mapping,
đồng thời lưu tọa độ theo biểu diễn
\((x_{\min},y_{\min},x_{\max},y_{\max})\).
Cấu trúc này bảo đảm mối liên kết nhất quán giữa ảnh, class label và bounding
box trước khi thực hiện chuyển đổi sang COCO.

\subsubsection{Xây dựng biểu diễn ảnh tham chiếu từ DICOM}
\label{subsubsec:dicom_reference_representation}

Để tạo đầu vào tương thích với MMDetection, nghiên cứu chuyển đổi ảnh DICOM thành ảnh xám 8-bit và mã hóa dưới định dạng JPEG thông qua \textit{DICOM metadata-aware, standard-aligned reference representation pipeline}, sau đây gọi tắt là \textit{metadata-aware pipeline}. Đây là pipeline xây dựng biểu diễn ảnh tham chiếu cho nghiên cứu, không được xem là một thuật toán xử lý ảnh mới, một trình hiển thị đạt chuẩn lâm sàng hoặc một phương pháp tiền xử lý tối ưu.

Trong quy trình tổng thể ở Hình~\ref{fig:data_standardization_coco_jpeg}, pipeline này tương ứng với nhánh biểu diễn ảnh. Pipeline chỉ làm thay đổi cách biểu diễn cường độ điểm ảnh; không thực hiện biến đổi hình học đối với ma trận ảnh.

Chuỗi xử lý được xây dựng dựa trên mô hình khái niệm của chuỗi biến đổi ảnh xám trong DICOM, gồm modality transformation, VOI transformation và presentation/polarity normalization \cite{dicom2026ps34}. Việc khai thác siêu dữ liệu DICOM cũng phù hợp với tiền lệ sử dụng các thuộc tính như
\texttt{WindowCenter}, \texttt{WindowWidth} và
\texttt{BitsStored} trong xử lý ảnh X-quang ngực cho bài toán phát hiện đối tượng \cite{cheng2024deep}.

Sau khi giải mã dữ liệu điểm ảnh, các pixel padding được nhận diện bằng thuộc tính
\texttt{PixelPaddingValue} và, khi hiện diện,
\texttt{PixelPaddingRangeLimit}. Các pixel này được loại khỏi quá trình xác định miền cường độ và được ánh xạ về nền sau khi hoàn tất chuỗi biến đổi ảnh.

Ở bước modality transformation, nếu
\texttt{ModalityLUTSequence} hiện diện thì Modality LUT được áp dụng. Trong trường hợp thuộc tính này không hiện diện, giá trị điểm ảnh được biến đổi bằng
\texttt{RescaleSlope} và \texttt{RescaleIntercept}:

\begin{equation}
I_{\mathrm{mod}}
=
I_{\mathrm{stored}}
\times
\mathrm{RescaleSlope}
+
\mathrm{RescaleIntercept}.
\end{equation}

Modality LUT và phép biến đổi rescale là hai nhánh thay thế; chúng không được áp dụng nối tiếp trên cùng một ảnh.

Sau modality transformation, VOI được xác định bằng
\texttt{VOILUTSequence} hoặc bằng cặp
\texttt{WindowCenter}/\texttt{WindowWidth}, có xét
\texttt{VOILUTFunction} khi thuộc tính này hiện diện
\cite{dicom2026ps34}. Khi một ảnh chứa nhiều cặp window, cặp hợp lệ đầu tiên được sử dụng theo giao thức đã được xác lập trước khi chuyển đổi toàn bộ dữ liệu.

Nếu cả VOI LUT và thông tin window đều không khả dụng, miền giá trị lý thuyết được xác định từ \texttt{BitsStored}, kiểu biểu diễn điểm ảnh và các tham số của modality transformation. Pipeline không sử dụng cực tiểu--cực đại quan sát riêng trên từng ảnh và không thực hiện cắt ngưỡng theo phân vị, qua đó tránh tạo ra quy tắc chuẩn hóa phụ thuộc vào phân bố cường độ riêng của mỗi ảnh.

Ở bước chuẩn hóa presentation/polarity, ảnh được đảo cường độ đúng một lần khi
\texttt{PresentationLUTShape} chỉ định \texttt{INVERSE}. Khi thuộc tính này không hiện diện, ảnh có
\texttt{PhotometricInterpretation} là \texttt{MONOCHROME1} được đảo cường độ. Đầu ra được chuẩn hóa về polarity tương đương
\texttt{MONOCHROME2}, trong đó giá trị thấp biểu diễn mức tối và giá trị cao biểu diễn mức sáng.

Giá trị sau VOI transformation được giới hạn theo miền đầu ra, ánh xạ về khoảng \([0,255]\), làm tròn và chuyển thành ảnh xám 8-bit một kênh. Ảnh này được sử dụng làm biểu diễn tham chiếu trước khi mã hóa JPEG.

Trong toàn bộ pipeline, các phép crop, resize, rotate, flip và transpose không được thực hiện. Do kích thước và thứ tự của ma trận điểm ảnh được giữ nguyên, tọa độ bounding box không cần co giãn hoặc biến đổi hình học. Mọi trường hợp xuất hiện thay đổi kích thước hoặc thứ tự ma trận đều được xem là lỗi chuyển đổi.

\paragraph{Lựa chọn cấu hình JPEG và kiểm tra độ trung thực.}

Ảnh xám 8-bit trước mã hóa JPEG được lưu ở định dạng không mất dữ liệu và sử dụng làm ảnh tham chiếu. Hai cấu hình JPEG với mức chất lượng 95 và 100 được đánh giá trên tập thí điểm gồm 64 ảnh, trong đó có 16 ảnh \textit{No Finding}.

Tập thí điểm được xây dựng theo nguyên tắc bao phủ xác định, gồm 14 lớp bất thường, các cực trị quan sát được về kích thước ảnh và bounding box, cùng các kiểu siêu dữ liệu DICOM xuất hiện trong phạm vi nghiên cứu. Do đây không phải là mẫu ngẫu nhiên đại diện cho toàn bộ dữ liệu, kết quả thí điểm chỉ được sử dụng để kiểm tra kỹ thuật và lựa chọn cấu hình mã hóa; không được dùng để suy rộng hiệu năng phát hiện cho toàn bộ tập dữ liệu.

Ảnh JPEG sau giải mã được so sánh với ảnh xám 8-bit trước mã hóa, thay vì so sánh trực tiếp với giá trị điểm ảnh DICOM ban đầu. Cách đối chiếu này giúp tách sai khác do mã hóa JPEG khỏi những thay đổi cường độ có chủ đích của chuỗi biến đổi DICOM.

Độ trung thực được đánh giá trên toàn ảnh và trong các vùng bounding box bằng sai số tuyệt đối trung bình (MAE), căn bậc hai của sai số bình phương trung bình (RMSE), tỷ số tín hiệu trên nhiễu đỉnh (PSNR) và chỉ số tương đồng cấu trúc (SSIM) \cite{wang2004image}. Quá trình kiểm tra còn bao gồm kích thước ảnh, tính bất biến của bounding box, ảnh sai khác và đánh giá trực quan ở cả toàn trường ảnh lẫn các vùng bên trong bounding box.

Các chỉ số đánh giá trên toàn ảnh được tính độc lập cho từng ảnh, sau đó lấy giá trị trung bình trên 64 ảnh thuộc tập thí điểm. Đối với đánh giá trong vùng bounding box, các chỉ số được tính trên từng vùng annotation và tổng hợp bằng giá trị trung bình trên toàn bộ 402 bounding box. Dung lượng tệp được xác định từ các tệp JPEG tương ứng của toàn bộ 64 ảnh trong tập thí điểm và được tổng hợp để đánh giá chi phí lưu trữ của từng cấu hình JPEG.

\begin{table}[htbp]
    \centering
    \caption{So sánh độ trung thực và dung lượng của JPEG chất lượng 95 và 100 trên tập thí điểm}
    \label{tab:jpeg_quality_comparison}
    \begin{tabular}{lcc}
        \toprule
        \textbf{Chỉ số} & \textbf{JPEG Q95} & \textbf{JPEG Q100} \\
        \midrule
        Số ảnh đánh giá
            & 64 & 64 \\
        MAE toàn ảnh, trung bình theo ảnh
            & 0,8733 & 0,0851 \\
        RMSE toàn ảnh, trung bình theo ảnh
            & 1,2354 & 0,2913 \\
        PSNR toàn ảnh, trung bình theo ảnh (dB)
            & 47,2414 & 58,8577 \\
        SSIM toàn ảnh, trung bình theo ảnh
            & 0,9812 & 0,9990 \\
        \midrule
        Số vùng bounding box đánh giá
            & 402 & 402 \\
        MAE trong vùng bounding box, trung bình theo vùng
            & 0,8480 & 0,0876 \\
        PSNR trong vùng bounding box, trung bình theo vùng (dB)
            & 47,9413 & 58,7247 \\
        SSIM trong vùng bounding box, trung bình theo vùng
            & 0,9966 & 0,9998 \\
        \midrule
        Tổng dung lượng JPEG (MiB)
            & 98,82 & 192,97 \\
        Dung lượng trung bình mỗi ảnh (MiB)
            & 1,54 & 3,02 \\
        Mức giảm tổng dung lượng so với Q100
            & 48,79\% & -- \\
        \bottomrule
    \end{tabular}
\end{table}

Như trình bày trong Bảng~\ref{tab:jpeg_quality_comparison}, JPEG Q100 đạt độ trung thực số học cao hơn Q95 trên tất cả các chỉ số toàn ảnh và vùng bounding box. Tuy nhiên, JPEG Q95 vẫn duy trì PSNR trung bình 47,24~dB và SSIM trung bình 0,9812 trên toàn ảnh. Trong các vùng bounding box, PSNR và SSIM trung bình lần lượt đạt 47,94~dB và 0,9966. Đồng thời, tổng dung lượng của 64 ảnh giảm từ 192,97~MiB ở Q100 xuống 98,82~MiB ở Q95, tương ứng mức giảm 48,79\%.

Kết hợp kết quả định lượng với kiểm tra kích thước ảnh, tính bất biến của bounding box và đánh giá trực quan, JPEG Q95 được lựa chọn và xác định làm cấu hình mã hóa cho toàn bộ dữ liệu. Quyết định này nhằm cân bằng độ trung thực của biểu diễn với chi phí lưu trữ và I/O. Kết quả trên không chứng minh rằng JPEG Q95 tối ưu về hiệu năng phát hiện, tương đương về giá trị chẩn đoán với ảnh DICOM hoặc vượt trội so với các biểu diễn ảnh khác.

Sau khi khóa cấu hình, toàn bộ 4.894 ảnh trong phạm vi nghiên cứu được chuyển đổi sang JPEG Q95. Kiểm tra sau chuyển đổi xác nhận không có tệp JPEG bị thiếu, không xuất hiện thay đổi kích thước hoặc thứ tự ma trận điểm ảnh và không phát sinh yêu cầu co giãn hay hiệu chỉnh tọa độ bounding box.

\subsubsection{Chuyển đổi annotation sang định dạng COCO}
\label{subsubsec:coco_conversion}

Từ canonical annotation schema đã được kiểm tra và chuẩn hóa ở bước trước, dữ liệu annotation được chuyển đổi sang định dạng COCO Detection JSON \cite{lin2014microsoft}. Đối với mỗi ground-truth bounding box, tọa độ ban đầu được biểu diễn dưới dạng
\((x_{\min},y_{\min},x_{\max},y_{\max})\) và được chuyển sang biểu diễn bounding box của COCO theo dạng \([x,y,w,h]\), trong đó \(x\) và \(y\) là tọa độ góc trên bên trái, còn \(w\) và \(h\) lần lượt là chiều rộng và chiều cao của bounding box:

\begin{equation}
[x,y,w,h]
=
[x_{\min},y_{\min},
x_{\max}-x_{\min},
y_{\max}-y_{\min}].
\end{equation}

Diện tích của mỗi bounding box được tính từ chiều rộng và chiều cao sau chuyển đổi:

\begin{equation}
A = w \times h.
\end{equation}

Tệp COCO Detection JSON được tổ chức từ ba thành phần chính gồm
\texttt{images}, \texttt{annotations} và \texttt{categories}. Thành phần
\texttt{images} lưu thông tin của từng ảnh; \texttt{annotations} lưu các
ground-truth bounding box và liên kết mỗi bounding box với ảnh và lớp bất thường
tương ứng; \texttt{categories} xác lập class mapping cho 14 lớp phát hiện.
Trong quá trình xây dựng, các trường định danh \texttt{image\_id},
\texttt{annotation\_id} và \texttt{category\_id} được kiểm tra về tính duy nhất
và referential integrity. Các \texttt{category\_id} được thiết lập liên tục từ
1 đến 14 và duy trì ánh xạ cố định với 14 lớp bất thường trong canonical
annotation schema.

Các ảnh \textit{No Finding} vẫn được khai báo trong thành phần
\texttt{images}, nhưng không có ground-truth bounding box tương ứng trong
\texttt{annotations}. Vì vậy, trong COCO Detection JSON, các ảnh này được biểu
diễn dưới dạng ảnh \textit{zero-GT} thay vì được định nghĩa thành một detection class
riêng. Cách biểu diễn này duy trì các ảnh không chứa đối tượng đích trong tập dữ
liệu mà không làm thay đổi không gian phát hiện gồm 14 lớp bất thường.

Kết quả của bước chuyển đổi tạo ra tệp COCO master gồm 4.894 ảnh, 36.096
ground-truth bounding box và 14 lớp bất thường; trong đó có 500 ảnh
\textit{No Finding} với số lượng ground-truth bounding box bằng không. COCO
master được sử dụng làm nguồn annotation chuẩn cho các bước xử lý tiếp theo và
được xây dựng độc lập với quyết định về định dạng biểu diễn ảnh đầu vào. Tại
thời điểm này, dữ liệu ảnh và annotation đã được chuẩn hóa về cấu trúc và quan
hệ tham chiếu, nhưng chưa thực hiện training/validation/test split và chưa xây dựng
các \textit{labeled subset}/\textit{unlabeled subset} phục vụ thí nghiệm semi-supervised learning.

Sau khi nhánh chuyển đổi biểu diễn ảnh DICOM hoàn tất, bộ dữ liệu COCO--JPEG
được tạo từ COCO master bằng cách ánh xạ đường dẫn của từng ảnh DICOM sang tệp
JPEG tương ứng thông qua cùng image identifier. Quá trình này chỉ thay đổi tham
chiếu đến định dạng ảnh đầu vào; các image identifier, kích thước ảnh, class
mapping và tọa độ bounding box được giữ nguyên. Do quá trình chuyển đổi
DICOM--JPEG không thực hiện crop, resize, rotate, flip hoặc các phép biến đổi làm
thay đổi geometry của ma trận điểm ảnh, hệ tọa độ bounding box trong COCO vẫn
tương ứng trực tiếp với ảnh JPEG. Đây là điểm hợp nhất giữa nhánh chuẩn hóa
annotation và nhánh biểu diễn ảnh trong
Hình~\ref{fig:data_standardization_coco_jpeg}.

Sau khi hợp nhất, automated quality assurance được thực hiện trên toàn bộ bộ dữ
liệu COCO--JPEG. Kết quả xác nhận đầy đủ 4.894 ảnh JPEG có thể được tham chiếu,
36.096 ground-truth bounding box thuộc 14 lớp bất thường và 500 ảnh \textit{zero-GT}.
Không phát hiện tệp JPEG bị thiếu, liên kết image--annotation--category không hợp
lệ hoặc thay đổi kích thước ma trận điểm ảnh. Đồng thời, các bounding box được
kiểm tra lại theo kích thước ảnh JPEG tương ứng nhằm xác nhận rằng quá trình
chuyển đổi và hợp nhất không làm phát sinh sai lệch geometry giữa ảnh và hệ tọa
độ bounding box.

\paragraph{Kiểm tra chất lượng sau chuyển đổi COCO.}

Sau khi hoàn tất chuyển đổi và liên kết với ảnh JPEG, quy trình kiểm tra chất lượng được thực hiện ở cả mức tự động và trực quan. Ở mức tự động, toàn bộ 4.894 ảnh và 36.096 bounding box được đối chiếu
lại về số lượng, kích thước ảnh, tính duy nhất và toàn vẹn tham chiếu của
các định danh, liên kết ảnh, chú giải lớp, giới hạn tọa độ, tính hợp lệ
của bounding box và tính nhất quán giữa geometry của ảnh với hệ tọa độ
bounding box. Các kiểm tra trùng lặp được đối chiếu với kết quả sàng lọc
trước chuyển đổi nhằm xác nhận quá trình chuyển đổi không làm phát sinh
annotation trùng mới.

Bên cạnh kiểm tra tự động, các bounding box trong tệp COCO sau chuyển đổi
được vẽ chồng (\textit{overlay}) lên ảnh JPEG để kiểm tra trực quan trên
16 ảnh \textit{representative/stress}. Tập kiểm tra này bao phủ 14 lớp bất
thường và bao gồm 4 ảnh không chứa ground-truth bounding box. Việc đối chiếu
trực quan không phát hiện sai lệch không gian rõ rệt giữa vị trí bounding box
và ảnh sau chuyển đổi. Đây là kiểm tra hỗ trợ nhằm phát hiện các lỗi ánh xạ
hoặc biến đổi hình học có thể không biểu hiện qua kiểm tra cấu trúc JSON đơn
thuần; nó không được diễn giải là bằng chứng bảo toàn tuyệt đối toàn bộ thông
tin pixel hoặc thông tin lâm sàng của DICOM gốc.

Sai khác do bước mã hóa JPEG được đánh giá riêng bằng cách so sánh JPEG sau giải mã với biểu diễn xám 8-bit không mất dữ liệu trước JPEG, sử dụng các chỉ số MAE, RMSE, PSNR và SSIM như mô tả ở trên. Vì vậy, kiểm tra geometry/annotation sau COCO và kiểm tra độ trung thực của mã hóa JPEG được xem là hai lớp kiểm soát chất lượng bổ sung cho nhau: lớp thứ nhất xác nhận tính nhất quán không gian của dữ liệu phát hiện, còn lớp thứ hai định lượng sai khác cường độ phát sinh do mã hóa JPEG Q95.

Để tổng hợp tính nhất quán của dữ liệu qua toàn bộ quá trình chuẩn hóa và
chuyển đổi, các đặc trưng chính của ảnh và chú giải được đối chiếu trước
và sau khi tạo bộ dữ liệu COCO sử dụng ảnh JPEG. Kết quả được trình bày
trong Bảng~\ref{tab:pre_post_coco_consistency}. Việc đối chiếu tập trung
vào tính toàn vẹn của phạm vi dữ liệu, tính hợp lệ của bounding box,
tính nhất quán về kích thước và geometry, cũng như khả năng duy trì
các ảnh không chứa đối tượng đích.

\begin{table*}[htbp]
    \centering
    \caption{Kết quả kiểm tra tính nhất quán dữ liệu trước và sau quá trình
    chuẩn hóa và chuyển đổi sang COCO--JPEG}
    \label{tab:pre_post_coco_consistency}
    \small
    \setlength{\tabcolsep}{5pt}
    \renewcommand{\arraystretch}{1.15}

    \begin{tabularx}{\textwidth}{
        p{3.3cm}
        >{\raggedright\arraybackslash}X
        >{\raggedright\arraybackslash}X}
        \toprule
        \textbf{Tiêu chí} &
        \textbf{Trước chuyển đổi} &
        \textbf{Sau chuyển đổi COCO--JPEG} \\
        \midrule

        Số ảnh
        & 4.894
        & 4.894 \\

        Ảnh bất thường
        & 4.394
        & 4.394 \\

        Ảnh \textit{No Finding}/zero-GT
        & 500
        & 500 \\

        Số lớp bất thường
        & 14
        & 14 \\

        Số bounding box
        & 36.096
        & 36.096 \\

        Bounding box không hợp lệ
        & 0
        & 0; không phát sinh loại bỏ \\

        Exact duplicate
        & 0
        & Không phát hiện phát sinh \\

        Near-duplicate candidates
        & 147 bản ghi, IoU $\geq 0{,}95$; giữ lại
        & Annotation nguồn được duy trì \\

        Kích thước và geometry
        & Đối chiếu với ma trận điểm ảnh DICOM
        & JPEG được đối chiếu với
        \texttt{width}, \texttt{height} trong COCO \\

        Hệ tọa độ bounding box
        & $(x_{\min},y_{\min},x_{\max},y_{\max})$;
        kiểm tra theo biên ảnh
        & Chuyển sang $[x,y,w,h]$;
        kiểm tra lại theo kích thước ảnh \\

        Biến đổi hình học
        & Không crop, resize, rotate, flip hoặc transpose
        & Không thay đổi kích thước hoặc thứ tự ma trận;
        không cần hiệu chỉnh bounding box \\

        Mã hóa JPEG
        & Biểu diễn xám 8-bit trước JPEG
        & JPEG Q95 được đánh giá bằng
        MAE, RMSE, PSNR và SSIM \\

        Kiểm tra trực quan
        & --
        & 16 ảnh representative/stress;
        14/14 lớp và 4 zero-GT;
        không phát hiện sai lệch không gian rõ rệt \\

        Nạp bằng MMDetection
        & Chưa áp dụng
        & 4.894 ảnh, 36.096 bounding box;
        \texttt{filter\_empty\_gt=false}; nạp thành công toàn bộ phạm vi dữ liệu \\

        \bottomrule
    \end{tabularx}
\end{table*}

Bảng~\ref{tab:pre_post_coco_consistency} cho thấy số lượng ảnh, số lớp
và số bounding box được duy trì nhất quán qua quá trình chuẩn hóa và
chuyển đổi. Không phát hiện bounding box không hợp lệ nên không phát sinh
trường hợp phải loại bỏ. Các near-duplicate candidates được nhận diện
trước chuyển đổi được giữ lại có chủ đích thay vì tự động hợp nhất hoặc
loại bỏ, do đặc điểm chú giải đa bác sĩ của dữ liệu nguồn.

Bên cạnh các kiểm tra tự động, Hình~\ref{fig:post_coco_overlay} minh họa
một số trường hợp trong tập kiểm tra trực quan hậu chuyển đổi. Bounding
box được đọc trực tiếp từ tệp COCO và overlay lên ảnh JPEG tương ứng,
qua đó hỗ trợ kiểm tra tính nhất quán giữa tọa độ annotation và geometry
của ảnh sau chuyển đổi.

\begin{figure*}[!ht] \centering \includegraphics[width=\textwidth] {z_contact_sheet_post_coco_overlay.png} \caption{Minh họa kiểm tra trực quan bounding box sau chuyển đổi COCO--JPEG. Các bounding box được đọc từ tệp COCO và vẽ chồng lên ảnh JPEG tương ứng. Tập kiểm tra trực quan gồm 16 ảnh representative/stress, bao phủ 14 lớp bất thường và 4 ảnh \textit{zero-GT}. Không phát hiện sai lệch không gian rõ rệt trên các mẫu được kiểm tra.} \label{fig:post_coco_overlay} \end{figure*}

Các kết quả trên cung cấp bằng chứng về tính toàn vẹn của annotation,
tính nhất quán của geometry và hệ tọa độ bounding box cần thiết cho bài
toán phát hiện đối tượng. Tuy nhiên, các kiểm tra này không được diễn giải
là bằng chứng cho việc bảo toàn tuyệt đối toàn bộ thông tin pixel hoặc
thông tin lâm sàng của DICOM gốc; sai khác do mã hóa JPEG Q95 được đánh
giá riêng trên biểu diễn xám 8-bit trước JPEG như đã trình bày ở trên.

\subsubsection{Kiểm tra khả năng nạp dữ liệu bằng MMDetection}
\label{subsubsec:mmdetection_loading_validation}

Khả năng sử dụng bộ dữ liệu COCO với ảnh JPEG được kiểm tra bằng
\texttt{CocoDataset} của MMDetection phiên bản 3.3.0. MMDetection là một toolbox phát hiện đối tượng mã nguồn mở dựa trên PyTorch, cung cấp các thành phần phục vụ cả phát hiện đối tượng có giám sát và bán giám sát \cite{mmdetection}.

Hai cấu hình \texttt{filter\_empty\_gt=false} và
\texttt{filter\_empty\_gt=true} được xây dựng độc lập để kiểm tra chính sách xử lý ảnh không chứa bounding box. Với
\texttt{filter\_empty\_gt=false}, tập dữ liệu giữ đủ 4.894 ảnh, bao gồm 500 ảnh \textit{No Finding} có mảng bounding box và nhãn rỗng hợp lệ. Khi đặt
\texttt{filter\_empty\_gt=true}, đúng 500 ảnh này bị loại và kích thước tập dữ liệu giảm còn 4.394 ảnh.

Do nghiên cứu cần duy trì các ảnh không chứa đối tượng đích trong những thí nghiệm tiếp theo, cấu hình dữ liệu được thiết lập với
\texttt{filter\_empty\_gt=false}. Ánh xạ nhãn nội bộ của MMDetection cũng được kiểm tra để bảo đảm tương ứng với các
\texttt{category\_id} liên tục từ 1 đến 14 trong tệp COCO.

Việc kiểm tra này chỉ xác nhận khả năng tương thích kỹ thuật của bộ dữ liệu COCO sử dụng ảnh JPEG, tính toàn vẹn của chú giải và chính sách giữ lại ảnh không chứa đối tượng đích. Nó không bao gồm phân chia tập dữ liệu, xây dựng các \textit{labeled subset}/\textit{unlabeled subset}, huấn luyện mô hình, sinh pseudo-label hoặc đánh giá hiệu năng phát hiện.

\subsection{Phân chia tập huấn luyện, xác thực và kiểm tra}
\label{subsec:fixed_data_split}

\subsubsection{Cơ sở lựa chọn tỷ lệ phân chia}

Phạm vi thực nghiệm gồm 4.894 ảnh, được phân chia thành \textit{training
set}, \textit{validation set} và \textit{test set} theo tỷ lệ mục tiêu xấp
xỉ 70\%/15\%/15\%. Trong thiết kế thực nghiệm, \textit{training set} được
sử dụng để huấn luyện các mô hình \textit{supervised} và làm nguồn dữ liệu
để xây dựng các \textit{labeled subset} và \textit{unlabeled subset} trong
các thí nghiệm \textit{semi-supervised learning}. \textit{Validation set}
được dành cho quá trình phát triển mô hình, bao gồm lựa chọn
\textit{checkpoint} và các quyết định cấu hình được quy định trong giao
thức thực nghiệm. Ngược lại, \textit{test set} được tách khỏi quá trình
phát triển mô hình và chỉ được sử dụng để đánh giá cuối cùng. Cách tổ chức
này nhằm hạn chế nguy cơ thông tin từ dữ liệu đánh giá cuối cùng ảnh hưởng
đến quá trình lựa chọn hoặc điều chỉnh mô hình.

Tỷ lệ 70\%/15\%/15\% được lựa chọn như một quyết định thiết kế nhằm cân
bằng hai yêu cầu của nghiên cứu. Thứ nhất, các thí nghiệm
\textit{semi-supervised learning} cần duy trì phần lớn dữ liệu trong
\textit{training set}, vì tập này tiếp tục được phân chia thành
\textit{labeled subset} và \textit{unlabeled subset} tại nhiều
ngân sách nhãn. Thứ hai, \textit{validation set} và
\textit{test set} cần có quy mô đủ lớn để giảm nguy cơ các lớp hiếm chỉ
xuất hiện trên một số lượng ảnh quá nhỏ, từ đó hạn chế khả năng đánh giá
hiệu năng ở cấp lớp.

Sự đánh đổi này đặc biệt quan trọng đối với lớp \textit{Pneumothorax}, là
lớp có số ảnh dương tính thấp nhất trong phạm vi thực nghiệm, với 96 trong
tổng số 4.894 ảnh. Nếu áp dụng tỷ lệ 80\%/10\%/10\%, số ảnh
\textit{Pneumothorax} kỳ vọng trong mỗi \textit{validation set} và
\textit{test set} chỉ là

\begin{equation}
96 \times 0{,}10 = 9{,}6.
\end{equation}

Với tỷ lệ 70\%/15\%/15\%, số ảnh kỳ vọng của lớp này trong mỗi tập đánh giá
tăng lên

\begin{equation}
96 \times 0{,}15 = 14{,}4,
\end{equation}

trong khi khoảng 70\% tổng số ảnh vẫn được giữ lại cho
\textit{training set}. Ngược lại, phương án 60\%/20\%/20\% có thể làm tăng
số ảnh \textit{Pneumothorax} kỳ vọng trong mỗi tập đánh giá lên 19,2 ảnh,
nhưng đồng thời làm giảm quy mô \textit{training set} xuống khoảng 2.936
ảnh. Sự suy giảm này thu hẹp nguồn dữ liệu dùng để xây dựng
\textit{labeled subset} và \textit{unlabeled subset}, đặc biệt tại các
ngân sách nhãn thấp.

Việc phân bổ đồng đều 15\% cho \textit{validation set} và 15\% cho
\textit{test set} còn tạo ra hai tập đánh giá có quy mô tương đương, qua đó
tránh để một trong hai tập trở nên quá nhỏ đối với các lớp hiếm. Trên cơ sở
đó, tỷ lệ 70\%/15\%/15\% được lựa chọn như một sự đánh đổi thực dụng giữa
quy mô dữ liệu huấn luyện, mức độ mất cân bằng lớp và yêu cầu đánh giá của
thiết kế thực nghiệm \textit{semi-supervised learning}. Tuy nhiên, tỷ lệ
này không được xem là một tiêu chuẩn phổ quát hoặc một phân hoạch tối ưu về
mặt thống kê. Số lượng kỳ vọng nêu trên cũng chỉ được sử dụng để phân tích
tính khả thi của các phương án tỷ lệ, không phải là số lượng quan sát thực
tế sau khi phân chia.

Với tổng số 4.894 ảnh, kích thước mục tiêu được xác định là 3.426 ảnh cho
\textit{training set}, 734 ảnh cho \textit{validation set} và 734 ảnh cho
\textit{test set}. Các giá trị nguyên này tương ứng xấp xỉ với tỷ lệ mục
tiêu 70\%/15\%/15\% và bảo đảm toàn bộ ảnh trong phạm vi thực nghiệm được
phân bổ đúng một lần.

\subsubsection{Phương pháp xây dựng phân hoạch}

Quá trình phân chia dữ liệu được thực hiện ở cấp ảnh, với
\texttt{image\_id} là đơn vị phân hoạch. Do một ảnh X-quang có thể đồng thời
chứa nhiều lớp bất thường, \textit{single-label stratification} không phản ánh
đầy đủ cấu trúc đa nhãn của dữ liệu. Vì vậy, nghiên cứu sử dụng
\textit{iterative multilabel stratification}, được triển khai thông qua
\texttt{MultilabelStratifiedShuffleSplit}, nhằm duy trì gần đúng tỷ lệ ảnh
chứa từng lớp giữa \textit{training set}, \textit{validation set} và
\textit{test set}.

Đầu vào của quá trình phân tầng là ma trận nhị phân

\begin{equation}
\mathbf{Y}\in\{0,1\}^{N\times 15},
\qquad N=4.894,
\label{eq:split_stratification_matrix}
\end{equation}

trong đó mỗi hàng của \(\mathbf{Y}\) tương ứng với một
\texttt{image\_id}, còn mỗi cột biểu diễn một chỉ báo ở cấp ảnh. Đối với 14
lớp bất thường, các phần tử của ma trận được xác định như sau:

\begin{equation}
Y_{i,k}
=
\begin{cases}
1, & \text{nếu ảnh } i \text{ chứa ít nhất một bounding box thuộc lớp } k,\\
0, & \text{nếu ảnh } i \text{ không chứa bounding box thuộc lớp } k,
\end{cases}
\qquad k=1,\ldots,14.
\label{eq:split_class_presence_matrix}
\end{equation}

Như vậy, 14 cột đầu tiên mã hóa \textit{class presence} của từng lớp bất
thường trên mỗi ảnh, không phụ thuộc vào số lượng \textit{bounding box} của
lớp đó. Một ảnh có một hoặc nhiều \textit{bounding box} thuộc cùng một lớp
đều nhận giá trị 1 tại cột tương ứng. Do đó, ma trận phân tầng phản ánh sự
hiện diện của lớp ở cấp ảnh, thay vì tần suất \textit{bounding box}.

Cột thứ 15 biểu diễn trạng thái \textit{No Finding} và được xác định bởi

\begin{equation}
Y_{i,15}
=
\begin{cases}
1, & \text{nếu ảnh } i \text{ không có ground-truth bounding box},\\
0, & \text{nếu ảnh } i \text{ có ít nhất một ground-truth bounding box}.
\end{cases}
\label{eq:split_no_finding_indicator}
\end{equation}

Chỉ báo này được sử dụng để hỗ trợ phân bổ các ảnh \textit{zero-GT} giữa
ba tập. \textit{No Finding} không được xem là \textit{detection class} thứ
15 và không được bổ sung vào hệ thống 14 lớp của \textit{COCO annotation
schema}. Trong \textit{COCO Detection JSON}, ảnh \textit{No Finding} được
biểu diễn dưới dạng ảnh không có \textit{ground-truth annotation}, thay vì
được gán một \textit{bounding box} hoặc một \textit{category} riêng.

Quá trình \textit{iterative multilabel stratification} được thực hiện qua
hai tầng. Ở tầng thứ nhất, toàn bộ \(N=4.894\) ảnh cùng ma trận chỉ báo đa nhãn
\(\mathbf{Y}\) được phân chia bằng
\texttt{MultilabelStratifiedShuffleSplit} với
\texttt{train\_size=0.70}, tạo ra một \textit{training candidate} có tỷ
lệ mục tiêu 70\% và một \textit{holdout candidate} chứa phần dữ liệu còn
lại.

Ở tầng thứ hai, chỉ các ảnh thuộc \textit{holdout set} và các hàng tương ứng
của \(\mathbf{Y}\) được tiếp tục phân tầng với
\texttt{test\_size=0.50}. Bước này chia \textit{holdout set} thành
\textit{validation set} và \textit{test set} theo tỷ lệ 50\%/50\%, tương
ứng với tỷ lệ mục tiêu 15\%/15\% khi tính trên toàn bộ phạm vi thực nghiệm.
Quy trình hai tầng do đó hướng đến phân hoạch cuối cùng theo tỷ lệ
70\%/15\%/15\%.

Cả hai tầng đều sử dụng 15 chỉ báo nhị phân đã định nghĩa, bao gồm
\textit{class presence} của 14 lớp bất thường và trạng thái
\textit{No Finding}. Sau tầng thứ nhất, mỗi ảnh được phân vào
\textit{training set} hoặc \textit{holdout set}. Ở tầng thứ hai, chỉ các ảnh
thuộc \textit{holdout set} được tiếp tục phân chia vào
\textit{validation set} hoặc \textit{test set}; các ảnh đã được phân vào
\textit{training set} không tham gia bước phân chia này. Kết quả của hai tầng
xác định tập mà mỗi ảnh tạm thời thuộc về. Sau đó, một số ảnh có thể được
chuyển từ tập này sang tập khác trong các bước \textit{deterministic repair}
nhằm đưa số lượng ảnh của từng tập và số lượng ảnh \textit{No Finding} về
đúng các giá trị mục tiêu. Các bước hiệu chỉnh này chỉ thay đổi tập chứa ảnh,
không làm thay đổi nội dung ảnh, \texttt{image\_id} hoặc
\textit{ground-truth annotation} tương ứng.

Do một hàng của \(\mathbf{Y}\) có thể đồng thời nhận giá trị 1 tại nhiều
cột, biểu diễn này duy trì cấu trúc đa nhãn ở mức \textit{class presence}.
Tuy nhiên, ma trận không trực tiếp mã hóa số lượng \textit{bounding box}, số
tổn thương, kích thước hoặc vị trí \textit{bounding box}, cũng như toàn bộ
cấu trúc đồng xuất hiện giữa các lớp. Vì vậy, quá trình phân tầng chỉ hướng
đến việc duy trì gần đúng \textit{marginal image-level prevalence} của từng
lớp và tỷ lệ ảnh \textit{zero-GT}; phương pháp này không bảo đảm rằng mọi
đặc trưng của phân bố chú giải đều được cân bằng giữa ba tập.

\textit{Partition seed} 42 được ấn định trước khi thực hiện quá trình phân chia
và được sử dụng thống nhất ở cả hai tầng \textit{iterative multilabel
stratification}. Việc cố định \textit{partition seed} hỗ trợ khả năng tái lập
thành phần phân hoạch khi quy trình được thực thi lại với cùng dữ liệu đầu
vào, cùng phiên bản triển khai và cùng môi trường thực thi. Giá trị 42 chỉ là
tham số kỹ thuật dùng để kiểm soát thành phần giả ngẫu nhiên của quy trình; nghiên cứu không giả định rằng giá trị này có ưu thế thống kê so với các giá trị \textit{partition seed} khác. Đồng thời, không thực hiện \textit{seed search}
để lựa chọn phân hoạch dựa trên kết quả các thí nghiệm huấn luyện và đánh giá mô hình hoặc dựa trên
các thống kê thuận lợi quan sát được sau khi phân chia.

Sau hai tầng phân tầng, phân hoạch ban đầu gồm 3.449 ảnh trong training set,
713 ảnh trong validation set và 732 ảnh trong test set. Kết quả này chưa đạt
kích thước mục tiêu 3.426/734/734 ảnh, tương ứng với tỷ lệ
70\%/15\%/15\% trên tổng số 4.894 ảnh. Sự chênh lệch phát sinh do
\textit{iterative multilabel stratification} ưu tiên duy trì phân bố của
các chỉ báo đa nhãn giữa các tập nhưng không áp đặt ràng buộc để kích thước
của từng tập phải bằng chính xác kích thước nguyên mục tiêu. Cụ thể, so với
mục tiêu, training set dư 23 ảnh, trong khi validation set thiếu 21 ảnh và
test set thiếu 2 ảnh. Vì vậy, sau bước phân tầng ban đầu, một bước
\textit{deterministic constrained greedy repair} được áp dụng để hiệu chỉnh thực hiện tổng cộng 23 phép chuyển ảnh, bao gồm 21 ảnh từ
\textit{training set} sang \textit{validation set} và 2 ảnh từ
\textit{training set} sang \textit{test set}. Sau bước này, kích thước của
ba tập đạt chính xác 3.426/734/734 ảnh.

Việc lựa chọn ảnh cần chuyển tiếp tục dựa trên biểu diễn nhị phân 15 chiều
\(\mathbf{Y}\). Gọi \(\mathbf{p}\) là vector tỷ lệ hiện diện của 15 chỉ báo
trên toàn bộ phạm vi thực nghiệm và \(\mathbf{p}_s\) là vector tỷ lệ tương ứng trong tập
\(s\in\{\mathrm{tr},\mathrm{va},\mathrm{te}\}\),

trong đó \(\mathrm{tr}\), \(\mathrm{va}\) và \(\mathrm{te}\) lần lượt ký
hiệu \textit{training set}, \textit{validation set} và \textit{test set}.
Đối với mỗi trạng thái phân hoạch ứng viên, sai lệch cực đại được xác định
bởi

\begin{equation}
J_{\max}
=
\max_{s,k}
\left|
p_{s,k}-p_k
\right|,
\label{eq:size_repair_max_deviation}
\end{equation}

và tổng bình phương sai lệch được xác định bởi

\begin{equation}
J_{\mathrm{sq}}
=
\sum_s\sum_{k=1}^{15}
\left(
p_{s,k}-p_k
\right)^2,
\label{eq:size_repair_squared_deviation}
\end{equation}

trong đó \(k\) biểu thị một trong 15 chỉ báo được sử dụng trong quá trình
phân tầng.

Quá trình hiệu chỉnh được thực hiện tuần tự. Tại mỗi vòng lặp, thuật toán xác
định tập đang vượt và tập đang thiếu kích thước mục tiêu, sau đó đánh giá
toàn bộ phép chuyển khả thi của một ảnh từ tập nguồn sang tập đích. Phép
chuyển được lựa chọn bằng cách tối thiểu hóa theo thứ tự từ điển

\begin{equation}
\left(
J_{\max},
J_{\mathrm{sq}},
r_{\mathrm{src}},
r_{\mathrm{dst}},
\operatorname{key}(\texttt{image\_id})
\right),
\label{eq:size_repair_lexicographic_objective}
\end{equation}

trong đó \(r_{\mathrm{src}}\) và \(r_{\mathrm{dst}}\) biểu thị thứ tự cố
định của tập nguồn và tập đích, còn
\(\operatorname{key}(\texttt{image\_id})\) là khóa định danh ảnh được sử
dụng để xử lý các trường hợp đồng hạng theo cách xác định. Thuật toán trước
hết ưu tiên phép chuyển tạo ra \(J_{\max}\) nhỏ hơn; nếu \(J_{\max}\) bằng
nhau, phép chuyển tạo ra \(J_{\mathrm{sq}}\) nhỏ hơn được lựa chọn. Các thành
phần còn lại trong hàm mục tiêu bảo đảm rằng kết quả vẫn được xác định khi
nhiều phép chuyển có cùng giá trị của hai thành phần mục tiêu chính.

Sau mỗi vòng lặp, phép chuyển được lựa chọn được áp dụng và kích thước của ba
tập được đánh giá lại. Quá trình kết thúc khi đạt chính xác
3.426/734/734 ảnh. Như vậy, các ảnh được điều chuyển theo một hàm mục tiêu và
quy tắc xử lý đồng hạng đã xác định, thay vì được lựa chọn thủ công hoặc tùy
ý.

Tuy nhiên, \textit{deterministic constrained greedy repair} là một
\textit{greedy algorithm}, trong đó phương án phù hợp nhất được lựa chọn tại
từng vòng lặp. Thuật toán không giải đồng thời toàn bộ bài toán tổ hợp và
không được diễn giải là đã tìm thấy nghiệm tối ưu toàn cục.

Bên cạnh các ràng buộc về kích thước và phân bố đa nhãn, nghiên cứu kiểm soát
riêng việc phân bổ 500 ảnh \textit{No Finding}, với số lượng mục tiêu là 350
ảnh trong \textit{training set}, 75 ảnh trong \textit{validation set} và 75
ảnh trong \textit{test set}. Sau bước hiệu chỉnh kích thước, phân bố này
được đưa về chính xác mục tiêu thông qua \textit{exact No Finding swap
repair}. Mỗi phép hoán đổi chuyển một ảnh \textit{No Finding} từ tập đang
thừa sang tập đang thiếu, đồng thời chuyển một ảnh bất thường theo chiều
ngược lại. Do đó, kích thước của hai tập tham gia hoán đổi được giữ nguyên.

Hàm mục tiêu của bước hoán đổi được tính trên 14 lớp bất thường. Gọi
\(\mathbf{q}\) là vector \textit{image-level prevalence} của 14 lớp trên
toàn bộ phạm vi thực nghiệm, \(\mathbf{q}_s\) là vector tỷ lệ tương ứng trong
tập \(s\), còn \(LC_s\) và \(LC_{\mathrm{full}}\) lần lượt là
\textit{label cardinality} của tập \(s\) và của toàn bộ phạm vi thực nghiệm.
Ba thành phần mục tiêu được xác định như sau:

\begin{equation}
K_{\max}
=
\max_{s,c}
\left|
q_{s,c}-q_c
\right|,
\label{eq:nofinding_repair_max_deviation}
\end{equation}

\begin{equation}
K_{\mathrm{sq}}
=
\sum_s\sum_{c=1}^{14}
\left(
q_{s,c}-q_c
\right)^2,
\label{eq:nofinding_repair_squared_deviation}
\end{equation}

và

\begin{equation}
K_{\mathrm{LC}}
=
\max_s
\left|
LC_s-LC_{\mathrm{full}}
\right|.
\label{eq:nofinding_repair_cardinality_deviation}
\end{equation}

Một phép hoán đổi chỉ được xem là khả thi nếu sau khi thực hiện, cả 14 lớp
bất thường vẫn hiện diện trong mỗi tập, đồng thời \textit{validation set} và
\textit{test set} đều chứa ít nhất 10 ảnh \textit{Pneumothorax}. Trong số
các phép hoán đổi khả thi, thuật toán lựa chọn phương án tối thiểu hóa theo
thứ tự từ điển

\begin{equation}
\left(
K_{\max},
K_{\mathrm{sq}},
K_{\mathrm{LC}},
r_{\mathrm{src}},
r_{\mathrm{dst}},
\operatorname{key}(\texttt{image\_id}_{\mathrm{NF}}),
\operatorname{key}(\texttt{image\_id}_{\mathrm{abn}})
\right),
\label{eq:nofinding_repair_lexicographic_objective}
\end{equation}

trong đó hai khóa định danh cuối cùng lần lượt tương ứng với ảnh
\textit{No Finding} và ảnh bất thường, được sử dụng để xử lý các trường hợp
đồng hạng theo cách xác định. Tổng cộng ba phép hoán đổi cặp được thực hiện
để đạt chính xác phân bổ 350/75/75 ảnh \textit{No Finding}, đồng thời duy
trì kích thước 3.426/734/734 của ba tập.

Các phép chuyển và hoán đổi chỉ làm thay đổi quan hệ thành viên của ảnh giữa
các tập. Không có ảnh nào bị loại khỏi phạm vi thực nghiệm và nội dung
\textit{ground-truth annotation} của từng ảnh không bị sửa đổi. Thuật ngữ
``\textit{exact}'' trong \textit{exact No Finding swap repair} chỉ biểu thị
rằng kích thước tập và số lượng ảnh \textit{No Finding} đạt chính xác các
mục tiêu đã xác định; thuật ngữ này không hàm ý rằng thuật toán đã tìm được
nghiệm tối ưu toàn cục.

Sau khi phân hoạch đáp ứng toàn bộ \textit{acceptance criteria}, thành phần
của \textit{training set}, \textit{validation set} và \textit{test set}
được khóa để làm đầu vào thống nhất cho các giai đoạn tiếp theo. Theo giao
thức nghiên cứu, các thí nghiệm huấn luyện và đánh giá mô hình sẽ sử dụng phân hoạch đã
khóa này, thay vì thực hiện lại quá trình phân chia cho từng thí nghiệm.
Việc khóa phân hoạch không đồng nghĩa với việc các \textit{labeled subset},
\textit{unlabeled subset} hoặc quá trình huấn luyện mô hình đã được thực
hiện.


\subsubsection{Kiểm định tính toàn vẹn và các đặc trưng phân phối của phân hoạch}

Sau khi hoàn tất quá trình xây dựng và hiệu chỉnh, phân hoạch
\textit{training/validation/test} được kiểm định trước khi cố định để sử dụng
trong các giai đoạn thực nghiệm tiếp theo. Quá trình kiểm định tập trung vào
bốn khía cạnh: (i) tính toàn vẹn cấu trúc của phân hoạch; (ii) sự tách biệt
giữa các tập trong phạm vi các định danh khả dụng; (iii) mức độ duy trì các
đặc trưng phân phối được lựa chọn để kiểm tra; và (iv)
\textit{reproducibility} của quá trình phân chia.

Về tính toàn vẹn cấu trúc, \textit{training set},
\textit{validation set} và \textit{test set} lần lượt chứa 3.426, 734 và
734 ảnh, tương ứng với toàn bộ 4.894 ảnh thuộc phạm vi thực nghiệm. Kiểm tra
tại cấp \texttt{image\_id} cho thấy giao của từng cặp tập bằng rỗng, trong
khi hợp của ba tập bao phủ toàn bộ phạm vi thực nghiệm. Do đó, trong phạm vi
định danh ảnh khả dụng, mỗi ảnh thuộc đúng một trong ba tập; không có ảnh nào
bị bỏ sót hoặc đồng thời xuất hiện trong nhiều tập.

Mỗi \textit{annotation} được phân bổ vào tập chứa ảnh mà nó tham chiếu. Quá
trình tạo các tệp \textit{COCO Detection JSON} giữ nguyên giá trị
\texttt{id} của ảnh, \textit{annotation} và lớp từ tệp COCO master, không
thực hiện đánh lại số định danh. Kết quả kiểm định không phát hiện giá trị
\texttt{id} của các phần tử trong \texttt{annotations} bị trùng giữa các
tập. Đồng thời, không có \textit{annotation} nào có trường
\texttt{image\_id} tham chiếu đến ảnh nằm ngoài tập tương ứng, và hợp của
các \textit{annotation} trong ba tập bao phủ đủ 36.096
\textit{annotation} của tệp COCO master. Những kết quả này xác nhận tính
nhất quán của quan hệ ảnh--\textit{annotation} sau quá trình phân hoạch.

Nhóm 500 ảnh \textit{No Finding} được phân bổ thành 350 ảnh trong
\textit{training set}, 75 ảnh trong \textit{validation set} và 75 ảnh trong
\textit{test set}. Trong \textit{COCO Detection JSON}, các ảnh này được biểu
diễn dưới dạng ảnh \textit{zero-GT}, tức không chứa
\textit{ground-truth bounding box}. Việc kiểm soát riêng nhóm
\textit{No Finding} giúp hạn chế sự khác biệt không chủ đích về tỷ lệ ảnh
không chứa bất thường được chú giải giữa ba tập, đặc biệt khi các ảnh này liên
quan trực tiếp đến việc đánh giá hành vi dự đoán \textit{false positive} của
mô hình.

Bộ dữ liệu VinDr-CXR đã trải qua quá trình khử định danh trước khi được công
bố. Nguyen et al.~\cite{nguyen2022vindr} cho biết các thông tin có khả năng
nhận dạng cá nhân gắn với ảnh đã được loại bỏ hoặc thay thế bằng giá trị ngẫu
nhiên. Các DICOM tag chứa thông tin sức khỏe được bảo vệ bị loại bỏ bao gồm
tên bệnh nhân, ngày sinh và \texttt{PatientID}. Các tác giả cũng cho biết
toàn bộ DICOM metadata đã được phân tích và rà soát thủ công để kiểm tra việc
loại bỏ thông tin có khả năng nhận dạng bệnh nhân. Trong dữ liệu được công
bố, mỗi ảnh được gán một định danh ẩn danh duy nhất, được mã hóa từ
\texttt{SOPInstanceUID}~\cite{nguyen2022vindr}. Vì vậy,
\texttt{image\_id} là định danh ảnh và không thể được diễn giải như định danh
bệnh nhân.

Dữ liệu khả dụng không cung cấp định danh nhóm bệnh nhân phù hợp để xác định
những ảnh thuộc cùng một bệnh nhân sau quá trình khử định danh. Đồng thời,
bài báo mô tả VinDr-CXR không khẳng định rằng 18.000 ảnh được công bố tương
ứng với 18.000 bệnh nhân duy nhất~\cite{nguyen2022vindr}. Vì vậy, nghiên cứu
không thể thực hiện phân chia hoặc kiểm tra giao nhau ở cấp bệnh nhân. Kết quả
không phát hiện giao nhau tại cấp \texttt{image\_id} do đó không được suy
rộng thành bằng chứng về sự vắng mặt tuyệt đối của \textit{data leakage} tại
cấp bệnh nhân. Sự tách biệt giữa các tập trong nghiên cứu này chỉ được xác
nhận trong phạm vi \texttt{image\_id}, định danh của \textit{annotation} và
các chú giải khả dụng. Đây là một giới hạn của quá trình kiểm soát phân hoạch
dữ liệu.

Bên cạnh tính toàn vẹn cấu trúc, phân hoạch được kiểm tra về một số đặc trưng
phân phối ở cấp ảnh. Trước hết, cả 14 lớp bất thường đều hiện diện trong
\textit{training set}, \textit{validation set} và \textit{test set}. Đối với
mỗi lớp \(c\), \textit{image-level prevalence} được xác định là tỷ lệ ảnh có
ít nhất một \textit{ground-truth bounding box} thuộc lớp đó. Sai lệch
\textit{image-level prevalence} giữa một tập \(s\) và toàn bộ phạm vi thực
nghiệm được tính bởi

\begin{equation}
\Delta_{s,c}
=
\left|
\frac{n_{s,c}}{N_s}\times 100
-
\frac{n_{\mathrm{full},c}}{N_{\mathrm{full}}}\times 100
\right|,
\label{eq:split_prevalence_deviation}
\end{equation}

\noindent
trong đó \(n_{s,c}\) là số ảnh chứa lớp \(c\) trong tập \(s\),
\(N_s\) là tổng số ảnh của tập \(s\), \(n_{\mathrm{full},c}\) là số ảnh chứa
lớp \(c\) trong toàn bộ phạm vi thực nghiệm và
\(N_{\mathrm{full}}=4.894\). Do một ảnh có thể chứa đồng thời nhiều lớp,
\textit{image-level prevalence} được tính độc lập cho từng lớp và tổng tỷ lệ
của 14 lớp không bị ràng buộc bằng 100\%.

Sai lệch \textit{image-level prevalence} lớn nhất quan sát được trên toàn bộ
các cặp tập--lớp là 0,3026 điểm phần trăm, xảy ra đối với lớp
\textit{Pleural effusion} trong \textit{validation set}. Giá trị này thấp hơn
\textit{internal acceptance criterion} 1 điểm phần trăm được xác định trước
để phát hiện những sai lệch rõ rệt về tỷ lệ ảnh chứa từng lớp sau phân hoạch.

Tuy nhiên, việc kiểm tra riêng \textit{image-level prevalence} của từng lớp
chưa mô tả đầy đủ cấu trúc đa nhãn của dữ liệu. Vì vậy, nghiên cứu kiểm tra
thêm \textit{label cardinality}, được định nghĩa là số lớp bất thường trung
bình hiện diện trên mỗi ảnh trong một tập dữ liệu \(D\):

\begin{equation}
LC(D)
=
\frac{1}{|D|}
\sum_{i=1}^{|D|}
|\mathcal{L}_i|,
\label{eq:label_cardinality}
\end{equation}

\noindent
trong đó \(\mathcal{L}_i\) là tập các lớp bất thường hiện diện trên ảnh
\(i\). \textit{Label cardinality} của toàn bộ phạm vi thực nghiệm là
3,139559. Các giá trị tương ứng của \textit{training set},
\textit{validation set} và \textit{test set} lần lượt là 3,131057,
3,170300 và 3,148501. Sai lệch tuyệt đối lớn nhất so với toàn bộ phạm vi thực
nghiệm là 0,030741, quan sát tại \textit{validation set}, và thấp hơn
\textit{internal acceptance criterion} 0,10.

Lớp \textit{Pneumothorax} được kiểm tra riêng do đây là lớp có số ảnh dương
tính thấp nhất trong phạm vi thực nghiệm, với tổng cộng 96 ảnh. Số ảnh chứa
\textit{Pneumothorax} trong \textit{training set},
\textit{validation set} và \textit{test set} lần lượt là 66, 15 và 15.
Nếu phân bổ hoàn toàn theo tỷ lệ danh nghĩa 70\%/15\%/15\%, số lượng kỳ vọng
tương ứng là 67,2, 14,4 và 14,4 ảnh. Như vậy, cả
\textit{validation set} và \textit{test set} đều chứa 15 ảnh
\textit{Pneumothorax}, đồng thời đáp ứng
\textit{internal acceptance criterion} tối thiểu 10 ảnh được sử dụng trong
quá trình xây dựng và kiểm định phân hoạch.

Tiêu chí trên chỉ bảo đảm rằng lớp hiếm nhất không bị mất hoàn toàn hoặc có
số lượng thấp hơn mức đã xác định trong \textit{validation set} và
\textit{test set}. Việc có 15 ảnh \textit{Pneumothorax} trong mỗi tập không
đồng nghĩa với việc \textit{AP} của lớp này có độ ổn định thống kê cao. Vì
vậy, kết quả đánh giá theo lớp, đặc biệt đối với các lớp hiếm, cần được diễn
giải thận trọng và xem xét cùng với phân tích bất định ở giai đoạn đánh giá
cuối cùng.

Bảng~\ref{tab:fixed_split_validation} tóm tắt các kết quả kiểm định chính của
phân hoạch cuối cùng.

\begin{table*}[htbp]
    \centering
    \caption{Đặc trưng và kết quả kiểm định phân hoạch
    \textit{training/validation/test}}
    \label{tab:fixed_split_validation}
    \begin{tabular}{p{5.0cm}p{9.0cm}}
        \toprule
        \textbf{Tiêu chí} & \textbf{Kết quả} \\
        \midrule

        Kích thước \textit{training/validation/test}
            & 3.426/734/734 ảnh \\

        Bao phủ phạm vi thực nghiệm
            & 4.894/4.894 ảnh và 36.096/36.096
              \textit{annotation} \\

        Sự tách biệt và toàn vẹn giữa các tập
            & Không phát hiện \texttt{image\_id} hoặc định danh của
              \textit{annotation} trùng giữa các tập; không có
              \textit{annotation} tham chiếu đến ảnh ngoài tập tương ứng;
              quan hệ ảnh--\textit{annotation} được bảo toàn \\

        Phân bổ \textit{No Finding}
            & 350/75/75 ảnh \\

        Sự hiện diện của các lớp
            & Cả 14 lớp bất thường hiện diện trong cả ba tập \\

        Sai lệch \textit{image-level prevalence} lớn nhất
            & 0,3026 điểm phần trăm
              (\textit{Pleural effusion}, \textit{validation set}) \\

        \textit{Label cardinality}
            & Toàn bộ: 3,139559; \textit{training set}: 3,131057;
              \textit{validation set}: 3,170300;
              \textit{test set}: 3,148501 \\

        Sai lệch \textit{label cardinality} lớn nhất
            & 0,030741 \\

        Phân bổ \textit{Pneumothorax}
            & 66/15/15 ảnh \\

        \bottomrule
    \end{tabular}
\end{table*}

Ba tiêu chí gồm sai lệch \textit{image-level prevalence} không vượt quá
1 điểm phần trăm, sai lệch \textit{label cardinality} không vượt quá 0,10 và
ít nhất 10 ảnh của lớp hiếm nhất trong mỗi tập giữ lại được sử dụng như các tiêu chí về sai lệch \textit{image-level prevalence}, \textit{label cardinality} và số ảnh của lớp hiếm nhất được xác định trước trong giao thức nghiên cứu nhằm phát hiện các phân hoạch có sai lệch đáng
kể đối với những đặc trưng đang được kiểm soát. Để phát hiện những phân hoạch có sai lệch rõ rệt đối với các đặc trưng đang được kiểm soát. Các giá trị này là ngưỡng kiểm soát được xác định trong giao thức nghiên cứu, không phải các ngưỡng chuẩn do \textit{iterative multilabel stratification} quy định và
cũng không phải tiêu chuẩn chung của bài toán \textit{object detection}.
Do đó, việc thỏa mãn các tiêu chí trên không được diễn giải là bằng chứng
rằng phân hoạch đạt tối ưu thống kê.

\textit{Reproducibility} của phân hoạch được kiểm tra ở cấp thành phần ảnh.
\textit{Partition seed} 42 được ấn định trước quá trình phân chia và giữ cố định
ở cả hai tầng \textit{iterative multilabel stratification}. Các bước
\textit{deterministic constrained greedy repair} và
\textit{exact No Finding swap repair} sử dụng các hàm mục tiêu và quy tắc xử
lý đồng hạng xác định. Khi toàn bộ quy trình được thực thi lại trên cùng dữ
liệu đầu vào, cùng phiên bản triển khai và cùng môi trường thực thi, thành
phần \texttt{image\_id} của từng tập nhất quán với phân hoạch đã kiểm định.
Việc cố định \textit{partition seed} nhằm kiểm soát thành phần giả ngẫu nhiên
của quá trình phân chia và hỗ trợ \textit{reproducibility}; giá trị này
không được xem là bằng chứng về tính đại diện hoặc tính tối ưu thống kê của
phân hoạch.

Tổng hợp các kiểm tra cho thấy phân hoạch cuối cùng đạt kích thước
3.426/734/734, bao phủ toàn bộ 4.894 ảnh và 36.096
\textit{annotation}, đồng thời không phát hiện giao nhau tại cấp
\texttt{image\_id} hoặc định danh của \textit{annotation}. Không có
\textit{annotation} nào có trường \texttt{image\_id} tham chiếu đến ảnh
ngoài tập tương ứng. Phân hoạch duy trì đúng phân bổ
\textit{No Finding}, bảo toàn sự hiện diện của cả 14 lớp và đáp ứng các
các tiêu chí kiểm định được xác định trước trong giao thức nghiên cứu đối với
\textit{image-level prevalence}, \textit{label cardinality} và lớp hiếm
nhất. Kiểm tra quan hệ ảnh--\textit{annotation} đồng thời xác nhận rằng mỗi
\textit{ground-truth annotation} vẫn được gắn với đúng tập chứa ảnh tương
ứng.

Tuy nhiên, các kết quả trên chỉ xác nhận những thuộc tính đã được kiểm tra và
không chứng minh rằng toàn bộ phân phối dữ liệu giữa ba tập là tương đương.
Cụ thể, quá trình kiểm định hiện tại không trực tiếp đánh giá đầy đủ cấu trúc
đồng xuất hiện giữa các lớp, phân phối số lượng \textit{bounding box} trên
mỗi ảnh, kích thước và vị trí tổn thương hoặc các đặc trưng liên quan đến
thiết bị và quá trình thu nhận ảnh. Tương tự, việc không phát hiện trùng lặp
tại cấp \texttt{image\_id} và định danh của \textit{annotation} không chứng
minh tính độc lập thống kê giữa ba tập và không loại trừ hoàn toàn khả năng
\textit{data leakage} ở cấp bệnh nhân khi định danh cần thiết cho kiểm tra
này không khả dụng.

Sau khi đáp ứng các tiêu chí kiểm định được xác định trước trong giao thức nghiên cứu, thành phần
\textit{training set}, \textit{validation set} và \textit{test set} được cố
định trước khi tiến hành các thí nghiệm huấn luyện và đánh giá mô hình. Phân hoạch này
được giữ nguyên giữa các thí nghiệm để tránh thay đổi thành phần
\textit{validation/test} giữa các cấu hình. Theo giao thức thực nghiệm,
\textit{validation set} được sử dụng cho lựa chọn \textit{checkpoint}, lựa
chọn mô hình và các quyết định về \textit{hyperparameter}.
\textit{Test set} không được sử dụng cho lựa chọn mô hình, điều chỉnh
\textit{hyperparameter}, lựa chọn \textit{pseudo-label threshold},
\textit{early stopping} hoặc lựa chọn cấu hình \textit{ablation}; tập này
chỉ được sử dụng cho đánh giá cuối cùng.

Ở bước xây dựng dữ liệu huấn luyện tiếp theo, các \textit{labeled subset} và
\textit{unlabeled subset} theo các mức ngân sách nhãn chỉ được
tạo từ 3.426 ảnh thuộc \textit{training set}. Quá trình này không làm thay
đổi thành phần của \textit{validation set} hoặc \textit{test set}. Tại thời
điểm phân hoạch \textit{training/validation/test} được cố định, các
\textit{labeled subset} và \textit{unlabeled subset} này được xây dựng ở mục tiếp theo.

\subsubsection{Xây dựng các tập con có nhãn và chưa gán nhãn}
\label{subsubsec:labeled_unlabeled_construction}

Sau khi thành phần của \textit{training set}, \textit{validation set} và
\textit{test set} được xác định và cố định theo quy trình phân chia trình
bày ở các tiểu mục trước, các tập con phục vụ thí nghiệm
\textit{semi-supervised learning} được tiếp tục xây dựng. Quá trình này chỉ
sử dụng 3.426 ảnh thuộc \textit{training set}; thành phần của
\textit{validation set} và \textit{test set} được giữ nguyên và không tham
gia vào quá trình \textit{sampling}, hiệu chỉnh hoặc lựa chọn
\textit{membership}. Mục tiêu là tạo ra các mức ngân sách dữ liệu có nhãn 1\%, 5\%,
10\% và 20\%, trong đó tỷ lệ ngân sách được xác định theo số lượng ảnh
thuộc \textit{training set}, không theo số lượng \textit{bounding box}.
Các mức ngân sách này được sử dụng thống nhất cho các cấu hình
\textit{supervised} và \textit{semi-supervised} tại cùng một mức ngân sách.

\paragraph{Lựa chọn các mức ngân sách nhãn.}

Trong nghiên cứu này, \textit{labeled-data budget} được định nghĩa ở cấp độ ảnh:
mức ngân sách \(b\) biểu thị tỷ lệ số ảnh của tập huấn luyện cố định được đưa vào
\textit{labeled subset}; mỗi ảnh được chọn giữ toàn bộ các chú giải
\textit{bounding box} tương ứng. Các mức ngân sách \(1\%\), \(5\%\), \(10\%\) và \(20\%\) được lựa chọn
nhằm khảo sát mô hình trong một dải điều kiện thiếu hụt chú giải, từ mức
cực kỳ hạn chế đến mức có lượng dữ liệu có nhãn tương đối lớn hơn. Việc
đánh giá \textit{semi-supervised object detection} tại nhiều tỷ lệ dữ
liệu có nhãn thấp đã được sử dụng trong các giao thức thực nghiệm trước
đây. Chẳng hạn, STAC đánh giá mô hình trên MS-COCO tại các mức \(1\%\),
\(2\%\), \(5\%\) và \(10\%\), trong khi Soft Teacher báo cáo kết quả tại
các mức \(1\%\), \(5\%\) và \(10\%\)
\cite{sohn2020simple,xu2021end}. Unbiased Teacher cũng khảo sát hiệu quả
của \textit{semi-supervised object detection} trong nhiều điều kiện dữ
liệu có nhãn hạn chế \cite{liu2021unbiased}. Những công trình này cung
cấp cơ sở thực nghiệm cho việc sử dụng các mức \(1\%\), \(5\%\) và
\(10\%\) để biểu diễn các chế độ thiếu nhãn khác nhau và tạo điều kiện
đối chiếu với các nghiên cứu trước.

Trong lĩnh vực X-quang ngực, Ji và cộng sự khảo sát bài toán định vị bất
thường theo thiết lập \textit{weakly semi-supervised} trên RSNA và
VinDr-CXR với nhiều tỷ lệ ảnh có chú giải \textit{bounding box}, bao gồm
các mức \(5\%\), \(10\%\) và \(20\%\)
\cite{ji2022benchmark}. Mặc dù thiết lập chú giải yếu của nghiên cứu này
khác với bài toán \textit{semi-supervised object detection} sử dụng dữ liệu chưa gán nhãn trong luận văn, công trình này cho thấy việc khảo
sát nhiều ngân sách nhãn, bao gồm mức \(20\%\), đã được sử dụng trong
bối cảnh định vị bất thường trên ảnh X-quang ngực.

Trên cơ sở đó, nghiên cứu hiện tại sử dụng \(1\%\) làm điều kiện chú giải
cực kỳ hạn chế; \(5\%\) và \(10\%\) làm các mức trung gian; và \(20\%\)
làm điểm quan sát có ngân sách nhãn lớn hơn. Mức \(20\%\) cho phép khảo
sát liệu lợi ích của việc khai thác dữ liệu chưa gán nhãn còn được duy trì khi \textit{labeled subset} đã được mở rộng đáng kể. Bốn mức ngân sách này không được xem là các ngưỡng tối ưu hoặc có tính phổ quát, mà là các điểm thực nghiệm được xác định trước nhằm mô tả sự thay đổi của kết quả theo lượng dữ liệu có nhãn.

Cụ thể, tại ngân sách \(1\%\), \textit{labeled subset} chỉ gồm 34 ảnh,
trong khi nghiệm cuối vẫn phải thỏa mãn ràng buộc bao phủ đủ 14 lớp bất
thường. Khi ngân sách tăng lên \(5\%\), \(10\%\) và \(20\%\), số ảnh có
nhãn tương ứng tăng lên 171, 343 và 685 ảnh. Để kiểm soát sự thay đổi về \textit{membership} giữa các mức ngân sách, các \textit{labeled subset} được xây dựng theo thiết kế \textit{nested}. Do đó, ngân sách lớn hơn giữ lại toàn bộ ảnh đã được lựa chọn ở ngân sách nhỏ hơn và chỉ bổ sung ảnh mới.

Gọi \(T\) là tập huấn luyện cố định, \(L_b\) là \textit{labeled subset} tại
ngân sách \(b\), và \(U_b\) là \textit{unlabeled subset} tương ứng. Với mỗi
ngân sách, \textit{unlabeled subset} được xác định là phần bù của \textit{labeled subset} trong \(T\):

\begin{equation}
U_b=T\setminus L_b,
\qquad
L_b\cap U_b=\varnothing,
\qquad
L_b\cup U_b=T.
\label{eq:labeled_unlabeled_complement}
\end{equation}

Gọi \(\operatorname{RHU}(\cdot)\) là phép làm tròn
\textit{round-half-up}. Với \(N=\lvert T\rvert=3.426\) là tổng số ảnh
trong tập huấn luyện cố định, số ảnh mục tiêu của
\textit{labeled subset} tại ngân sách \(b\) được xác định bởi

\begin{equation}
n_b
=
\operatorname{RHU}
\left(
bN
\right),
\label{eq:labeled_budget_size}
\end{equation}

trong đó \(b\in\{0{,}01;0{,}05;0{,}10;0{,}20\}\). Tập huấn luyện chứa
350 ảnh \textit{No Finding}; do đó, số ảnh \textit{No Finding} mục tiêu
trong \(L_b\) được xác định theo tỷ lệ tương ứng trong \(T\):

\begin{equation}
n_b^{\mathrm{NF}}
=
\operatorname{RHU}
\left(
n_b\frac{350}{3.426}
\right).
\label{eq:labeled_budget_no_finding_size}
\end{equation}

Quy mô cuối cùng của các cặp \textit{labeled subset}/\textit{unlabeled subset} được
trình bày trong Bảng~\ref{tab:labeled_unlabeled_budgets}.

\begin{table*}[htbp]
    \centering
    \small
    \setlength{\tabcolsep}{8pt}
    \caption{Quy mô các tập con có nhãn và chưa gán nhãn trong tập huấn luyện}
    \label{tab:labeled_unlabeled_budgets}
    \begin{tabular}{lccc}
        \toprule
        \shortstack{\textbf{Ngân sách}\\\textbf{nhãn}} &
        \shortstack{\textbf{Số ảnh}\\\textbf{có nhãn}} &
        \shortstack{\textbf{Số ảnh}\\\textbf{chưa gán nhãn}} &
        \shortstack{\textbf{\textit{No Finding}}\\
                    \textbf{trong tập có nhãn}} \\
        \midrule
        1\%  & 34  & 3.392 & 3  \\
        5\%  & 171 & 3.255 & 17 \\
        10\% & 343 & 3.083 & 35 \\
        20\% & 685 & 2.741 & 70 \\
        \bottomrule
    \end{tabular}
\end{table*}

Các \textit{labeled subset} được xây dựng theo thiết kế
\textit{nested}:

\begin{equation}
L_{1\%}
\subsetneq
L_{5\%}
\subsetneq
L_{10\%}
\subsetneq
L_{20\%}
\subsetneq T.
\label{eq:nested_labeled_subsets}
\end{equation}

Quá trình xây dựng được thực hiện tuần tự từ ngân sách nhỏ đến ngân sách
lớn. Khi chuyển sang ngân sách lớn hơn, toàn bộ ảnh đã thuộc
\textit{labeled subset} ở ngân sách trước được giữ cố định và không được phép loại ra trong các bước hiệu chỉnh tiếp theo. Ngân sách lớn vì vậy chỉ bổ sung ảnh mới vào thành phần đã có của ngân sách nhỏ hơn. Thiết kế này giúp kiểm soát sự thay đổi \textit{membership} khi khảo sát ảnh hưởng của lượng dữ
liệu có nhãn trên cùng một chuỗi tập con. Tuy nhiên, thiết kế
\textit{nested} không được giả định là luôn tốt hơn
\textit{independent sampling}; đây là một quyết định nhằm kiểm soát \textit{membership} giữa các mức ngân sách trong thiết kế thực nghiệm hiện tại.

Mỗi ảnh \(i\) trong \(T\) được biểu diễn bằng vector chỉ báo nhị phân

\begin{equation}
\mathbf{y}_i
=
\left(
y_{i,1},y_{i,2},\ldots,y_{i,14}
\right)
\in\{0,1\}^{14},
\label{eq:labeled_subset_class_vector}
\end{equation}

trong đó

\begin{equation}
y_{i,c}
=
\begin{cases}
1, & \text{nếu ảnh }i\text{ chứa ít nhất một bounding box thuộc lớp }c,\\
0, & \text{nếu ảnh }i\text{ không chứa bounding box thuộc lớp }c,
\end{cases}
\qquad c=1,\ldots,14.
\label{eq:labeled_subset_class_indicator}
\end{equation}

Biểu diễn này mã hóa \textit{class presence} ở cấp ảnh, không phụ thuộc
vào số lượng \textit{bounding box} của lớp đó. Vì một ảnh có thể đồng thời
chứa nhiều bất thường, vector \(\mathbf{y}_i\) có thể có nhiều phần tử bằng
1. Ảnh \textit{No Finding} không chứa \textit{ground-truth bounding box},
do đó có vector \(\mathbf{y}_i=\mathbf{0}\).

Bên cạnh 14 chỉ báo lớp, một chỉ báo riêng \(z_i\) được sử dụng để biểu diễn
trạng thái \textit{zero-GT}:

\begin{equation}
z_i
=
\begin{cases}
1, & \text{nếu ảnh }i\text{ không có ground-truth bounding box},\\
0, & \text{nếu ảnh }i\text{ có ít nhất một ground-truth bounding box}.
\end{cases}
\label{eq:labeled_subset_zero_gt_indicator}
\end{equation}

Ma trận sử dụng trong bước \textit{sampling} vì vậy gồm 14 chỉ báo
\textit{class presence} và một chỉ báo \textit{zero-GT}. Chỉ báo \(z_i\)
được sử dụng để hỗ trợ phân bổ ảnh \textit{No Finding}; nó không được đưa
vào không gian 14 lớp phát hiện và không làm cho \textit{No Finding} trở
thành một \textit{detection class} thứ 15.

\paragraph{Lựa chọn phương pháp sampling.}

Các phương pháp có thể được xem xét để tạo \textit{initial candidate} gồm
\textit{random sampling}, \textit{single-label stratification},
\textit{labelset-based stratification}, \textit{iterative multilabel
stratification}, \textit{second-order iterative stratification} và các
phương pháp \textit{evolutionary optimization}
\cite{sechidis2011stratification,szymanski2017network,
florezrevuelta2021evosplit}. Những phương pháp này kiểm soát các thuộc tính
phân phối khác nhau; không có phương pháp nào được xem là tối ưu phổ quát
cho mọi bộ dữ liệu đa nhãn.

\textit{Random sampling} không sử dụng thông tin nhãn trong quá trình lựa
chọn ảnh. Ưu điểm của phương pháp này là đơn giản và không áp đặt cấu trúc
phân tầng. Tuy nhiên, khi ngân sách nhãn rất nhỏ, \textit{random sampling}
có thể làm cho một lớp hiếm không xuất hiện trong \textit{labeled subset} hoặc
tạo ra sai lệch lớn về tỷ lệ hiện diện lớp. Rủi ro này đặc biệt đáng chú ý
trong nghiên cứu hiện tại vì \textit{labeled subset} nhỏ nhất chỉ có 34 ảnh
nhưng phải bao phủ 14 lớp bất thường. Vì vậy, quy trình xây dựng
\textit{labeled subset} trong nghiên cứu không sử dụng
\textit{random sampling} thuần túy, mà sử dụng thông tin hiện diện lớp ở
cấp ảnh trong quá trình tạo tập con.

\textit{Single-label stratification} yêu cầu mỗi ảnh chỉ thuộc một nhóm
nhãn. Để áp dụng phương pháp này cho dữ liệu đa nhãn, cần chọn một nhãn đại
diện hoặc nhãn chính cho từng ảnh. Việc quy giản như vậy làm mất thông tin
về các lớp đồng xuất hiện: tỷ lệ của nhãn đại diện có thể được duy trì,
trong khi tỷ lệ của các nhãn còn lại vẫn bị sai lệch. Do đó,
\textit{single-label stratification} không phù hợp với cấu trúc của dữ liệu
X-quang trong nghiên cứu, nơi một ảnh có thể đồng thời chứa nhiều bất
thường.

\textit{Labelset-based stratification} xem mỗi vector nhãn khác nhau như
một nhóm riêng. Phương pháp này có khả năng phản ánh cấu trúc đồng xuất hiện
đầy đủ hơn so với việc xem xét từng lớp độc lập. Tuy nhiên, với 14 chỉ báo
nhị phân, không gian lý thuyết có thể chứa tới \(2^{14}\) tổ hợp nhãn. Khi
nhiều tổ hợp chỉ xuất hiện trên một số ít ảnh, các nhóm phân tầng trở nên
thưa và khó được duy trì đồng thời, đặc biệt ở ngân sách 1\%.

\textit{Second-order iterative stratification} mở rộng
\textit{iterative multilabel stratification} bằng cách kiểm soát thêm phân
phối của các cặp nhãn \cite{szymanski2017network}. Phương pháp này phù hợp
khi mục tiêu chính là duy trì cấu trúc đồng xuất hiện bậc hai. Tuy nhiên,
các thuộc tính phân phối được kiểm soát trong nghiên cứu hiện tại là tỷ lệ
hiện diện của từng lớp và \textit{label cardinality}, không phải toàn bộ
phân phối cặp nhãn. Việc bổ sung các ràng buộc theo cặp, đặc biệt tại ngân
sách 1\%, sẽ làm tăng số điều kiện rời rạc phải thỏa mãn và thay đổi mục
tiêu thiết kế tập con.

Các phương pháp \textit{evolutionary optimization} hoặc
\textit{solver-based optimization} có thể tìm kiếm đồng thời nhiều thuộc
tính phân phối và mở rộng phạm vi tìm kiếm vượt ra ngoài một quy tắc
\textit{sampling} tham lam \cite{florezrevuelta2021evosplit}. Đổi lại,
chúng yêu cầu thêm các lựa chọn về \textit{fitness function},
\textit{population size}, điều kiện dừng, \textit{optimality gap} hoặc cấu
hình bộ giải. Những lựa chọn này làm tăng số \textit{hyperparameter} của
chính bước xây dựng dữ liệu và đòi hỏi một thiết kế đánh giá riêng.

Nghiên cứu lựa chọn \textit{iterative multilabel stratification} vì ba lý
do. Thứ nhất, phương pháp sử dụng trực tiếp vector hiện diện đa nhãn ở cấp
ảnh, do đó thông tin lớp được sử dụng trong bước xây dựng
\textit{initial candidate} thay vì thực hiện lấy mẫu mù nhãn. Thứ hai,
phương pháp hướng đến duy trì phân phối biên của từng nhãn, phù hợp trực
tiếp với mục tiêu kiểm soát tỷ lệ ảnh chứa từng lớp bất thường trong các
\textit{labeled subset} \cite{sechidis2011stratification}. Thứ ba, phương
pháp không phụ thuộc vào \textit{feature embedding}, kết quả dự đoán hoặc
hiệu năng phát hiện; vì vậy, \textit{membership} có thể được xác định trước
khi huấn luyện mô hình.

Sự lựa chọn này không hàm ý rằng \textit{iterative multilabel
stratification} luôn tốt hơn các phương pháp còn lại. Phương pháp được chọn
vì phù hợp với các thuộc tính được kiểm soát trong nghiên cứu, gồm phân phối
biên của 14 lớp, sự hiện diện của các lớp hiếm, tỷ lệ ảnh
\textit{No Finding} và khả năng tái lập \textit{membership} của tập con.

Tại mỗi ngân sách, \textit{iterative multilabel stratification} được áp
dụng trên phần dữ liệu chưa thuộc tập con đã cố định ở ngân sách trước để
tạo \textit{initial candidate}. \texttt{partition\_seed}=42 được xác định
trước và được truyền vào bộ chia \textit{multilabel stratified} để kiểm
soát thành phần nghiệm khởi đầu. Trong các bước hiệu chỉnh có nhiều
\textit{candidate move} tương đương theo tiêu chí ưu tiên phía trước,
\texttt{partition\_seed}=42 còn được sử dụng trong quy tắc
\textit{tie-break} có tính xác định. Seed này được giữ nguyên ở cả bốn mức
ngân sách; không thực hiện \textit{seed search}, không thay đổi seed giữa
các ngân sách và không lựa chọn seed dựa trên thống kê phân phối sau
\textit{sampling} hoặc kết quả của mô hình.

Do nghiên cứu mô phỏng hồi cứu các mức thiếu hụt chú giải, nhãn đầy đủ của
tập huấn luyện được sử dụng trong bước xây dựng phân hoạch để tạo các chỉ
báo \textit{class presence}, thực hiện \textit{sampling} và kiểm tra các
ràng buộc phân phối. Vì thông tin lớp được sử dụng trong quá trình xây dựng
và nghiệm cuối phải thỏa mãn yêu cầu bao phủ đủ 14 lớp bất thường, quy
trình này là lấy mẫu có sử dụng thông tin nhãn và có ràng buộc bao phủ lớp
(\textit{label-aware, coverage-constrained subset construction}), không
phải \textit{random sampling} thuần túy. Sau khi \textit{membership} được
cố định, các chú giải của \(U_b\) bị loại khỏi tệp cung cấp cho nhánh
\textit{unlabeled}. Vì vậy, các chú giải này chỉ phục vụ việc xây dựng và
kiểm định phân hoạch ngoại tuyến, không được sử dụng làm tín hiệu giám sát
trong quá trình huấn luyện \textit{semi-supervised}.

\paragraph{Deterministic constrained repair và local search.}

\textit{Iterative multilabel stratification} chỉ tạo
\textit{initial candidate} có phân phối nhãn gần với tập huấn luyện; bước
này không tự nó được xem là cơ chế bảo đảm tất cả các ràng buộc khả thi của
nghiên cứu. Do ảnh là đơn vị không thể chia nhỏ, \textit{initial candidate}
có thể chưa đạt chính xác kích thước ngân sách, số ảnh
\textit{No Finding} mục tiêu hoặc yêu cầu bao phủ đủ 14 lớp bất thường.
Trong giao thức xây dựng tập con, sự hiện diện của cả 14 lớp được xem là
một ràng buộc khả thi bắt buộc đối với nghiệm cuối. Vì vậy,
\textit{deterministic constrained repair} được thực hiện sau
\textit{sampling} để đưa nghiệm khởi đầu về một nghiệm khả thi và sau đó
tiếp tục cải thiện các đặc trưng phân phối.

Cần nhấn mạnh rằng \textit{sampling} và \textit{local search} đảm nhiệm hai
vai trò khác nhau. \textit{Iterative multilabel stratification} xây dựng
nghiệm khởi đầu từ tập ứng viên; \textit{local search} nhận nghiệm đó làm
điểm xuất phát và khảo sát các nghiệm lân cận. Vì vậy,
\textit{one-for-one local search} không thay thế
\textit{iterative multilabel stratification}, mà là bước refinement được
thực hiện sau \textit{sampling}.

Các lựa chọn có thể được xem xét cho bước refinement gồm
\textit{single-image add/remove}, \textit{one-for-one swap},
\textit{two-for-two swap}, \textit{\(k\)-for-\(k\) swap},
\textit{random restart}, \textit{simulated annealing},
\textit{tabu search}, \textit{mixed-integer programming} và
\textit{constraint programming}. Các phương pháp này khác nhau về kích
thước lân cận, khả năng thoát khỏi \textit{local optimum}, số
\textit{hyperparameter}, chi phí tìm kiếm và mức độ có thể tái lập.

\textit{Single-image add/remove} phù hợp với bước hiệu chỉnh kích thước,
nhưng không phù hợp với giai đoạn cải thiện phân phối sau khi kích thước đã
được cố định, vì mỗi phép thêm hoặc loại đơn lẻ làm thay đổi số ảnh của tập.
\textit{Two-for-two swap} và \textit{\(k\)-for-\(k\) swap} tạo ra lân cận
rộng hơn và có thể phát hiện một cải thiện không thể đạt được bằng một phép
\textit{one-for-one swap}. Tuy nhiên, số tổ hợp cần xem xét tăng nhanh theo
\(k\), và việc mở rộng lân cận không tự động cung cấp bảo đảm
\textit{global optimum}.

Các phương pháp như \textit{simulated annealing} hoặc
\textit{tabu search} có thể thoát khỏi một số \textit{local optimum} bằng
cách chấp nhận tạm thời trạng thái không cải thiện hoặc sử dụng bộ nhớ tìm
kiếm. Đổi lại, chúng yêu cầu thêm các lựa chọn như
\textit{temperature schedule}, xác suất chấp nhận, \textit{tabu tenure} và
điều kiện dừng. Tương tự, \textit{random restart} có thể tạo ra nhiều
nghiệm khởi đầu, nhưng việc chọn nghiệm tốt nhất sau khi quan sát các kết
quả sẽ biến seed thành một tham số được tìm kiếm và làm tăng nguy cơ
\textit{selection bias}.

Các phương pháp \textit{mixed-integer programming} hoặc
\textit{constraint programming} có thể biểu diễn trực tiếp các ràng buộc
về kích thước, bao phủ lớp, số ảnh \textit{No Finding} và quan hệ
\textit{nested}. Trong một số trường hợp, bộ giải còn có thể cung cấp
\textit{optimality bound} hoặc \textit{optimality certificate}. Tuy nhiên,
cách tiếp cận này đòi hỏi xây dựng một mô hình tối ưu riêng, lựa chọn bộ
giải và quy định rõ \textit{optimality gap}, giới hạn tài nguyên và cách xử
lý trường hợp bộ giải không hoàn tất. Đây là một thiết kế phương pháp khác,
không phải sự mở rộng trực tiếp của quy trình \textit{sampling} được sử
dụng trong nghiên cứu.

\textit{One-for-one swap} được lựa chọn làm lân cận của
\textit{local search} vì đây là thay đổi nhỏ nhất có thể điều chỉnh thành
phần tập con mà vẫn giữ nguyên kích thước theo cấu trúc: một ảnh được loại
ra và một ảnh khác được bổ sung. Khi hai ảnh có cùng trạng thái
\textit{zero-GT}, phép hoán đổi còn giữ nguyên số ảnh
\textit{No Finding}. Việc chỉ cho phép loại ảnh khỏi phần mới được bổ sung
cũng bảo toàn \textit{membership} đã cố định của các ngân sách nhỏ hơn.
Do đó, cùng một phép biến đổi có thể duy trì đồng thời kích thước, số ảnh
\textit{No Finding} và quan hệ \textit{nested}, trong khi vẫn thay đổi
vector \textit{class presence} để cải thiện phân phối.

Một lý do khác là toàn bộ lân cận \textit{one-for-one} khả dụng có thể
được duyệt tại mỗi vòng lặp mà không cần lấy mẫu ngẫu nhiên các phép chuyển
hoặc đặt thêm giới hạn về số ứng viên. Sau mỗi phép hoán đổi được chấp
nhận, lân cận được xây dựng lại trên nghiệm mới và tiếp tục được duyệt cho
đến khi không còn phép hoán đổi khả dụng nào cải thiện mục tiêu. Cơ chế
này tạo ra một điều kiện dừng tường minh và có thể kiểm tra được.

Như vậy, quy trình xây dựng mỗi tập \(L_b\) gồm hai giai đoạn bổ sung cho
nhau. \textit{Iterative multilabel stratification} tạo
\textit{initial candidate} có định hướng phân phối; sau đó,
\textit{deterministic constrained repair} và
\textit{one-for-one local search} đưa nghiệm khởi đầu về đúng các ràng buộc
và cải thiện mục tiêu phân phối. Sự kết hợp này được lựa chọn nhằm cân bằng khả năng duy trì phân phối đa nhãn, bảo toàn các ràng buộc, tính xác định và khả năng duyệt đầy đủ lân cận được lựa chọn.

Toàn bộ quy trình xây dựng các tập con có nhãn và chưa gán nhãn được khái
quát trong Hình~\ref{fig:labeled_unlabeled_construction}.

\begin{figure*}[t]
    \centering
    \includegraphics[width=\textwidth]
    {z_quy-trinh-xay-dung-labeled-unlabeled.png}
    \caption{Quy trình xây dựng các tập con có nhãn và chưa gán nhãn từ
    tập huấn luyện cố định. \textit{Iterative multilabel stratification}
    tạo \textit{initial candidate} tại từng mức ngân sách;
    \textit{deterministic constrained repair} tiếp tục hiệu chỉnh kích
    thước, số ảnh \textit{No Finding}, mức bao phủ lớp và các đặc trưng
    phân phối. Các tập có nhãn được xây dựng theo quan hệ \textit{nested},
    trong khi tập chưa gán nhãn được xác định là phần bù của tập có nhãn
    trong tập huấn luyện.}
    \label{fig:labeled_unlabeled_construction}
\end{figure*}

\paragraph{Hàm mục tiêu và thứ tự hiệu chỉnh.}

Với mỗi lớp \(c\), gọi \(p_{L_b,c}\) và \(p_{T,c}\) lần lượt là
\textit{image-level prevalence} của lớp đó trong \(L_b\) và \(T\). Ba chỉ
số sai lệch được sử dụng gồm

\begin{equation}
D_{\max}
=
\max_{c=1,\ldots,14}
\left|
p_{L_b,c}-p_{T,c}
\right|,
\label{eq:labeled_subset_dmax}
\end{equation}

\begin{equation}
D_{\mathrm{mean}}
=
\frac{1}{14}
\sum_{c=1}^{14}
\left|
p_{L_b,c}-p_{T,c}
\right|,
\label{eq:labeled_subset_dmean}
\end{equation}

và

\begin{equation}
D_{\mathrm{LC}}
=
\left|
LC(L_b)-LC(T)
\right|,
\label{eq:labeled_subset_dlc}
\end{equation}

trong đó \(LC(\cdot)\) là \textit{label cardinality}, tức số lớp bất
thường trung bình hiện diện trên mỗi ảnh.

Các mục tiêu được ưu tiên theo thứ tự từ điển

\begin{equation}
D_{\max}
\;\longrightarrow\;
D_{\mathrm{mean}}
\;\longrightarrow\;
D_{\mathrm{LC}}.
\label{eq:labeled_subset_objective_order}
\end{equation}

Thứ tự này có nghĩa là một phương án có \(D_{\max}\) nhỏ hơn luôn được ưu
tiên; \(D_{\mathrm{mean}}\) chỉ được xét khi \(D_{\max}\) bằng nhau, và
\(D_{\mathrm{LC}}\) chỉ được xét khi hai thành phần trước cùng bằng nhau.

Để tránh sai số số thực trong quá trình lựa chọn, việc so sánh các nghiệm
được thực hiện bằng một bộ ba số nguyên tương đương. Gọi
\(n_b=\lvert L_b\rvert\), \(N=\lvert T\rvert\), \(k_{b,c}\) là số ảnh chứa
lớp \(c\) trong \(L_b\), \(K_c\) là số ảnh chứa lớp \(c\) trong \(T\),
\(S_{L_b}=\sum_{c=1}^{14}k_{b,c}\), và
\(S_T=\sum_{c=1}^{14}K_c\). Khi đó,

\begin{equation}
\begin{aligned}
E_{\max}
&=
\max_{c=1,\ldots,14}
\left|
k_{b,c}N-K_cn_b
\right|,\\
E_{\mathrm{mean}}
&=
\sum_{c=1}^{14}
\left|
k_{b,c}N-K_cn_b
\right|,\\
E_{\mathrm{LC}}
&=
\left|
S_{L_b}N-S_Tn_b
\right|.
\end{aligned}
\label{eq:labeled_subset_integer_objective}
\end{equation}

Trong cùng một vòng lặp, \(n_b\) và \(N\) không thay đổi giữa các
\textit{candidate move}; vì vậy, phép so sánh theo

\begin{equation}
\left(
E_{\max},
E_{\mathrm{mean}},
E_{\mathrm{LC}}
\right)
\label{eq:labeled_subset_integer_objective_tuple}
\end{equation}

bảo toàn thứ tự từ điển của ba mục tiêu ban đầu mà không cần thực hiện phép
chia số thực.

Quy trình \textit{repair} được thực hiện theo bốn bước theo thứ tự cố định.
Thứ nhất, \textit{exact-size repair} sử dụng các phép
\textit{add/remove} để đưa số ảnh về đúng \(n_b\). Thứ hai,
\textit{exact-No-Finding repair} sử dụng \textit{one-for-one swap} giữa
hai ảnh có trạng thái \textit{zero-GT} đối nghịch để đưa số ảnh
\textit{No Finding} về đúng \(n_b^{\mathrm{NF}}\) trong khi vẫn giữ
nguyên kích thước. Thứ ba, \textit{minimum-class-coverage repair} duyệt
các phép \textit{one-for-one swap} khả dụng và chỉ chấp nhận những phép
chuyển làm giảm nghiêm ngặt số lớp còn thiếu, đồng thời giữ nguyên kích
thước, số ảnh \textit{No Finding} và \textit{nested membership}. Khi có
nhiều phép chuyển thỏa điều kiện, phương án được ưu tiên theo thứ tự:
số lớp còn thiếu sau phép chuyển, bộ ba mục tiêu nguyên, quy tắc
\textit{tie-break} dựa trên \texttt{partition\_seed}, và cuối cùng là
thứ tự định danh số học của phép chuyển. Thứ tư,
\textit{objective repair} áp dụng \textit{one-for-one local search} để
cải thiện nghiêm ngặt bộ ba mục tiêu nguyên trong khi bảo toàn toàn bộ các
ràng buộc đã đạt được.

Quy trình này không sử dụng chiến lược \textit{greedy class-by-class}
trong đó các ảnh được chọn tuần tự để lần lượt ``phủ'' từng lớp. Nghiệm
khởi đầu được tạo đồng thời bằng \textit{iterative multilabel
stratification}; các quyết định tuần tự chỉ xuất hiện trong các bước
\textit{repair}, nơi mỗi phép chuyển được đánh giá theo các ràng buộc và
thứ tự ưu tiên đã xác định trước.

Trong \textit{objective repair}, một \textit{candidate move} chỉ được xem
là khả dụng khi giữ nguyên kích thước, số ảnh \textit{No Finding}, bao phủ
đủ 14 lớp và không loại ảnh thuộc tập con đã cố định ở ngân sách nhỏ hơn.
Một phép chuyển có mục tiêu bằng với nghiệm hiện tại không được chấp nhận.
Khi nhiều phép chuyển có cùng giá trị mục tiêu, một quy tắc
\textit{tie-break} xác định dựa trên \texttt{partition\_seed} và
\texttt{image\_id} được sử dụng để lựa chọn duy nhất một phương án.

\paragraph{Phạm vi của kết luận local optimum.}

\textit{Objective repair} kết thúc khi toàn bộ lân cận
\textit{one-for-one} khả dụng đã được duyệt mà không còn
\textit{candidate move} nào cải thiện nghiêm ngặt bộ ba mục tiêu. Nghiệm
cuối vì vậy đạt \textit{one-for-one local optimum}.

Kết luận này chỉ áp dụng đối với lân cận được định nghĩa bởi một phép
\textit{one-for-one swap}. Nó không phải bằng chứng về
\textit{global optimum} trong toàn bộ không gian tập con khả thi. Cụ thể,
một \textit{two-for-two swap}, một \textit{\(k\)-for-\(k\) swap} hoặc một
chuỗi phép chuyển đi qua trạng thái tạm thời không cải thiện vẫn có thể dẫn
đến một nghiệm tốt hơn. Do đó, kết quả không được diễn giải là tập con tối
ưu toàn cục hoặc tốt nhất trong toàn bộ không gian tổ hợp.

\paragraph{Kết quả phân phối.}

Ảnh \textit{No Finding} tiếp tục được giữ dưới dạng ảnh
\textit{zero-GT}, không được chuyển thành một \textit{detection class}
riêng. Số ảnh \textit{No Finding} trong bốn tập \(L_b\) lần lượt là 3, 17,
35 và 70. Cả 14 lớp bất thường đều hiện diện trong mỗi
\textit{labeled subset}. Riêng tại ngân sách \(1\%\), bước
\textit{minimum-class-coverage repair} không cần thực hiện phép hoán đổi
nào, cho thấy nghiệm sau các bước trước đó đã thỏa mãn yêu cầu bao phủ
14 lớp; tuy nhiên, ràng buộc bao phủ vẫn được kiểm tra và giữ bắt buộc
trong giao thức xây dựng tập con. Sự hiện diện đầy đủ của các lớp không
đồng nghĩa phân phối lớp được bảo toàn hoàn hảo.

Mức độ duy trì phân phối được mô tả bằng \(D_{\max}\),
\(D_{\mathrm{mean}}\) và \(D_{\mathrm{LC}}\). Cả ba chỉ số đều không âm;
giá trị nhỏ hơn biểu thị \textit{labeled subset} gần với tập huấn luyện hơn
theo thuộc tính tương ứng. Kết quả tại bốn mức ngân sách được trình bày
trong Bảng~\ref{tab:labeled_distribution_deviation}.

\begin{table*}[htbp]
    \centering
    \small
    \setlength{\tabcolsep}{6pt}
    \caption{Sai lệch phân phối của các tập có nhãn so với tập huấn luyện}
    \label{tab:labeled_distribution_deviation}
    \begin{tabular}{lccclc}
        \toprule
        \shortstack{\textbf{Ngân sách}\\\textbf{nhãn}} &
        \(D_{\max}\) &
        \(D_{\mathrm{mean}}\) &
        \(D_{\mathrm{LC}}\) &
        \shortstack{\textbf{Lớp có sai lệch}\\
                    \textbf{tỷ lệ lớn nhất}} &
        \shortstack{\textbf{Tổng số phép}\\
                    \textbf{hiệu chỉnh}} \\
        \midrule
        1\%
            & 0,0130 & 0,0070 & 0,0134
            & \textit{Consolidation}      & 3  \\
        5\%
            & 0,0029 & 0,0015 & 0,0141
            & \textit{Atelectasis}        & 7  \\
        10\%
            & 0,0013 & 0,0007 & 0,0031
            & \textit{Pleural thickening} & 10 \\
        20\%
            & 0,0007 & 0,0003 & 0,0018
            & \textit{Lung Opacity}       & 26 \\
        \bottomrule
    \end{tabular}
\end{table*}

Các giá trị trong Bảng~\ref{tab:labeled_distribution_deviation} được biểu
diễn dưới dạng tỷ lệ. Theo đơn vị điểm phần trăm, \(D_{\max}\) lần lượt
tương ứng với 1,2980; 0,2857; 0,1278 và 0,0663 điểm phần trăm tại các ngân
sách 1\%, 5\%, 10\% và 20\%. Kết quả cho thấy sai lệch lớn nhất về tỷ lệ
hiện diện lớp giảm khi ngân sách nhãn tăng. Khi số ảnh có nhãn lớn hơn, số
lượng nghiệm khả thi có khả năng biểu diễn gần đúng phân phối của tập huấn
luyện cũng tăng lên.

Tuy nhiên, \(D_{\mathrm{LC}}\) không giảm đơn điệu giữa ngân sách 1\% và
5\%; giá trị tăng từ 0,0134 lên 0,0141 trước khi giảm ở hai ngân sách lớn
hơn. Kết quả này cho thấy việc giảm sai lệch tỷ lệ hiện diện của từng lớp
không nhất thiết đồng thời làm giảm sai lệch về số lớp trung bình trên mỗi
ảnh. Hiện tượng này phù hợp với thứ tự từ điển của hàm mục tiêu, trong đó
\(D_{\max}\) được ưu tiên trước \(D_{\mathrm{mean}}\) và
\(D_{\mathrm{LC}}\), đồng thời nghiệm phải thỏa mãn các ràng buộc về kích
thước, số ảnh \textit{No Finding}, bao phủ lớp và quan hệ
\textit{nested}.

Cột tổng số phép hiệu chỉnh bao gồm tất cả thay đổi được chấp nhận trong
\textit{exact-size repair}, \textit{exact-No-Finding repair},
\textit{minimum-class-coverage repair} và \textit{objective repair}. Đại
lượng này mô tả mức độ thay đổi từ \textit{initial candidate} đến tập cuối
cùng; nó không phải thước đo chất lượng phân phối, chi phí tính toán hoặc
số \textit{one-for-one swap} riêng của \textit{objective repair}. Do quy
mô phần dữ liệu mới được bổ sung khác nhau giữa các ngân sách, số phép hiệu
chỉnh cũng không nên được sử dụng như một chỉ số hiệu quả độc lập.

Giới hạn rời rạc thể hiện rõ nhất tại ngân sách 1\%, nơi chỉ có 34 ảnh
nhưng vẫn phải đồng thời bao phủ 14 lớp và duy trì số ảnh
\textit{No Finding} mục tiêu. Vì ảnh là đơn vị không thể chia nhỏ, một ảnh
có thể chứa nhiều lớp và một số lớp có tần suất thấp hoặc đồng xuất hiện
với các lớp khác, không phải mọi tỷ lệ lớp trong tập huấn luyện đều có thể
được tái tạo chính xác. Các chỉ số trên vì vậy chỉ mô tả mức độ duy trì
phân phối theo các thuộc tính được đo, không chứng minh rằng toàn bộ phân
phối chú giải đã được bảo toàn.

\paragraph{Kiểm định và khả năng tái lập.}

Các tệp dữ liệu sau khi được ghi được đọc lại độc lập trước khi cố định.
Quá trình kiểm định xác nhận: (i) kích thước của các \textit{labeled subset} và \textit{unlabeled subset} đạt đúng mục tiêu; (ii) mọi
\texttt{image\_id} đều thuộc tập huấn luyện cố định; (iii) \(L_b\) và
\(U_b\) không giao nhau và hợp của chúng khôi phục toàn bộ \(T\);
(iv) quan hệ \textit{nested} của các \textit{labeled subset} và của riêng các
ảnh \textit{No Finding} được duy trì; (v) ánh xạ
\texttt{category\_id}--tên lớp nhất quán với tập huấn luyện; và
(vi) các tham chiếu giữa ảnh, \textit{annotation} và lớp trong các tệp COCO có nhãn vẫn hợp lệ.

Các kết quả kiểm định định lượng theo từng mức ngân sách được trình bày
trong Bảng~\ref{tab:labeled_unlabeled_integrity_validation}.

\begin{table*}[htbp]
    \centering
    \footnotesize
    \setlength{\tabcolsep}{3pt}
    \renewcommand{\arraystretch}{1.15}
    \caption{Kết quả kiểm định quy mô và tính toàn vẹn của các tập con có
    nhãn và chưa gán nhãn}
    \label{tab:labeled_unlabeled_integrity_validation}

    \begin{tabularx}{\textwidth}{
        @{}
        >{\raggedright\arraybackslash}X
        *{4}{>{\centering\arraybackslash}p{0.115\textwidth}}
        @{}
    }
        \toprule
        \textbf{Nội dung kiểm định} &
        \textbf{1\%} &
        \textbf{5\%} &
        \textbf{10\%} &
        \textbf{20\%} \\
        \midrule

        \(\lvert L_b\rvert\), quan sát/mục tiêu
            & \(34/34\)
            & \(171/171\)
            & \(343/343\)
            & \(685/685\) \\

        \(\lvert U_b\rvert\), quan sát/mục tiêu
            & \(3.392/3.392\)
            & \(3.255/3.255\)
            & \(3.083/3.083\)
            & \(2.741/2.741\) \\

        \textit{No Finding} trong \(L_b\), quan sát/mục tiêu
            & \(3/3\)
            & \(17/17\)
            & \(35/35\)
            & \(70/70\) \\

        Số lớp hiện diện trong \(L_b\)
            & \(14/14\)
            & \(14/14\)
            & \(14/14\)
            & \(14/14\) \\

        \(\lvert L_b\cap U_b\rvert\)
            & 0
            & 0
            & 0
            & 0 \\

        \(\lvert L_b\cup U_b\rvert\)
            & 3.426
            & 3.426
            & 3.426
            & 3.426 \\

        Số \texttt{image\_id} ngoài \(T\)
            & 0
            & 0
            & 0
            & 0 \\

        Tái tạo \textit{membership}
            & Trùng khớp
            & Trùng khớp
            & Trùng khớp
            & Trùng khớp \\

        \bottomrule
    \end{tabularx}
\end{table*}

Kết quả cho thấy quy mô quan sát của \(L_b\), \(U_b\) và số ảnh
\textit{No Finding} trùng với các giá trị mục tiêu tại cả bốn mức ngân
sách. Mỗi tập \(L_b\) chứa đủ 14 lớp bất thường. Đồng thời,
\(\lvert L_b\cap U_b\rvert=0\) và
\(\lvert L_b\cup U_b\rvert=3.426\), xác nhận rằng hai tập không giao nhau
và hợp của chúng khôi phục toàn bộ tập huấn luyện. Không ghi nhận
\texttt{image\_id} nằm ngoài \(T\). Khi quy trình được thực hiện lại với
cùng dữ liệu đầu vào, cấu hình và \texttt{partition\_seed}, kết quả tái
tạo \textit{membership} trùng khớp tại cả \(4/4\) mức ngân sách.

Ba trên ba quan hệ bao hàm giữa các mức ngân sách kế tiếp đều được duy trì
đối với \textit{labeled subset}:

\begin{equation}
L_{1\%}\subsetneq L_{5\%},\qquad
L_{5\%}\subsetneq L_{10\%},\qquad
L_{10\%}\subsetneq L_{20\%}.
\label{eq:validated_strict_nested_labeled_subsets}
\end{equation}

Tương tự, cả \(3/3\) quan hệ \textit{nested} tương ứng của riêng các ảnh
\textit{No Finding} cũng được duy trì.

Đối với các tệp COCO \textit{unlabeled}, trường
\texttt{annotations} được đặt thành mảng rỗng. Các trường ở cấp ảnh được
suy ra trực tiếp từ \textit{ground truth}, chẳng hạn trạng thái ảnh âm tính
hoặc vector \textit{class presence}, cũng được loại khỏi tệp cung cấp cho
nhánh \textit{unlabeled}. Quá trình đọc lại không ghi nhận
\textit{annotation} hoặc trường dẫn xuất từ \textit{ground truth} còn tồn
tại trong các tệp này.

Không phát hiện \texttt{image\_id} của \textit{validation set} hoặc
\textit{test set} trong các tập \(L_b\) và \(U_b\). Kết quả này chỉ xác
nhận sự tách biệt trong phạm vi \texttt{image\_id}, cấu trúc COCO và các
trường đã được kiểm tra; nó không được suy rộng thành bằng chứng loại trừ
\textit{patient-level leakage} khi định danh bệnh nhân không khả dụng.

Sau khi hoàn tất quá trình xây dựng và kiểm định, thành phần
\texttt{image\_id} của từng tập con được cố định để sử dụng thống nhất
trong các thí nghiệm tiếp theo. Khi quy trình được thực hiện lại với cùng
dữ liệu đầu vào, cùng cấu hình và cùng \texttt{partition\_seed}, thành
phần ảnh của cả bốn mức ngân sách được tái tạo nhất quán.

Hai loại seed được sử dụng cho hai mục đích độc lập.
\texttt{partition\_seed} được sử dụng trong quá trình lấy mẫu và hiệu
chỉnh có tính xác định để xây dựng các tập dữ liệu, trong khi
\texttt{training\_seed} được sử dụng để kiểm soát các nguồn ngẫu nhiên
phát sinh khi huấn luyện mô hình. Cụ thể,
\texttt{partition\_seed}=42 được ấn định trước, không được lựa chọn thông
qua \textit{seed search}, và được giữ cố định để xây dựng cũng như tái lập
thành phần ảnh của \(T\), \(L_b\) và \(U_b\). Sau khi các tập này được
khóa, việc thay đổi \texttt{training\_seed} chỉ tác động đến những thành
phần ngẫu nhiên của từng lần huấn luyện, chẳng hạn
\textit{model initialization}, thứ tự lấy mẫu của
\textit{data loader} và \textit{stochastic augmentation}; nó không kích
hoạt lại quá trình lấy mẫu dữ liệu và không làm thay đổi thành phần ảnh,
quan hệ \textit{nested} hoặc các ràng buộc phân hoạch đã xác lập. Giao
thức sử dụng nhiều \texttt{training\_seed} cho các lần huấn luyện được
trình bày trong Mục~\ref{subsec:experimental_protocol}.

Tổng hợp các kiểm định cho thấy bốn cặp
\textit{labeled subset}/\textit{unlabeled subset} đạt đúng kích thước mục tiêu, chỉ chứa
các ảnh thuộc tập huấn luyện cố định, duy trì quan hệ \textit{nested} và
thỏa mãn quan hệ phần bù tại từng mức ngân sách. Bước xây dựng dữ liệu này
không bao gồm huấn luyện mô hình, sinh \textit{pseudo-label}, lựa chọn
ngưỡng dự đoán hoặc đánh giá trên \textit{test set}. Các tập con đã cố định
được sử dụng làm đầu vào thống nhất cho các thí nghiệm
\textit{supervised} và \textit{semi-supervised} được trình bày ở các phần
tiếp theo.

\subsection{Chẩn đoán dữ liệu trước huấn luyện}
\label{subsec:pretraining_dataset_diagnostics}

\paragraph{Mục tiêu, phạm vi và nguyên tắc sử dụng dữ liệu.}

Sau khi thành phần của \textit{training set}, \textit{validation set},
\textit{test set} và các cặp
\textit{labeled subset}/\textit{unlabeled subset} đã được xác định theo
quy trình trình bày tại Mục~\ref{subsec:fixed_data_split} và
Mục~\ref{subsubsec:labeled_unlabeled_construction}, nghiên cứu thực hiện
bước chẩn đoán dữ liệu trước huấn luyện
(\textit{pre-training dataset diagnostics}). Mục tiêu của bước này là mô
tả có hệ thống cấu trúc của dữ liệu đầu vào trước khi mô hình được huấn
luyện, nhận diện các đặc trưng dữ liệu cần được lưu ý khi diễn giải kết
quả thực nghiệm và thiết lập một tập thống kê mô tả nhất quán để đối chiếu
với hiệu năng mô hình ở các giai đoạn sau.

Các phép chẩn đoán trong mục này có bản chất mô tả dữ liệu, không phải
thí nghiệm đánh giá mô hình. Vì vậy, quá trình này không bao gồm huấn
luyện mô hình, sinh \textit{pseudo-label}, lựa chọn
\textit{checkpoint}, tính AP/mAP, tìm kiếm siêu tham số, lựa chọn
\textit{confidence threshold}, điều chỉnh augmentation, thay đổi hàm mất
mát hoặc lựa chọn kiến trúc. Một đặc trưng được phát hiện trong dữ liệu
chỉ được ghi nhận như một yếu tố cần theo dõi; nó không tự động dẫn đến
việc thay đổi cấu hình huấn luyện.

Phân tích nội dung chi tiết được thực hiện trên tập huấn luyện \(T\) đã
được xác định ở bước phân hoạch trước đó. Bốn tập có nhãn
\(L_{1\%}\), \(L_{5\%}\), \(L_{10\%}\) và \(L_{20\%}\) cũng được phân
tích nhằm đánh giá mức bao phủ lớp và mức độ đại diện của các ngân sách
nhãn đã xây dựng. Thành phần ảnh của \(T\) và \(L_b\) được giữ nguyên
trong toàn bộ quá trình chẩn đoán; kết quả chẩn đoán không được sử dụng
để lấy mẫu lại hoặc thay đổi \textit{membership} của các tập này.

Đối với \textit{validation set} và \textit{test set}, nghiên cứu chỉ thực
hiện các kiểm tra cấu trúc và tính toàn vẹn cần thiết để xác nhận rằng
các phân hoạch sử dụng trong quá trình nghiên cứu vẫn tương ứng với các
phân hoạch đã được xác lập trước. Các phân tích nội dung chi tiết như
phân bố lớp, phân bố kích thước bounding box, phân bố vị trí,
\textit{label cardinality}, \textit{class co-occurrence} và phân tích độ
nhạy không được thực hiện trên \textit{test set} để phục vụ thiết kế hoặc
điều chỉnh phương pháp.

Một giới hạn đặc biệt được áp dụng đối với các
\textit{unlabeled subset} \(U_b\). Mặc dù các ảnh trong \(U_b\) có nguồn
gốc từ tập huấn luyện \(T\), các ground-truth annotations tương ứng được
xem là thông tin bị ẩn trong thiết lập học bán giám sát. Do đó, nghiên
cứu không sử dụng \texttt{image\_id} của \(U_b\) để truy ngược vào
ground truth của \(T\), khôi phục lớp hoặc bounding box rồi sử dụng các
thông tin này cho chẩn đoán. Việc làm như vậy sẽ cung cấp cho quá trình
thiết kế thông tin mà nhánh \textit{unlabeled} về nguyên tắc không được
quan sát và có thể tạo ra \textit{oracle leakage}. Đối với \(U_b\), chỉ
các thuộc tính cấu trúc không tiết lộ ground truth, chẳng hạn số ảnh,
\textit{membership}, quan hệ phần bù

\begin{equation}
U_b=T\setminus L_b
\end{equation}

và tính toàn vẹn của tệp dữ liệu được phép kiểm tra.

Các phép chẩn đoán dữ liệu cũng không sử dụng
\texttt{training\_seed}. Các \texttt{training\_seed} chỉ kiểm soát những
nguồn ngẫu nhiên phát sinh trong quá trình huấn luyện mô hình, trong khi
các thống kê ở mục này được tính trên các tập dữ liệu có thành phần đã
được xác định trước.

\paragraph{Phân bố lớp và mức độ mất cân bằng lớp.}

Mức hỗ trợ của từng lớp bất thường được mô tả chủ yếu ở cấp độ ảnh. Với
lớp \(c\), gọi \(N_c^{\mathrm{img}}\) là số ảnh duy nhất trong tập
huấn luyện \(T\) có ít nhất một ground-truth bounding box thuộc lớp đó:

\begin{equation}
N_c^{\mathrm{img}}
=
\left|
\left\{
i\in T:
c\ \text{xuất hiện trong ảnh }i
\right\}
\right|.
\label{eq:diag_class_image_support}
\end{equation}

Một ảnh chỉ được tính một lần đối với một lớp, bất kể ảnh đó chứa một hay
nhiều bounding box của cùng lớp. Tỷ lệ hiện diện ở cấp độ ảnh của lớp
\(c\) được xác định bởi

\begin{equation}
P_c^{\mathrm{img}}
=
\frac{N_c^{\mathrm{img}}}{|T|}.
\label{eq:diag_class_image_prevalence}
\end{equation}

Việc sử dụng \textit{image-level support} làm đại lượng chính nhằm phân
biệt giữa số ảnh độc lập cung cấp ví dụ của một lớp với số bản ghi
bounding box của lớp đó. Một lớp có thể xuất hiện trên ít ảnh nhưng có
nhiều bounding box trên mỗi ảnh; trong trường hợp này, số bounding box
không phản ánh trực tiếp số lượng ảnh độc lập mà mô hình quan sát được.

Mức độ không đồng đều giữa các lớp được mô tả bằng toàn bộ vector
\(\{N_c^{\mathrm{img}}\}_{c=1}^{14}\),

cùng các thống kê như giá trị lớn nhất, nhỏ nhất và trung vị. Tỷ số giữa
mức hỗ trợ lớn nhất và nhỏ nhất được xác định bởi

\begin{equation}
R_{\max/\min}
=
\frac{
\max_c N_c^{\mathrm{img}}
}{
\min_c N_c^{\mathrm{img}}
}.
\label{eq:diag_imbalance_ratio}
\end{equation}

Bên cạnh đó, hệ số biến thiên của mức hỗ trợ giữa các lớp được sử dụng
như một thống kê mô tả bổ sung:

\begin{equation}
CV_{\mathrm{img}}
=
\frac{
s\!\left(N_c^{\mathrm{img}}\right)
}{
\overline{N}^{\mathrm{img}}
},
\label{eq:diag_class_support_cv}
\end{equation}

trong đó \(s(N_c^{\mathrm{img}})\) là độ lệch chuẩn mẫu của mức hỗ trợ
giữa 14 lớp và \(\overline{N}^{\mathrm{img}}\) là giá trị trung bình
tương ứng. Các đại lượng này chỉ được sử dụng để mô tả mức độ không đồng
đều của dữ liệu. Nghiên cứu không tự thiết lập các mức định tính như
``mất cân bằng nhẹ'', ``trung bình'' hoặc ``nghiêm trọng'' khi không có
ngưỡng định lượng tương ứng được xác định trước.

\paragraph{Tiêu chí xác định lớp có mức hỗ trợ thấp.}

Để tránh sử dụng thuật ngữ ``lớp hiếm'' theo nghĩa định tính hoặc thay
đổi tùy theo cách diễn giải, nghiên cứu xác định một tiêu chí vận hành ở
cấp độ ảnh. Một lớp \(c\) được xác định là có mức hỗ trợ thấp trong tập
huấn luyện nếu

\begin{equation}
\mathrm{Rare}(c)
=
\mathbf{1}
\left[
P_c^{\mathrm{img}}<0{,}05
\right].
\label{eq:diag_rare_definition}
\end{equation}

Ngưỡng 5\% được xác định trước khi tính và trực quan hóa phân bố lớp của tập huấn luyện. Vì tập huấn luyện có \(|T|=3426\) ảnh, ta có

\begin{equation}
0{,}05\times3426=171{,}3.
\end{equation}

Do phép so sánh sử dụng dấu nhỏ hơn nghiêm ngặt, tiêu chí trên tương
đương với

\begin{equation}
N_c^{\mathrm{img}}\le171
\Rightarrow \mathrm{Rare},
\qquad
N_c^{\mathrm{img}}\ge172
\Rightarrow \mathrm{Non\text{-}rare}.
\label{eq:diag_rare_count_boundary}
\end{equation}

Ngưỡng 5\% được sử dụng như một ranh giới vận hành trung gian để nhận diện
các lớp có mức hỗ trợ thấp cần được lưu ý trong phạm vi dữ liệu nghiên
cứu. Giá trị này không được lựa chọn bằng cách quan sát histogram rồi tìm
một điểm cắt thuận lợi và cũng không được tuyên bố là một tiêu chuẩn phổ
quát của object detection hoặc một định nghĩa y khoa về bệnh hiếm.

Số bounding box của một lớp không được sử dụng để xác định hoặc thay đổi
trạng thái rare/non-rare. Quy tắc này tránh trường hợp một lớp xuất hiện
trên ít ảnh nhưng có nhiều annotation trên mỗi ảnh bị đánh giá là có mức
hỗ trợ cao chỉ vì số bounding box lớn.

Thuật ngữ ``rare'' trong nghiên cứu vì vậy chỉ mang nghĩa
\textit{dataset-level low support} trong tập huấn luyện \(T\) theo
Phương trình~\ref{eq:diag_rare_definition}. Nó không đồng nghĩa với bệnh
hiếm, bất thường hiếm trong quần thể hoặc tần suất mắc bệnh trong thực
hành lâm sàng.

\paragraph{Phân bố số lượng bounding box.}

Song song với mức hỗ trợ ở cấp ảnh, nghiên cứu mô tả mật độ annotation ở
cấp bounding box. Với lớp \(c\), số bản ghi bounding box được ký hiệu là

\begin{equation}
N_c^{\mathrm{bbox}}
=
\text{số ground-truth bounding-box annotations thuộc lớp }c.
\label{eq:diag_bbox_class_count}
\end{equation}

Tỷ trọng bounding box của lớp \(c\) trên toàn bộ annotation của tập
huấn luyện được xác định bởi

\begin{equation}
P_c^{\mathrm{bbox}}
=
\frac{
N_c^{\mathrm{bbox}}
}{
\sum_{k=1}^{14}
N_k^{\mathrm{bbox}}
}.
\label{eq:diag_bbox_class_share}
\end{equation}

Để mô tả mật độ annotation trên các ảnh thực sự chứa lớp \(c\), số
bounding box trung bình trên mỗi ảnh có lớp được tính bằng

\begin{equation}
R_c^{\mathrm{bbox/img}}
=
\frac{
N_c^{\mathrm{bbox}}
}{
N_c^{\mathrm{img}}
}.
\label{eq:diag_bbox_per_class_positive_image}
\end{equation}

Ngoài phân tích theo lớp, số bounding box trên từng ảnh cũng được khảo
sát. Với ảnh \(i\), đặt

\begin{equation}
B_i
=
\text{số ground-truth bounding-box annotations trên ảnh }i.
\label{eq:diag_bbox_count_per_image}
\end{equation}

Phân bố \(B_i\) được mô tả riêng cho hai phạm vi: toàn bộ tập huấn luyện
và nhóm ảnh có ít nhất một ground-truth bounding box. Việc tách hai phạm
vi này nhằm tránh để các ảnh \textit{No Finding} với \(B_i=0\) che khuất
mật độ annotation trên các ảnh có bất thường. Các thống kê được tính gồm
trung bình, độ lệch chuẩn mẫu, trung vị, phân vị thứ 95 và giá trị lớn
nhất.

Các đại lượng \(N_c^{\mathrm{bbox}}\), \(R_c^{\mathrm{bbox/img}}\) và
\(B_i\) phản ánh số bản ghi annotation. Chúng không được mặc nhiên diễn
giải là số tổn thương lâm sàng độc lập, đặc biệt trong bối cảnh bộ dữ
liệu có nguồn annotation từ nhiều bác sĩ.

\paragraph{Phân bố hình học và kích thước bounding box.}

Đối với bounding box thứ \(j\) trên ảnh \(i\), annotation COCO được biểu diễn dưới dạng
\([x_{ij},y_{ij},w_{ij},h_{ij}]\),

trong đó \(x_{ij}\) và \(y_{ij}\) là tọa độ góc trên bên trái,
\(w_{ij}\) và \(h_{ij}\) lần lượt là chiều rộng và chiều cao bounding
box; \(W_i\) và \(H_i\) là chiều rộng và chiều cao của ảnh.

Các đặc trưng hình học tuyệt đối gồm chiều rộng \(w_{ij}\), chiều cao
\(h_{ij}\) và diện tích

\begin{equation}
A_{ij}
=
w_{ij}h_{ij}.
\label{eq:diag_bbox_raw_area}
\end{equation}

Do các ảnh X-quang trong dữ liệu có thể có kích thước pixel khác nhau,
nghiên cứu đồng thời tính các đặc trưng đã chuẩn hóa:

\begin{equation}
w_{n,ij}
=
\frac{w_{ij}}{W_i},
\qquad
h_{n,ij}
=
\frac{h_{ij}}{H_i},
\label{eq:diag_bbox_normalized_dimensions}
\end{equation}

diện tích chuẩn hóa

\begin{equation}
a_{n,ij}
=
\frac{
w_{ij}h_{ij}
}{
W_iH_i
},
\label{eq:diag_bbox_normalized_area}
\end{equation}

và tỷ lệ khung hình

\begin{equation}
AR_{ij}
=
\frac{w_{ij}}{h_{ij}}.
\label{eq:diag_bbox_aspect_ratio}
\end{equation}

Trong các đại lượng trên, \(a_n\) được sử dụng làm đại lượng chính để
mô tả quy mô tương đối của bounding box. Diện tích tuyệt đối theo pixel
phụ thuộc trực tiếp vào độ phân giải ảnh; cùng một diện tích pixel có thể
chiếm tỷ lệ rất khác nhau trên hai ảnh có kích thước khác nhau. Ngược lại,

\[
a_n
=
\frac{wh}{WH}
\]

là đại lượng không thứ nguyên và biểu diễn trực tiếp tỷ lệ diện tích ảnh
được bounding box bao phủ.

Trước khi tính các đặc trưng hình học, mỗi bounding box phải thỏa mãn

\begin{equation}
w>0,
\qquad
h>0,
\qquad
x\ge0,
\qquad
y\ge0,
\end{equation}

\begin{equation}
x+w\le W,
\qquad
y+h\le H,
\qquad
0<a_n\le1.
\label{eq:diag_bbox_validity}
\end{equation}

Nếu phát hiện annotation vi phạm các điều kiện này, trường hợp đó được
ghi nhận như một vấn đề về tính toàn vẹn dữ liệu. Bước chẩn đoán không tự
động clamp, dịch chuyển, resize, xóa hoặc sửa bounding box để làm cho
annotation hợp lệ.

Để hỗ trợ diễn giải kích thước tương đối, bounding box được phân thành ba
nhóm dựa trên diện tích chuẩn hóa:

\begin{equation}
\mathrm{Small}
\iff
a_n<0{,}01,
\label{eq:diag_bbox_small}
\end{equation}

\begin{equation}
\mathrm{Medium}
\iff
0{,}01\le a_n<0{,}10,
\label{eq:diag_bbox_medium}
\end{equation}

và

\begin{equation}
\mathrm{Large}
\iff
a_n\ge0{,}10.
\label{eq:diag_bbox_large}
\end{equation}

Theo định nghĩa này, bounding box có \(a_n=0{,}01\) thuộc nhóm Medium,
còn bounding box có \(a_n=0{,}10\) thuộc nhóm Large. Hai ranh giới 1\%
và 10\% được xác định trước khi phân bố \(a_n\) của dữ liệu được tính và
trực quan hóa. Chúng tạo ra ba khoảng dễ diễn giải theo tỷ lệ diện tích
ảnh: dưới một phần trăm, từ một phần trăm đến dưới một phần mười, và từ
một phần mười diện tích ảnh trở lên.

Các nhóm này được sử dụng như
\textit{normalized-area-based bbox size categories}. Chúng không phải
các nhóm Small, Medium và Large được định nghĩa theo diện tích pixel tuyệt
đối trong giao thức COCO \cite{lin2014microsoft}. Vì vậy, các thuật ngữ
``COCO Small'', ``COCO Medium'' và ``COCO Large'' không được sử dụng cho
các nhóm kích thước trong nghiên cứu này.

Nghiên cứu cũng không sử dụng các phân vị của chính dữ liệu, chẳng hạn
P33 và P66, để xác định ba nhóm kích thước chính. Cách phân chia dựa trên
phân vị sẽ làm cho số lượng quan sát trong mỗi nhóm phần lớn trở thành hệ
quả của cách đặt ranh giới, thay vì một đặc trưng cần được quan sát từ dữ
liệu. Việc sử dụng các ranh giới được xác định trước cho phép tỷ lệ
Small, Medium và Large trở thành kết quả mô tả của dữ liệu.

Bên cạnh các nhóm rời rạc, các phân bố liên tục của \(w\), \(h\), \(A\),
\(w_n\), \(h_n\), \(a_n\) và \(AR\) vẫn được giữ lại. Các thống kê mô tả
bao gồm số quan sát, trung bình, độ lệch chuẩn mẫu, giá trị nhỏ nhất, các
phân vị 5\%, 25\%, 50\%, 75\%, 95\% và giá trị lớn nhất. Cách tổ chức này
tránh việc ba nhóm kích thước thay thế hoàn toàn thông tin liên tục về
hình học bounding box.

\paragraph{Phân bố không gian của bounding box.}

Để khảo sát vị trí tương đối của các annotations trong khi loại bỏ ảnh
hưởng trực tiếp của kích thước ảnh, tâm của mỗi bounding box được chuẩn
hóa về miền \([0,1]\times[0,1]\):

\begin{equation}
x_{c,n}
=
\frac{x+w/2}{W},
\qquad
y_{c,n}
=
\frac{y+h/2}{H}.
\label{eq:diag_bbox_normalized_center}
\end{equation}

Hệ tọa độ sử dụng góc trên bên trái làm gốc; \(x_{c,n}\) tăng từ trái
sang phải và \(y_{c,n}\) tăng từ trên xuống dưới. Việc chuẩn hóa tọa độ
cho phép tổng hợp vị trí bounding box giữa các ảnh có kích thước pixel
khác nhau.

Miền chuẩn hóa được chia thành lưới \(50\times50\). Với mỗi bounding box,
tọa độ tâm xác định một ô trong lưới và ô đó được tăng một đơn vị đếm.
Mỗi bounding box vì vậy đóng góp đúng một quan sát, nên heatmap thu được
có bản chất \textit{annotation-weighted bbox-center distribution}.

Độ phân giải \(50\times50\) được xác định trước và sử dụng thống nhất cho
toàn bộ phép phân tích. Lưới không được thay đổi sau khi quan sát dữ liệu
nhằm tạo ra hoặc làm nổi bật một pattern cụ thể.

Heatmap chỉ mô tả phân bố của tâm các bounding-box annotations. Một vùng
có mật độ cao chỉ có nghĩa là nhiều tâm annotation nằm trong vùng tương
ứng; kết quả không được mặc nhiên diễn giải là mật độ của các tổn thương
lâm sàng độc lập, không chứng minh quan hệ giải phẫu nhân quả và không
cho phép suy ra trước hiệu năng phát hiện của mô hình tại vùng đó.

\paragraph{Phân bố ảnh âm tính.}

Ảnh \textit{No Finding} tiếp tục được định nghĩa như ở các bước xây dựng
dữ liệu trước đó: đây là ảnh âm tính không có ground-truth bounding box
(\textit{zero-GT negative image}), không phải lớp phát hiện thứ 15.

Với \(N_{\mathrm{neg}}\) là số ảnh \textit{No Finding} trong tập huấn
luyện, tỷ lệ ảnh âm tính được tính bởi

\begin{equation}
P_{\mathrm{neg}}
=
\frac{
N_{\mathrm{neg}}
}{
|T|
}.
\label{eq:diag_negative_prevalence}
\end{equation}

Số ảnh có ít nhất một ground-truth bounding box tương ứng là

\begin{equation}
N_{\mathrm{pos}}
=
|T|-N_{\mathrm{neg}}.
\label{eq:diag_positive_count}
\end{equation}

Việc mô tả riêng tỷ lệ ảnh âm tính cho phép đặc trưng hóa thành phần
zero-GT của dữ liệu trước khi đánh giá hành vi false positive của mô hình
ở các giai đoạn sau. Tuy nhiên, trong bước chẩn đoán hiện tại,
\(P_{\mathrm{neg}}\) chỉ được xem là thuộc tính dữ liệu, không phải metric
hiệu năng.

Cần phân biệt rõ ảnh âm tính và ảnh chưa gán nhãn. Một ảnh
\textit{No Finding} là ảnh mà ground truth xác định không có đối tượng
đích, trong khi một ảnh thuộc \(U_b\) chỉ có nghĩa annotation của ảnh bị
ẩn khỏi quá trình học bán giám sát. Do đó,

\begin{equation}
\textit{unlabeled image}
\not\Rightarrow
\textit{negative image}.
\label{eq:diag_unlabeled_not_negative}
\end{equation}

Việc đồng nhất hai khái niệm này sẽ làm sai bản chất của thiết lập
\textit{semi-supervised object detection}.

\paragraph{Chẩn đoán các tập dữ liệu có nhãn theo ngân sách.}

Bốn tập \(L_{1\%}\), \(L_{5\%}\), \(L_{10\%}\) và \(L_{20\%}\) đã được
xây dựng trước bước chẩn đoán và được giữ nguyên thành phần ảnh. Mục tiêu
ở đây không phải lựa chọn lại ảnh cho các tập có nhãn mà là định lượng
mức độ mà phân bố của từng \(L_b\) phản ánh phân bố của tập huấn luyện
\(T\).

Với ngân sách \(b\) và lớp \(c\), gọi
\(N_{b,c}^{\mathrm{img}}\) là số ảnh trong \(L_b\) có ít nhất một
ground-truth bounding box thuộc lớp \(c\). Tỷ lệ hiện diện của lớp đó
trong tập có nhãn được xác định bởi

\begin{equation}
P_{b,c}^{\mathrm{img}}
=
\frac{
N_{b,c}^{\mathrm{img}}
}{
|L_b|
}.
\label{eq:diag_budget_class_prevalence}
\end{equation}

Độ lệch tuyệt đối của tỷ lệ lớp so với tập huấn luyện được tính bằng

\begin{equation}
D_{b,c}
=
\left|
P_{b,c}^{\mathrm{img}}
-
P_c^{\mathrm{img}}
\right|.
\label{eq:diag_budget_class_deviation}
\end{equation}

Để tóm tắt độ lệch trên 14 lớp, hai đại lượng được sử dụng:

\begin{equation}
D_{\max,b}
=
\max_{c=1,\ldots,14}
D_{b,c},
\label{eq:diag_budget_max_deviation}
\end{equation}

và

\begin{equation}
D_{\mathrm{mean},b}
=
\frac{1}{14}
\sum_{c=1}^{14}
D_{b,c}.
\label{eq:diag_budget_mean_deviation}
\end{equation}

Trong đó, \(D_{\max,b}\) phản ánh sai lệch lớn nhất của một lớp tại ngân
sách \(b\), còn \(D_{\mathrm{mean},b}\) phản ánh mức sai lệch trung bình
trên toàn bộ 14 lớp. Hai đại lượng chỉ được sử dụng để mô tả mức độ đại
diện của các tập có nhãn; chúng không được sử dụng để tái tối ưu
\textit{membership} của \(L_b\).

Mức bao phủ lớp của mỗi ngân sách được xác định bởi

\begin{equation}
C_b
=
\sum_{c=1}^{14}
\mathbf{1}
\left[
N_{b,c}^{\mathrm{img}}>0
\right],
\label{eq:diag_budget_class_coverage}
\end{equation}

với \(0\le C_b\le14\).

Ngoài ra, số ảnh \textit{No Finding}, tỷ lệ ảnh âm tính, mức hỗ trợ của
từng lớp và mức hỗ trợ của các lớp đã được xác định là rare từ tập
huấn luyện \(T\) cũng được ghi nhận cho từng \(L_b\).

Trạng thái rare/non-rare không được định nghĩa lại riêng trong từng
\(L_b\). Một lớp chỉ có một trạng thái chính thức được xác định từ toàn
bộ tập huấn luyện theo
Phương trình~\ref{eq:diag_rare_definition}; tại từng ngân sách, nghiên
cứu chỉ khảo sát mức hỗ trợ mà lớp đó nhận được trong tập có nhãn. Cách
xử lý này tránh việc cùng một lớp thay đổi ý nghĩa ``rare'' chỉ do mẫu số
của phép tính chuyển từ \(T\) sang một tập con nhỏ hơn.

Tất cả các phép chẩn đoán trên \(L_b\) chỉ sử dụng ground truth thuộc
chính \(L_b\). Ground truth bị ẩn của \(U_b\) không được khôi phục để
bổ sung class support, bounding-box statistics hoặc bất kỳ thống kê nội
dung nào.

\paragraph{Các đặc trưng đa nhãn bổ sung.}

Do một ảnh X-quang có thể đồng thời chứa nhiều lớp bất thường, nghiên cứu
bổ sung hai phép chẩn đoán ở cấp độ đa nhãn:
\textit{label cardinality} và \textit{class co-occurrence}.

Với ảnh \(i\), \textit{label cardinality} được định nghĩa là số lớp phát
hiện duy nhất hiện diện trên ảnh:

\begin{equation}
LC_i
=
\left|
\left\{
c:
c\ \text{xuất hiện trong ảnh }i
\right\}
\right|.
\label{eq:diag_label_cardinality}
\end{equation}

Một lớp chỉ được tính một lần trong \(LC_i\), bất kể ảnh có bao nhiêu
bounding box thuộc lớp đó. Đối với ảnh \textit{No Finding},

\begin{equation}
LC_i=0.
\end{equation}

Phân bố \(\{LC_i\}\) được mô tả bằng số ảnh tại từng mức cardinality,
trung bình, độ lệch chuẩn mẫu, trung vị, phân vị thứ 95 và giá trị lớn
nhất. \textit{Label cardinality} phản ánh số loại bất thường đồng thời
xuất hiện trên ảnh và khác về bản chất với \(B_i\), là tổng số bounding
box annotations trên ảnh.

Đối với hai lớp khác nhau \(c\) và \(d\), số ảnh đồng xuất hiện được
định nghĩa bởi

\begin{equation}
N_{cd}
=
\left|
\left\{
i\in T:
c\ \text{và}\ d
\ \text{cùng xuất hiện trong ảnh }i
\right\}
\right|.
\label{eq:diag_class_cooccurrence}
\end{equation}

Với 14 lớp phát hiện, số cặp lớp không thứ tự cần khảo sát là

\begin{equation}
\binom{14}{2}
=
91.
\label{eq:diag_number_of_class_pairs}
\end{equation}

Mỗi ảnh chỉ đóng góp một lần cho một cặp lớp, bất kể trên ảnh có bao
nhiêu bounding box thuộc lớp \(c\) hoặc lớp \(d\). Vì vậy,
\(N_{cd}\) phản ánh số ảnh duy nhất trong tập huấn luyện có sự hiện diện
đồng thời của hai lớp, thay vì phản ánh số lượng bounding box của hai
lớp đó.

Phân tích \textit{class co-occurrence} chỉ nhằm mô tả cấu trúc đa nhãn
quan sát được trong dữ liệu. Số ảnh đồng xuất hiện lớn cho biết hai lớp
cùng xuất hiện trên nhiều ảnh hơn trong phạm vi tập dữ liệu nghiên cứu,
nhưng không chứng minh rằng một lớp gây ra lớp kia, không xác lập cơ chế
bệnh sinh và không chứng minh mối liên hệ lâm sàng giữa hai bất thường.
Kết quả đồng xuất hiện cũng không được sử dụng như một tiêu chí tự động
để thay đổi sampling, hàm mất mát, kiến trúc hoặc cấu hình huấn luyện.

\paragraph{Phân tích độ nhạy của các ngưỡng chẩn đoán.}

Ngưỡng 5\% cho lớp có mức hỗ trợ thấp và các ranh giới 1\%/10\% cho
kích thước bounding box là các định nghĩa vận hành được xác định trước,
không phải các hằng số tự nhiên hoặc tiêu chuẩn lâm sàng. Vì vậy, nghiên
cứu bổ sung một phân tích độ nhạy nhằm đánh giá mức độ phụ thuộc của các
kết luận mô tả vào việc lựa chọn các ranh giới này.

Các ngưỡng chính và các giá trị khảo sát được trình bày tại
Bảng~\ref{tab:diagnostic_threshold_sensitivity}.

\begin{table*}[!ht]
    \centering
    \caption{Ngưỡng chính và các giá trị sử dụng trong phân tích độ nhạy}
    \label{tab:diagnostic_threshold_sensitivity}
    \small
    \setlength{\tabcolsep}{3pt}
    \renewcommand{\arraystretch}{1.15}

    \begin{tabularx}{\textwidth}{
        >{\raggedright\arraybackslash}p{2.5cm}
        >{\centering\arraybackslash}p{2.0cm}
        >{\centering\arraybackslash}p{1.8cm}
        >{\centering\arraybackslash}p{2.7cm}
        >{\raggedright\arraybackslash}X
    }
        \toprule
        \textbf{Thành phần} &
        \textbf{Đại lượng} &
        \textbf{Ngưỡng chính} &
        \textbf{Các giá trị khảo sát} &
        \textbf{Điều kiện giữ cố định} \\
        \midrule

        Lớp có mức hỗ trợ thấp &
        \(P_c^{\mathrm{img}}\) &
        \(0{,}05\) &
        \(0{,}01;\;0{,}05;\;0{,}10\) &
        Mức hỗ trợ luôn được xác định ở cấp ảnh; số bounding box không
        thay thế tiêu chí này. \\

        \addlinespace[2pt]

        Ranh giới Small &
        \(a_n\) &
        \(0{,}01\) &
        \(0{,}005;\;0{,}01;\;0{,}02\) &
        Ranh giới Large được giữ cố định tại \(0{,}10\). \\

        \addlinespace[2pt]

        Ranh giới Large &
        \(a_n\) &
        \(0{,}10\) &
        \(0{,}05;\;0{,}10;\;0{,}20\) &
        Ranh giới Small được giữ cố định tại \(0{,}01\). \\

        \bottomrule
    \end{tabularx}
\end{table*}

Đối với tiêu chí lớp có mức hỗ trợ thấp, cùng một vector
\(P_c^{\mathrm{img}}\) được lần lượt đánh giá tại các ranh giới 1\%,
5\% và 10\%. Mục tiêu của phép khảo sát không phải tìm ngưỡng tạo ra một
số lượng lớp rare mong muốn, mà là đánh giá trạng thái rare/non-rare có
nhạy với định nghĩa vận hành hay không. Ngưỡng 5\% tiếp tục được sử dụng
làm ngưỡng chính bất kể kết quả của phép khảo sát bổ sung.

Đối với kích thước bounding box, phân tích độ nhạy được thực hiện theo
nguyên tắc thay đổi từng yếu tố một
(\textit{one-factor-at-a-time}). Khi khảo sát ranh giới Small, ranh giới này lần lượt nhận các giá trị
\(0{,}005\), \(0{,}01\) và \(0{,}02\), trong khi ranh giới Large được giữ
cố định tại \(0{,}10\). Ngược lại, khi khảo sát ranh giới Large, các giá trị
được sử dụng là \(0{,}05\), \(0{,}10\) và \(0{,}20\), trong khi ranh giới
Small được giữ cố định tại \(0{,}01\).

Việc chỉ thay đổi một ranh giới trong mỗi phép khảo sát giúp tách ảnh
hưởng của ranh giới đang được xem xét. Nghiên cứu không tạo toàn bộ tích
Descartes giữa các ranh giới Small và Large, bởi mục tiêu của bước này
không phải tìm kiếm tổ hợp ngưỡng tối ưu.

Các ngưỡng chính được xác định trước khi các thống kê chẩn đoán được tính
và được giữ nguyên sau khi quan sát kết quả. Do đó,

\begin{equation}
\text{kết quả phân tích độ nhạy}
\not\Rightarrow
\text{lựa chọn lại ngưỡng chính}.
\label{eq:diag_sensitivity_no_reselection}
\end{equation}

Phân tích độ nhạy ở đây cũng không phải
\textit{Ablation Study}. Phân tích này chỉ thay đổi định nghĩa sử dụng
để mô tả dữ liệu và không bao gồm huấn luyện mô hình hoặc đánh giá AP/mAP.
Ngược lại, \textit{Ablation Study} ở các thí nghiệm mô hình liên quan
đến việc thay đổi một thành phần của phương pháp và đo ảnh hưởng của thay
đổi đó lên hiệu năng. Vì vậy, các ngưỡng rare và Small/Medium/Large không
được xem là các biến ablation của mô hình.

\paragraph{Kiểm tra cấu trúc các phân hoạch và giới hạn sử dụng dữ liệu.}

Các kiểm tra cấu trúc của \textit{training set},
\textit{validation set} và \textit{test set} kế thừa quy trình kiểm định
đã trình bày tại Mục~\ref{subsec:fixed_data_split}. Trong bước chẩn đoán
trước huấn luyện, các thuộc tính cấu trúc được đối chiếu nhằm xác nhận
rằng các tập dữ liệu sử dụng vẫn tương ứng với các phân hoạch đã được
xác lập trước, bao gồm số ảnh, số annotation, thành phần
\texttt{image\_id}, sự rời nhau giữa các tập, tính đầy đủ của hợp các
phân hoạch, ánh xạ category và các thông tin kiểm tra tính toàn vẹn có
liên quan.

Mục tiêu của các kiểm tra này là phát hiện sự thay đổi ngoài ý muốn của
dữ liệu giữa các bước xử lý, không phải mở lại quá trình lựa chọn phân
hoạch. Vì vậy, \textit{validation set} và \textit{test set} không được
sử dụng để xác định lại class distribution, tiêu chí rare, phân bố kích
thước bounding box, \textit{label cardinality},
\textit{class co-occurrence} hoặc các ranh giới chẩn đoán.

Đặc biệt, \textit{test set} không được sử dụng để lựa chọn các ngưỡng
chẩn đoán, kiến trúc mô hình, augmentation, confidence threshold,
\textit{pseudo-label filtering} hoặc các quyết định thiết kế khác. Các
phân tích nội dung chi tiết mới trên \textit{test set} không được thực
hiện trong bước này.

Tương tự, ground truth bị ẩn của \(U_b\) không được sử dụng để tính phân
bố lớp, phân bố kích thước bounding box, mức hỗ trợ của lớp rare,
\textit{label cardinality}, \textit{class co-occurrence} hoặc heatmap
cho phần dữ liệu chưa gán nhãn. Quy định này bảo đảm rằng việc xem một
ảnh là \textit{unlabeled} được duy trì không chỉ trong quá trình huấn
luyện mà cả trong quá trình phân tích dữ liệu trước huấn luyện.

\paragraph{Ý nghĩa và giới hạn của chẩn đoán trước huấn luyện.}

Các phép chẩn đoán trên tạo ra một mô tả có cấu trúc về dữ liệu đầu vào,
bao gồm mức hỗ trợ lớp, mức độ không đồng đều giữa các lớp, mật độ
annotation, hình học và vị trí bounding box, tỷ lệ ảnh âm tính, mức độ
đại diện của các ngân sách nhãn và cấu trúc đa nhãn. Những thông tin này
được sử dụng để xác định các đặc trưng dữ liệu cần được theo dõi khi diễn
giải kết quả \textit{supervised} và \textit{semi-supervised} ở các phần
sau.

Tuy nhiên, các phép chẩn đoán dữ liệu không phải bằng chứng về hành vi
của mô hình. Việc một lớp có mức hỗ trợ thấp không chứng minh rằng lớp đó
sẽ có AP thấp. Một bounding box thuộc nhóm Small không chứng minh rằng
detector sẽ thất bại trên bounding box đó. Tương tự, mức mất cân bằng lớp
không tự chứng minh nguyên nhân của một sai số mô hình, và một vùng có
mật độ tâm bounding box cao không cho phép suy ra trước chất lượng phát
hiện tại vùng đó.

Các kết quả đồng xuất hiện giữa các lớp chỉ phản ánh cấu trúc nhãn quan
sát được trong dữ liệu. Số ảnh đồng xuất hiện lớn không chứng minh quan
hệ nhân quả, cơ chế bệnh sinh hoặc mối liên hệ lâm sàng giữa các bất
thường. Số bounding-box annotations cũng không được mặc nhiên diễn giải
thành số tổn thương lâm sàng độc lập.

Bước chẩn đoán trước huấn luyện không cung cấp bằng chứng về chất lượng
\textit{pseudo-label}, độ ổn định hội tụ, lợi ích của EMA, hiệu quả của
augmentation, ảnh hưởng của confidence threshold hoặc ưu thế của học bán
giám sát so với học có giám sát. Những kết luận này chỉ có thể được đánh
giá sau khi các mô hình được huấn luyện và được đo lường theo giao thức
thực nghiệm đã xác định.

Do đó, vai trò của mục này được giới hạn ở việc xác định rõ cách mô tả và
kiểm tra dữ liệu trước huấn luyện. Các giá trị quan sát cụ thể, bao gồm
lớp nào thỏa tiêu chí rare, mức mất cân bằng thực tế, phân bố
Small/Medium/Large, phân bố tâm bounding box, tỷ lệ ảnh âm tính, mức độ
đại diện thực tế của từng ngân sách nhãn, phân bố
\textit{label cardinality}, số ảnh đồng xuất hiện giữa các cặp lớp và
kết quả phân tích độ nhạy, được trình bày và diễn giải trong phần kết quả.
\subsection{Supervised Baseline}
\label{subsec:Supervised_Baseline}

Supervised baseline được xây dựng nhằm cung cấp mốc tham chiếu có kiểm soát cho
các phép so sánh với semi-supervised learning tại cùng kiến trúc, cùng ngân
sách nhãn và cùng điều kiện thực nghiệm. Nghiên cứu không xem supervised
baseline chỉ là một cấu hình huấn luyện phụ trợ, mà là điều kiện đối chứng cần
thiết để định lượng phần cải thiện gắn với việc khai thác thêm dữ liệu chưa gán
nhãn.

Để hạn chế ảnh hưởng của khác biệt từ họ detector, nghiên cứu giữ cùng cấu trúc
Faster R-CNN và FPN cho cả hai nhánh kiến trúc. Thành phần kiến trúc được thay
đổi có chủ đích nằm ở backbone, gồm một kiến trúc CNN truyền thống dựa trên
ResNet-50 và một kiến trúc Vision Transformer phân cấp dựa trên Swin-T. Thiết
kế này giúp tập trung phép so sánh vào sự khác biệt về backbone trong cùng một
họ detector thay vì trộn lẫn nhiều thay đổi kiến trúc đồng thời.

\subsubsection{Lựa chọn họ detector và hai kiến trúc thực nghiệm}
\label{subsubsec:supervised_detector_architecture}

Nghiên cứu lựa chọn Faster R-CNN làm họ detector chung cho cả hai kiến trúc
thực nghiệm. Faster R-CNN là một detector hai giai đoạn, trong đó Region
Proposal Network tạo các vùng đề xuất và RoI Head thực hiện phân loại cùng hồi
quy bounding box trên các vùng này \cite{ren2016faster}. Việc sử dụng cùng họ
Faster R-CNN giúp giữ nhất quán cấu trúc phát hiện giữa các điều kiện và giảm
khả năng kết quả bị chi phối bởi khác biệt giữa các detector family.

Feature Pyramid Network (FPN) được giữ chung cho cả hai nhánh kiến trúc nhằm
cung cấp biểu diễn đặc trưng đa tỉ lệ cho bài toán phát hiện đối tượng
\cite{lin2017feature}. Vì vậy, hai cấu hình chính được xác định là

\begin{equation}
\mathcal{A}_{R50}
=
Faster\ R\text{-}CNN
+
ResNet\text{-}50
+
FPN,
\label{eq:supervised_arch_r50}
\end{equation}

và

\begin{equation}
\mathcal{A}_{Swin}
=
Faster\ R\text{-}CNN
+
Swin\text{-}T
+
FPN.
\label{eq:supervised_arch_swin}
\end{equation}

Cách tổ chức này giữ chung các thành phần chính gồm Faster R-CNN, FPN, RPN và
họ RoI Head; khác biệt kiến trúc chính được giới hạn ở backbone ResNet-50 so
với Swin-T. Do đó, mục tiêu không phải xác định detector family nào tốt hơn, mà
là khảo sát liệu mức lợi ích của semi-supervised learning có phụ thuộc vào kiểu
kiến trúc backbone hay không. Đây cũng là cơ sở để diễn giải RQ7 trong phạm vi
các giao thức huấn luyện đặc thù theo kiến trúc đã được xác định trước.

\paragraph{Lý do lựa chọn ResNet-50.}

ResNet-50 được sử dụng như kiến trúc CNN truyền thống tham chiếu
(\textit{ResNet-50-based conventional CNN architecture}). Kiến trúc ResNet
được xây dựng trên cơ chế residual learning, trong đó các kết nối tắt
(\textit{skip connections}) hỗ trợ tối ưu các mạng sâu và đã trở thành một
backbone nền tảng trong nhiều bài toán thị giác máy tính
\cite{he2016deep}.

Trong nghiên cứu này, ResNet-50 được lựa chọn không chỉ vì tính phổ biến mà
còn vì nó tạo ra một mốc CNN trưởng thành, có cấu hình Faster R-CNN + FPN
được sử dụng rộng rãi trong object detection và được hỗ trợ trực tiếp trong
MMDetection. Faster R-CNN cung cấp khung phát hiện hai giai đoạn ổn định,
trong khi FPN hỗ trợ biểu diễn đặc trưng đa tỉ lệ cho object detection
\cite{ren2016faster,lin2017feature}. Sự kết hợp này vì vậy tạo thành một
baseline CNN phù hợp để đối chiếu với kiến trúc Vision Transformer phân cấp
mà không cần thay đổi đồng thời họ detector.

Ngoài ra, ResNet-50 có mức độ phức tạp tính toán phù hợp với thiết kế thực
nghiệm lặp lại nhiều training seeds. Đây là một yêu cầu thực tế quan trọng
của nghiên cứu do mỗi điều kiện phải được huấn luyện lặp lại trên cùng danh
sách 10 seed. Vì vậy, ResNet-50 được chọn như một kiến trúc CNN tham chiếu
trưởng thành, có cơ sở rõ ràng trong hệ sinh thái Faster R-CNN/FPN và khả thi
về tài nguyên cho thiết kế thực nghiệm nhiều lần lặp; lựa chọn này không hàm
ý ResNet-50 là backbone CNN tối ưu cho dữ liệu X-quang phổi.
\paragraph{Lý do lựa chọn Swin-T.}

Nhánh kiến trúc thứ hai sử dụng Swin-T, tương ứng với một kiến trúc
Vision Transformer phân cấp
(\textit{Swin-T-based hierarchical Vision Transformer architecture}).
Swin Transformer xây dựng biểu diễn đặc trưng phân cấp và sử dụng cơ chế
\textit{shifted-window self-attention}. Thiết kế attention theo các cửa sổ
cục bộ và cơ chế dịch cửa sổ giúp hạn chế chi phí tính toán so với attention
toàn cục, đồng thời duy trì khả năng trao đổi thông tin giữa các vùng ảnh.
Kiến trúc phân cấp này cho phép Swin Transformer được tích hợp thuận lợi vào
các bài toán dense prediction như object detection
\cite{liu2021swin}.

Swin-T được lựa chọn vì nó đại diện cho một hướng kiến trúc khác biệt rõ ràng
so với CNN truyền thống của ResNet-50, trong khi vẫn có thể tích hợp với FPN
và cùng họ Faster R-CNN. Nhờ đó, nghiên cứu có thể khảo sát sự khác biệt về
mức lợi ích của semi-supervised learning giữa một backbone CNN truyền thống
và một backbone Vision Transformer phân cấp mà không đồng thời thay đổi họ
detector.

Việc lựa chọn biến thể Tiny thay vì các biến thể Swin lớn hơn còn nhằm giữ
chi phí tính toán ở mức phù hợp với thiết kế nhiều điều kiện và nhiều
training seeds. Đồng thời, Swin-T đã có cấu hình tích hợp trong MMDetection,
giúp việc triển khai cùng Faster R-CNN/FPN có cơ sở kỹ thuật rõ ràng và hạn
chế việc phải xây dựng một detector riêng cho nhánh Transformer.

Lựa chọn này được xác định trước khi huấn luyện chính thức và không xuất phát
từ quá trình thử nhiều backbone rồi chọn mô hình có kết quả tốt nhất. Vì vậy,
nghiên cứu không đặt giả thuyết rằng Swin-T phải tốt hơn ResNet-50; mục tiêu
là đánh giá liệu mức lợi ích của semi-supervised learning có khác nhau giữa
hai kiểu kiến trúc hay không.
\subsubsection{Khởi tạo backbone và biểu diễn đầu vào}
\label{subsubsec:supervised_pretraining_input}

Để hạn chế khác biệt về nguồn thông tin tiền huấn luyện giữa hai kiến trúc,
cả ResNet-50 và Swin-T đều sử dụng backbone pretrained trên ImageNet-1K. Nghiên
cứu không sử dụng detector-level initialization pretrained trên COCO.

Nguyên tắc khởi tạo được biểu diễn bởi

\begin{equation}
Pretrain_{R50}
=
Pretrain_{Swin}
=
ImageNet\text{-}1K\ backbone,
\label{eq:supervised_pretraining_fairness}
\end{equation}

trong khi

\begin{equation}
COCO\ detector\text{-}level\ initialization
=
\varnothing.
\label{eq:no_coco_detector_pretraining}
\end{equation}

Thiết kế này nhằm làm cho nguồn external supervision giữa hai backbone tương
đương về loại. Nếu một kiến trúc được khởi tạo từ detector đã học trên COCO còn
kiến trúc kia chỉ dùng backbone pretrained trên ImageNet, phần chênh lệch về
hiệu năng có thể bị trộn lẫn với lượng thông tin giám sát bổ sung từ dữ liệu
detection bên ngoài.

Ảnh đầu vào của nghiên cứu được lưu dưới dạng JPEG grayscale một kênh. Để tương
thích với các backbone pretrained trên ảnh ba kênh, mỗi ảnh một kênh được ánh
xạ thành ba kênh có giá trị giống nhau:

\begin{equation}
I_{1ch}
\rightarrow
I_{3ch},
\qquad
R=G=B.
\label{eq:grayscale_to_three_channels}
\end{equation}

Phép chuyển đổi này không tạo thêm thông tin màu. Ba kênh chỉ là ba bản sao của
cùng tín hiệu grayscale nhằm bảo đảm khả năng sử dụng trực tiếp các trọng số
backbone pretrained.

Trong cả train, validation và test, ảnh được resize với cấu hình

\begin{equation}
scale=(1333,800),
\qquad
keep\_ratio=True,
\label{eq:supervised_input_scale}
\end{equation}

và được padding theo bội số

\begin{equation}
pad\_size\_divisor=32.
\label{eq:supervised_padding}
\end{equation}

Tỷ lệ khung hình được giữ trong quá trình resize nhằm tránh biến dạng hình học
không cần thiết của ảnh và bounding box.

\subsubsection{Cấu hình huấn luyện theo kiến trúc}
\label{subsubsec:supervised_training_configuration}

Luồng dữ liệu có nhãn trong cấu hình supervised chính được giữ theo hướng bảo
thủ. Quy trình cơ bản gồm nạp dữ liệu, resize và đóng gói đầu vào cho mô hình.
Nghiên cứu không mặc định sử dụng flip, crop, rotation, mosaic hoặc các phép
biến đổi quang học mạnh trong supervised baseline chính.

Lựa chọn này nhằm giảm số nguồn biến thiên không cần thiết trong phép so sánh
SUP--SSL và tránh việc augmentation trở thành một yếu tố gây nhiễu trước khi
ảnh hưởng của augmentation được khảo sát riêng trong thí nghiệm phân tích thành
phần.

Effective labeled batch size được cố định ở mức

\begin{equation}
B^{L}_{\mathrm{eff}}
=
4.
\label{eq:effective_labeled_batch}
\end{equation}

Nếu giới hạn phần cứng yêu cầu sử dụng microbatch hoặc gradient accumulation,
đó chỉ được xem là cơ chế triển khai để đạt cùng effective batch size, không
được xem là thay đổi thiết kế khoa học.

\paragraph{Cấu hình tối ưu hóa cho ResNet-50.}

Nhánh Faster R-CNN + ResNet-50 + FPN sử dụng SGD với các tham số:

\begin{equation}
\begin{aligned}
lr_{R50} &= 0.005,\\
momentum &= 0.9,\\
weight\ decay &= 0.0001.
\end{aligned}
\label{eq:r50_optimizer}
\end{equation}

Cấu hình này phù hợp với họ training recipe phổ biến của Faster R-CNN với
ResNet trong MMDetection. Learning rate được điều chỉnh theo effective batch
của nghiên cứu thay vì sao chép nguyên giá trị từ một reference batch lớn hơn.

\paragraph{Cấu hình tối ưu hóa cho Swin-T.}

Nhánh Faster R-CNN + Swin-T + FPN sử dụng AdamW với các tham số:

\begin{equation}
\begin{aligned}
lr_{Swin} &= 0.000025,\\
\beta_1 &= 0.9,\\
\beta_2 &= 0.999,\\
weight\ decay &= 0.05.
\end{aligned}
\label{eq:swin_optimizer}
\end{equation}

Các quy tắc no-decay đặc thù của kiến trúc Swin được giữ theo training recipe
gốc của kiến trúc.

Việc R50 và Swin-T sử dụng optimizer khác nhau không được xem là vi phạm nguyên
tắc kiểm soát. Hai backbone có đặc trưng tối ưu hóa khác nhau và được huấn luyện
bằng các giao thức đặc thù theo kiến trúc đã xác định trước. Fairness trong
nghiên cứu được áp dụng bằng cách giữ cùng recipe giữa SUP và SSL trong từng
kiến trúc:

\begin{equation}
Recipe^{SUP}_{a}
=
Recipe^{SSL}_{a},
\qquad
a\in\{R50,Swin\text{-}T\}.
\label{eq:supervised_architecture_specific_matching}
\end{equation}

Nghiên cứu không yêu cầu

\begin{equation}
Recipe_{R50}
=
Recipe_{Swin\text{-}T}.
\label{eq:supervised_no_cross_arch_recipe_equality}
\end{equation}

Do đó, mọi so sánh về ảnh hưởng của kiến trúc phải được diễn giải trong bối
cảnh các training recipe đặc thù theo kiến trúc đã được xác định trước, không
được xem là bằng chứng về hiệu ứng nhân quả thuần túy của backbone.

Ngoài ra, AMP được bật trong huấn luyện nhằm giảm chi phí bộ nhớ và tăng hiệu
quả tính toán. Toàn bộ các stage của backbone được phép cập nhật tham số trong
quá trình huấn luyện. Gradient clipping không được sử dụng trong giao thức thực nghiệm chính.

Các ảnh \textit{No Finding} có zero ground-truth bounding box vẫn được giữ
trong quá trình huấn luyện. Các ảnh này không được loại bỏ chỉ vì không chứa
annotation, nhằm bảo toàn phân bố dữ liệu âm tính đã được xác định trong tập
huấn luyện.

\subsubsection{Ngân sách cập nhật và lựa chọn checkpoint}
\label{subsubsec:supervised_update_checkpoint}

Ngân sách huấn luyện được định nghĩa theo số lần cập nhật optimizer thay vì số
epoch. Lựa chọn này nhằm tạo một đơn vị ngân sách huấn luyện trực tiếp gắn với
số lần cập nhật tham số, đặc biệt khi kích thước của các tập có nhãn thay đổi
rất lớn giữa các mức ngân sách \(1\%\), \(5\%\), \(10\%\) và \(20\%\).

Nếu chỉ sử dụng cùng số epoch cho mọi ngân sách, số lần cập nhật optimizer có
thể khác đáng kể giữa các điều kiện do số lượng ảnh có nhãn khác nhau. Vì vậy,
nghiên cứu xác định trước ngân sách huấn luyện theo optimizer updates:

\begin{equation}
U_{\mathrm{opt}}(b)
=
\begin{cases}
1032, & b=1\%,\\
2064, & b=5\%,\\
2064, & b=10\%,\\
2064, & b=20\%.
\end{cases}
\label{eq:supervised_optimizer_update_budget}
\end{equation}

Đối với supervised baseline sử dụng toàn bộ tập huấn luyện có nhãn:

\begin{equation}
U_{\mathrm{opt}}(100\%)
=
10284.
\label{eq:supervised_full_label_update_budget}
\end{equation}

Các giá trị trong
Phương trình~\eqref{eq:supervised_optimizer_update_budget} và
Phương trình~\eqref{eq:supervised_full_label_update_budget} là các quyết định
được xác định trước của nghiên cứu, không phải các giá trị bắt buộc được kế
thừa từ một paper bên ngoài.

Để duy trì cùng cấu trúc scheduler giữa các điều kiện có ngân sách huấn luyện
khác nhau, nghiên cứu xác định hai mốc giảm learning rate tại các vị trí tương
đối bằng \(2/3\) và \(11/12\) của tổng ngân sách actual optimizer updates.
Theo quy tắc này, đối với các điều kiện \(5\%\), \(10\%\) và \(20\%\), mỗi điều
kiện có 2064 optimizer updates và các mốc giảm learning rate lần lượt là update
1376 và 1892. Với ngân sách \(1\%\) gồm 1032 optimizer updates, các mốc tương
ứng là update 688 và 946. Đối với điều kiện supervised tham chiếu sử dụng
\(100\%\) dữ liệu có nhãn với ngân sách 10284 optimizer updates, các mốc tương
ứng được đặt tại update 6856 và 9427. Quy tắc này được xác định theo vị trí
tương đối trong tổng ngân sách cập nhật nhằm giữ nhất quán cấu trúc scheduler
giữa các điều kiện huấn luyện có độ dài khác nhau.

Validation được thực hiện sau mỗi 172 optimizer updates. Khoảng cách này được
giữ như một quy tắc thống nhất để theo dõi tiến trình huấn luyện và hỗ trợ lựa
chọn checkpoint theo cùng một giao thức.

Checkpoint BEST được xác định bởi bbox mAP@[0.50:0.95] trên tập validation:

\begin{equation}
Checkpoint_{\mathrm{BEST}}
=
\arg\max_{t}
mAP^{val}_{t}@[0.50:0.95].
\label{eq:supervised_best_checkpoint}
\end{equation}

Đồng thời, nghiên cứu giữ cả checkpoint BEST và checkpoint LAST:

\begin{equation}
\mathcal{C}_{\mathrm{retain}}
=
\left\{
BEST,\,
LAST
\right\}.
\label{eq:supervised_checkpoint_retention}
\end{equation}

BEST được sử dụng để đại diện cho mô hình supervised được lựa chọn theo hiệu
năng validation, trong khi LAST được giữ để phục vụ truy vết trạng thái huấn
luyện và bảo đảm khả năng đối chiếu nhất quán với các nhánh phân tích khác khi
cần thiết.

Tóm lại, supervised baseline được xây dựng trên cùng họ Faster R-CNN + FPN với
hai backbone ResNet-50 và Swin-T. Hai kiến trúc được khởi tạo từ backbone
pretrained trên ImageNet-1K, sử dụng cùng biểu diễn ảnh đầu vào và cùng nguyên
tắc đánh giá nhưng giữ các training recipe đặc thù theo kiến trúc. Ngân sách
huấn luyện được kiểm soát theo optimizer updates thay vì epoch, và checkpoint
được lựa chọn trên tập validation theo bbox mAP@[0.50:0.95]. Thiết kế này cung
cấp các mốc supervised có kiểm soát cho các phép so sánh với semi-supervised
learning ở các phần tiếp theo.

\subsection{Semi-supervised Object Detection Framework}
\label{subsec:Semi_supervised_Object_Detection_Framework}

Nghiên cứu sử dụng
\textbf{khung semi-supervised object detection teacher--student dựa trên
pseudo-labeling, triển khai theo SoftTeacher}. Việc lựa chọn cơ chế
teacher--student pseudo-labeling được đặt trong tiến trình phát triển của
semi-supervised object detection, trong đó các công trình như STAC,
Unbiased Teacher và Soft Teacher đã cho thấy khả năng khai thác dữ liệu
chưa gán nhãn thông qua pseudo-label trong điều kiện dữ liệu gán nhãn hạn
chế \cite{sohn2020simple,liu2021unbiased,xu2021end}. Các công trình này
cung cấp cơ sở học thuật cho việc sử dụng teacher--student và
pseudo-labeling trong bài toán object detection; tuy nhiên, việc lựa chọn
SoftTeacher làm cấu hình SSL chính của nghiên cứu là một quyết định thiết kế
được xác định trước, không phải kết luận rằng đây là phương pháp tốt nhất
cho mọi bài toán.

Trong số các hướng teacher--student pseudo-labeling nói trên, nghiên cứu lựa
chọn triển khai theo SoftTeacher vì khung này tích hợp trực tiếp các thành
phần phù hợp với mục tiêu thực nghiệm của đề tài, bao gồm pseudo-label trực
tuyến, EMA Teacher, lọc pseudo-label cho nhánh phân loại, đánh giá độ tin
cậy của pseudo-label cho hồi quy bounding box và cơ chế weak/strong
augmentation \cite{xu2021end}. Các thành phần này tạo thành một quy trình
semi-supervised object detection end-to-end và đồng thời cung cấp các cơ chế
có thể được khảo sát riêng trong thí nghiệm phân tích thành phần.

Một lý do quan trọng khác là SoftTeacher được xây dựng trong cùng họ detector
hai giai đoạn Faster R-CNN với supervised baseline của nghiên cứu
\cite{ren2016faster,xu2021end}. Nhờ đó, khi chuyển từ supervised learning
sang semi-supervised learning, nghiên cứu có thể giữ nguyên họ detector,
backbone tương ứng, FPN và các điều kiện huấn luyện được kiểm soát, trong khi
khác biệt có chủ đích chủ yếu nằm ở việc SSL khai thác thêm tập chưa gán nhãn
\(U_b\) thông qua Teacher và pseudo-label. Thiết kế này hạn chế việc trộn lẫn
ảnh hưởng của semi-supervised learning với ảnh hưởng do thay đổi detector
family.

Ngoài ra, các thành phần confidence filtering, strong augmentation và
regression pseudo-label reliability trong SoftTeacher có liên hệ trực tiếp
với mục tiêu phân tích cơ chế của nghiên cứu \cite{xu2021end}. Chúng cho phép
xây dựng các thí nghiệm phân tích thành phần đã xác định trước mà không cần
chuyển sang một framework SSL khác. Vì vậy, SoftTeacher được sử dụng như một
khung SSOD đại diện để đánh giá khả năng khai thác dữ liệu chưa gán nhãn trong
một thiết kế thực nghiệm có kiểm soát; nghiên cứu không đề xuất một thuật toán
SSOD mới và không giả định SoftTeacher là phương pháp tối ưu cho dữ liệu
X-quang phổi.

Việc lựa chọn SoftTeacher phù hợp với mục tiêu của nghiên cứu vì ba lý do
chính. Thứ nhất, kiến trúc teacher--student và cơ chế pseudo-label trực tuyến
cho phép khai thác trực tiếp tập dữ liệu chưa gán nhãn \(U_b\) mà không thay
đổi cấu trúc nhãn của bài toán object detection. Thứ hai, SoftTeacher được xây
dựng trên họ detector hai giai đoạn Faster R-CNN, phù hợp với supervised
baseline đã được lựa chọn ở Mục~\ref{subsec:Supervised_Baseline}, nhờ đó giảm
việc thay đổi đồng thời detector family khi chuyển từ SUP sang SSL. Thứ ba,
khung này cung cấp các thành phần có thể phân tích riêng như confidence
filtering, strong augmentation và regression uncertainty, phù hợp với mục
tiêu phân tích cơ chế ở thí nghiệm ablation phía sau. Việc lựa chọn
SoftTeacher vì vậy phục vụ mục tiêu đánh giá một phương pháp SSOD đại diện
trong controlled experiment, không hàm ý đây là phương pháp tốt nhất cho mọi
bài toán X-quang phổi.

\subsubsection{Cấu trúc Teacher--Student}
\label{subsubsec:teacher_student_formulation}

Trong mỗi điều kiện SSL, Student và Teacher có cùng kiến trúc detector. Nếu
kiến trúc đang xét là R50 thì cả hai sử dụng Faster R-CNN + ResNet-50 + FPN;
nếu kiến trúc đang xét là Swin-T thì cả hai sử dụng Faster R-CNN + Swin-T +
FPN. Vì vậy,

\begin{equation}
\mathcal{A}_{Teacher}
=
\mathcal{A}_{Student}.
\label{eq:teacher_student_same_architecture}
\end{equation}

Student là mô hình được tối ưu trực tiếp bằng gradient descent. Teacher không
được tối ưu bằng một optimizer độc lập mà được cập nhật từ trạng thái của
Student thông qua cơ chế exponential moving average (EMA).

Tại thời điểm khởi tạo \(t_0\), Teacher và Student có cùng trạng thái tham số:

\begin{equation}
\theta_T^{(t_0)}
=
\theta_S^{(t_0)},
\label{eq:teacher_student_initialization}
\end{equation}

trong đó \(\theta_T\) và \(\theta_S\) lần lượt biểu diễn tham số của Teacher
và Student.

Nghiên cứu không sử dụng một giai đoạn supervised burn-in riêng để huấn luyện
Student trước rồi mới khởi tạo Teacher từ một supervised checkpoint. Teacher
được tạo đồng thời với Student ngay từ thời điểm bắt đầu huấn luyện SSL. Cách
khởi tạo này được giữ nhất quán trong toàn bộ các điều kiện SSL.

Khi một lần huấn luyện được tiếp tục từ checkpoint sau gián đoạn kỹ thuật,
trạng thái Teacher đã lưu trong checkpoint phải được khôi phục cùng Student;
Teacher không được đồng bộ lại từ Student tại thời điểm resume. Quy tắc này
giúp duy trì đúng lịch sử EMA của mô hình.

\subsubsection{Cập nhật Teacher bằng EMA}
\label{subsubsec:teacher_ema_update}

Teacher được cập nhật một lần sau mỗi lần cập nhật optimizer thực tế của
Student. Với quy ước EMA đã xác định trước, tham số Teacher tại bước \(t\)
được tính bởi

\begin{equation}
\theta_T^{(t)}
=
0.999\,
\theta_T^{(t-1)}
+
0.001\,
\theta_S^{(t)}.
\label{eq:teacher_ema_update}
\end{equation}

Trong Phương trình~\eqref{eq:teacher_ema_update}, hệ số \(0.999\) giữ phần lớn
thông tin tích lũy của Teacher từ các bước trước, trong khi hệ số \(0.001\)
đưa trạng thái mới của Student vào Teacher. Nhờ đó, Teacher thay đổi chậm hơn
Student và đóng vai trò mô hình tham chiếu ổn định hơn để sinh pseudo-label.

Trong triển khai, quy tắc này tương ứng với
\texttt{MeanTeacherHook} có \texttt{momentum}=0.001 và
\texttt{skip\_buffers=True}. EMA chỉ được cập nhật sau một optimizer update
thực sự của Student. Nếu một optimizer step bị bỏ qua do lỗi số học trong
quá trình huấn luyện với mixed precision, Teacher không được cập nhật cho
bước đó.

\subsubsection{Luồng dữ liệu có nhãn và chưa gán nhãn}
\label{subsubsec:ssl_labeled_unlabeled_streams}

Tại mỗi ngân sách nhãn \(b\), semi-supervised learning sử dụng cùng tập có
nhãn \(L_b\) như supervised baseline và bổ sung tập chưa gán nhãn \(U_b\).
Hai nguồn dữ liệu có vai trò khác nhau trong quá trình huấn luyện.

Đối với dữ liệu có nhãn, ảnh và ground-truth annotation được đưa trực tiếp
vào Student để tính supervised loss:

\begin{equation}
L_b
\longrightarrow
Student
\longrightarrow
\mathcal{L}_{sup}.
\label{eq:labeled_stream_student}
\end{equation}

Đối với dữ liệu chưa gán nhãn, cùng một ảnh \(x_u\in U_b\) được tạo thành hai
view khác nhau. Weak view được đưa vào Teacher để sinh pseudo-label, trong khi
strong view của cùng ảnh được đưa vào Student:

\begin{equation}
x_u^{weak}
\longrightarrow
Teacher
\longrightarrow
\hat{y}_u,
\label{eq:unlabeled_weak_teacher}
\end{equation}

và

\begin{equation}
\left(
x_u^{strong},
\hat{y}_u
\right)
\longrightarrow
Student
\longrightarrow
\mathcal{L}_{unsup},
\label{eq:unlabeled_strong_student}
\end{equation}

trong đó \(\hat{y}_u\) là tập pseudo-label được Teacher sinh ra sau các bước
lọc đã xác định trước.

Ground truth bị ẩn của \(U_b\) không được sử dụng trong quá trình huấn luyện,
sinh pseudo-label, lọc pseudo-label hoặc lựa chọn threshold. Quy tắc này duy
trì đúng vai trò chưa gán nhãn của \(U_b\) trong toàn bộ SSL protocol.

\subsubsection{Weak và strong augmentation cho dữ liệu chưa gán nhãn}
\label{subsubsec:ssl_weak_strong_augmentation}

Weak view của dữ liệu chưa gán nhãn được giữ bảo thủ nhằm tạo đầu vào tương
đối ổn định cho Teacher. Luồng weak augmentation gồm resize theo cùng quy tắc
hình học của supervised baseline, sau đó áp dụng preprocessing và padding
chuẩn. Nghiên cứu không sử dụng flip, crop, rotation hoặc photometric
augmentation trên weak view.

Strong view được tạo từ cùng ảnh chưa gán nhãn nhưng bổ sung hai phép biến đổi
quang học grayscale gồm brightness và contrast. Hệ số biến đổi được lấy trong
các khoảng đã xác định trước:

\begin{equation}
f_{\mathrm{brightness}}
\sim
\mathcal{U}(0.90,1.10),
\qquad
f_{\mathrm{contrast}}
\sim
\mathcal{U}(0.90,1.10).
\label{eq:ssl_strong_photometric_range}
\end{equation}

Do dữ liệu X-quang được biểu diễn dưới dạng grayscale và ánh xạ thành ba kênh
giống nhau, các phép biến đổi strong view phải giữ

\begin{equation}
R=G=B.
\label{eq:ssl_strong_grayscale_constraint}
\end{equation}

Nhờ đó, strong augmentation làm thay đổi cường độ ảnh nhưng không tạo thông
tin màu giả.

Nghiên cứu không sử dụng geometric augmentation, random erasing hoặc mosaic
trong strong view chính. Việc giới hạn augmentation theo hướng này giúp bảo
toàn geometry của ảnh và bounding box, đồng thời tạo một cấu hình Main rõ
ràng để sau đó đánh giá riêng ảnh hưởng của brightness và contrast trong
thí nghiệm phân tích thành phần.

\subsubsection{Hàm mục tiêu semi-supervised}
\label{subsubsec:ssl_training_objective}

Student được tối ưu bằng tổng của supervised loss trên \(L_b\) và unsupervised
loss sinh từ pseudo-label trên \(U_b\). Hàm mục tiêu chính được xác định bởi

\begin{equation}
\mathcal{L}_{total}
=
\mathcal{L}_{sup}
+
4\mathcal{L}_{unsup}.
\label{eq:ssl_total_loss}
\end{equation}

Trong Phương trình~\eqref{eq:ssl_total_loss},
\(\mathcal{L}_{sup}\) là supervised Faster R-CNN loss được tính từ
ground-truth annotation của \(L_b\), còn
\(\mathcal{L}_{unsup}\) là loss từ nhánh dữ liệu chưa gán nhãn sau khi
pseudo-label của Teacher được chuyển sang Student.

Trọng số supervised được cố định bằng \(1\), trong khi trọng số của
unsupervised loss được cố định bằng

\begin{equation}
\lambda_u
=
4.
\label{eq:ssl_unsupervised_weight}
\end{equation}

Nghiên cứu không sử dụng burn-in hoặc ramp-up cho
\(\lambda_u\). Vì vậy,

\begin{equation}
\lambda_u(t)
=
4,
\qquad
\forall t,
\label{eq:ssl_constant_unsupervised_weight}
\end{equation}

trong toàn bộ quá trình huấn luyện SSL.

Việc giữ \(\lambda_u\) cố định giúp cấu hình Main được xác định trước và
tránh việc bổ sung thêm một quá trình điều chỉnh trọng số không giám sát dựa
trên kết quả validation sau khi huấn luyện.

\subsubsection{Tỷ lệ lấy mẫu labeled--unlabeled}
\label{subsubsec:ssl_sampling_ratio}

Tại mỗi optimizer update, dữ liệu có nhãn và chưa gán nhãn được lấy theo tỷ
lệ

\begin{equation}
L:U
=
1:1.
\label{eq:ssl_labeled_unlabeled_ratio}
\end{equation}

Với effective labeled batch size bằng \(4\), mỗi optimizer update của nhánh
SSL sử dụng hiệu dụng:

\begin{equation}
B^{L}_{eff}
=
4,
\qquad
B^{U}_{eff}
=
4.
\label{eq:ssl_effective_batches}
\end{equation}

Do đó, Student nhận tín hiệu có giám sát từ cùng mức labeled exposure đã được
kiểm soát trong supervised baseline, trong khi SSL được bổ sung tín hiệu từ
dữ liệu chưa gán nhãn. Quy tắc này giúp hạn chế cách diễn giải rằng khác biệt
SUP--SSL xuất hiện chỉ vì SSL nhận nhiều lần cập nhật có giám sát hơn.

\subsubsection{Mô hình Teacher dùng cho đánh giá}
\label{subsubsec:ssl_official_evaluation_model}

Trong các điều kiện SSL, EMA Teacher là mô hình chính thức được sử dụng để
đánh giá. Student được giữ như một mô hình phụ trợ và không thay thế Teacher
trong endpoint đánh giá chính.

Checkpoint BEST của EMA Teacher được xác định bằng bbox
mAP@[0.50:0.95] trên tập validation:

\begin{equation}
Teacher_{\mathrm{BEST}}
=
\arg\max_t
mAP^{val}_{Teacher,t}@[0.50:0.95].
\label{eq:ssl_best_ema_teacher}
\end{equation}

BEST EMA Teacher là mô hình SSL được sử dụng cho đánh giá phát hiện cuối cùng.
Quy tắc này tương ứng trực tiếp với supervised baseline, trong đó checkpoint
BEST cũng được lựa chọn theo bbox mAP@[0.50:0.95] trên cùng tập validation,
nhưng việc lựa chọn trong SSL được thực hiện trên Teacher thay vì Student.

Bên cạnh BEST EMA Teacher, checkpoint LAST EMA Teacher cũng được giữ lại.
Hai checkpoint này có vai trò khác nhau:

\begin{equation}
\begin{aligned}
Teacher_{\mathrm{BEST}}
&\rightarrow
\text{đánh giá hiệu năng phát hiện cuối cùng},\\
Teacher_{\mathrm{LAST}}
&\rightarrow
Q_{\mathrm{pseudo}}\text{ trên validation}.
\end{aligned}
\label{eq:ssl_best_last_teacher_roles}
\end{equation}

Việc sử dụng LAST EMA Teacher để tính \(Q_{\mathrm{pseudo}}\) được xác định
trước và không được thay bằng BEST Teacher sau khi quan sát kết quả. Quy tắc
này tách biệt lựa chọn checkpoint tối ưu cho hiệu năng phát hiện khỏi phép đo
chất lượng pseudo-label dùng trong phân tích cơ chế.

Tóm lại, SSL framework của nghiên cứu giữ cùng họ Faster R-CNN và cùng
backbone tương ứng với supervised baseline, nhưng bổ sung Teacher được cập
nhật bằng EMA và khai thác tập chưa gán nhãn thông qua pseudo-label trực
tuyến. Weak view của \(U_b\) được sử dụng để Teacher sinh pseudo-label,
strong view được đưa vào Student, và Student được tối ưu bởi tổng supervised
và unsupervised loss với trọng số cố định. EMA Teacher là mô hình chính thức
cho đánh giá SSL, trong đó BEST Teacher phục vụ đánh giá hiệu năng phát hiện
và LAST Teacher phục vụ phép đo \(Q_{\mathrm{pseudo}}\).

\subsection{Pseudo-label Generation and Filtering}
\label{subsec:Pseudo_label_Generation_and_Filtering}

Pseudo-label là thành phần trung tâm kết nối EMA Teacher và Student trong
khung semi-supervised object detection của nghiên cứu. Đối với mỗi ảnh chưa
gán nhãn \(x_u\in U_b\), Teacher thực hiện suy luận trên weak view để sinh các
dự đoán detection, sau đó các dự đoán này được lọc theo các tiêu chí đã xác
định trước trước khi được sử dụng như tín hiệu huấn luyện cho Student trên
strong view.

Cơ chế này được triển khai theo SoftTeacher, một khung teacher--student cho
semi-supervised object detection kết hợp pseudo-label trực tuyến, EMA Teacher,
confidence filtering, regression uncertainty filtering và weak/strong
augmentation \cite{xu2021end}. Việc lựa chọn cách tổ chức này phù hợp với mục
tiêu của nghiên cứu là đánh giá một phương pháp SSOD đại diện trong một thiết
kế thực nghiệm có kiểm soát, không đề xuất một thuật toán SSOD mới.

\subsubsection{Nguồn và cơ chế sinh pseudo-label}
\label{subsubsec:pseudo_label_source}

Trong cấu hình SSL chính, pseudo-label được sinh \textbf{trực tuyến}
(\textit{online}) từ các detection của EMA Teacher thay vì được xây dựng trước
như một tập nhãn cố định. Lựa chọn này được giữ theo thiết kế SoftTeacher đã
được xác định cho nghiên cứu, trong đó Teacher, Student và quá trình
pseudo-labeling được tích hợp trong cùng một quy trình huấn luyện end-to-end
\cite{xu2021end}.

Khi trạng thái EMA Teacher thay đổi trong quá trình huấn luyện, pseudo-label
sinh ra từ Teacher cũng được cập nhật tương ứng. Vì vậy, tín hiệu không giám
sát truyền tới Student phản ánh trạng thái hiện hành của Teacher thay vì bị cố
định tại một thời điểm trước huấn luyện.

Trong nghiên cứu này, cơ chế online pseudo-labeling không được diễn giải như
một phương án đã được chứng minh vượt trội so với offline pseudo-labeling.
Đây là một thành phần của khung SoftTeacher được xác định trước và được giữ
nhất quán trong toàn bộ các điều kiện SSL.

Đối với một ảnh chưa gán nhãn \(x_u\), Teacher tạo tập detection trên weak
view:

\begin{equation}
\mathcal{D}_{T}(x_u^{weak})
=
\left\{
\left(
\hat{b}_j,\,
\hat{c}_j,\,
\hat{s}_j
\right)
\right\}_{j=1}^{N_u},
\label{eq:teacher_raw_detections}
\end{equation}

trong đó \(\hat{b}_j\) là bounding box dự đoán, \(\hat{c}_j\) là lớp dự đoán
và \(\hat{s}_j\) là confidence score của detection thứ \(j\).

Pseudo-label cuối cùng được xây dựng từ các detection của RCNN Teacher, không
phải trực tiếp từ proposal của RPN. Một pseudo-label có dạng

\begin{equation}
\hat{y}_j
=
\left(
\hat{b}_j,\,
\hat{c}_j,\,
\hat{s}_j
\right).
\label{eq:pseudo_label_structure}
\end{equation}

RPN chỉ tạo các vùng đề xuất ứng viên trong quy trình hai giai đoạn của Faster
R-CNN. Các proposal này chưa đồng thời chứa đầy đủ quyết định lớp, bounding box
sau tinh chỉnh và confidence score ở mức detection. Vì vậy, RCNN detections
của Teacher được sử dụng làm nguồn pseudo-label cuối cùng. Cách lựa chọn này
cũng giữ pseudo-label ở dạng thông tin tương ứng trực tiếp với bài toán
detection cần truyền cho Student.

\subsubsection{Confidence filtering cho pseudo-label}
\label{subsubsec:pseudo_label_confidence_thresholds}

Pseudo-label là nhãn do mô hình tự sinh nên có thể chứa sai số. Nếu mọi
detection của Teacher đều được chuyển trực tiếp sang Student, các dự đoán có
độ tin cậy thấp cũng có thể trở thành tín hiệu huấn luyện và làm khuếch đại
sai số pseudo-label. Vì vậy, cấu hình Main sử dụng confidence filtering để
giới hạn tập pseudo-label được chấp nhận cho các nhánh tương ứng. Cơ chế này
phù hợp với thiết kế SoftTeacher \cite{xu2021end}.

Quy trình sử dụng nhiều ngưỡng với các vai trò khác nhau. Trước hết, detector
giữ các detection ứng viên với score floor

\begin{equation}
\tau_{\mathrm{det}}
=
0.05.
\label{eq:pseudo_detector_score_floor}
\end{equation}

Ngưỡng initial score được xác định là

\begin{equation}
\tau_{\mathrm{init}}
=
0.50.
\label{eq:pseudo_initial_score_threshold}
\end{equation}

Đối với proposal dùng trong nhánh không giám sát của RPN, ngưỡng confidence
được đặt ở mức

\begin{equation}
\tau_{\mathrm{RPN}}
=
0.90.
\label{eq:rpn_pseudo_threshold}
\end{equation}

Đối với pseudo-label sử dụng cho nhánh phân loại RCNN, ngưỡng chính được xác
định là

\begin{equation}
\tau_{\mathrm{cls}}
=
0.90.
\label{eq:rcnn_classification_pseudo_threshold}
\end{equation}

Do đó, tập pseudo-label được chấp nhận cho nhánh classification được biểu diễn
bởi

\begin{equation}
\widehat{\mathcal{Y}}^{cls}_{u}
=
\left\{
\hat{y}_j
\in
\mathcal{D}_{T}(x_u^{weak})
:
\hat{s}_j
>
\tau_{\mathrm{cls}}
\right\}.
\label{eq:accepted_classification_pseudo_set}
\end{equation}

Các giá trị
\(\tau_{\mathrm{det}}\),
\(\tau_{\mathrm{init}}\),
\(\tau_{\mathrm{RPN}}\) và
\(\tau_{\mathrm{cls}}\)
là các giá trị được nghiên cứu xác định trước cho cấu hình SSL chính. Chúng
không được trình bày như các giá trị bắt buộc do SoftTeacher paper quy định
và không được điều chỉnh hậu nghiệm sau khi quan sát kết quả test.

Đặc biệt, ảnh hưởng của
\(\tau_{\mathrm{cls}}\)
được khảo sát riêng trong ablation study. Kết quả ablation không được sử dụng
để thay đổi ngược cấu hình Main sau khi đã quan sát validation performance.

\subsubsection{Non-maximum suppression và giới hạn số proposal}
\label{subsubsec:pseudo_label_nms}

Non-maximum suppression (NMS) được sử dụng để giảm các proposal hoặc detection
trùng lặp mạnh trước khi hình thành tập pseudo-label cuối cùng.

Đối với RPN, ngưỡng IoU của NMS được xác định là

\begin{equation}
\tau_{\mathrm{NMS}}^{RPN}
=
0.70,
\label{eq:rpn_nms_threshold}
\end{equation}

với số proposal tối đa

\begin{equation}
N_{\max}^{RPN}
=
1000.
\label{eq:rpn_max_proposals}
\end{equation}

Đối với RCNN, ngưỡng NMS được xác định là

\begin{equation}
\tau_{\mathrm{NMS}}^{RCNN}
=
0.50,
\label{eq:rcnn_nms_threshold}
\end{equation}

và số detection tối đa sau NMS là

\begin{equation}
N_{\max}^{RCNN}
=
100.
\label{eq:rcnn_max_detections}
\end{equation}

Các thiết lập này được giữ cố định trong cấu hình Main và không thay đổi giữa
các ngân sách nhãn hoặc training seeds.

\subsubsection{Độ tin cậy của pseudo-label cho hồi quy bounding box}
\label{subsubsec:pseudo_regression_reliability}

Confidence của lớp và độ tin cậy của tọa độ bounding box không phản ánh cùng
một dạng bất định. Một detection có confidence phân loại cao chưa bảo đảm rằng
tọa độ bounding box của nó ổn định; ngược lại, một bounding box có vị trí ổn
định không nhất thiết phải vượt một ngưỡng phân loại nghiêm ngặt.

Vì vậy, nghiên cứu không sử dụng một confidence threshold duy nhất cho cả
classification và regression. Nhánh classification được kiểm soát chủ yếu
bằng confidence score, trong khi nhánh regression sử dụng initial score kết
hợp với regression uncertainty. Cách tách này kế thừa cơ chế xử lý
pseudo-label của SoftTeacher và cho phép kiểm soát riêng sai số về lớp và sai
số về định vị \cite{xu2021end}.

Bounding-box regression là một bài toán liên tục; confidence score của lớp
không trực tiếp đo độ ổn định của tọa độ bounding box. Vì vậy, nghiên cứu sử
dụng regression uncertainty như một chỉ báo bổ sung về độ tin cậy của
pseudo-label cho nhánh hồi quy.

Đối với mỗi bounding box ứng viên, Teacher thực hiện jitter nhiều lần quanh
bounding box để đánh giá mức biến thiên của dự đoán hồi quy. Số lần jitter
được xác định là

\begin{equation}
N_{\mathrm{jitter}}
=
10,
\label{eq:regression_jitter_times}
\end{equation}

với biên độ jitter

\begin{equation}
\gamma_{\mathrm{jitter}}
=
0.06.
\label{eq:regression_jitter_scale}
\end{equation}

Nếu dự đoán tọa độ thay đổi mạnh khi bounding box đầu vào chỉ bị nhiễu nhẹ,
pseudo-label được xem là kém ổn định hơn. Ngược lại, mức biến thiên thấp cho
thấy vị trí bounding box ổn định hơn. Cơ chế này phù hợp với regression
uncertainty filtering trong SoftTeacher \cite{xu2021end}.

Ngưỡng regression uncertainty được xác định là

\begin{equation}
\tau_{\mathrm{reg}}
=
0.02.
\label{eq:regression_uncertainty_threshold}
\end{equation}

Đồng thời, pseudo bounding box phải thỏa điều kiện kích thước tối thiểu

\begin{equation}
\left(
w_{\min},
h_{\min}
\right)
=
(0.01,0.01).
\label{eq:min_pseudo_bbox_wh}
\end{equation}

Trong cấu hình Main, tập pseudo-label dùng cho nhánh hồi quy được xác định bởi

\begin{equation}
\widehat{\mathcal{Y}}^{reg}_{u}
=
\left\{
\hat{y}_j
:
\hat{s}^{initial}_j>0.50
\;\land\;
u_{\mathrm{reg},j}<0.02
\right\}.
\label{eq:accepted_regression_pseudo_set}
\end{equation}

Một điểm cần phân biệt rõ là pseudo-label dùng cho regression
\textbf{không bắt buộc} phải vượt ngưỡng classification
\(\tau_{\mathrm{cls}}=0.90\). Cụ thể,

\begin{equation}
\hat{y}_j
\in
\widehat{\mathcal{Y}}^{reg}_{u}
\not\Rightarrow
\hat{s}_j
>
\tau_{\mathrm{cls}}.
\label{eq:regression_not_require_cls_threshold}
\end{equation}

Quy tắc này tránh việc đồng nhất confidence của lớp với độ ổn định của tọa độ
bounding box, hai thuộc tính có vai trò khác nhau trong pseudo-label quality.

\subsubsection{Luồng pseudo-label từ Teacher sang Student}
\label{subsubsec:pseudo_label_teacher_student_flow}

Sau các bước lọc, pseudo-label được truyền từ weak view của Teacher sang strong
view tương ứng của Student. Luồng tổng quát được biểu diễn bởi

\begin{equation}
x_u^{weak}
\xrightarrow{Teacher}
\mathcal{D}_{T}
\xrightarrow{\mathrm{filtering}}
\widehat{\mathcal{Y}}_{u}
\xrightarrow{}
Student(x_u^{strong}).
\label{eq:pseudo_label_generation_pipeline}
\end{equation}

Trong đó,
\(\widehat{\mathcal{Y}}_{u}\)
không nhất thiết là một tập duy nhất được sử dụng giống nhau cho mọi thành
phần loss. Nhánh classification và nhánh regression áp dụng các tiêu chí
reliability khác nhau như đã trình bày tại
Phương trình~\eqref{eq:accepted_classification_pseudo_set}
và
Phương trình~\eqref{eq:accepted_regression_pseudo_set}.

Cách tổ chức này cho phép Teacher lựa chọn pseudo-label phù hợp với từng nhiệm
vụ con của detector thay vì sử dụng một confidence threshold duy nhất như đại
diện cho toàn bộ chất lượng pseudo-label.

\subsubsection{Định nghĩa và vai trò của chất lượng pseudo-label}
\label{subsubsec:pseudo_label_quality}

Ngoài việc sử dụng pseudo-label làm tín hiệu huấn luyện, nghiên cứu còn định
nghĩa một chỉ số riêng để định lượng chất lượng pseudo-label phục vụ RQ4 và
phân tích cơ chế.

Việc lựa chọn metric xuất phát từ chính cấu trúc của pseudo-label trong object
detection. Một pseudo-label không chỉ chứa confidence mà còn bao gồm quyết
định lớp và bounding box. Vì vậy, các chỉ số như mean confidence, số lượng
pseudo-label hoặc retention rate chỉ phản ánh một phần hành vi của
pseudo-label và không đánh giá đồng thời chất lượng classification và
localization.

Nghiên cứu vì vậy sử dụng

\begin{equation}
Q_{\mathrm{pseudo}}
=
PL\text{-}mAP@[0.50:0.95]
\label{eq:pseudo_quality_metric}
\end{equation}

làm chỉ số chất lượng pseudo-label chính.

PL-mAP@[0.50:0.95] đánh giá pseudo-label trên nhiều ngưỡng IoU và do đó phản
ánh đồng thời khả năng phân loại và độ chính xác định vị của pseudo-label.
Lựa chọn này phù hợp hơn việc chỉ sử dụng mean confidence, pseudo-label count
hoặc retention rate để đại diện cho toàn bộ chất lượng pseudo-label.

Các chỉ số như PL-AP50, PL-AP75, precision, recall, F1, mean matched IoU,
class-wise pseudo-label AP, pseudo-label count, retention rate hoặc regression
uncertainty có thể được sử dụng như các chỉ số chẩn đoán cơ chế bổ sung, nhưng
không thay thế
\(Q_{\mathrm{pseudo}}\).

\(Q_{\mathrm{pseudo}}\) cũng không thay thế chỉ số hiệu năng phát hiện của mô
hình. Hai loại chỉ số trả lời hai câu hỏi khác nhau: bbox mAP phản ánh thay
đổi ở detector, trong khi \(Q_{\mathrm{pseudo}}\) phản ánh thay đổi ở tín hiệu
pseudo-label trung gian.

Chuỗi phân tích cơ chế của nghiên cứu được biểu diễn bởi

\begin{equation}
\text{thành phần SSL}
\longrightarrow
\text{hành vi/chất lượng pseudo-label}
\longrightarrow
\text{hiệu năng phát hiện}.
\label{eq:pseudo_mechanistic_chain}
\end{equation}

Nếu chỉ sử dụng bbox mAP, nghiên cứu có thể biết hiệu năng detector thay đổi
nhưng khó xác định thay đổi đó có đi kèm với thay đổi chất lượng pseudo-label
hay không. Ngược lại, nếu chỉ sử dụng \(Q_{\mathrm{pseudo}}\), nghiên cứu biết
pseudo-label thay đổi nhưng chưa thể xác định sự thay đổi đó có chuyển thành
thay đổi hiệu năng phát hiện hay không.

Việc sử dụng đồng thời hai mức đo vì vậy hỗ trợ phân tích cơ chế mà không biến
\(Q_{\mathrm{pseudo}}\) thành một endpoint chính thay thế cho chỉ số hiệu năng
phát hiện.

\subsubsection{Checkpoint dùng để đánh giá chất lượng pseudo-label}
\label{subsubsec:qpseudo_checkpoint}

\(Q_{\mathrm{pseudo}}\) được tính trên tập validation cố định bằng
\textbf{LAST EMA Teacher} và trên tập RCNN classification pseudo-label đã được
chấp nhận theo giao thức chính:

\begin{equation}
Teacher_{Q_{\mathrm{pseudo}}}
=
Teacher_{\mathrm{LAST}}.
\label{eq:qpseudo_last_teacher}
\end{equation}

Việc cố định LAST EMA Teacher tạo một quy tắc đo thống nhất giữa các điều kiện
và tách phép đo chất lượng pseudo-label khỏi quy trình lựa chọn BEST checkpoint
theo validation bbox mAP.

Hai checkpoint vì vậy có vai trò khác nhau:

\begin{equation}
\begin{aligned}
Teacher_{\mathrm{BEST}}
&\rightarrow
\text{đánh giá hiệu năng phát hiện},\\
Teacher_{\mathrm{LAST}}
&\rightarrow
Q_{\mathrm{pseudo}}.
\end{aligned}
\label{eq:best_last_teacher_separation}
\end{equation}

BEST EMA Teacher không được thay thế LAST EMA Teacher để tính
\(Q_{\mathrm{pseudo}}\) sau khi quan sát kết quả. Quy tắc này được xác định
trước để tách lựa chọn checkpoint tốt nhất về detection performance khỏi phép
đo pseudo-label quality dùng cho phân tích cơ chế.

\subsubsection{Firewall đối với ground truth của dữ liệu chưa gán nhãn}
\label{subsubsec:pseudo_label_hidden_gt_firewall}

Ground truth bị ẩn của \(U_b\) không được sử dụng trong bất kỳ bước nào liên
quan đến xây dựng, điều chỉnh hoặc đánh giá pseudo-label trong quá trình
training.

Nguyên tắc này xuất phát từ vai trò của \(U_b\) như nguồn dữ liệu chưa gán
nhãn trong thiết kế semi-supervised learning. Nếu ground truth của \(U_b\)
được sử dụng để lựa chọn threshold, lọc pseudo-label, điều chỉnh EMA, lựa chọn
augmentation, lựa chọn checkpoint hoặc trực tiếp đánh giá chất lượng
pseudo-label, thông tin nhãn thật của tập chưa gán nhãn sẽ quay trở lại quá
trình xây dựng phương pháp. Khi đó, điều kiện thực nghiệm không còn phản ánh
đúng bài toán semi-supervised learning đã xác định.

Firewall được biểu diễn bởi

\begin{equation}
GT(U_b)
\not\rightarrow
\left\{
\begin{array}{l}
\text{training},\\
\text{pseudo-label filtering},\\
\text{threshold tuning},\\
\text{EMA tuning},\\
\text{augmentation selection},\\
\text{checkpoint selection},\\
Q_{\mathrm{pseudo}}.
\end{array}
\right.
\label{eq:hidden_u_gt_firewall}
\end{equation}

Chất lượng pseudo-label chỉ được đánh giá trên tập validation cố định, nơi
ground truth được phép sử dụng cho mục đích đánh giá đã xác định trước. Hidden
ground truth của \(U_b\) không được sử dụng để tính
\(Q_{\mathrm{pseudo}}\).

Tập test cũng không được sử dụng để tính \(Q_{\mathrm{pseudo}}\), lựa chọn
confidence threshold, lựa chọn filtering strategy hoặc điều chỉnh
pseudo-label pipeline. Việc giữ test ngoài quá trình này nhằm bảo toàn vai trò
đánh giá cuối cùng độc lập của tập test.

Tóm lại, pseudo-label trong nghiên cứu được sinh trực tuyến từ RCNN detections
của EMA Teacher và được xử lý bằng các tiêu chí riêng cho classification và
bounding-box regression. Classification sử dụng confidence filtering, trong
khi regression kết hợp initial score với regression uncertainty. Chất lượng
pseudo-label được định lượng bằng
\(PL\text{-}mAP@[0.50:0.95]\) trên fixed validation sử dụng LAST EMA Teacher.
Thiết kế này vừa cung cấp tín hiệu không giám sát cho Student, vừa cho phép
phân tích cơ chế pseudo-label mà không sử dụng ground truth bị ẩn của
\(U_b\).

\subsection{Giao thức thực nghiệm}
\label{subsec:experimental_protocol}

\subsubsection{Giả thuyết nghiên cứu}
\label{subsubsec:research_hypotheses}

Các giả thuyết nghiên cứu được xác định trước khi thực hiện các thí nghiệm
huấn luyện chính thức nhằm liên kết tường minh giữa các câu hỏi nghiên cứu,
thiết kế thực nghiệm, chỉ số đánh giá và phương pháp phân tích thống kê.
Không phải tất cả các câu hỏi nghiên cứu đều được chuyển thành một giả thuyết
kiểm định chính thức. Cụ thể, RQ1 có vai trò mô tả và RQ5 có vai trò khám phá,
trong khi RQ2, RQ3, RQ4, RQ6 và RQ7 được gắn với các giả thuyết nghiên cứu
chính thức với vai trò và hướng kiểm định được xác định trước.

Chỉ số hiệu năng phát hiện chính được sử dụng cho giả thuyết trung tâm của
nghiên cứu là bbox mAP@[0.50:0.95]. Với một kiến trúc \(a\), ngân sách nhãn
\(b\) và training seed \(s\), mức lợi ích của semi-supervised learning so với
supervised baseline được định nghĩa bởi hiệu ứng ghép cặp

\begin{equation}
\Delta M_{a,b,s}
=
M^{SSL}_{a,b,s}
-
M^{SUP}_{a,b,s},
\label{eq:paired_metric_effect}
\end{equation}

trong đó \(M\) biểu diễn chỉ số đánh giá được xem xét. Đối với giả thuyết chính
RQ2, \(M\) được xác định là bbox mAP@[0.50:0.95]. Việc ghép cặp được thực hiện
trên cùng kiến trúc, cùng ngân sách nhãn và cùng training seed nhằm giảm ảnh
hưởng của biến thiên thực nghiệm không liên quan đến phương thức học.

\paragraph{RQ1 -- Hiệu năng supervised baseline theo ngân sách nhãn.}

RQ1 có vai trò mô tả và không gắn với một giả thuyết thống kê khẳng định
riêng. Nghiên cứu không giả định trước rằng hiệu năng supervised baseline bắt
buộc phải tăng đơn điệu theo thứ tự

\begin{equation}
M_{1\%}
<
M_{5\%}
<
M_{10\%}
<
M_{20\%}.
\label{eq:rq1_monotonic_not_assumed}
\end{equation}

Quan hệ trong Phương trình~\eqref{eq:rq1_monotonic_not_assumed} chỉ biểu diễn
một xu hướng đơn điệu có thể được kỳ vọng trực giác và \textbf{không} được
đặt làm giả thuyết của nghiên cứu. Thay vào đó, hiệu năng của supervised
baseline tại bốn mức ngân sách nhãn được báo cáo nhằm mô tả mức độ phụ thuộc
của detector vào lượng bounding-box supervision và tạo mốc tham chiếu cho các
phép so sánh với semi-supervised learning.

\paragraph{RQ2 -- Lợi ích tổng thể của semi-supervised learning.}

RQ2 là giả thuyết nghiên cứu chính (\textit{primary hypothesis}) và có hướng
dương. Nghiên cứu giả thuyết rằng, khi các điều kiện thực nghiệm được ghép cặp,
semi-supervised learning khai thác thêm dữ liệu chưa gán nhãn sẽ tạo ra mức cải
thiện dương về bbox mAP@[0.50:0.95] so với supervised baseline.

Với bbox mAP@[0.50:0.95], hiệu ứng ghép cặp tại từng tổ hợp kiến trúc, ngân
sách nhãn và training seed được xác định bởi

\begin{equation}
\Delta_{a,b,s}
=
mAP^{SSL}_{a,b,s}
-
mAP^{SUP}_{a,b,s}.
\label{eq:rq2_paired_map_gain}
\end{equation}

Mức lợi ích tổng hợp trên tám tổ hợp
\(2\) kiến trúc \(\times\) \(4\) ngân sách nhãn đối với seed \(s\) được định
nghĩa là

\begin{equation}
G_s
=
\frac{1}{8}
\sum_{a,b}
\Delta_{a,b,s}.
\label{eq:rq2_overall_gain}
\end{equation}

Giả thuyết thống kê tương ứng là

\begin{equation}
H_{0,\mathrm{RQ2}}:
\mu_G \leq 0,
\label{eq:rq2_h0}
\end{equation}

và

\begin{equation}
H_{1,\mathrm{RQ2}}:
\mu_G > 0.
\label{eq:rq2_h1}
\end{equation}

Do đó, giả thuyết chính của nghiên cứu chỉ có một phát biểu tổng thể, không
được tách thành bốn giả thuyết chính theo ngân sách nhãn hoặc tám giả thuyết
chính theo từng tổ hợp kiến trúc--ngân sách.

\paragraph{RQ3 -- Sự thay đổi của mức lợi ích SSL theo ngân sách nhãn.}

RQ3 là giả thuyết thứ cấp, chính thức và không định hướng. Nghiên cứu kiểm tra
liệu mức lợi ích của semi-supervised learning so với supervised baseline có
thay đổi giữa bốn mức ngân sách nhãn hay không, nhưng không giả định trước
rằng lợi ích này phải tăng hoặc giảm đơn điệu theo lượng dữ liệu có nhãn.

Để cân bằng ảnh hưởng của hai kiến trúc, đối với mỗi ngân sách \(b\) và seed
\(s\), mức lợi ích được xác định bởi

\begin{equation}
B_{b,s}
=
\frac{
\Delta_{R50,b,s}
+
\Delta_{Swin,b,s}
}{2}.
\label{eq:rq3_budget_gain}
\end{equation}

Giả thuyết không của RQ3 phát biểu rằng mức cải thiện trung bình của
semi-supervised learning không khác giữa bốn ngân sách nhãn:

\begin{equation}
H_{0,\mathrm{RQ3}}:
\mu_{B_{1\%}}
=
\mu_{B_{5\%}}
=
\mu_{B_{10\%}}
=
\mu_{B_{20\%}}.
\label{eq:rq3_h0}
\end{equation}

Giả thuyết đối là

\begin{equation}
H_{1,\mathrm{RQ3}}:
\text{ít nhất một } \mu_{B_b}
\text{ khác với các mức còn lại}.
\label{eq:rq3_h1}
\end{equation}

Nghiên cứu không đặt trước giả thuyết tuyến tính hoặc đơn điệu như

\begin{equation}
\Delta_{1\%}
>
\Delta_{5\%}
>
\Delta_{10\%}
>
\Delta_{20\%}.
\label{eq:rq3_monotonic_not_assumed}
\end{equation}

Quan hệ trong Phương trình~\eqref{eq:rq3_monotonic_not_assumed} vì vậy chỉ là
một dạng xu hướng có thể quan sát được sau thực nghiệm, không phải giả thuyết
được xác định trước.

\paragraph{RQ4 -- Mối liên hệ giữa chất lượng pseudo-label và mức cải thiện hiệu năng phát hiện.}

RQ4 là giả thuyết thứ cấp mang tính cơ chế và có hướng dương. Chất lượng
pseudo-label được xác định trước bởi

\begin{equation}
Q_{\mathrm{pseudo}}
=
PL\text{-}mAP@[0.50:0.95],
\label{eq:rq4_qpseudo}
\end{equation}

được tính trên tập validation cố định bằng EMA Teacher tại checkpoint LAST.
Ground truth ẩn của các \textit{unlabeled subsets} không được sử dụng để tính
hoặc hiệu chỉnh \(Q_{\mathrm{pseudo}}\).

Đối với mỗi tổ hợp kiến trúc \(a\), ngân sách nhãn \(b\) và seed \(s\), mức
cải thiện hiệu năng phát hiện của semi-supervised learning so với supervised
baseline trên tập test cuối cùng được xác định bởi

\begin{equation}
Y_{a,b,s}
=
\Delta mAP^{test}_{a,b,s}
=
mAP^{SSL,test}_{a,b,s}
-
mAP^{SUP,test}_{a,b,s}.
\label{eq:rq4_test_map_gain}
\end{equation}

Nghiên cứu giả thuyết rằng chất lượng pseudo-label cao hơn trên tập validation
có mối liên hệ dương với mức cải thiện bbox mAP@[0.50:0.95] của
semi-supervised learning trên tập test cuối cùng.

Với hệ số hồi quy \(\beta_Q\) biểu diễn mối liên hệ giữa
\(Q_{\mathrm{pseudo}}\) trong Phương trình~\eqref{eq:rq4_qpseudo} và
\(\Delta mAP^{test}\) trong Phương trình~\eqref{eq:rq4_test_map_gain},
giả thuyết thống kê được xác định là

\begin{equation}
H_{0,\mathrm{RQ4}}:
\beta_Q \leq 0,
\label{eq:rq4_h0}
\end{equation}

và

\begin{equation}
H_{1,\mathrm{RQ4}}:
\beta_Q > 0.
\label{eq:rq4_h1}
\end{equation}

Giả thuyết này chỉ mô tả \textit{association} giữa chất lượng pseudo-label
được đo trên validation và mức cải thiện hiệu năng được đo trên test, không
được diễn giải như một quan hệ nhân quả. Chất lượng pseudo-label không phải là
một can thiệp ngẫu nhiên độc lập trong thiết kế thực nghiệm. Tập test chỉ được
sử dụng để xác định kết quả hiệu năng cuối cùng sau khi toàn bộ giao thức nghiên
cứu, mô hình và phương pháp phân tích đã được cố định trước.

\paragraph{RQ5 -- Lợi ích đối với các lớp hiếm.}

RQ5 có vai trò khám phá (\textit{exploratory}) và không gắn với một phép kiểm
định giả thuyết xác nhận bắt buộc. Hai lớp được xác định trước là lớp hiếm từ
phân tích dữ liệu huấn luyện gồm \textit{Atelectasis} và
\textit{Pneumothorax}. Đối với mỗi lớp \(c\), mức thay đổi AP được xác định bởi

\begin{equation}
\Delta AP_{c,a,b,s}
=
AP^{SSL}_{c,a,b,s}
-
AP^{SUP}_{c,a,b,s}.
\label{eq:rq5_class_ap_gain}
\end{equation}

Các hiệu ứng trong Phương trình~\eqref{eq:rq5_class_ap_gain} được tổng hợp và
báo cáo bằng ước lượng, độ lệch chuẩn và khoảng tin cậy nhằm mô tả mức lợi ích
hoặc suy giảm của SSL đối với từng lớp. Nghiên cứu không tạo thêm một chỉ số
``rare-mAP'' và không chuyển các kết quả khám phá này thành giả thuyết xác nhận
sau khi quan sát dữ liệu.

\paragraph{RQ6 -- False positive trên ảnh \textit{No Finding}.}

RQ6 là giả thuyết thứ cấp, chính thức và không định hướng. Nghiên cứu kiểm tra
liệu semi-supervised learning có làm thay đổi số false positive trung bình trên
các ảnh \textit{No Finding} so với supervised baseline hay không, nhưng không
giả định trước rằng SSL nhất thiết phải làm giảm hoặc làm tăng false positive.

Đối với mỗi kiến trúc \(a\), ngân sách \(b\) và seed \(s\), hiệu ứng được định
nghĩa là

\begin{equation}
\Delta F_{a,b,s}
=
FP^{SSL}_{neg,a,b,s}
-
FP^{SUP}_{neg,a,b,s},
\label{eq:rq6_negative_fp_effect}
\end{equation}

trong đó \(FP_{neg}\) là số false positive trung bình trên mỗi ảnh
\textit{No Finding}. Hiệu ứng tổng hợp theo seed được xác định bởi

\begin{equation}
H_s
=
\frac{1}{8}
\sum_{a,b}
\Delta F_{a,b,s}.
\label{eq:rq6_overall_fp_effect}
\end{equation}

Giả thuyết thống kê tương ứng là

\begin{equation}
H_{0,\mathrm{RQ6}}:
\mu_H = 0,
\label{eq:rq6_h0}
\end{equation}

và

\begin{equation}
H_{1,\mathrm{RQ6}}:
\mu_H \neq 0.
\label{eq:rq6_h1}
\end{equation}

Theo quy ước trong Phương trình~\eqref{eq:rq6_negative_fp_effect}, giá trị âm
biểu thị SSL tạo ít false positive hơn supervised baseline, trong khi giá trị
dương biểu thị số false positive trung bình trên ảnh \textit{No Finding} cao
hơn.

\paragraph{RQ7 -- Sự phụ thuộc của mức lợi ích SSL vào kiến trúc.}

RQ7 là giả thuyết thứ cấp, chính thức và không định hướng. Mục tiêu không phải
kiểm định liệu Swin-T có hiệu năng tuyệt đối cao hơn ResNet-50 hay không, mà
kiểm tra liệu mức cải thiện do semi-supervised learning mang lại so với
supervised baseline có khác giữa hai kiến trúc dưới các giao thức huấn luyện
đặc thù theo kiến trúc đã được xác định trước.

Đối với kiến trúc \(a\) và seed \(s\), mức cải thiện trung bình của
semi-supervised learning qua bốn ngân sách nhãn được xác định bởi

\begin{equation}
A_{a,s}
=
\frac{1}{4}
\sum_b
\Delta_{a,b,s}.
\label{eq:rq7_architecture_gain}
\end{equation}

Hiệu ứng khác biệt giữa hai kiến trúc được định nghĩa là

\begin{equation}
D_s
=
A_{Swin,s}
-
A_{R50,s}.
\label{eq:rq7_architecture_difference}
\end{equation}

Giả thuyết thống kê tương ứng là

\begin{equation}
H_{0,\mathrm{RQ7}}:
\mu_D = 0,
\label{eq:rq7_h0}
\end{equation}

và

\begin{equation}
H_{1,\mathrm{RQ7}}:
\mu_D \neq 0.
\label{eq:rq7_h1}
\end{equation}

Do hai kiến trúc sử dụng các giao thức tối ưu hóa đặc thù theo kiến trúc,
kết quả RQ7 chỉ được diễn giải là sự khác biệt về mức lợi ích của
semi-supervised learning giữa hai kiến trúc dưới các giao thức huấn luyện đã
được xác định trước. Kết quả này không được xem là bằng chứng cho thấy Swin-T
có ưu thế nhân quả tổng quát so với ResNet-50.

Tóm lại, hệ giả thuyết của nghiên cứu được tổ chức theo một thứ bậc xác định
trước: RQ2 là giả thuyết chính; RQ3, RQ4, RQ6 và RQ7 là các giả thuyết thứ cấp
chính thức; RQ1 có vai trò mô tả; và RQ5 có vai trò khám phá. Cách tổ chức này
giúp phân biệt rõ giữa mục tiêu xác nhận, mục tiêu phân tích cơ chế và mục tiêu
khám phá, đồng thời hạn chế việc hình thành giả thuyết hậu nghiệm sau khi quan
sát kết quả thực nghiệm. Các phép kiểm định cụ thể, mức ý nghĩa, khoảng tin
cậy, kiểm soát đa kiểm định và phân tích độ nhạy được trình bày tại phần
\textit{Statistical Analysis}.


\subsubsection{Biến nghiên cứu và nguyên tắc kiểm soát}
\label{subsubsec:research_variables}

Thiết kế thực nghiệm được xây dựng nhằm phân biệt rõ các yếu tố được thay đổi
có chủ đích, các đại lượng được đo lường và các yếu tố cần được giữ ổn định
giữa các điều kiện so sánh. Theo đó, nghiên cứu xác định ba biến độc lập chính,
các biến phụ thuộc dùng để đánh giá hiệu năng phát hiện, một biến dự báo mang
tính cơ chế liên quan đến chất lượng pseudo-label, cùng các yếu tố kiểm soát và
yếu tố ghép cặp được xác định trước.

\paragraph{Biến độc lập.}

Biến độc lập thứ nhất là \textit{training paradigm}, ký hiệu \(X_1\), gồm hai
mức supervised learning và semi-supervised learning:

\begin{equation}
X_1
\in
\left\{
SUP,\,
SSL
\right\}.
\label{eq:research_variable_method}
\end{equation}

Đây là yếu tố thực nghiệm chính của nghiên cứu. Trong mỗi phép so sánh ghép
cặp, supervised baseline chỉ sử dụng tập dữ liệu có nhãn, trong khi thiết lập
semi-supervised sử dụng cùng tập dữ liệu có nhãn và khai thác thêm tập dữ liệu
chưa gán nhãn thông qua cơ chế teacher--student pseudo-labeling.

Biến độc lập thứ hai là ngân sách dữ liệu có nhãn
(\textit{labeled-data budget}), ký hiệu \(X_2\), được xác định theo tỷ lệ
số ảnh của tập huấn luyện cố định được đưa vào \textit{labeled subset}
và gồm bốn mức:

\begin{equation}
X_2
\in
\left\{
1\%,\,
5\%,\,
10\%,\,
20\%
\right\}.
\label{eq:research_variable_budget}
\end{equation}

Biến này được xem là một yếu tố phân loại có thứ tự
(\textit{ordered categorical factor}). Tuy nhiên, trong phân tích thống kê
chính, nghiên cứu không giả định trước quan hệ tuyến tính hoặc đơn điệu giữa
ngân sách nhãn và hiệu quả của semi-supervised learning.

Biến độc lập thứ ba là kiến trúc mô hình (\textit{architecture}), ký hiệu
\(X_3\), gồm hai mức:

\begin{equation}
X_3
\in
\left\{
R50,\,
Swin\text{-}T
\right\}.
\label{eq:research_variable_architecture}
\end{equation}

Trong đó, \(R50\) biểu diễn Faster R-CNN sử dụng backbone ResNet-50 kết hợp
FPN, còn \(Swin\text{-}T\) biểu diễn Faster R-CNN sử dụng backbone Swin-T kết
hợp FPN. ResNet-50, Faster R-CNN và FPN lần lượt được mô tả chi tiết tại phần
Supervised Baseline \cite{ren2016faster,he2016deep,lin2017feature}. Swin-T được
sử dụng như kiến trúc Vision Transformer phân cấp trong nhánh kiến trúc thứ hai.

Ba biến độc lập trên tạo thành ma trận thực nghiệm chính:

\begin{equation}
2
\times
4
\times
2
=
16
\text{ experimental cells}.
\label{eq:main_experimental_cells}
\end{equation}

Trong đó, hai mức của \(X_1\) tương ứng với SUP và SSL, bốn mức của \(X_2\)
tương ứng với bốn ngân sách nhãn, và hai mức của \(X_3\) tương ứng với hai
kiến trúc.

\paragraph{Biến phụ thuộc.}

Biến phụ thuộc chính của nghiên cứu là bbox mAP@[0.50:0.95], được sử dụng để
đánh giá hiệu năng phát hiện tổng thể và làm chỉ số trung tâm cho giả thuyết
chính RQ2.

Ngoài biến phụ thuộc chính, nghiên cứu sử dụng các biến phụ thuộc bổ sung gồm:

\begin{itemize}
    \item AP50;
    \item AP75;
    \item AP theo từng lớp;
    \item recall;
    \item số false positive trung bình trên mỗi ảnh;
    \item số false positive trung bình trên mỗi ảnh \textit{No Finding}.
\end{itemize}

Các chỉ số này không có cùng vai trò suy luận. Bbox mAP@[0.50:0.95] là chỉ số
chính ở cấp toàn nghiên cứu, trong khi AP50, AP75, recall và AP theo lớp đóng
vai trò bổ sung. Số false positive trung bình trên ảnh \textit{No Finding} là
biến kết quả chính trong phân tích RQ6.

Đối với các phép so sánh SUP--SSL, biến kết quả thường được biểu diễn dưới dạng
hiệu ứng ghép cặp:

\begin{equation}
\Delta M_{a,b,s}
=
M^{SSL}_{a,b,s}
-
M^{SUP}_{a,b,s},
\label{eq:research_variable_paired_effect}
\end{equation}

trong đó \(M\) là chỉ số đánh giá được xem xét, \(a\) là kiến trúc, \(b\) là
ngân sách nhãn và \(s\) là training seed. Cách biểu diễn trong
Phương trình~\eqref{eq:research_variable_paired_effect} cho phép phân tích trực
tiếp phần thay đổi gắn với semi-supervised learning trong từng điều kiện matched.

\paragraph{Biến dự báo mang tính cơ chế.}

Chất lượng pseudo-label không được xem là một biến độc lập được thao tác trực
tiếp trong ma trận thực nghiệm chính. Thay vào đó, nó được sử dụng như một biến
dự báo mang tính cơ chế trong RQ4.

Biến này được ký hiệu là \(Q_{\mathrm{pseudo}}\) và được định nghĩa bởi:

\begin{equation}
Q_{\mathrm{pseudo}}
=
PL\text{-}mAP@[0.50:0.95].
\label{eq:research_variable_qpseudo}
\end{equation}

\(Q_{\mathrm{pseudo}}\) được tính trên tập validation cố định bằng EMA Teacher
tại checkpoint LAST. Ground truth ẩn của các ảnh thuộc
\textit{unlabeled subsets} không được sử dụng để tính, hiệu chỉnh hoặc lựa chọn
\(Q_{\mathrm{pseudo}}\).

Việc phân loại \(Q_{\mathrm{pseudo}}\) như một biến dự báo thay vì một biến độc
lập là cần thiết vì chất lượng pseudo-label không được randomize hoặc điều khiển
như một treatment độc lập. Do đó, các kết quả liên quan đến RQ4 chỉ được diễn
giải dưới dạng mối liên hệ thống kê, không phải quan hệ nhân quả.

\paragraph{Yếu tố ghép cặp và chặn.}

\textit{Training seed} được sử dụng như một yếu tố ghép cặp và chặn
(\textit{blocking/pairing factor}), không phải một biến độc lập chính của
nghiên cứu.

Đối với mỗi tổ hợp kiến trúc và ngân sách nhãn, kết quả SUP và SSL được ghép
cặp theo cùng training seed. Khóa ghép cặp được xác định bởi:

\begin{equation}
P
=
\left(
a,\,
b,\,
s
\right),
\label{eq:research_pairing_key}
\end{equation}

trong đó \(a\) là kiến trúc, \(b\) là ngân sách nhãn và \(s\) là training
seed.

Trong triển khai thực nghiệm, khóa này được ánh xạ tường minh tới
\texttt{architecture}, \texttt{budget}, \texttt{seed\_index} và
\texttt{training\_seed}. Việc ghép cặp không được thực hiện dựa trên thứ tự
dòng kết quả, hiệu năng mô hình hoặc lựa chọn seed có kết quả tốt hơn.

Nghiên cứu sử dụng cùng một danh sách có thứ tự gồm 10 training seeds cho SUP
và SSL. Seed không được tìm kiếm hoặc thay thế dựa trên hiệu năng. Nếu một lần
huấn luyện gặp lỗi kỹ thuật hợp lệ và phải thực hiện lại, lần chạy thay thế phải
sử dụng lại đúng training seed ban đầu.

\paragraph{Phân biệt partition seed và training seed.}

Seed dùng để xây dựng phân hoạch dữ liệu và seed dùng cho quá trình huấn luyện
được xem là hai khái niệm khác nhau.

Seed xây dựng phân hoạch dữ liệu được cố định tại:

\begin{equation}
s_{\mathrm{partition}}
=
42.
\label{eq:partition_seed}
\end{equation}

Trong khi đó, training seeds được dùng để biểu diễn biến thiên ngẫu nhiên của
quá trình huấn luyện mô hình. Vì vậy:

\begin{equation}
s_{\mathrm{partition}}
\neq
s_{\mathrm{training}}.
\label{eq:partition_training_seed_distinction}
\end{equation}

Việc thay đổi training seed không làm thay đổi thành phần của tập train,
validation, test hoặc các \textit{labeled/unlabeled subsets} đã được cố định
trước.

\paragraph{Các yếu tố được kiểm soát trong so sánh SUP--SSL.}

Để một hiệu ứng ghép cặp có thể được diễn giải là sự khác biệt giữa supervised
và semi-supervised learning trong phạm vi thiết kế đã xác định, các yếu tố sau
được giữ nhất quán trong từng cặp so sánh:

\begin{itemize}
    \item phạm vi dữ liệu nghiên cứu;
    \item phân chia train/validation/test cố định;
    \item cùng \textit{labeled subset} \(L_b\);
    \item cùng kiến trúc;
    \item cùng training seed;
    \item cùng biểu diễn ảnh đầu vào và quy trình tiền xử lý;
    \item cùng họ detector Faster R-CNN;
    \item cùng giao thức đánh giá;
    \item cùng tổng ngân sách cập nhật tham số;
    \item cùng mức tiếp xúc với dữ liệu có nhãn.
\end{itemize}

Trong từng kiến trúc, optimizer và các thông số huấn luyện đặc thù theo kiến
trúc được giữ giống nhau giữa SUP và SSL. Tuy nhiên, hai kiến trúc không bắt
buộc sử dụng cùng optimizer. ResNet-50 sử dụng giao thức tối ưu hóa được xác
định riêng cho nhánh R50, trong khi Swin-T sử dụng giao thức tối ưu hóa đặc thù
cho kiến trúc Transformer. Do đó, nghiên cứu kiểm soát SUP--SSL
\textbf{within architecture}, thay vì áp đặt một optimizer duy nhất cho cả
hai kiến trúc.

Khác biệt có chủ đích trong mỗi cặp SUP--SSL là khả năng sử dụng dữ liệu chưa
gán nhãn. Với ngân sách nhãn \(b\), supervised baseline sử dụng:

\begin{equation}
\mathcal{D}^{SUP}_{b}
=
L_b,
\label{eq:sup_training_data}
\end{equation}

trong khi semi-supervised learning sử dụng:

\begin{equation}
\mathcal{D}^{SSL}_{b}
=
L_b
\cup
U_b,
\label{eq:ssl_training_data}
\end{equation}

trong đó chỉ \(L_b\) cung cấp ground-truth bounding box, còn \(U_b\) được khai
thác thông qua pseudo-label do Teacher sinh ra.

Hai tập này thỏa mãn:

\begin{equation}
L_b
\cap
U_b
=
\varnothing,
\label{eq:labeled_unlabeled_disjoint}
\end{equation}

và

\begin{equation}
L_b
\cup
U_b
=
T,
\label{eq:labeled_unlabeled_union}
\end{equation}

với \(T\) là tập huấn luyện cố định.

\paragraph{Kiểm soát việc sử dụng validation và test.}

Tập validation được sử dụng cho các mục đích đã được xác định trước như lựa
chọn checkpoint, tính các chỉ số validation cần thiết, đánh giá
\(Q_{\mathrm{pseudo}}\) và thực hiện các phân tích ablation đã định nghĩa trước.

Ngược lại, tập test được giữ ngoài quá trình phát triển mô hình và không được
sử dụng để:

\begin{itemize}
    \item lựa chọn kiến trúc;
    \item lựa chọn checkpoint;
    \item điều chỉnh confidence threshold;
    \item lựa chọn chiến lược augmentation;
    \item lựa chọn pseudo-label filtering;
    \item lựa chọn metric;
    \item lựa chọn mô hình thống kê;
    \item lựa chọn cấu hình ablation.
\end{itemize}

Tập test chỉ được sử dụng ở giai đoạn đánh giá cuối cùng sau khi giao thức mô
hình, metric và phương pháp phân tích đã được cố định.

\paragraph{Kiểm soát ground truth của dữ liệu chưa gán nhãn.}

Mặc dù các ảnh thuộc \(U_b\) có nguồn gốc từ bộ dữ liệu đã được chú giải,
ground truth của chúng được xem là ẩn đối với quá trình semi-supervised
learning. Ground truth này không được sử dụng để:

\begin{itemize}
    \item huấn luyện mô hình;
    \item lọc pseudo-label;
    \item điều chỉnh confidence threshold;
    \item lựa chọn EMA;
    \item lựa chọn augmentation;
    \item lựa chọn checkpoint;
    \item tính chất lượng pseudo-label.
\end{itemize}

Quy tắc này bảo đảm rằng \(U_b\) thực sự đóng vai trò dữ liệu chưa gán nhãn
trong giao thức thực nghiệm, thay vì trở thành nguồn thông tin giám sát gián
tiếp.

\paragraph{Vai trò của các lớp hiếm và ảnh \textit{No Finding}.}

Trạng thái lớp hiếm không phải là một biến độc lập của ma trận thực nghiệm mà
là một đặc trưng dữ liệu được xác định trước để phục vụ phân tích khám phá RQ5.
Theo kết quả chẩn đoán dữ liệu trên tập huấn luyện, hai lớp được xác định là
lớp hiếm gồm \textit{Atelectasis} và \textit{Pneumothorax}.

Tương tự, \textit{No Finding} không được xem là một detection class hoặc một
biến độc lập. Các ảnh \textit{No Finding} được giữ như một nhóm ảnh âm tính
zero-GT để đánh giá riêng false positive trong RQ6.

Tóm lại, cấu trúc biến của nghiên cứu gồm ba biến độc lập chính là training
paradigm, labeled-data budget và architecture; các biến phụ thuộc phản ánh hiệu
năng phát hiện và false-positive behavior; \(Q_{\mathrm{pseudo}}\) đóng vai trò
biến dự báo mang tính cơ chế; training seed là yếu tố ghép cặp/chặn; còn các
điều kiện dữ liệu, kiến trúc, biểu diễn đầu vào, ngân sách cập nhật và giao thức
đánh giá được kiểm soát theo từng phép so sánh. Cấu trúc này tạo cơ sở trực
tiếp cho thiết kế các điều kiện thực nghiệm và các phép phân tích thống kê được
trình bày ở các tiểu mục tiếp theo.

\subsubsection{Thiết kế các điều kiện thực nghiệm}
\label{subsubsec:experimental_conditions}

Các điều kiện thực nghiệm được xây dựng từ ba yếu tố đã xác định trước gồm
phương thức học, ngân sách dữ liệu có nhãn và kiến trúc mô hình. Phương thức
học gồm supervised learning (SUP) và semi-supervised learning (SSL); ngân sách
nhãn gồm \(1\%\), \(5\%\), \(10\%\) và \(20\%\) số ảnh của tập huấn luyện cố định;
kiến trúc gồm Faster R-CNN với ResNet-50 và FPN (R50) và Faster R-CNN với
Swin-T và FPN (Swin-T).

Ma trận thực nghiệm chính trong chế độ hạn chế nhãn được xác định bởi

\begin{equation}
N_{\mathrm{cells}}
=
2_{\mathrm{method}}
\times
4_{\mathrm{budget}}
\times
2_{\mathrm{architecture}}
=
16.
\label{eq:main_experimental_matrix}
\end{equation}

Mỗi điều kiện được thực hiện với cùng danh sách có thứ tự gồm 10
\texttt{training\_seed} đã xác định trước. Do đó, số lần huấn luyện trong ma
trận thực nghiệm chính là

\begin{equation}
N_{\mathrm{low\text{-}label}}
=
16
\times
10
=
160.
\label{eq:low_label_training_runs}
\end{equation}

Bên cạnh các điều kiện hạn chế nhãn, supervised baseline được huấn luyện thêm
với \(100\%\) tập huấn luyện có nhãn cho cả hai kiến trúc. Mỗi kiến trúc tiếp
tục sử dụng cùng 10 \texttt{training\_seed}, tạo ra

\begin{equation}
N_{\mathrm{100\%\text{-}SUP}}
=
2
\times
10
=
20
\label{eq:full_supervised_reference_runs}
\end{equation}

lần huấn luyện tham chiếu. Tổng số lần huấn luyện thuộc ma trận chính và các
điều kiện supervised tham chiếu vì vậy là

\begin{equation}
N_{\mathrm{main+reference}}
=
160
+
20
=
180.
\label{eq:main_reference_total_runs}
\end{equation}

Các điều kiện \(100\%\)-SUP chỉ được sử dụng làm mốc tham chiếu cho hiệu năng
supervised khi toàn bộ tập huấn luyện có nhãn được sử dụng. Các điều kiện này
không thuộc các phép so sánh ghép cặp chính giữa SUP và SSL tại các ngân sách
\(1\%\), \(5\%\), \(10\%\) và \(20\%\).

Tại mỗi ngân sách \(b\), supervised baseline và phương pháp semi-supervised sử
dụng cùng tập có nhãn \(L_b\) đã được cố định. Supervised baseline chỉ sử dụng
\(L_b\) để huấn luyện, trong khi phương pháp semi-supervised sử dụng cùng
\(L_b\) kết hợp với tập chưa gán nhãn tương ứng \(U_b\). Thành phần của
\(L_b\), \(U_b\), tập validation và tập test không thay đổi giữa các lần huấn
luyện có cùng ngân sách. Cách xây dựng quan hệ lồng nhau giữa các
\textit{labeled subsets} và nguyên tắc cô lập validation/test đã được trình
bày ở các tiểu mục dữ liệu trước đó và được giữ nguyên trong toàn bộ ma trận
thực nghiệm.

\paragraph{Kiểm soát ngân sách huấn luyện trong các phép so sánh SUP--SSL.}

Trong mỗi cặp so sánh có cùng kiến trúc, ngân sách nhãn và
\texttt{training\_seed}, SUP và SSL được kiểm soát theo cùng ngân sách cập nhật
optimizer và cùng mức tiếp xúc với dữ liệu có nhãn. Ngân sách cập nhật optimizer
được xác định trước theo ngân sách nhãn như sau:

\begin{equation}
U_{\mathrm{opt}}(b)
=
\begin{cases}
1032, & b=1\%,\\
2064, & b=5\%,\\
2064, & b=10\%,\\
2064, & b=20\%.
\end{cases}
\label{eq:optimizer_update_budgets_low_label}
\end{equation}

Đối với điều kiện supervised tham chiếu sử dụng \(100\%\) dữ liệu có nhãn,
ngân sách được xác định là

\begin{equation}
U_{\mathrm{opt}}(100\%)
=
10284.
\label{eq:optimizer_update_budget_full_label}
\end{equation}

Các ngân sách trong
Phương trình~\eqref{eq:optimizer_update_budgets_low_label} và
Phương trình~\eqref{eq:optimizer_update_budget_full_label} được định nghĩa theo
số lần cập nhật optimizer, không phải theo số epoch. Trong mỗi cặp SUP--SSL,
SSL được phép khai thác thêm \(U_b\), nhưng không được tăng tổng ngân sách cập
nhật hoặc số lần tiếp xúc với tín hiệu giám sát từ \(L_b\) so với supervised
baseline tương ứng.

Các tham số tối ưu hóa được xác định riêng cho từng kiến trúc. Vì vậy, nguyên
tắc matching được áp dụng giữa SUP và SSL \textbf{trong cùng kiến trúc}, thay
vì yêu cầu R50 và Swin-T sử dụng một cấu hình tối ưu hóa giống nhau. Với kiến
trúc \(a\), nguyên tắc này được biểu diễn bởi

\begin{equation}
Recipe^{SUP}_{a}
=
Recipe^{SSL}_{a},
\qquad
a\in\{R50,Swin\text{-}T\}.
\label{eq:within_architecture_recipe_matching}
\end{equation}

Do đó, sự khác biệt giữa R50 và Swin-T được phân tích trong bối cảnh các giao
thức huấn luyện đặc thù theo kiến trúc đã được xác định trước, không được diễn
giải như một phép kiểm định thuần túy về ưu thế nhân quả của backbone.

\paragraph{Giao thức lặp lại và ghép cặp theo seed.}

Nghiên cứu xác định trước một danh sách có thứ tự gồm 10
\texttt{training\_seed} khác nhau trước khi thực hiện các thí nghiệm huấn
luyện chính thức. Cùng một danh sách seed được sử dụng cho supervised
learning và semi-supervised learning trong toàn bộ ma trận thực nghiệm.
Số lượng, giá trị và thứ tự của các seed được giữ cố định trong suốt quá
trình thực nghiệm và không được thay đổi dựa trên kết quả huấn luyện,
validation hoặc test.

Mục đích của việc sử dụng 10 \texttt{training\_seed} là mô tả biến thiên
giữa các lần huấn luyện, giảm sự phụ thuộc của kết luận vào một lần chạy
duy nhất và hỗ trợ phép so sánh ghép cặp
(\textit{paired-seed comparison}) giữa supervised learning và
semi-supervised learning. Số lượng 10 seed được xác định trước trong giao
thức nghiên cứu; nghiên cứu không xem đây là số lượng tối ưu về mặt thống
kê và không lựa chọn số seed dựa trên kết quả quan sát sau huấn luyện.

Tại mỗi tổ hợp kiến trúc \(a\), ngân sách nhãn \(b\) và chỉ số seed \(k\),
supervised baseline và phương pháp semi-supervised được thực hiện với cùng
giá trị training seed. Khóa ghép cặp được xác định bởi

\begin{equation}
P
=
\left(
a,\,
b,\,
k,\,
s_k
\right),
\label{eq:experimental_pairing_key}
\end{equation}

trong đó \(a\) là kiến trúc, \(b\) là ngân sách nhãn, \(k\) là chỉ số của
seed trong danh sách đã xác định trước và \(s_k\) là giá trị training seed
tương ứng.

Cụ thể, lần huấn luyện supervised tại chỉ số seed \(k\) được so sánh với
lần huấn luyện semi-supervised có cùng chỉ số \(k\). Hai lần chạy sử dụng
cùng training seed, cùng phân hoạch dữ liệu cố định, cùng kiến trúc và cùng
tập có nhãn \(L_b\); tập \(U_b\) chỉ được phương pháp semi-supervised khai
thác và cũng được giữ cố định tại ngân sách tương ứng. Việc sử dụng cùng
seed trong mỗi cặp nhằm giảm ảnh hưởng của biến thiên ngẫu nhiên lên phép
so sánh giữa hai phương thức học; điều này không hàm ý rằng hai phương thức
phải tạo ra kết quả giống nhau.

Một lần huấn luyện chỉ được thực hiện lại khi xảy ra lỗi kỹ thuật có thể
xác định và ghi nhận, chẳng hạn tiến trình bị gián đoạn, hết bộ nhớ, lỗi
phần cứng, lỗi hạ tầng hoặc lỗi I/O. Lần chạy lại phải sử dụng cùng training
seed với lần chạy ban đầu. Lỗi kỹ thuật và lần thực hiện lại được lưu vết để
bảo đảm khả năng truy xuất thí nghiệm.

Ngược lại, một lần chạy hoàn thành hợp lệ nhưng có hiệu năng thấp hoặc hội
tụ kém không được xem là lỗi kỹ thuật và không phải là căn cứ để thay đổi
seed, loại bỏ kết quả hoặc thực hiện lại thí nghiệm. Quy tắc này nhằm tránh
việc lựa chọn kết quả thuận lợi sau khi đã quan sát hiệu năng mô hình.
\subsubsection{Thiết kế thí nghiệm phân tích thành phần}
\label{subsubsec:ablation_design}

Bên cạnh ma trận thực nghiệm chính, nghiên cứu thực hiện thí nghiệm phân tích
thành phần (\textit{ablation study}) nhằm làm rõ ảnh hưởng của một số thành
phần đã được xác định trước trong cấu hình semi-supervised learning chính.
Mục tiêu của ablation không phải tìm kiếm cấu hình có hiệu năng cao nhất,
lựa chọn lại mô hình hoặc tối ưu hậu nghiệm giao thức SSL, mà là phân tích cơ
chế để trả lời câu hỏi: khi một thành phần cụ thể của cấu hình chính thay đổi,
hành vi pseudo-label và hiệu năng phát hiện thay đổi như thế nào.

Vì vậy, ablation được xác định là một phân tích
\textbf{thứ cấp, mang tính cơ chế và theo nguyên tắc one-factor-at-a-time
(OFAT)}. Cấu hình SSL Main đã được xác định trước vẫn là mốc tham chiếu cố
định; kết quả ablation, kể cả khi một biến thể cho kết quả tốt hơn Main trên
validation, không được sử dụng để sửa ngược threshold, augmentation hoặc
filtering strategy của giao thức chính. Cách tổ chức này giúp hạn chế
hyperparameter search hậu nghiệm, selective reporting và sự gia tăng không
cần thiết của số phép so sánh.

\paragraph{Nguyên tắc one-factor-at-a-time.}

Trong mỗi nhóm ablation, chỉ một yếu tố thuộc nhóm đang khảo sát được thay
đổi, trong khi tất cả các yếu tố khoa học còn lại được giữ giống cấu hình
Main. Nguyên tắc này được biểu diễn bởi

\begin{equation}
\mathcal{C}_{v}
=
\mathcal{C}_{\mathrm{Main}}
\setminus
\{f_{\mathrm{Main}}\}
\cup
\{f_v\},
\label{eq:ablation_ofat}
\end{equation}

trong đó \(\mathcal{C}_{\mathrm{Main}}\) là cấu hình SSL chính,
\(f_{\mathrm{Main}}\) là giá trị của thành phần đang khảo sát trong Main và
\(f_v\) là giá trị tương ứng của biến thể ablation \(v\).

Theo nguyên tắc này, khi khảo sát confidence threshold thì augmentation,
regression uncertainty, EMA, optimizer, lịch học và các điều kiện khác được
giữ nguyên; khi khảo sát augmentation thì các threshold và filtering rule
được giữ nguyên; và khi khảo sát regression filtering thì confidence
threshold, augmentation và các thành phần khác tiếp tục được giữ cố định.

Việc lựa chọn OFAT thay vì một thiết kế full factorial nhằm giữ khả năng diễn
giải ảnh hưởng của từng thành phần và kiểm soát quy mô thực nghiệm. Nếu kết
hợp đầy đủ 5 mức confidence, 4 cấu hình augmentation và 3 filtering strategy,
thiết kế sẽ tạo ra

\begin{equation}
5
\times
4
\times
3
=
60
\label{eq:ablation_full_factorial_explosion}
\end{equation}

tổ hợp trước khi nhân với số training seed. Điều này vừa làm tăng đáng kể chi
phí tính toán, vừa khiến mục tiêu phân tích cơ chế của từng thành phần bị trộn
lẫn với các tương tác giữa nhiều yếu tố.

\paragraph{Điều kiện tham chiếu của ablation.}

Toàn bộ ablation được thực hiện trên một mốc tham chiếu duy nhất gồm Faster
R-CNN với ResNet-50 và FPN, ngân sách nhãn \(10\%\), cùng 10 training seed đã
được xác định trước và cấu hình SSL Main.

Phạm vi được biểu diễn bởi

\begin{equation}
\mathcal{A}_{\mathrm{abl}}
=
\left\{
R50,\,
10\%,\,
10\ \mathrm{training\ seeds}
\right\}.
\label{eq:ablation_scope}
\end{equation}

ResNet-50 được chọn làm kiến trúc anchor vì đây là nhánh gần với đường triển
khai tham chiếu Faster R-CNN của SoftTeacher \cite{ren2016faster,xu2021end},
đồng thời có chi phí tính toán phù hợp cho phân tích lặp nhiều seed. Quan
trọng hơn, mục tiêu của ablation không phải khảo sát tương tác giữa kiến trúc
và từng thành phần SSL; ảnh hưởng của kiến trúc đối với mức lợi ích của
semi-supervised learning đã được đặt thành một câu hỏi nghiên cứu riêng trong
RQ7. Giới hạn ablation ở R50 vì vậy giúp giữ phân tích tập trung vào các thành
phần SSL thay vì lặp lại phép phân tích kiến trúc.

Ngân sách \(10\%\) được chọn vì vẫn thuộc chế độ hạn chế nhãn nhưng ít cực
đoan hơn mức \(1\%\). Ở mức này, tập có nhãn cung cấp mức hỗ trợ ổn định hơn
cho việc phân tích cơ chế, qua đó giảm nguy cơ kết luận về một thành phần SSL
chủ yếu phản ánh sự bất ổn do số lượng dữ liệu có nhãn quá thấp. Vì vậy,
\(10\%\) được sử dụng như một điều kiện trung gian phù hợp giữa tính
\textit{low-label} và độ ổn định cần thiết cho component analysis.

Việc sử dụng một anchor duy nhất còn giúp giảm confounding, cho phép so sánh
ghép cặp theo seed, tiết kiệm tài nguyên tính toán và giữ diễn giải ở mức
ảnh hưởng của từng thành phần cụ thể.

Nếu 10 lần huấn luyện Main tại R50--\(10\%\) trong ma trận thực nghiệm chính
đã được thực hiện hợp lệ, các kết quả này được tái sử dụng làm điều kiện Main
của ablation thay vì huấn luyện lại một mốc tham chiếu mới.

\paragraph{Ablation ngưỡng confidence của pseudo-label.}

Nhóm thứ nhất khảo sát ảnh hưởng của ngưỡng confidence dùng để chấp nhận
pseudo-label cho nhánh classification RCNN. Confidence filtering là một
thành phần trực tiếp của cơ chế pseudo-labeling theo SoftTeacher
\cite{xu2021end}, đồng thời liên quan đến nguy cơ pseudo-label noise đã được
trình bày ở phần trước. Vì vậy, đây là một thành phần phù hợp để kiểm tra liệu
mức độ nghiêm ngặt của việc chấp nhận pseudo-label có liên hệ với thay đổi
chất lượng pseudo-label và hiệu năng phát hiện hay không.

Các mức được xác định trước là

\begin{equation}
\tau_{\mathrm{cls}}
\in
\left\{
0.50,\,
0.60,\,
0.70,\,
0.80,\,
0.90
\right\}.
\label{eq:ablation_confidence_thresholds}
\end{equation}

Các biến thể được ký hiệu:

\begin{itemize}
    \item C0: \(\tau_{\mathrm{cls}}=0.50\);
    \item C1: \(\tau_{\mathrm{cls}}=0.60\);
    \item C2: \(\tau_{\mathrm{cls}}=0.70\);
    \item C3: \(\tau_{\mathrm{cls}}=0.80\);
    \item C4: \(\tau_{\mathrm{cls}}=0.90\), tương ứng cấu hình Main.
\end{itemize}

Trong nhóm này, chỉ
\(\tau_{\mathrm{cls}}\)
được thay đổi; strong augmentation, regression reliability filtering, EMA,
tỷ lệ \(L:U\), trọng số loss, optimizer, lịch học, ngân sách cập nhật và các
thành phần còn lại được giữ giống Main.

\paragraph{Ablation strong photometric augmentation.}

Nhóm thứ hai khảo sát vai trò của hai phép biến đổi quang học được sử dụng
trong strong view của dữ liệu chưa gán nhãn. Strong augmentation tạo sự khác
biệt giữa đầu vào của Teacher và Student và là một thành phần trực tiếp của
khung teacher--student được triển khai theo SoftTeacher \cite{xu2021end}.
Trong nghiên cứu này, brightness và contrast được lựa chọn vì chúng thay đổi
đặc tính cường độ của ảnh mà không thay đổi geometry của ảnh hoặc bounding
box, phù hợp với nguyên tắc bảo toàn cấu trúc không gian của dữ liệu X-quang.

Bốn cấu hình được xác định:

\begin{itemize}
    \item A0: không brightness và không contrast;
    \item A1: chỉ brightness;
    \item A2: chỉ contrast;
    \item A3: brightness kết hợp contrast, tương ứng cấu hình Main.
\end{itemize}

Khi được sử dụng, hai phép biến đổi giữ cùng phạm vi đã định nghĩa trong cấu
hình Main:

\begin{equation}
f_{\mathrm{brightness}},
f_{\mathrm{contrast}}
\in
[0.90,1.10].
\label{eq:ablation_photometric_ranges}
\end{equation}

Các phép biến đổi tiếp tục duy trì điều kiện grayscale

\begin{equation}
R=G=B.
\label{eq:ablation_grayscale_constraint}
\end{equation}

Không bổ sung geometric augmentation, random erasing hoặc mosaic trong nhóm
ablation này. Ngoài sự hiện diện hoặc vắng mặt của brightness và contrast,
các thành phần còn lại của cấu hình Main được giữ nguyên.

\paragraph{Ablation chiến lược lọc pseudo-label cho hồi quy bounding box.}

Nhóm thứ ba khảo sát chiến lược đánh giá độ tin cậy của pseudo-label dùng cho
bbox regression. Lý do lựa chọn thành phần này xuất phát từ việc confidence
của lớp và độ ổn định của tọa độ bounding box phản ánh hai dạng bất định khác
nhau. SoftTeacher sử dụng regression uncertainty để kiểm soát riêng độ tin
cậy của pseudo-label cho nhánh hồi quy \cite{xu2021end}; vì vậy, đây là một
thành phần cơ chế quan trọng cần được phân tích độc lập.

Ba cấu hình được xác định trước:

\begin{itemize}
    \item F0: giữ pseudo-label có initial score lớn hơn \(0.50\), không áp dụng
    thêm tiêu chí độ tin cậy cho bbox regression;

    \item F1: sử dụng confidence-based filtering với score lớn hơn \(0.90\);

    \item F2: giữ pseudo-label có initial score lớn hơn \(0.50\) đồng thời có
    regression uncertainty nhỏ hơn \(0.02\), tương ứng cấu hình Main.
\end{itemize}

Điều kiện Main của nhóm này là

\begin{equation}
F_{\mathrm{Main}}
=
\left(
s_{\mathrm{initial}}>0.50
\right)
\land
\left(
u_{\mathrm{reg}}<0.02
\right).
\label{eq:ablation_filtering_main}
\end{equation}

Các tham số của cơ chế regression uncertainty, bao gồm số lần jitter và biên
độ jitter, được giữ nguyên như đã định nghĩa trong
Mục~\ref{subsubsec:pseudo_regression_reliability}. Như vậy, ablation này chỉ
thay đổi chiến lược quyết định pseudo-label nào được chấp nhận cho nhánh hồi
quy, không thay đổi cơ chế sinh pseudo-label hoặc các thành phần SSL khác.

\paragraph{Các thành phần không đưa vào ablation.}

Nghiên cứu không thực hiện ablation đối với EMA hoặc burn-in/ramp-up. Việc
loại các thành phần này khỏi phạm vi không có nghĩa chúng không quan trọng,
mà phản ánh giới hạn có chủ đích của thiết kế component analysis.

Mục tiêu chính của ablation đã được tập trung vào ba nhóm thành phần liên quan
trực tiếp đến quá trình hình thành và sử dụng pseudo-label gồm confidence,
strong photometric augmentation và bbox-regression reliability filtering.
Bổ sung thêm EMA hoặc burn-in sẽ làm tăng đáng kể số lượng lần huấn luyện,
trong khi không cần thiết để trả lời mục tiêu cơ chế đã xác định. Phạm vi
giới hạn này giúp giữ ablation khả thi về tài nguyên và tránh biến nghiên cứu
thành một quá trình khảo sát toàn bộ không gian cấu hình SSL.

\paragraph{Số lượng cấu hình và số lần huấn luyện.}

Ba nhóm ablation lần lượt gồm 5 cấu hình confidence, 4 cấu hình augmentation
và 3 cấu hình filtering. Tuy nhiên, C4, A3 và F2 cùng quy chiếu về một cấu
hình Main duy nhất. Do đó số cấu hình khác nhau là

\begin{equation}
N_{\mathrm{abl,config}}
=
5+4+3-2
=
10.
\label{eq:ablation_unique_configs}
\end{equation}

Với 10 training seed cho mỗi cấu hình, tổng số quan sát
configuration--seed là

\begin{equation}
N_{\mathrm{abl,obs}}
=
10
\times
10
=
100.
\label{eq:ablation_total_observations}
\end{equation}

Nếu 10 lần huấn luyện Main R50--\(10\%\) đã có trong ma trận thực nghiệm chính
được tái sử dụng, số lần huấn luyện bổ sung cần thực hiện cho ablation là

\begin{equation}
N_{\mathrm{abl,additional}}
=
100
-
10
=
90.
\label{eq:ablation_additional_runs}
\end{equation}

\paragraph{Chỉ số đánh giá và logic phân tích cơ chế.}

Chỉ số chính của ablation là bbox mAP@[0.50:0.95] trên tập validation cố định.
AP50 và AP75 được sử dụng như các chỉ số hiệu năng phát hiện bổ sung.

Chất lượng pseudo-label được sử dụng như một chỉ số cơ chế thứ cấp quan trọng,
với định nghĩa đã trình bày trước đó:

\begin{equation}
Q_{\mathrm{pseudo}}
=
PL\text{-}mAP@[0.50:0.95].
\label{eq:ablation_qpseudo}
\end{equation}

Việc sử dụng đồng thời hai tầng chỉ số nhằm hỗ trợ chuỗi diễn giải

\begin{equation}
\text{thành phần SSL}
\longrightarrow
\text{hành vi/chất lượng pseudo-label}
\longrightarrow
\text{hiệu năng phát hiện}.
\label{eq:ablation_mechanistic_chain}
\end{equation}

Nếu chỉ sử dụng bbox mAP, nghiên cứu có thể xác định hiệu năng thay đổi nhưng
khó giải thích cơ chế trung gian. Ngược lại, nếu chỉ sử dụng
\(Q_{\mathrm{pseudo}}\), nghiên cứu có thể quan sát thay đổi ở pseudo-label
nhưng chưa biết thay đổi đó có chuyển thành thay đổi hiệu năng detector hay
không. Vì vậy, bbox mAP@[0.50:0.95] trên validation là chỉ số ablation chính,
còn \(Q_{\mathrm{pseudo}}\) giữ vai trò chỉ số cơ chế thứ cấp; hai chỉ số
không được xem là co-primary endpoint.

Các chỉ số pseudo-label bổ sung như PL-AP50, PL-AP75, precision, recall, F1,
mean matched IoU, pseudo-label count, retention rate, pseudo false positive
trên ảnh validation âm tính hoặc regression uncertainty có thể được báo cáo
để hỗ trợ diễn giải nhưng không thay thế \(Q_{\mathrm{pseudo}}\).

\paragraph{Phân tích ghép cặp theo training seed.}

Ablation được phân tích theo hướng ước lượng
(\textit{estimation-focused}) thay vì xây dựng thêm một hệ thống kiểm định
giả thuyết chính thức.

Với một chỉ số \(M\), biến thể \(v\) và training seed \(s\), hiệu ứng ghép cặp
so với Main được định nghĩa bởi

\begin{equation}
\Delta M_{v,s}
=
M_{v,s}
-
M_{\mathrm{Main},s}.
\label{eq:ablation_paired_metric_effect}
\end{equation}

Đối với chất lượng pseudo-label,

\begin{equation}
\Delta Q_{v,s}
=
Q_{v,s}
-
Q_{\mathrm{Main},s}.
\label{eq:ablation_paired_qpseudo_effect}
\end{equation}

Mỗi biến thể được ghép với Main theo cùng training seed. Đối với từng variant,
nghiên cứu báo cáo tối thiểu giá trị trung bình, độ lệch chuẩn mẫu, giá trị
trung bình của Main, chênh lệch trung bình ghép cặp, khoảng tin cậy 95\% hai
phía của chênh lệch và các giá trị ghép cặp theo seed.

Không thực hiện formal hypothesis test riêng cho từng biến thể ablation và
không tạo thêm một family điều chỉnh đa kiểm định. Lý do là ablation có vai
trò thứ cấp và cơ chế, không phải một confirmatory hypothesis family. Phân
tích theo effect size, khoảng tin cậy và tính nhất quán giữa các seed cung cấp
thông tin về hướng, độ lớn và độ bất định của ảnh hưởng mà không làm phình
hệ thống multiplicity đã xác định cho các câu hỏi nghiên cứu chính thức.

Kết quả âm, gần bằng không hoặc bất lợi so với Main vẫn được giữ lại và báo
cáo. Không sử dụng khái niệm ``ablation winner'' để thay cấu hình Main.

\paragraph{Giới hạn sử dụng validation và test.}

Các biến thể ablation khác Main chỉ được đánh giá trên tập validation cố định:

\begin{equation}
v\neq Main
\quad\Rightarrow\quad
Evaluation(v)=Validation\ only.
\label{eq:ablation_validation_only}
\end{equation}

Các biến thể này không được đưa sang tập test cuối cùng. Tập test không được
sử dụng để lựa chọn confidence threshold, augmentation, filtering strategy
hoặc xác định biến thể ``tốt nhất''.

Cấu hình SSL Main đã được xác định trước vẫn là cấu hình duy nhất của nhánh
thực nghiệm chính được đưa vào quy trình đánh giá test cuối cùng. Quy tắc này
giúp giữ test như một tập đánh giá độc lập và ngăn ablation trở thành một
quá trình lựa chọn mô hình dựa trên kết quả đánh giá.

Tóm lại, ablation study được giới hạn tại Faster R-CNN + ResNet-50 + FPN,
ngân sách nhãn \(10\%\) và cùng 10 training seed. Theo nguyên tắc OFAT, nghiên
cứu khảo sát ba nhóm thành phần gồm confidence threshold của pseudo-label,
strong photometric augmentation và chiến lược bbox-regression pseudo-label
reliability filtering. Thiết kế này nhằm giải thích hướng và độ lớn của ảnh
hưởng từng thành phần lên chất lượng pseudo-label và hiệu năng phát hiện, đồng
thời giữ cấu hình SSL Main như một mốc tham chiếu được xác định trước.

\subsection{Evaluation Metrics}
\label{subsec:Evaluation_Metrics}

Hệ thống chỉ số đánh giá được xác định trước nhằm phản ánh đồng thời chất
lượng phân loại, độ chính xác định vị bounding box, khả năng thu hồi tổn
thương, hành vi trên từng lớp bất thường, gánh nặng false positive và chất
lượng pseudo-label. Việc sử dụng nhiều nhóm chỉ số là cần thiết vì một giá
trị mAP tổng hợp, mặc dù phù hợp làm endpoint chính, không mô tả đầy đủ các
khía cạnh có liên quan đến các câu hỏi nghiên cứu về lớp hiếm, ảnh âm tính
\textit{No Finding} và cơ chế pseudo-label.

Các chỉ số detection chính được tính theo giao thức COCO cho bài toán bounding
box \cite{lin2014microsoft}. Cùng một định nghĩa metric được sử dụng cho
supervised learning và semi-supervised learning; sự khác biệt giữa hai phương
thức học không làm thay đổi evaluator hoặc quy tắc tính metric.

\subsubsection{Chỉ số hiệu năng phát hiện chính}
\label{subsubsec:primary_detection_metric}

Chỉ số hiệu năng phát hiện chính của toàn nghiên cứu là
bbox mAP@[0.50:0.95], tương ứng với trung bình Average Precision trên mười
ngưỡng Intersection over Union (IoU)

\begin{equation}
\mathcal{T}_{IoU}
=
\left\{
0.50,\,
0.55,\,
0.60,\,
\ldots,\,
0.95
\right\}.
\label{eq:coco_iou_threshold_set}
\end{equation}

Với một bounding box dự đoán \(B_p\) và một ground-truth bounding box
\(B_g\), IoU được xác định bởi

\begin{equation}
IoU(B_p,B_g)
=
\frac{
|B_p\cap B_g|
}{
|B_p\cup B_g|
}.
\label{eq:detection_iou}
\end{equation}

Tại mỗi ngưỡng IoU \(t\in\mathcal{T}_{IoU}\), Average Precision được tính từ
đường precision--recall của detector theo giao thức COCO. Chỉ số chính được
xác định bởi

\begin{equation}
mAP@[0.50:0.95]
=
\frac{1}{10}
\sum_{t\in\mathcal{T}_{IoU}}
AP_t.
\label{eq:primary_bbox_map}
\end{equation}

bbox mAP@[0.50:0.95] được lựa chọn làm endpoint chính vì nó đánh giá detector
trên nhiều mức yêu cầu về độ chồng lấp giữa bounding box dự đoán và ground
truth. Do đó, metric này phản ánh đồng thời khả năng phân loại đúng đối tượng
và mức chính xác của việc định vị, nghiêm ngặt hơn việc chỉ đánh giá tại một
ngưỡng IoU duy nhất như AP50 \cite{lin2014microsoft}.

Đây là endpoint chính dùng để đánh giá hiệu năng phát hiện trong RQ2 và là
metric nền tảng cho các hiệu ứng SUP--SSL đã được định nghĩa trong phần giả
thuyết nghiên cứu. Việc lựa chọn endpoint này được xác định trước và không
thay đổi theo kết quả thực nghiệm.

\subsubsection{AP50 và AP75}
\label{subsubsec:ap50_ap75_metrics}

Bên cạnh endpoint chính, nghiên cứu báo cáo AP50 và AP75 như các chỉ số
detection thứ cấp:

\begin{equation}
AP50
=
AP_{IoU=0.50},
\label{eq:ap50_metric}
\end{equation}

và

\begin{equation}
AP75
=
AP_{IoU=0.75}.
\label{eq:ap75_metric}
\end{equation}

AP50 sử dụng tiêu chuẩn chồng lấp ít nghiêm ngặt hơn và phản ánh khả năng phát
hiện đối tượng khi bounding box dự đoán đạt mức tương hợp cơ bản với ground
truth. Ngược lại, AP75 đặt yêu cầu định vị chặt hơn và nhạy hơn với sai lệch
tọa độ bounding box.

Việc báo cáo đồng thời AP50, AP75 và bbox mAP@[0.50:0.95] giúp phân biệt một
tình huống trong đó mô hình phát hiện đúng vùng bất thường ở mức tương đối
nhưng chưa định vị chính xác với một tình huống trong đó bounding box dự đoán
bám sát ground truth hơn. Tuy nhiên, AP50 và AP75 giữ vai trò thứ cấp và không
thay thế endpoint chính của nghiên cứu.

\subsubsection{Giao thức COCO cho các chỉ số Average Precision}
\label{subsubsec:coco_evaluation_contract}

Các chỉ số AP được tính bằng evaluator dành cho bounding box theo giao thức
COCO \cite{lin2014microsoft}. Các thiết lập đánh giá chính được cố định như
sau:

\begin{equation}
\mathcal{T}_{IoU}
=
\{0.50,0.55,\ldots,0.95\},
\qquad
\mathcal{T}_{recall}
=
\{0.00,0.01,\ldots,1.00\}.
\label{eq:coco_eval_grids}
\end{equation}

Đánh giá được thực hiện theo từng category với

\begin{equation}
useCats
=
1,
\qquad
maxDets
=
100.
\label{eq:coco_eval_category_maxdet}
\end{equation}

Không áp dụng thêm một score threshold cố định trước khi tính AP. Các dự đoán
được xếp hạng theo confidence score trong evaluator, và precision--recall được
xây dựng từ thứ hạng đó theo giao thức COCO. Quy tắc này cần được phân biệt
với ngưỡng vận hành
\(\tau_{\mathrm{eval}}=0.50\)
được sử dụng cho các chỉ số như Recall tại operating point và FP/image ở các tiểu mục phía sau.

Nghiên cứu cũng không báo cáo các chỉ số COCO
\(AP_S\), \(AP_M\) và \(AP_L\) như các endpoint chính thức. Lý do là trong
phân tích đặc điểm dữ liệu trước huấn luyện, các nhóm Small, Medium và Large
được xác định theo diện tích bounding box chuẩn hóa theo kích thước ảnh,
không phải theo các ngưỡng diện tích pixel tuyệt đối của giao thức COCO
\cite{lin2014microsoft}. Vì vậy, việc sử dụng các ký hiệu
\(AP_S\), \(AP_M\) hoặc \(AP_L\) có thể gây nhầm lẫn giữa định nghĩa kích
thước riêng của nghiên cứu và định nghĩa kích thước chuẩn của COCO.

\subsubsection{Average Recall và Recall tại điểm vận hành cố định}
\label{subsubsec:recall_metrics}

Khả năng thu hồi tổn thương được đánh giá trước hết bằng COCO Average Recall
trên toàn bộ khoảng IoU từ 0.50 đến 0.95, với tối đa 100 detection trên mỗi
ảnh và không giới hạn theo nhóm diện tích:

\begin{equation}
AR
=
AR@[0.50:0.95],
\qquad
maxDets=100,
\qquad
area=all.
\label{eq:coco_average_recall}
\end{equation}

Average Recall bổ sung cho Average Precision bằng cách nhấn mạnh khả năng mô
hình tìm lại các ground-truth objects khi thay đổi yêu cầu IoU.

Ngoài COCO AR, nghiên cứu sử dụng Recall tại một điểm vận hành cố định như một
chỉ số chẩn đoán bổ sung. Điểm vận hành toàn cục được xác định trước bởi

\begin{equation}
\tau_{\mathrm{eval}}
=
0.50.
\label{eq:fixed_evaluation_threshold}
\end{equation}

Sau detector-native NMS và giới hạn tối đa 100 detection trên mỗi ảnh, chỉ
các detection có confidence score thỏa

\begin{equation}
s_j
\ge
\tau_{\mathrm{eval}}
\label{eq:operating_score_filter}
\end{equation}

mới được giữ lại để tính các chỉ số tại operating point.

Để phép matching tại operating point có tính xác định và tái lập, các
detection còn lại sau bước lọc theo
\(\tau_{\mathrm{eval}}\)
được duyệt theo confidence score giảm dần. Khi hai detection có cùng
confidence score, thứ tự đầu ra của detector sau detector-native
post-processing được giữ nguyên.

Với mỗi detection theo thứ tự này, chỉ các ground-truth bounding box chưa
được ghép và thuộc cùng lớp mới được xem là ứng viên. Detection được ghép
với ground-truth bounding box có IoU lớn nhất trong tập ứng viên; cặp ghép
chỉ được chấp nhận khi

\begin{equation}
IoU(B_p,B_g)
\ge
0.50.
\label{eq:operating_point_matching_iou}
\end{equation}

Nếu nhiều ground-truth bounding box có cùng giá trị IoU lớn nhất, ground
truth xuất hiện trước theo thứ tự annotation cố định trong tệp COCO của tập
đánh giá được chọn. Sau khi một ground truth được ghép thành công, ground
truth đó không còn khả dụng cho các detection tiếp theo. Quy trình này tạo
thành deterministic score-descending greedy, category-aware, one-to-one
matching. Từ số true positive \(TP\) và false negative \(FN\), Recall tại
operating point được tính bởi

\begin{equation}
Recall_{\tau_{\mathrm{eval}}}
=
\frac{TP}{TP+FN}.
\label{eq:operating_point_recall}
\end{equation}

Việc cố định cùng một operating point cho mọi phương pháp, kiến trúc, ngân
sách nhãn và training seed giúp các chỉ số phụ thuộc threshold có thể được
so sánh trực tiếp mà không lựa chọn một threshold riêng có lợi cho từng mô
hình.

\subsubsection{Average Precision theo từng lớp}
\label{subsubsec:classwise_ap}

Để tránh việc metric tổng hợp che khuất sự khác biệt giữa các lớp bất thường,
nghiên cứu báo cáo Average Precision riêng cho toàn bộ 14 detection classes.
Với lớp \(c\), chỉ số được ký hiệu

\begin{equation}
AP_c
=
AP_c@[0.50:0.95],
\qquad
c=1,\ldots,14.
\label{eq:classwise_ap_metric}
\end{equation}

Class-wise AP đặc biệt quan trọng trong bộ dữ liệu có phân phối không cân bằng,
vì một giá trị mAP tổng hợp có thể không phản ánh rõ hành vi của các lớp có
support thấp.

Các class-wise AP được tính bằng cùng evaluator và cùng checkpoint đã dùng
cho đánh giá tổng thể. Nghiên cứu không lựa chọn checkpoint riêng cho từng lớp
và không thay đổi evaluator hoặc threshold nhằm tối đa hóa kết quả của một lớp
cụ thể.

\subsubsection{Đánh giá các lớp hiếm}
\label{subsubsec:rare_class_metrics}

Các lớp hiếm được xác định từ chẩn đoán dữ liệu trước huấn luyện trên tập train
cố định, sử dụng tiêu chí image-level prevalence nhỏ hơn \(5\%\). Theo định
nghĩa đã xác định trước khi huấn luyện, hai lớp thuộc nhóm này là
\textit{Atelectasis} và \textit{Pneumothorax}.

Đối với mỗi lớp hiếm \(c\), metric chính được giữ ở mức class-specific AP:

\begin{equation}
AP^{rare}_{c}
=
AP_c@[0.50:0.95],
\qquad
c
\in
\{
\text{Atelectasis},
\text{Pneumothorax}
\}.
\label{eq:rare_class_ap}
\end{equation}

Nghiên cứu không gộp hai lớp này thành một chỉ số ``rare-mAP''. Lý do là hai
lớp có support và mức độ khó khác nhau; gộp chúng thành một trung bình mới có
thể che khuất hành vi riêng của từng lớp và tạo thêm một endpoint không được
xác định trong giao thức chính.

Phân tích lớp hiếm vì vậy giữ nguyên ở mức từng lớp và phục vụ RQ5. Các kết
quả này có vai trò khám phá; cách tổng hợp qua seed và cách biểu diễn độ bất
định được trình bày tại Mục~\ref{subsec:Statistical_Analysis}.

\subsubsection{False positive trên toàn bộ ảnh}
\label{subsubsec:false_positive_per_image}

Ngoài precision--recall metrics, nghiên cứu đánh giá gánh nặng false positive
ở điểm vận hành cố định
\(\tau_{\mathrm{eval}}=0.50\).
Sau detector-native NMS, giới hạn tối đa 100 detection và matching
category-aware one-to-one với ground truth ở
\(IoU\ge0.50\), một detection không được ghép hợp lệ được tính là false
positive.

Với \(FP_i\) là số false positive trên ảnh \(i\) và \(N\) là tổng số ảnh trong
tập đánh giá, số false positive trung bình trên mỗi ảnh được xác định bởi

\begin{equation}
FP/image
=
\frac{
\sum_{i=1}^{N}FP_i
}{
N
}.
\label{eq:fp_per_image}
\end{equation}

Metric này bổ sung thông tin mà AP không thể hiện trực tiếp: một detector có
thể đạt AP tương đối tốt nhưng vẫn tạo ra nhiều dự đoán sai tại một operating
point thực tế. Việc báo cáo FP/image do đó giúp mô tả thêm mức độ ``quá phát
hiện'' của mô hình.

\subsubsection{False positive trên ảnh \textit{No Finding}}
\label{subsubsec:no_finding_metrics}

Trong nghiên cứu này, \textit{No Finding} không phải là một detection class.
Các ảnh này được xem là một stratum đánh giá âm tính có zero ground-truth
bounding box. Do không có đối tượng bất thường thật, mọi detection được giữ
lại sau detector-native NMS và ngưỡng
\(\tau_{\mathrm{eval}}=0.50\)
trên một ảnh \textit{No Finding} đều được tính là false positive.

Gọi \(\mathcal{N}\) là tập các ảnh \textit{No Finding} trong tập đánh giá.
Metric chính phục vụ RQ6 là số false positive trung bình trên mỗi ảnh âm tính:

\begin{equation}
FP/negative
=
\frac{
\sum_{i\in\mathcal{N}}FP_i
}{
|\mathcal{N}|
}.
\label{eq:fp_per_negative_image}
\end{equation}

Tập validation và test cố định mỗi tập chứa 75 ảnh
\textit{No Finding}; membership của các ảnh này không thay đổi giữa các mô
hình được so sánh.

Bên cạnh số false positive trung bình, nghiên cứu báo cáo
\textit{false-alarm rate} trên ảnh âm tính như một chỉ số thứ cấp:

\begin{equation}
FAR_{\mathrm{neg}}
=
\frac{
\#\left\{
i\in\mathcal{N}:FP_i\ge1
\right\}
}{
|\mathcal{N}|
}.
\label{eq:false_alarm_rate_negative}
\end{equation}

\(FP/negative\) trả lời câu hỏi một ảnh âm tính trung bình chứa bao nhiêu
detection sai, trong khi \(FAR_{\mathrm{neg}}\) phản ánh tỷ lệ ảnh âm tính có
ít nhất một false positive. Hai đại lượng vì vậy mô tả hai khía cạnh khác nhau
của false-positive burden.

Việc phân tích riêng nhóm \textit{No Finding} là cần thiết vì mAP được tính
trên các detection classes không trực tiếp cho biết detector có xu hướng sinh
bao nhiêu cảnh báo giả trên những ảnh không chứa bất thường. Điều này đặc biệt
liên quan đến RQ6 và rủi ro pseudo-label sai trên ảnh âm tính đã được thảo
luận ở phần cơ sở nghiên cứu.

\subsubsection{Chất lượng pseudo-label}
\label{subsubsec:evaluation_pseudo_label_quality}

Đối với các điều kiện semi-supervised, nghiên cứu bổ sung phép đánh giá chất
lượng pseudo-label nhằm phục vụ RQ4 và các phân tích cơ chế. Chỉ số này đã
được định nghĩa tại
Mục~\ref{subsubsec:pseudo_label_quality} là
\(Q_{\mathrm{pseudo}}\), với

\begin{equation}
Q_{\mathrm{pseudo}}
=
PL\text{-}mAP@[0.50:0.95].
\label{eq:evaluation_qpseudo}
\end{equation}

\(Q_{\mathrm{pseudo}}\) được tính trên tập validation cố định bằng LAST EMA
Teacher và chỉ sử dụng tập RCNN classification pseudo-label đã được chấp nhận
theo giao thức Main. Hidden ground truth của các tập \(U_b\) không được sử
dụng để tính metric này, và \(Q_{\mathrm{pseudo}}\) không được tính trên tập
test.

Việc sử dụng PL-mAP@[0.50:0.95] xuất phát từ cấu trúc của pseudo-label trong
object detection: mỗi pseudo-label chứa đồng thời quyết định lớp và vị trí
bounding box. Do đó, chỉ số chất lượng cần phản ánh cả classification và
localization trên nhiều mức IoU, thay vì chỉ dựa trên mean confidence, số
lượng pseudo-label hoặc tỷ lệ pseudo-label được giữ lại.

Các metric như PL-AP50, PL-AP75, precision, recall, F1, mean matched IoU,
class-wise pseudo-label AP, pseudo-label count, retention rate, pseudo false
positive trên ảnh validation âm tính và regression uncertainty có thể được
sử dụng như các chỉ số cơ chế bổ sung. Tuy nhiên, chúng không thay thế
\(Q_{\mathrm{pseudo}}\).

\(Q_{\mathrm{pseudo}}\) cũng không phải là một endpoint đồng cấp với bbox
mAP@[0.50:0.95] của toàn nghiên cứu. bbox mAP đánh giá hiệu năng detector,
trong khi \(Q_{\mathrm{pseudo}}\) đánh giá chất lượng tín hiệu pseudo-label
trung gian. Hai metric được sử dụng đồng thời để hỗ trợ phân tích mối liên hệ
giữa chất lượng pseudo-label và mức lợi ích của semi-supervised learning.

\subsubsection{Phạm vi sử dụng validation và test}
\label{subsubsec:evaluation_validation_test_roles}

Các metric nêu trên được tính theo cùng định nghĩa trên các tập dữ liệu mà
giao thức cho phép sử dụng, nhưng validation và test có vai trò khác nhau.

Tập validation được phép sử dụng cho lựa chọn checkpoint theo bbox
mAP@[0.50:0.95], tính \(Q_{\mathrm{pseudo}}\) và đánh giá các biến thể ablation
đã xác định trước. Ngược lại, tập test được giữ cho chiến dịch đánh giá cuối
cùng của các điều kiện Main và supervised reference theo giao thức đã xác định trước.

Đặc biệt,

\begin{equation}
Q_{\mathrm{pseudo}}^{test}
=
\varnothing,
\label{eq:no_test_qpseudo}
\end{equation}

và các biến thể ablation khác Main không được đưa sang test.

Việc tách vai trò validation và test giúp ngăn các metric đánh giá trở thành
công cụ điều chỉnh phương pháp sau khi đã quan sát kết quả test, đồng thời giữ
tập test như nguồn bằng chứng cuối cùng độc lập cho các phép so sánh chính.

Tóm lại, nghiên cứu sử dụng bbox mAP@[0.50:0.95] làm endpoint phát hiện chính;
AP50, AP75, COCO AR, Recall tại operating point, FP/image và class-wise AP làm
các chỉ số bổ sung; AP riêng cho Atelectasis và Pneumothorax phục vụ phân tích
lớp hiếm; FP/negative và \(FAR_{\mathrm{neg}}\) phục vụ đánh giá trên ảnh
\textit{No Finding}; và \(Q_{\mathrm{pseudo}}\) phục vụ phân tích chất lượng
pseudo-label. Hệ metric này cho phép đánh giá đồng thời hiệu năng tổng thể,
độ chính xác định vị, khả năng thu hồi, hành vi theo lớp, false-positive burden
và cơ chế pseudo-label mà không thay đổi endpoint chính của nghiên cứu.

\subsection{Statistical Analysis}
\label{subsec:Statistical_Analysis}

Phân tích thống kê được xác định trước nhằm lượng hóa mức độ khác biệt giữa
supervised learning và semi-supervised learning, đánh giá sự phụ thuộc của
mức lợi ích này vào ngân sách nhãn và kiến trúc mô hình, đồng thời phân tích
mối liên hệ giữa chất lượng pseudo-label và hiệu năng phát hiện.

Đơn vị suy luận thống kê của nghiên cứu là một lần huấn luyện hợp lệ ứng với
một \texttt{training\_seed}. Ảnh, bounding box và detection không được xem là
các lần lặp độc lập của phương pháp huấn luyện, vì chúng cùng được sinh ra từ
một mô hình đã huấn luyện. Với mỗi điều kiện thực nghiệm, số lần lặp chính
được xác định là

\begin{equation}
n
=
10
\quad
\text{training seeds}.
\label{eq:statistical_inferential_unit}
\end{equation}

Mức ý nghĩa thống kê chung được xác định trước là

\begin{equation}
\alpha
=
0.05.
\label{eq:statistical_alpha}
\end{equation}

Việc diễn giải kết quả ưu tiên độ lớn hiệu ứng và độ bất định trước giá trị
\(p\). Theo đó, thứ tự báo cáo chính được tổ chức theo hướng

\begin{equation}
\text{effect estimate}
\longrightarrow
95\%\,CI
\longrightarrow
p\text{-value}.
\label{eq:statistical_reporting_hierarchy}
\end{equation}

Cách tiếp cận này nhằm tránh việc diễn giải kết quả chỉ dựa trên việc một
\(p\)-value nằm ở phía nào của ngưỡng \(0.05\)
\cite{wasserstein2016asa}.

\subsubsection{Tổng hợp kết quả qua nhiều lần huấn luyện}
\label{subsubsec:seed_level_summary}

Đối với mỗi tổ hợp giữa phương thức học, kiến trúc và ngân sách nhãn, kết quả
được lưu riêng cho từng training seed. Với \(x_k\) là giá trị của một metric
tại seed thứ \(k\), giá trị trung bình trên \(n=10\) seed được xác định bởi

\begin{equation}
\bar{x}
=
\frac{1}{n}
\sum_{k=1}^{n}
x_k,
\qquad
n=10.
\label{eq:seed_mean}
\end{equation}

Độ lệch chuẩn mẫu được tính bởi

\begin{equation}
s_x
=
\sqrt{
\frac{1}{n-1}
\sum_{k=1}^{n}
\left(
x_k-\bar{x}
\right)^2
}.
\label{eq:seed_sample_sd}
\end{equation}

Mẫu số \(n-1\) tương ứng với cách tính sample standard deviation với
\texttt{ddof}=1.

Kết quả của từng seed được giữ lại và trình bày cùng giá trị trung bình và độ
lệch chuẩn mẫu. Một lần chạy hợp lệ nhưng có hiệu năng thấp không bị loại bỏ,
không được thay bằng seed khác và không được thực hiện lại chỉ vì kết quả
không thuận lợi. Không lựa chọn seed có hiệu năng cao nhất để đại diện cho một
điều kiện thực nghiệm.

\subsubsection{Hiệu ứng ghép cặp SUP--SSL ở mức seed}
\label{subsubsec:paired_seed_effect}

Các phép so sánh giữa supervised learning và semi-supervised learning được
thực hiện theo cặp trên cùng kiến trúc, ngân sách nhãn và training seed.

Với metric \(M\), kiến trúc \(a\), ngân sách nhãn \(b\) và seed \(s\), hiệu
ứng ghép cặp cơ bản được xác định bởi

\begin{equation}
\Delta M_{a,b,s}
=
M^{SSL}_{a,b,s}
-
M^{SUP}_{a,b,s}.
\label{eq:paired_seed_metric_effect}
\end{equation}

Đối với bbox mAP@[0.50:0.95], giá trị dương của
\(\Delta M_{a,b,s}\) biểu thị hiệu năng SSL cao hơn supervised baseline ở
cặp tương ứng.

Độ lớn hiệu ứng chính thức được báo cáo dưới dạng chênh lệch tuyệt đối của
metric. Đối với mAP và AP, các chênh lệch được trình bày theo đơn vị
\textit{AP points}. Tỷ lệ thay đổi tương đối theo phần trăm, nếu được báo cáo,
chỉ có vai trò mô tả bổ sung và không thay thế hiệu ứng tuyệt đối.

\subsubsection{Phân tích chính: SSL so với supervised baseline}
\label{subsubsec:statistical_rq2}

Phân tích chính đánh giá liệu semi-supervised learning có tạo ra mức cải thiện
bbox mAP@[0.50:0.95] dương so với supervised baseline trên toàn bộ ma trận
hạn chế nhãn hay không.

Với mỗi training seed \(s\), mức cải thiện tổng hợp được tính bằng trung bình
đều của tám tổ hợp gồm hai kiến trúc và bốn ngân sách nhãn:

\begin{equation}
G_s
=
\frac{1}{8}
\sum_{a}
\sum_{b}
\Delta mAP_{a,b,s},
\label{eq:statistical_rq2_gain}
\end{equation}

trong đó

\begin{equation}
a
\in
\{
R50,
Swin\text{-}T
\},
\qquad
b
\in
\{
1\%,
5\%,
10\%,
20\%
\}.
\label{eq:statistical_rq2_cells}
\end{equation}

Phân tích chính sử dụng one-sample \(t\)-test một phía trên
\(\{G_s\}_{s=1}^{10}\), tương ứng với giả thuyết có hướng đã xác định trước.
Với \(n=10\), bậc tự do là

\begin{equation}
df
=
n-1
=
9.
\label{eq:rq2_degrees_freedom}
\end{equation}

Mặc dù phép kiểm định chính là một phía, độ bất định của hiệu ứng vẫn được
báo cáo bằng khoảng tin cậy 95\% hai phía. Các hiệu ứng tại từng tổ hợp riêng
lẻ giữa kiến trúc và ngân sách được báo cáo mô tả bằng mean, sample SD và
95\% CI, nhưng không tạo thêm các phép kiểm định chính thức cho từng cell.

\subsubsection{Ảnh hưởng của ngân sách nhãn đến mức lợi ích SSL}
\label{subsubsec:statistical_rq3}

Để đánh giá liệu mức lợi ích của SSL thay đổi theo ngân sách nhãn hay không,
hiệu ứng tại mỗi ngân sách được trung bình trên hai kiến trúc trong cùng
training seed:

\begin{equation}
B_{b,s}
=
\frac{
\Delta mAP_{R50,b,s}
+
\Delta mAP_{Swin\text{-}T,b,s}
}{
2
}.
\label{eq:statistical_budget_gain}
\end{equation}

Bốn mức

\begin{equation}
b
\in
\{
1\%,
5\%,
10\%,
20\%
\}
\label{eq:statistical_budget_levels}
\end{equation}

được xem là bốn điều kiện lặp lại trên cùng 10 training seeds.

Phân tích chính sử dụng one-way repeated-measures ANOVA. Hiệu chỉnh
Greenhouse--Geisser được áp dụng nhất quán để giảm phụ thuộc vào giả định
sphericity \cite{greenhouse1959profile}. Kết quả báo cáo gồm thống kê \(F\),
hệ số hiệu chỉnh

\begin{equation}
\epsilon_{GG},
\label{eq:greenhouse_geisser_epsilon}
\end{equation}

bậc tự do đã hiệu chỉnh và \(p_{GG}\).

Nghiên cứu không đặt trước giả thuyết tuyến tính hoặc đơn điệu theo ngân sách
nhãn. Nếu cần xác định cụ thể các mức ngân sách khác nhau ở đâu, sáu tương
phản ghép cặp đã xác định trước được xem xét:

\begin{equation}
\mathcal{C}_{budget}
=
\left\{
\begin{array}{l}
5\%-1\%,\\
10\%-1\%,\\
20\%-1\%,\\
10\%-5\%,\\
20\%-5\%,\\
20\%-10\%
\end{array}
\right\}.
\label{eq:rq3_budget_contrasts}
\end{equation}

Đối với mỗi tương phản \(c=(b_2-b_1)\in\mathcal{C}_{budget}\) và mỗi
training seed \(s\), hiệu ứng ghép cặp giữa hai mức ngân sách được xác định bởi

\begin{equation}
D_{c,s}
=
B_{b_2,s}
-
B_{b_1,s}.
\label{eq:rq3_budget_contrast_difference}
\end{equation}

Mỗi tương phản được kiểm định bằng two-sided one-sample \(t\)-test trên
\(\{D_{c,s}\}_{s=1}^{10}\) so với giá trị không, tương đương với paired
\(t\)-test giữa hai mức ngân sách trên cùng các training seeds. Với \(n=10\),
bậc tự do của mỗi phép kiểm định là

\begin{equation}
df
=
n-1
=
9.
\label{eq:rq3_budget_contrast_df}
\end{equation}

Đối với mỗi tương phản, nghiên cứu báo cáo mean paired difference, sample
standard deviation với \(ddof=1\), khoảng tin cậy 95\% hai phía cho hiệu ứng
riêng và raw two-sided \(p\)-value. Sáu raw \(p\)-value này tạo thành
multiplicity family riêng của RQ3 và được điều chỉnh bằng thủ tục Holm như
trình bày ở Mục~\ref{subsubsec:statistical_multiplicity}.

\subsubsection{Sự khác biệt về mức lợi ích SSL giữa hai kiến trúc}
\label{subsubsec:statistical_rq7}

Đối với mỗi kiến trúc \(a\) và training seed \(s\), mức lợi ích SSL trung bình
qua bốn ngân sách được xác định bởi

\begin{equation}
A_{a,s}
=
\frac{1}{4}
\sum_b
\Delta mAP_{a,b,s}.
\label{eq:statistical_architecture_gain}
\end{equation}

Sự khác biệt về mức lợi ích giữa Swin-T và ResNet-50 trong cùng seed được tính
bởi

\begin{equation}
D_s
=
A_{Swin\text{-}T,s}
-
A_{R50,s}.
\label{eq:statistical_architecture_difference}
\end{equation}

Phân tích sử dụng two-sided one-sample \(t\)-test trên
\(\{D_s\}_{s=1}^{10}\), với \(df=9\).

Phân tích này kiểm tra sự khác biệt về \textbf{mức lợi ích của SSL} giữa hai
kiến trúc dưới các giao thức huấn luyện đặc thù theo kiến trúc đã được xác
định trước. Nó không được diễn giải như một kiểm định về ưu thế tổng quát hay
hiệu ứng nhân quả thuần túy của Swin-T so với ResNet-50.

\subsubsection{Mối liên hệ giữa chất lượng pseudo-label và mức cải thiện hiệu năng}
\label{subsubsec:statistical_rq4}

Phân tích cơ chế đánh giá mối liên hệ giữa chất lượng pseudo-label trên tập
validation và mức cải thiện bbox mAP@[0.50:0.95] của SSL so với supervised
baseline trên tập test cuối cùng.

Biến kết quả tại kiến trúc \(a\), ngân sách \(b\) và seed \(s\) được xác định
bởi

\begin{equation}
Y_{a,b,s}
=
\Delta mAP^{test}_{a,b,s}
=
mAP^{SSL,test}_{a,b,s}
-
mAP^{SUP,test}_{a,b,s}.
\label{eq:statistical_rq4_outcome}
\end{equation}

Biến dự báo là chất lượng pseudo-label

\begin{equation}
Q_{a,b,s}
=
PL\text{-}mAP@[0.50:0.95]_{a,b,s},
\label{eq:statistical_rq4_predictor}
\end{equation}

được tính trên tập validation cố định bằng LAST EMA Teacher như đã trình bày
trong phần đánh giá pseudo-label.

Để tách biến thiên của pseudo-label quality bên trong từng tổ hợp kiến trúc và
ngân sách khỏi khác biệt trung bình giữa các tổ hợp, biến dự báo được
centered trong từng cell \cite{enders2007centering}:

\begin{equation}
Q^{WC}_{a,b,s}
=
Q_{a,b,s}
-
\bar{Q}_{a,b},
\label{eq:rq4_within_cell_centering}
\end{equation}

trong đó

\begin{equation}
\bar{Q}_{a,b}
=
\frac{1}{10}
\sum_{s=1}^{10}
Q_{a,b,s}.
\label{eq:rq4_cell_mean_qpseudo}
\end{equation}

Để hệ số hồi quy có diễn giải thuận tiện theo AP points, biến centered được
chuẩn hóa về đơn vị một pseudo AP point:

\begin{equation}
X_{a,b,s}
=
\frac{
Q^{WC}_{a,b,s}
}{
0.01
}.
\label{eq:rq4_pseudo_ap_point_scale}
\end{equation}

Mối liên hệ được ước lượng bằng linear mixed-effects model:

\begin{equation}
Y_{a,b,s}
=
\beta_0
+
\beta_Q X_{a,b,s}
+
\gamma_{a,b}
+
u_s
+
\varepsilon_{a,b,s},
\label{eq:rq4_linear_mixed_model}
\end{equation}

trong đó \(\gamma_{a,b}\) biểu diễn fixed effect của tổ hợp
architecture--budget, còn

\begin{equation}
u_s
\sim
\mathcal{N}
\left(
0,\sigma_u^2
\right)
\label{eq:rq4_seed_random_intercept}
\end{equation}

là random intercept của training seed.

Mô hình được ước lượng bằng restricted maximum likelihood (REML). Suy luận
đối với fixed effect sử dụng hiệu chỉnh Kenward--Roger cho bậc tự do và sai
số chuẩn \cite{kenward1997small}. Kiểm định đối với
\(\beta_Q\) là một phía theo giả thuyết có hướng dương đã xác định trước.

Mặc dù mô hình có tối đa 80 quan sát
architecture--budget--seed, các quan sát này không được xem là 80 training
replications độc lập. Random intercept của training seed được sử dụng để phản
ánh cấu trúc lặp lại này.

Hệ số \(\beta_Q\) được diễn giải là mức thay đổi trung bình của
\(\Delta mAP^{test}\), tính theo AP points, khi chất lượng pseudo-label tăng
thêm một pseudo AP point trong cùng architecture--budget cell. Phân tích này
chỉ hỗ trợ diễn giải về \textbf{mối liên hệ}, không được diễn giải như bằng
chứng nhân quả.

\subsubsection{Phân tích các lớp hiếm}
\label{subsubsec:statistical_rq5}

Phân tích đối với các lớp hiếm có vai trò khám phá và được thực hiện riêng cho
\textit{Atelectasis} và \textit{Pneumothorax}.

Với lớp \(c\), hiệu ứng class-specific tại kiến trúc \(a\), ngân sách \(b\) và
seed \(s\) là

\begin{equation}
\Delta AP_{c,a,b,s}
=
AP^{SSL}_{c,a,b,s}
-
AP^{SUP}_{c,a,b,s}.
\label{eq:statistical_rare_class_effect}
\end{equation}

Hiệu ứng tổng hợp trong mỗi seed được xác định bởi

\begin{equation}
G_{c,s}
=
\frac{1}{8}
\sum_{a,b}
\Delta AP_{c,a,b,s}.
\label{eq:statistical_rare_class_seed_gain}
\end{equation}

Đối với mỗi lớp, nghiên cứu báo cáo mean paired gain, sample SD và khoảng tin
cậy 95\% hai phía. Không yêu cầu một phép kiểm định giả thuyết chính thức hoặc
\(p\)-value cho RQ5.

Nếu một lớp không có ground-truth support trong phạm vi cần đánh giá khiến AP
không xác định được, giá trị đó được ghi nhận là không xác định chứ không được
gán bằng 0.

\subsubsection{False positive trên ảnh \textit{No Finding}}
\label{subsubsec:statistical_rq6}

Đối với mỗi kiến trúc \(a\), ngân sách \(b\) và seed \(s\), hiệu ứng của SSL
lên số false positive trung bình trên ảnh âm tính được xác định bởi

\begin{equation}
\Delta F_{a,b,s}
=
FP^{SSL}_{neg,a,b,s}
-
FP^{SUP}_{neg,a,b,s}.
\label{eq:statistical_negative_fp_effect}
\end{equation}

Hiệu ứng tổng hợp trong mỗi seed là

\begin{equation}
H_s
=
\frac{1}{8}
\sum_{a,b}
\Delta F_{a,b,s}.
\label{eq:statistical_negative_fp_seed_effect}
\end{equation}

RQ6 sử dụng two-sided one-sample \(t\)-test trên
\(\{H_s\}_{s=1}^{10}\), với \(df=9\). Giá trị

\begin{equation}
H_s<0
\label{eq:negative_fp_interpretation}
\end{equation}

biểu thị ít false positive trên ảnh \textit{No Finding} hơn dưới SSL, trong
khi \(H_s>0\) biểu thị gánh nặng false positive cao hơn.

Chênh lệch về
\(FAR_{\mathrm{neg}}\)
được báo cáo theo hướng ước lượng với effect estimate và 95\% CI, nhưng không
được gắn với một phép kiểm định chính thức riêng.

\subsubsection{Kiểm soát đa kiểm định}
\label{subsubsec:statistical_multiplicity}

Cấu trúc kiểm soát đa kiểm định được xác định trước nhằm hạn chế family-wise
error rate khi nhiều câu hỏi nghiên cứu thứ cấp được phân tích.

Giả thuyết chính của nghiên cứu về lợi ích tổng thể của SSL được xem là một
family riêng chỉ chứa một phép kiểm định, vì vậy không cần điều chỉnh
multiplicity.

Bốn phân tích thứ cấp chính thức gồm:
RQ3, RQ4, RQ6 và RQ7 được đặt trong cùng một family. Bốn \(p\)-value đầu vào
gồm:

\begin{equation}
\mathcal{P}_{secondary}
=
\left\{
p_{\mathrm{RQ3}},
p_{\mathrm{RQ4}},
p_{\mathrm{RQ6}},
p_{\mathrm{RQ7}}
\right\}.
\label{eq:secondary_pvalue_family}
\end{equation}

Trong đó \(p_{\mathrm{RQ3}}\) là \(p\)-value đã hiệu chỉnh
Greenhouse--Geisser của omnibus repeated-measures ANOVA,
\(p_{\mathrm{RQ4}}\) là \(p\)-value một phía của hệ số
\(\beta_Q\), còn RQ6 và RQ7 cung cấp \(p\)-value hai phía.

Các giá trị này được điều chỉnh theo thủ tục Holm step-down để kiểm soát
family-wise error rate ở mức

\begin{equation}
FWER
=
0.05
\label{eq:holm_secondary_fwer}
\end{equation}

\cite{holm1979sequentially}.

Sáu phép so sánh cặp giữa bốn ngân sách nhãn của RQ3 tạo thành một
multiplicity family riêng, ký hiệu là \(F_3\). Các \(p\)-value đầu vào của
family này là sáu raw two-sided \(p\)-value thu được từ các phép kiểm định
ghép cặp đã xác định tại Phương trình~\eqref{eq:rq3_budget_contrast_difference}:

\begin{equation}
\mathcal{P}_{F_3}
=
\left\{
p_{5\%-1\%},
p_{10\%-1\%},
p_{20\%-1\%},
p_{10\%-5\%},
p_{20\%-5\%},
p_{20\%-10\%}
\right\}.
\label{eq:rq3_holm_f3_family}
\end{equation}

Family \(F_3\) có

\begin{equation}
m_{F_3}
=
6
\label{eq:rq3_holm_f3_size}
\end{equation}

phép kiểm định và được điều chỉnh bằng thủ tục Holm step-down để kiểm soát
family-wise error rate ở mức \(0.05\).

Một khác biệt cụ thể giữa hai ngân sách chỉ được diễn giải như một kết quả
chính thức của RQ3 khi cả omnibus analysis tương ứng và phép so sánh cặp sau
điều chỉnh đều đáp ứng tiêu chí đã xác định.

Không thu nhỏ số lượng phép kiểm định trong family sau khi quan sát kết quả
hoặc khi một kết quả không thuận lợi.

\subsubsection{Quy tắc báo cáo hiệu ứng và độ bất định}
\label{subsubsec:statistical_reporting}

Đối với các phân tích chính thức, nghiên cứu ưu tiên báo cáo hiệu ứng tuyệt đối
không chuẩn hóa. Cohen's \(d\), Hedges' \(g\) và \(\eta^2\) không được sử dụng
như effect size chính thức của nghiên cứu.

Đối với mAP và AP, các hiệu ứng từ thang \(0\)--\(1\) được chuyển sang AP
points và báo cáo với hai chữ số thập phân. Đối với số false positive, kết quả
được báo cáo với ba chữ số thập phân. Chênh lệch
\(FAR_{\mathrm{neg}}\)
được trình bày theo percentage points.

Mọi khoảng tin cậy chính thức là khoảng tin cậy 95\% hai phía cho từng hiệu
ứng riêng. Các khoảng tin cậy này không phải simultaneous confidence
intervals và không được điều chỉnh trực tiếp bằng Holm.

\(p\)-value được báo cáo với ba chữ số thập phân; đối với giá trị nhỏ hơn
0.001 sử dụng

\begin{equation}
p<0.001
\label{eq:pvalue_reporting_lower_bound}
\end{equation}

thay vì viết \(p=0.000\).

Các thuật ngữ như ``marginally significant'' không được sử dụng; diễn giải tập
trung vào hướng hiệu ứng, độ lớn, khoảng tin cậy và bối cảnh khoa học
\cite{wasserstein2016asa}.

\subsubsection{Phân tích độ bền và độ nhạy}
\label{subsubsec:statistical_robustness}

Các phân tích độ bền và độ nhạy được xác định trước như các phép kiểm tra bổ
sung cho phân tích chính. Chúng không được sử dụng để lựa chọn lại phương pháp
phân tích sau khi đã quan sát kết quả, cũng không thay thế các phép kiểm định
chính đã được xác định trước.

Trong nghiên cứu này, hai khái niệm được phân biệt như sau. Phân tích độ bền
(\textit{robustness analysis}) đánh giá liệu hướng và độ lớn của kết luận có
duy trì khi dữ liệu được xem xét dưới các phép kiểm tra thay thế hoặc khi từng
training seed được loại tạm thời. Phân tích độ nhạy
(\textit{sensitivity analysis}) đánh giá liệu kết luận có phụ thuộc mạnh vào
một lựa chọn phân tích thống kê cụ thể hay không, chẳng hạn giả định phân phối,
cách ước lượng phương sai hoặc phương pháp kiểm định.

\paragraph{Phân tích độ bền.}

Đối với RQ2, RQ6 và RQ7, độ bền của kết luận được kiểm tra bằng exact
sign-flip analysis trên 10 hiệu ứng ghép cặp. Với \(n=10\), số cấu hình đổi
dấu có thể được xét chính xác là

\begin{equation}
N_{\mathrm{sign\text{-}flip}}
=
2^{10}
=
1024.
\label{eq:exact_sign_flip_count}
\end{equation}

Phân tích này kiểm tra liệu kết luận về hiệu ứng ghép cặp có còn được hỗ trợ
khi không phụ thuộc hoàn toàn vào giả định phân phối chuẩn của one-sample
\(t\)-test.

Bên cạnh đó, ảnh hưởng của từng training seed được đánh giá bằng
leave-one-seed-out analysis. Với 10 seed, số lần phân tích là

\begin{equation}
N_{\mathrm{LOSO}}
=
10.
\label{eq:loso_count}
\end{equation}

Trong mỗi lần, một seed được loại tạm thời và hiệu ứng được tính lại trên các
seed còn lại. Mục tiêu là kiểm tra xem hướng và độ lớn của kết luận có bị chi
phối quá mức bởi một seed riêng lẻ hay không. LOSO không được sử dụng để loại
seed khỏi phân tích chính.

\paragraph{Phân tích độ nhạy.}

Đối với RQ3, repeated-measures ANOVA với hiệu chỉnh Greenhouse--Geisser vẫn là
phân tích chính. Friedman test được sử dụng như một phân tích độ nhạy phi tham
số nhằm đánh giá xem kết luận về sự khác biệt giữa các ngân sách nhãn có phụ
thuộc mạnh vào các giả định phân phối của repeated-measures ANOVA hay không.
Kết quả Friedman chỉ có vai trò hỗ trợ diễn giải và không thay thế kết luận
chính sau khi đã quan sát dữ liệu.

Đối với RQ4, mô hình mixed-effects chính được đánh giá trước hết thông qua các
chẩn đoán gồm khả năng hội tụ, singularity, residual-versus-fitted pattern,
Q--Q plot và mức ảnh hưởng của từng training seed. Các chẩn đoán này được sử
dụng để phát hiện vi phạm hoặc bất thường của mô hình, không phải để lựa chọn
một mô hình mới sau khi đã xem kết quả.

Ngoài mô hình mixed-effects chính, một phân tích độ nhạy sử dụng covariance
estimator CR2 với clustering theo training seed và Satterthwaite degrees of
freedom được thực hiện nhằm kiểm tra mức phụ thuộc của suy luận đối với cách
ước lượng sai số chuẩn \cite{pustejovsky2018small}. Nếu kết quả từ CR2 và mô
hình chính khác nhau, sự khác biệt được báo cáo như bằng chứng về độ nhạy của
kết luận đối với phương pháp suy luận, chứ không được sử dụng để chọn phương
pháp cho kết quả thuận lợi hơn.

Như vậy, exact sign-flip và leave-one-seed-out chủ yếu đánh giá độ bền của kết
luận trước biến thiên theo seed và trước một cách suy luận thay thế, trong khi
Friedman test và CR2 được sử dụng để đánh giá độ nhạy của kết luận đối với lựa
chọn phương pháp thống kê. Tất cả các phân tích này đều giữ vai trò bổ sung và
không làm thay đổi phương pháp phân tích chính đã được xác định trước.
\subsubsection{Nguyên tắc sử dụng tập test trong phân tích thống kê}
\label{subsubsec:statistical_test_firewall}

Tập test chỉ được sử dụng sau khi kiến trúc, checkpoint-selection rule,
threshold, metric, giả thuyết thống kê, mô hình phân tích và cấu trúc kiểm
soát đa kiểm định đã được xác định.

Trình tự phân tích được tổ chức theo nguyên tắc

\begin{equation}
\begin{aligned}
&\text{Hoàn tất giao thức}
\;\longrightarrow\;
\text{Cố định cấu hình}
\;\longrightarrow\;
\text{Đánh giá test} \\
&\qquad\longrightarrow\;
\text{Phân tích thống kê}
\;\longrightarrow\;
\text{Báo cáo}.
\end{aligned}
\label{eq:test_firewall_sequence}
\end{equation}

Sau khi kết quả test được mở để phân tích, không thay đổi kiến trúc, threshold,
quy tắc chọn checkpoint, metric, giả thuyết, statistical model hoặc phương
pháp điều chỉnh đa kiểm định dựa trên kết quả test. Nguyên tắc này nhằm hạn
chế adaptive reuse của tập test và bảo toàn vai trò của test như nguồn bằng
chứng đánh giá cuối cùng độc lập \cite{dwork2015reusable}.

Tóm lại, phân tích thống kê sử dụng training seed như đơn vị suy luận, tận
dụng cấu trúc ghép cặp giữa SUP và SSL, ưu tiên effect estimate và khoảng tin
cậy, đồng thời áp dụng các phép phân tích riêng phù hợp với từng câu hỏi
nghiên cứu. Các giả thuyết thứ cấp được kiểm soát đa kiểm định theo Holm, các
phân tích độ bền không thay thế phân tích chính, và tập test chỉ được sử dụng
sau khi toàn bộ giao thức phân tích đã được cố định.
\subsection{Reproducibility}
\label{subsec:Reproducibility}

Khả năng tái lập của nghiên cứu được xem xét ở hai mức: 
(i) các nguyên tắc tái lập được xác định trước trong thiết kế thực nghiệm; và
(ii) các thông tin thực thi được ghi nhận từ môi trường huấn luyện thực tế.
Các nguyên tắc thuộc thiết kế nghiên cứu được trình bày tại đây ngay từ trước
khi huấn luyện, trong khi các thông tin phụ thuộc implementation và runtime sẽ
được hoàn thiện sau khi toàn bộ pipeline thực nghiệm được triển khai.

\subsubsection{Kiểm soát ngẫu nhiên và training seed}
\label{subsubsec:reproducibility_seed}

Mỗi điều kiện thực nghiệm được thực hiện trên cùng một danh sách gồm 10
training seed đã xác định trước. Cùng một training seed được sử dụng cho
supervised baseline và semi-supervised learning trong mỗi phép so sánh ghép
cặp giữa cùng kiến trúc và ngân sách nhãn:

\begin{equation}
s^{SUP}_{a,b}
=
s^{SSL}_{a,b}.
\label{eq:reproducibility_paired_seed}
\end{equation}

Training seed đóng vai trò là yếu tố ghép cặp và blocking factor, không phải
một biến độc lập của nghiên cứu. Danh sách seed không được lựa chọn dựa trên
hiệu năng mô hình, không thay thế seed có kết quả thấp và không thực hiện lại
một thí nghiệm chỉ vì kết quả không thuận lợi. Nếu một lần huấn luyện thất bại
do lỗi kỹ thuật, lần thực hiện lại phải sử dụng chính seed ban đầu.

Sau khi implementation hoàn tất, phần này sẽ được bổ sung bằng thông tin thực
tế về cách training seed được truyền tới các nguồn ngẫu nhiên trong pipeline,
bao gồm:

\begin{itemize}
    \item Python random: \textbf{[BỔ SUNG SAU IMPLEMENTATION]};
    \item NumPy RNG: \textbf{[BỔ SUNG SAU IMPLEMENTATION]};
    \item PyTorch CPU RNG: \textbf{[BỔ SUNG SAU IMPLEMENTATION]};
    \item PyTorch CUDA RNG: \textbf{[BỔ SUNG SAU IMPLEMENTATION]};
    \item DataLoader workers: \textbf{[BỔ SUNG SAU IMPLEMENTATION]};
    \item sampler: \textbf{[BỔ SUNG SAU IMPLEMENTATION]};
    \item stochastic augmentation: \textbf{[BỔ SUNG SAU IMPLEMENTATION]}.
\end{itemize}

\subsubsection{Cố định thành phần dữ liệu và điều kiện so sánh}
\label{subsubsec:reproducibility_data}

Tất cả các thí nghiệm sử dụng cùng phân chia train, validation và test cố định
đã được xây dựng trước. Các tập dữ liệu có nhãn tại các ngân sách
\(1\%\), \(5\%\), \(10\%\) và \(20\%\) được giữ nguyên trong toàn bộ ma trận
thực nghiệm.

Đối với cùng kiến trúc \(a\) và ngân sách nhãn \(b\), supervised baseline và
semi-supervised learning sử dụng cùng tập dữ liệu có nhãn:

\begin{equation}
L^{SUP}_{a,b}
=
L^{SSL}_{a,b}
=
L_b.
\label{eq:reproducibility_same_labeled_set}
\end{equation}

Semi-supervised learning được phép sử dụng thêm tập chưa gán nhãn \(U_b\),
trong khi supervised baseline chỉ sử dụng \(L_b\). Ngoài khác biệt có chủ
đích này, các điều kiện nền tảng của phép so sánh được giữ nhất quán theo
giao thức thực nghiệm đã trình bày.

Sau khi triển khai pipeline huấn luyện, tính nhất quán của các thành phần dữ
liệu giữa các lần chạy sẽ được xác nhận từ các artifact thực tế:

\begin{itemize}
    \item train/validation/test membership:
    \textbf{[BỔ SUNG BẰNG CHỨNG SAU]};
    \item labeled subset membership:
    \textbf{[BỔ SUNG BẰNG CHỨNG SAU]};
    \item unlabeled subset membership:
    \textbf{[BỔ SUNG BẰNG CHỨNG SAU]};
    \item mapping giữa architecture, budget và seed:
    \textbf{[BỔ SUNG BẰNG CHỨNG SAU]}.
\end{itemize}

\subsubsection{Cấu hình huấn luyện}
\label{subsubsec:reproducibility_training_config}

Các thông số khoa học chính của supervised baseline và semi-supervised
learning được xác định trước trong phần thiết kế phương pháp. Khi triển khai,
mỗi lần huấn luyện phải sử dụng đúng cấu hình tương ứng với kiến trúc, ngân
sách nhãn và phương thức học đã định nghĩa.

Các thông tin thực tế cần được lưu cho mỗi lần huấn luyện gồm:

\begin{itemize}
    \item kiến trúc:
    \textbf{[R50/Swin-T]};
    \item phương thức học:
    \textbf{[SUP/SSL]};
    \item ngân sách nhãn:
    \textbf{[1\%/5\%/10\%/20\%/100\%]};
    \item training seed:
    \textbf{[BỔ SUNG TỪ RUN]};
    \item optimizer:
    \textbf{[BỔ SUNG/XÁC NHẬN TỪ CONFIG]};
    \item learning rate:
    \textbf{[BỔ SUNG/XÁC NHẬN TỪ CONFIG]};
    \item weight decay:
    \textbf{[BỔ SUNG/XÁC NHẬN TỪ CONFIG]};
    \item tổng số optimizer updates:
    \textbf{[BỔ SUNG/XÁC NHẬN TỪ CONFIG]};
    \item validation interval:
    \textbf{[BỔ SUNG/XÁC NHẬN TỪ CONFIG]};
    \item batch composition:
    \textbf{[BỔ SUNG/XÁC NHẬN TỪ CONFIG]};
    \item AMP:
    \textbf{[BỔ SUNG/XÁC NHẬN TỪ RUNTIME]}.
\end{itemize}

Đối với SSL, các thông số liên quan đến Teacher--Student, EMA, pseudo-label
threshold, strong/weak augmentation, regression uncertainty filtering,
\(L:U\) ratio và unsupervised loss weight cũng phải được ghi nhận từ cấu hình
thực tế và đối chiếu với phương pháp đã mô tả.

\subsubsection{Môi trường phần mềm và phần cứng}
\label{subsubsec:reproducibility_environment}

Thông tin môi trường thực thi chỉ được hoàn thiện sau khi các thí nghiệm thực
tế được triển khai. Các trường tối thiểu cần báo cáo gồm:

\begin{itemize}
    \item hệ điều hành:
    \textbf{[BỔ SUNG SAU]};
    \item Python:
    \textbf{[BỔ SUNG SAU]};
    \item PyTorch:
    \textbf{[BỔ SUNG SAU]};
    \item MMDetection:
    \textbf{[BỔ SUNG SAU]};
    \item MMCV/MMEngine:
    \textbf{[BỔ SUNG SAU]};
    \item CUDA:
    \textbf{[BỔ SUNG SAU]};
    \item cuDNN:
    \textbf{[BỔ SUNG SAU]};
    \item GPU:
    \textbf{[BỔ SUNG SAU]};
    \item GPU memory:
    \textbf{[BỔ SUNG SAU]};
    \item NVIDIA driver:
    \textbf{[BỔ SUNG SAU]}.
\end{itemize}

Sau khi có môi trường thực tế, nội dung này sẽ được chuyển từ dạng placeholder
sang mô tả đầy đủ, thay vì giả định trước phiên bản phần mềm hoặc phần cứng.

\subsubsection{Thiết lập deterministic và giới hạn tái lập số học}
\label{subsubsec:reproducibility_determinism}

Các thiết lập liên quan đến deterministic execution chỉ được mô tả sau khi
pipeline huấn luyện đã được kiểm tra thực tế. Các thông tin cần xác nhận gồm:

\begin{itemize}
    \item PyTorch deterministic algorithms:
    \textbf{[BỔ SUNG SAU]};
    \item cuDNN deterministic:
    \textbf{[BỔ SUNG SAU]};
    \item cuDNN benchmark:
    \textbf{[BỔ SUNG SAU]};
    \item worker seed initialization:
    \textbf{[BỔ SUNG SAU]};
    \item AMP skipped-step behavior:
    \textbf{[BỔ SUNG SAU]};
    \item EMA update behavior khi optimizer step bị bỏ qua:
    \textbf{[BỔ SUNG SAU]}.
\end{itemize}

Ngay cả khi training seed và cấu hình được cố định, kết quả deep learning có
thể không giống nhau từng bit giữa các môi trường phần cứng hoặc phần mềm khác
nhau. Vì vậy, nghiên cứu phân biệt giữa khả năng tái tạo giao thức thực nghiệm
và bitwise numerical reproducibility.

Sau khi có bằng chứng runtime, giới hạn tái lập thực tế sẽ được mô tả tại đây:

\begin{quote}
\textbf{[BỔ SUNG SAU THỰC NGHIỆM: mô tả mức deterministic đạt được và các
nguồn nondeterminism còn lại.]}
\end{quote}

\subsubsection{Lựa chọn checkpoint và khả năng tái lập đánh giá}
\label{subsubsec:reproducibility_evaluation}

Quy tắc lựa chọn checkpoint được xác định trước và không phụ thuộc vào kết quả
test.

Đối với supervised baseline, checkpoint BEST được lựa chọn theo bbox
mAP@[0.50:0.95] trên validation. Đối với semi-supervised learning, BEST EMA
Teacher được lựa chọn bằng cùng metric trên validation và được sử dụng cho
đánh giá phát hiện cuối cùng.

LAST EMA Teacher được giữ riêng để tính

\begin{equation}
Q_{\mathrm{pseudo}}
=
PL\text{-}mAP@[0.50:0.95]
\label{eq:reproducibility_qpseudo}
\end{equation}

trên validation.

Sau khi thực nghiệm hoàn tất, phần này sẽ được bổ sung bằng bằng chứng rằng
checkpoint selection và evaluator đã thực sự tuân theo các quy tắc trên:

\begin{itemize}
    \item checkpoint BEST của SUP:
    \textbf{[BỔ SUNG SAU]};
    \item checkpoint BEST EMA Teacher:
    \textbf{[BỔ SUNG SAU]};
    \item checkpoint LAST EMA Teacher:
    \textbf{[BỔ SUNG SAU]};
    \item evaluator/config đánh giá:
    \textbf{[BỔ SUNG SAU]}.
\end{itemize}

\subsubsection{Thông tin cần hoàn thiện sau thực nghiệm}
\label{subsubsec:reproducibility_post_execution}

Sau khi toàn bộ các thí nghiệm chính được triển khai, phần
Reproducibility sẽ được rà soát lại dựa trên bằng chứng thực tế. Các nội dung
cần được bổ sung tối thiểu gồm:

\begin{enumerate}
    \item môi trường phần mềm và phần cứng thực tế;
    \item cách khởi tạo và truyền training seed trong pipeline;
    \item deterministic settings thực tế;
    \item cấu hình DataLoader, sampler và augmentation;
    \item AMP và EMA runtime behavior;
    \item cấu hình huấn luyện thực tế của R50 và Swin-T;
    \item checkpoint-selection evidence;
    \item evaluator configuration;
    \item thông tin về các lần thực hiện lại do lỗi kỹ thuật, nếu có;
    \item các sai khác giữa thiết kế dự kiến và implementation thực tế, nếu có.
\end{enumerate}

Chỉ các thông tin được xác nhận từ implementation và các lần chạy thực tế mới
được chuyển thành phát biểu khẳng định trong phiên bản cuối của luận văn.

% ============================================================
% [PENDING_REPRODUCIBILITY_AFTER_EXPERIMENT]
% Revisit after implementation/training/evaluation are completed.
% Required evidence: environment, seed propagation, deterministic settings,
% DataLoader/sampler, AMP/EMA runtime behavior, checkpoints, evaluator, logs.
% ============================================================

\subsection{Threats to Validity}
\label{subsec:Threats_to_Validity}

Mặc dù nghiên cứu được xây dựng theo một giao thức thực nghiệm có kiểm soát,
một số giới hạn vẫn có thể ảnh hưởng đến cách diễn giải kết quả. Các mối đe
dọa đối với tính hợp lệ được xem xét chủ yếu theo bốn nhóm: giới hạn liên quan
đến dữ liệu và annotation, tính hợp lệ nội tại và các giả định của
semi-supervised learning, khả năng khái quát hóa, và tính hợp lệ của suy luận
thống kê cũng như quy trình đánh giá. Việc trình bày các giới hạn này nhằm xác
định rõ phạm vi mà các kết luận của nghiên cứu có thể được hỗ trợ và tránh
suy rộng quá mức.

\paragraph{Giới hạn liên quan đến dữ liệu và annotation.}

Nghiên cứu sử dụng VinDr-CXR làm nguồn dữ liệu trung tâm. Trong phạm vi thực
nghiệm có kiểm soát, nghiên cứu sử dụng 4.894 ảnh, gồm toàn bộ 4.394 ảnh có
bất thường và 500 ảnh \textit{No Finding}. Việc giới hạn số lượng ảnh
\textit{No Finding} giúp kiểm soát quy mô thực nghiệm và tránh để số lượng ảnh
âm tính lớn trong phân bố gốc chi phối quá mức tập dữ liệu nghiên cứu. Tuy
nhiên, hệ quả là phân bố của phạm vi thực nghiệm không tái tạo đầy đủ phân bố
của toàn bộ VinDr-CXR. Do đó, các kết quả, đặc biệt là false positive trên ảnh
âm tính, cần được diễn giải trong phạm vi tập dữ liệu đã xây dựng thay vì xem
là ước lượng trực tiếp cho toàn bộ phân bố VinDr-CXR hoặc thực hành lâm sàng.

Một giới hạn quan trọng khác là khả năng kiểm soát data leakage ở cấp bệnh
nhân. Trong dữ liệu công bố, \texttt{image\_id} là định danh ảnh và không thể
được diễn giải như định danh bệnh nhân \cite{nguyen2022vindr}. Dữ liệu khả
dụng không cung cấp định danh nhóm bệnh nhân phù hợp để xác định liệu nhiều ảnh
có thuộc cùng một bệnh nhân sau quá trình khử định danh hay không. Vì vậy,
nghiên cứu chỉ xác nhận sự tách biệt giữa train, validation và test ở cấp định
danh ảnh và annotation khả dụng, chứ không thể khẳng định tuyệt đối rằng
patient-level leakage đã được loại trừ.

Reference annotations cũng có thể chứa bất định giữa người đọc. Fleiss' kappa
được sử dụng để đặc trưng mức độ đồng thuận về trạng thái present/absent của
từng lớp ở cấp ảnh \cite{fleiss1971measuring}, nhưng không đánh giá trực tiếp
sự đồng thuận về vị trí, kích thước hoặc ranh giới bounding box giữa các bác
sĩ. Vì vậy, kết quả kappa không được xem là bằng chứng rằng annotation hoàn
toàn không có bất định. Các kết quả detection cần được diễn giải với lưu ý về
inter-reader variability của nguồn annotation \cite{le2023learning}.

\paragraph{Giới hạn về tính hợp lệ nội tại và các giả định của semi-supervised learning.}

Đối với các phép so sánh SUP--SSL, nghiên cứu giảm confounding bằng cách giữ
cùng tập dữ liệu có nhãn, cùng kiến trúc, cùng training seed, cùng biểu diễn
ảnh đầu vào, cùng họ detector, cùng ngân sách optimizer updates và cùng giao
thức đánh giá tại mỗi điều kiện ghép cặp. Khác biệt có chủ đích là SSL được
phép khai thác thêm tập dữ liệu chưa gán nhãn. Cấu trúc này giúp cô lập tương
đối ảnh hưởng của dữ liệu chưa gán nhãn, nhưng không loại bỏ hoàn toàn biến
thiên do quá trình tối ưu deep learning. Vì vậy, nghiên cứu sử dụng nhiều
training seeds và thực hiện so sánh ghép cặp theo seed thay vì dựa trên một
lần huấn luyện đơn lẻ.

Đối với hai kiến trúc được khảo sát, ResNet-50 và Swin-T sử dụng các giao thức
tối ưu hóa đặc thù theo kiến trúc. Do đó, phân tích architecture-dependent SSL
effect không thể được diễn giải như một phép can thiệp trong đó chỉ duy nhất
backbone thay đổi. Kết quả chỉ phản ánh sự khác biệt về mức lợi ích của SSL
giữa hai kiến trúc dưới các giao thức huấn luyện đặc thù theo kiến trúc đã được
xác định trước.

Semi-supervised learning cũng dựa trên các giả định về sự liên hệ giữa dữ liệu
có nhãn và chưa gán nhãn. Mặc dù \(L_b\) và \(U_b\) đều được lấy từ cùng nguồn
VinDr-CXR, cùng modality và cùng không gian lớp, điều này không bảo đảm hai
phân phối hoàn toàn đồng nhất. Khác biệt về mức độ bệnh, chất lượng ảnh, thiết
bị chụp hoặc tần suất lớp vẫn có thể tạo ra distribution shift.

Bên cạnh đó, teacher--student pseudo-labeling có thể chịu ảnh hưởng của
pseudo-label noise và confirmation bias. Teacher có thể sinh pseudo-label sai
về lớp, vị trí hoặc bỏ sót tổn thương, và Student có thể tiếp tục học từ các
tín hiệu này. Confidence filtering, regression-uncertainty filtering và EMA
Teacher được sử dụng để giảm rủi ro, nhưng không bảo đảm loại bỏ hoàn toàn
pseudo-label noise. Vì vậy, nếu SSL cải thiện hoặc suy giảm hiệu năng, kết quả
không được tự động quy cho một cơ chế duy nhất.

\paragraph{Giới hạn về khả năng khái quát hóa.}

Khả năng khái quát hóa của nghiên cứu bị giới hạn bởi phạm vi dữ liệu, kiến
trúc và phương pháp SSL được khảo sát. Nghiên cứu chỉ sử dụng một nguồn dữ liệu
chính là VinDr-CXR và không thực hiện external validation trên bệnh viện, thiết
bị chụp, quần thể bệnh nhân hoặc bộ dữ liệu CXR khác. Do đó, nghiên cứu không
trực tiếp kiểm chứng khả năng khái quát hóa dưới domain shift.

Hai kiến trúc được khảo sát đều thuộc cùng họ Faster R-CNN, gồm ResNet-50 +
FPN và Swin-T + FPN. Vì vậy, các kết luận về architecture-dependent SSL effect
không tự động áp dụng cho one-stage detectors, các backbone khác, foundation
models hoặc các detector paradigm khác.

Tương tự, phương pháp semi-supervised được triển khai theo một khung
teacher--student pseudo-labeling đại diện. Nghiên cứu không được thiết kế như
một benchmark toàn diện của mọi phương pháp semi-supervised object detection
hoặc mọi chiến lược label-efficient learning. Vì vậy, kết quả chỉ hỗ trợ kết
luận trong phạm vi giao thức SSL được khảo sát và không chứng minh rằng
teacher--student pseudo-labeling nói chung vượt trội hơn mọi hướng thay thế.

Khái niệm ``lớp hiếm'' trong nghiên cứu cũng chỉ phản ánh mức hỗ trợ thấp của
một lớp trong tập huấn luyện theo tiêu chí đã xác định trước. Thuật ngữ này
không được suy rộng thành kết luận rằng một bất thường là hiếm về mặt dịch tễ
hoặc lâm sàng.

\paragraph{Giới hạn của suy luận thống kê và quy trình đánh giá.}

Đơn vị suy luận của nghiên cứu là training seed/trained run, với 10 training
seeds cho mỗi điều kiện chính. Ảnh, bounding box và detection không được xem là
các lần lặp độc lập của phương pháp huấn luyện. Việc sử dụng 10 seeds và
paired-seed analysis giúp phản ánh training variability tốt hơn một single-run
comparison, nhưng số lượng này vẫn hữu hạn và có thể làm giảm độ chính xác của
ước lượng đối với các hiệu ứng nhỏ hoặc các mô hình thống kê phức tạp.

Các phép phân tích chính cũng dựa trên các giả định thống kê tương ứng.
Repeated-measures ANOVA sử dụng hiệu chỉnh Greenhouse--Geisser
\cite{greenhouse1959profile}; phân tích mối liên hệ giữa chất lượng pseudo-label
và SSL gain sử dụng linear mixed-effects model với REML và Kenward--Roger
inference \cite{kenward1997small}. Exact sign-flip, Friedman test,
leave-one-seed-out và CR2 được sử dụng như các kiểm tra độ bền hoặc độ nhạy,
không được sử dụng để lựa chọn hậu nghiệm phương pháp cho kết quả thuận lợi
hơn. Các giả thuyết thứ cấp chính thức được kiểm soát đa kiểm định bằng thủ tục
Holm \cite{holm1979sequentially}.

Đối với các lớp có mức hỗ trợ thấp, class-wise AP có thể biến thiên mạnh hơn do
số ground-truth instances hạn chế. Vì vậy, phân tích lớp hiếm được giữ ở vai
trò khám phá và tập trung vào effect estimate cùng khoảng tin cậy, thay vì xây
dựng thêm một hệ giả thuyết confirmatory riêng.

Ngoài ra, validation được sử dụng cho checkpoint selection,
\(Q_{\mathrm{pseudo}}\) và các biến thể ablation đã xác định trước, trong khi
test chỉ được sử dụng cho đánh giá cuối cùng của các điều kiện chính. Ground
truth ẩn của \(U_b\) không được sử dụng cho pseudo-label filtering, threshold
tuning, augmentation selection, EMA tuning, checkpoint selection hoặc tính
\(Q_{\mathrm{pseudo}}\). Cách tổ chức này giúp hạn chế adaptive evaluation và
bảo toàn vai trò của test như nguồn bằng chứng cuối cùng độc lập. Tuy nhiên,
do validation thực hiện nhiều vai trò đã xác định trước, các kết luận chính về
hiệu năng vẫn cần dựa trên test theo giao thức đã quy định.

Chất lượng pseudo-label được đánh giá bằng
\(Q_{\mathrm{pseudo}}=PL\text{-}mAP@[0.50:0.95]\) trên validation cố định.
Phép đo này không sử dụng ground truth ẩn của \(U_b\), vì vậy không được giả
định là đại diện hoàn toàn cho chất lượng pseudo-label trên toàn bộ dữ liệu
unlabeled đã dùng trong huấn luyện. Đây là một giới hạn được chấp nhận có chủ
đích để tránh vi phạm điều kiện unlabeled của thiết kế semi-supervised.

% TODO_IMPLEMENTATION_VALIDITY_AFTER_EXPERIMENT
% ============================================================
% [PENDING_IMPLEMENTATION_VALIDITY_AFTER_EXPERIMENT]
% Sau implementation/training/evaluation, rà lại:
% - seed propagation;
% - SUP/SSL data membership;
% - optimizer/schedule/update budget;
% - DataLoader/sampler;
% - augmentation;
% - EMA/AMP runtime behavior;
% - pseudo-label thresholds/filtering;
% - checkpoint selection;
% - evaluator configuration;
% - sai khác giữa Methodology và implementation thực tế.
% Chỉ bổ sung những gì có evidence thực tế.
% ============================================================

Tóm lại, thiết kế ghép cặp, nhiều training seeds, metric hierarchy, kiểm soát
đa kiểm định và quy tắc sử dụng validation/test giúp giảm nhiều nguồn đe dọa
đối với tính hợp lệ. Tuy nhiên, các kết luận của nghiên cứu vẫn cần được giới
hạn bởi phạm vi dữ liệu, khả năng kiểm soát leakage ở cấp bệnh nhân, bất định
của annotation, pseudo-label noise, số lượng kiến trúc và phương pháp SSL được
khảo sát, cũng như số lượng training seeds hữu hạn.
\section{Results}\label{sec:results}

\subsection{Kết quả chẩn đoán dữ liệu trước huấn luyện}
\label{subsec:results_pretraining_dataset_diagnostics}

Áp dụng các phép chẩn đoán được trình bày tại
Mục~\ref{subsec:pretraining_dataset_diagnostics} trên tập huấn luyện \(T\)
cho thấy dữ liệu gồm 3.426 ảnh, 25.260 bounding-box annotations và 14 lớp
bất thường. Trong số này, 3.076 ảnh có ít nhất một ground-truth bounding
box và 350 ảnh là \textit{No Finding}. Các kết quả trong mục này chỉ mô
tả những đặc trưng quan sát được của dữ liệu trước huấn luyện; mối liên hệ
giữa các đặc trưng này với hiệu năng của mô hình được đánh giá ở các thí
nghiệm tiếp theo.

\paragraph{Phân bố lớp, mức độ mất cân bằng và các lớp có mức hỗ trợ thấp.}

Mức hỗ trợ ở cấp độ ảnh khác nhau rõ rệt giữa 14 lớp. Số ảnh chứa từng
lớp dao động từ 66 đến 2.144 ảnh, với trung vị bằng 648 ảnh. Theo tỷ số mất cân bằng được định nghĩa tại
Phương trình~\ref{eq:diag_imbalance_ratio}, giá trị quan sát được là
\(R_{\max/\min}=2144/66=32{,}48\).

Hệ số biến thiên của mức hỗ trợ ở cấp ảnh đạt
\(CV_{\mathrm{img}}=0{,}803\).

Phân bố chi tiết của số ảnh và số bounding-box annotations theo lớp được
trình bày tại Bảng~\ref{tab:result_train_class_distribution}. Các lớp
được sắp xếp theo số ảnh hỗ trợ giảm dần để làm rõ mức độ không đồng đều
giữa các lớp.

\begin{table*}[htbp]
    \centering
    \caption{Phân bố ảnh và bounding-box annotations theo lớp trong tập
    huấn luyện}
    \label{tab:result_train_class_distribution}
    \small
    \setlength{\tabcolsep}{4pt}
    \renewcommand{\arraystretch}{1.10}

    \begin{tabular}{lrrrrl}
        \toprule
        \textbf{Lớp bất thường} &
        \textbf{Số ảnh} &
        \textbf{Tỷ lệ ảnh (\%)} &
        \textbf{Số bbox} &
        \textbf{Bbox/ảnh có lớp} &
        \textbf{Trạng thái} \\
        \midrule

        Aortic enlargement
            & 2.144 & 62,58 & 5.021 & 2,34 & Non-rare \\

        Cardiomegaly
            & 1.607 & 46,91 & 3.827 & 2,38 & Non-rare \\

        Pleural thickening
            & 1.384 & 40,40 & 3.415 & 2,47 & Non-rare \\

        Pulmonary fibrosis
            & 1.130 & 32,98 & 3.271 & 2,89 & Non-rare \\

        Lung Opacity
            & 923 & 26,94 & 1.728 & 1,87 & Non-rare \\

        Other lesion
            & 791 & 23,09 & 1.501 & 1,90 & Non-rare \\

        Pleural effusion
            & 720 & 21,02 & 1.779 & 2,47 & Non-rare \\

        Nodule/Mass
            & 576 & 16,81 & 1.790 & 3,11 & Non-rare \\

        Infiltration
            & 428 & 12,49 & 879 & 2,05 & Non-rare \\

        Calcification
            & 314 & 9,17 & 665 & 2,12 & Non-rare \\

        ILD
            & 268 & 7,82 & 658 & 2,46 & Non-rare \\

        Consolidation
            & 246 & 7,18 & 390 & 1,59 & Non-rare \\

        Atelectasis
            & 130 & 3,79 & 194 & 1,49 & Rare \\

        Pneumothorax
            & 66 & 1,93 & 142 & 2,15 & Rare \\

        \bottomrule
    \end{tabular}
\end{table*}

Theo tiêu chí \(P_c^{\mathrm{img}}<5\%\) đã được xác định tại
Phương trình~\ref{eq:diag_rare_definition}, có 2/14 lớp thỏa điều kiện
mức hỗ trợ thấp. \textit{Atelectasis} xuất hiện trên 130 ảnh, tương ứng
3,79\%, trong khi \textit{Pneumothorax} xuất hiện trên 66 ảnh, tương ứng
1,93\%. Hai lớp này cũng là hai lớp có số ảnh hỗ trợ thấp nhất trong tập
huấn luyện.

Hình~\ref{fig:result_class_distribution} trực quan hóa đồng thời mức hỗ
trợ ở cấp độ ảnh và số lượng bounding-box annotations. Sự khác biệt giữa
hai biểu diễn cho thấy số bounding box không hoàn toàn tỷ lệ với số ảnh
có lớp, do một ảnh có thể chứa nhiều annotations thuộc cùng một lớp.

\begin{figure*}[htbp]
    \centering
    \includegraphics[width=\textwidth]
    {z_class_distribution_training_set.png}
    \caption{Phân bố số ảnh và số bounding-box annotations theo 14 lớp
    bất thường trong tập huấn luyện.}
    \label{fig:result_class_distribution}
\end{figure*}

Trạng thái rare ở đây chỉ phản ánh mức hỗ trợ của lớp trong tập dữ liệu
nghiên cứu theo tiêu chí đã xác định trước; kết quả này không được diễn
giải là độ hiếm về mặt dịch tễ hoặc lâm sàng.

\paragraph{Mật độ annotation, hình học và kích thước bounding box.}

Trên toàn bộ 3.426 ảnh của tập huấn luyện, trung bình có 7,37
bounding-box annotations trên mỗi ảnh; trung vị bằng 6, phân vị thứ 95
bằng 17 và giá trị lớn nhất bằng 48. Khi chỉ xét 3.076 ảnh có ít nhất một
ground-truth bounding box, giá trị trung bình tăng lên 8,21,
trung vị bằng 7 và phân vị thứ 95 bằng 18, trong khi giá trị lớn nhất vẫn
bằng 48.

\begin{table}[htbp]
    \centering
    \caption{Phân bố số bounding-box annotations trên mỗi ảnh}
    \label{tab:result_bbox_count_per_image}
    \small

    \begin{tabular}{lrrrr}
        \toprule
        \textbf{Phạm vi} &
        \textbf{Mean} &
        \textbf{Median} &
        \textbf{P95} &
        \textbf{Max} \\
        \midrule

        Toàn bộ tập huấn luyện
            & 7,37 & 6 & 17 & 48 \\

        Ảnh có ít nhất một bbox
            & 8,21 & 7 & 18 & 48 \\

        \bottomrule
    \end{tabular}
\end{table}

Sự khác biệt giữa hai phạm vi chủ yếu xuất phát từ 350 ảnh
\textit{No Finding} có số bounding box bằng 0. Các giá trị trong
Bảng~\ref{tab:result_bbox_count_per_image} phản ánh số bản ghi annotation,
không phải số tổn thương lâm sàng độc lập trên mỗi ảnh.

Đối với 25.260 bounding-box annotations, diện tích chuẩn hóa có trung
bình \(0{,}0304\), trung vị \(0{,}0149\), phân vị thứ 95
\(0{,}1003\) và giá trị lớn nhất \(0{,}9384\). Tỷ lệ khung hình có
trung bình \(1{,}4099\), trung vị \(0{,}9853\), phân vị thứ 95
\(3{,}5423\) và giá trị lớn nhất bằng 45.

\begin{table}[htbp]
    \centering
    \caption{Một số thống kê hình học chính của bounding box}
    \label{tab:result_bbox_geometry}
    \small

    \begin{tabular}{lrrrr}
        \toprule
        \textbf{Đại lượng} &
        \textbf{Mean} &
        \textbf{Median} &
        \textbf{P95} &
        \textbf{Max} \\
        \midrule

        Diện tích chuẩn hóa
            & 0,0304 & 0,0149 & 0,1003 & 0,9384 \\

        Tỷ lệ khung hình
            & 1,4099 & 0,9853 & 3,5423 & 45,00 \\

        \bottomrule
    \end{tabular}
\end{table}

Theo các ranh giới diện tích chuẩn hóa đã xác định trước, 9.372
annotations thuộc nhóm Small, 14.612 thuộc nhóm Medium và 1.276 thuộc
nhóm Large. Tỷ lệ tương ứng là 37,10\%, 57,85\% và 5,05\%.

\begin{table}[htbp]
    \centering
    \caption{Phân bố bounding box theo các nhóm diện tích chuẩn hóa}
    \label{tab:result_bbox_size_categories}
    \small

    \begin{tabular}{lrr}
        \toprule
        \textbf{Nhóm} &
        \textbf{Số bbox} &
        \textbf{Tỷ lệ (\%)} \\
        \midrule

        Small  & 9.372  & 37,10 \\
        Medium & 14.612 & 57,85 \\
        Large  & 1.276  & 5,05 \\

        \bottomrule
    \end{tabular}
\end{table}

Hình~\ref{fig:result_bbox_distribution} cho thấy cả phân bố liên tục của
diện tích chuẩn hóa và kết quả phân nhóm theo các ranh giới chính. Ở
panel A, phần lớn annotations tập trung về phía các giá trị diện tích
chuẩn hóa thấp và phân bố có phần đuôi kéo dài tới các bounding box chiếm
tỷ lệ lớn của ảnh. Trục tung sử dụng thang logarithm để các vùng có số
lượng annotation khác nhau nhiều bậc vẫn có thể được quan sát đồng thời.

\begin{figure*}[htbp]
    \centering
    \includegraphics[width=\textwidth]
    {z_bbox_distribution.png}
    \caption{Phân bố kích thước bounding box trong tập huấn luyện.
    Panel A biểu diễn phân bố liên tục của diện tích bounding box chuẩn
    hóa, với hai ranh giới chính tại \(0{,}01\) và \(0{,}10\).
    Panel B biểu diễn số lượng và tỷ lệ bounding-box annotations thuộc
    các nhóm Small, Medium và Large. Các nhóm này dựa trên diện tích
    tương đối so với ảnh và không phải các nhóm Small/Medium/Large theo
    diện tích pixel tuyệt đối của COCO.}
    \label{fig:result_bbox_distribution}
\end{figure*}

Các tỷ lệ trên mô tả trực tiếp thành phần annotation theo định nghĩa diện
tích chuẩn hóa đã sử dụng trong nghiên cứu; chúng chưa cho phép kết luận
bounding box thuộc một nhóm kích thước cụ thể sẽ có hiệu năng phát hiện
thấp hơn hoặc cao hơn.

\paragraph{Phân bố không gian của bounding box và ảnh âm tính.}

Phân bố tâm của 25.260 bounding-box annotations sau khi chuẩn hóa tọa độ
được trình bày tại Hình~\ref{fig:result_bbox_location_heatmap}. Kết quả
cho thấy các tâm annotation không phân bố đồng đều trên toàn bộ miền ảnh
chuẩn hóa mà có mật độ khác nhau giữa các vùng.

\begin{figure}[htbp]
    \centering
    \includegraphics[width=\linewidth]
    {z_bbox_location_heatmap.png}
    \caption{Phân bố tâm bounding-box annotations trên lưới
    \(50\times50\) sau khi chuẩn hóa tọa độ ảnh. Mỗi bounding box đóng
    góp một quan sát tại vị trí tâm tương ứng.}
    \label{fig:result_bbox_location_heatmap}
\end{figure}

Heatmap trên chỉ phản ánh phân bố của tâm các bounding-box annotations.
Do đó, sự khác biệt mật độ giữa các vùng không được diễn giải thành mật
độ của các tổn thương lâm sàng độc lập hoặc bằng chứng về một quan hệ giải
phẫu nhân quả.

Đối với thành phần ảnh âm tính, tập huấn luyện chứa 350 ảnh
\textit{No Finding} trên tổng số 3.426 ảnh, tương ứng \(10{,}22\%\).
Phần còn lại gồm 3.076 ảnh có ít nhất một ground-truth bounding box,
tương ứng \(89{,}78\%\).

\begin{figure}[htbp]
    \centering
    \includegraphics[width=\linewidth]
    {z_negative_image_distribution.png}
    \caption{Phân bố ảnh có ít nhất một ground-truth bounding box và ảnh
    \textit{No Finding} trong tập huấn luyện.}
    \label{fig:result_negative_distribution}
\end{figure}

Kết quả này mô tả thành phần zero-GT của tập huấn luyện. Ảnh
\textit{No Finding} tiếp tục được xem là ảnh âm tính không có bounding
box, không phải lớp phát hiện thứ 15.

\paragraph{Mức độ đại diện của các tập có nhãn theo ngân sách.}

Bốn ngân sách nhãn \(1\%\), \(5\%\), \(10\%\) và \(20\%\) đều duy trì
sự hiện diện của đủ 14 lớp bất thường. Các chỉ số tổng hợp về độ lệch
phân bố so với toàn bộ tập huấn luyện được trình bày trong
Bảng~\ref{tab:result_labeled_budget_representativeness}.

\begin{table*}[htbp]
    \centering
    \caption{Mức bao phủ lớp và độ lệch phân bố của các tập có nhãn so
    với tập huấn luyện}
    \label{tab:result_labeled_budget_representativeness}
    \small
    \setlength{\tabcolsep}{6pt}
    \renewcommand{\arraystretch}{1.10}

    \begin{tabular}{lrrrrr}
        \toprule
        \textbf{Ngân sách} &
        \textbf{Số ảnh} &
        \textbf{No Finding} &
        \textbf{Bao phủ lớp} &
        \(\mathbf{D_{\max}}\) &
        \(\mathbf{D_{\mathrm{mean}}}\) \\
        \midrule

        1\%
            & 34  & 3  & 14/14 & 0,01298 & 0,00699 \\

        5\%
            & 171 & 17 & 14/14 & 0,00286 & 0,00154 \\

        10\%
            & 343 & 35 & 14/14 & 0,00128 & 0,00075 \\

        20\%
            & 685 & 70 & 14/14 & 0,00066 & 0,00034 \\

        \bottomrule
    \end{tabular}
\end{table*}

Cả \(D_{\max}\) và \(D_{\mathrm{mean}}\) đều giảm khi ngân sách nhãn
tăng. Cụ thể, \(D_{\max}\) giảm từ khoảng 0,01298 ở mức 1\% xuống
0,00066 ở mức 20\%, trong khi \(D_{\mathrm{mean}}\) giảm từ khoảng
0,00699 xuống 0,00034. Điều này cho thấy tỷ lệ hiện diện lớp trong các
tập có nhãn tiến gần hơn đến tỷ lệ của toàn bộ tập huấn luyện khi quy mô
\(L_b\) tăng.

Hình~\ref{fig:result_labeled_budget_representativeness} bổ sung thông tin
theo từng lớp bằng cách biểu diễn độ lệch tuyệt đối \(D_{b,c}\) theo đơn
vị điểm phần trăm.

\begin{figure*}[htbp]
    \centering
    \includegraphics[width=\textwidth]
    {z_labeled_budget_representativeness.png}
    \caption{Độ lệch tuyệt đối của tỷ lệ hiện diện từng lớp trong các tập
    có nhãn so với toàn bộ tập huấn luyện. Mỗi hàng tương ứng với một
    ngân sách nhãn, mỗi cột tương ứng với một lớp bất thường và cường độ
    màu biểu diễn \(100\times D_{b,c}\) theo đơn vị điểm phần trăm.}
    \label{fig:result_labeled_budget_representativeness}
\end{figure*}

Heatmap cho thấy độ lệch theo từng lớp lớn nhất tập trung ở ngân sách
1\%, trong khi các hàng tương ứng với 10\% và 20\% có giá trị nhìn chung
nhỏ hơn. Hình này bổ sung cho các chỉ số tổng hợp
\(D_{\max}\) và \(D_{\mathrm{mean}}\) bằng cách cho thấy sai lệch không
phân bố đồng đều giữa các lớp.

Một giới hạn quan trọng ở mức ngân sách thấp nhất là dù
\(L_{1\%}\) vẫn bao phủ đủ 14 lớp, hai lớp rare
\textit{Atelectasis} và \textit{Pneumothorax} mỗi lớp chỉ có một ảnh có
nhãn. Vì vậy, việc đạt 14/14 class coverage không đồng nghĩa với việc
các lớp nhận được mức hỗ trợ tuyệt đối tương đương nhau.

\paragraph{Cấu trúc đa nhãn của tập huấn luyện.}

\textit{Label cardinality} quan sát được dao động từ 0 đến 10 lớp trên
mỗi ảnh. Giá trị trung bình là 3,131 lớp/ảnh, độ lệch chuẩn mẫu là
1,932, trung vị bằng 3 và phân vị thứ 95 bằng 7.

Phân bố đầy đủ được trình bày tại
Bảng~\ref{tab:result_label_cardinality}.

\begin{table}[htbp]
    \centering
    \caption{Phân bố số lớp bất thường duy nhất trên mỗi ảnh}
    \label{tab:result_label_cardinality}
    \small

    \begin{tabular}{rrr}
        \toprule
        \textbf{Số lớp/ảnh} &
        \textbf{Số ảnh} &
        \textbf{Tỷ lệ (\%)} \\
        \midrule

        0  & 350 & 10,22 \\
        1  & 225 & 6,57 \\
        2  & 839 & 24,49 \\
        3  & 702 & 20,49 \\
        4  & 530 & 15,47 \\
        5  & 372 & 10,86 \\
        6  & 209 & 6,10 \\
        7  & 130 & 3,79 \\
        8  & 51  & 1,49 \\
        9  & 14  & 0,41 \\
        10 & 4   & 0,12 \\

        \bottomrule
    \end{tabular}
\end{table}

Hình~\ref{fig:result_label_cardinality_distribution} cho thấy
cardinality bằng 2 và 3 là hai giá trị xuất hiện nhiều nhất, lần lượt
với 839 ảnh (24,49\%) và 702 ảnh (20,49\%).

\begin{figure*}[htbp]
    \centering
    \includegraphics[width=0.92\textwidth]
    {z_label_cardinality_distribution.png}
    \caption{Phân bố \textit{label cardinality} trong tập huấn luyện.
    Trục hoành biểu diễn số lớp bất thường duy nhất xuất hiện trên mỗi
    ảnh; trục tung biểu diễn số ảnh tương ứng. Cardinality bằng 0 tương
    ứng với các ảnh \textit{No Finding}.}
    \label{fig:result_label_cardinality_distribution}
\end{figure*}

Ngoài các ảnh có 2--3 lớp, dữ liệu vẫn chứa các trường hợp có cardinality
cao hơn: 130 ảnh có 7 lớp, 51 ảnh có 8 lớp, 14 ảnh có 9 lớp và 4 ảnh có
10 lớp bất thường duy nhất. Kết quả này cho thấy cấu trúc nhãn ở cấp ảnh
không thể được mô tả đầy đủ như một bài toán mỗi ảnh chỉ có một lớp.

Cần phân biệt kết quả này với mật độ bounding box. Cardinality phản ánh
số lớp khác nhau trên ảnh, trong khi một lớp có thể có nhiều
bounding-box annotations.

Phân tích đồng xuất hiện giữa 14 lớp tạo ra 91 cặp lớp không thứ tự.
Theo định nghĩa tại
Phương trình~\ref{eq:diag_class_cooccurrence}, mỗi giá trị
\(N_{cd}\) là số ảnh duy nhất trong tập huấn luyện có đồng thời hai lớp
\(c\) và \(d\). Toàn bộ 91 giá trị được trực quan hóa tại
Hình~\ref{fig:result_class_cooccurrence_heatmap}.

\begin{figure*}[htbp]
    \centering
    \includegraphics[width=0.92\textwidth]
    {z_class_cooccurrence_heatmap.png}
    \caption{Ma trận đồng xuất hiện giữa 14 lớp bất thường trong tập
    huấn luyện. Mỗi ô ngoài đường chéo biểu diễn số ảnh duy nhất có đồng
    thời hai lớp tương ứng. Mỗi ảnh chỉ được tính một lần cho một cặp lớp,
    bất kể số bounding box của từng lớp trên ảnh. Đường chéo được để trống
    vì không xét đồng xuất hiện của một lớp với chính nó.}
    \label{fig:result_class_cooccurrence_heatmap}
\end{figure*}

Ma trận có tính đối xứng vì số ảnh chứa đồng thời lớp \(c\) và lớp \(d\)
không phụ thuộc vào thứ tự của hai lớp. Cặp
\textit{Aortic enlargement}--\textit{Cardiomegaly} có 1.373 ảnh đồng
xuất hiện và là cặp có số ảnh đồng xuất hiện nổi bật nhất trong ma trận.
Ngoài cặp này, heatmap cho thấy mức đồng xuất hiện giữa các cặp lớp có
sự khác biệt đáng kể trên tập huấn luyện, phản ánh cấu trúc đa nhãn không
đồng đều của dữ liệu.

Các giá trị \(N_{cd}\) chỉ phản ánh số ảnh có hai lớp cùng hiện diện trong
phạm vi tập dữ liệu nghiên cứu. Một giá trị đồng xuất hiện lớn không chứng
minh quan hệ nhân quả, cơ chế bệnh sinh hoặc mối liên hệ lâm sàng giữa hai
bất thường. Kết quả này cũng không được sử dụng để tự động thay đổi
sampling, hàm mất mát, kiến trúc hoặc cấu hình huấn luyện.

\paragraph{Độ nhạy đối với các ngưỡng chẩn đoán.}

Kết quả phân tích độ nhạy được tổng hợp tại
Hình~\ref{fig:result_threshold_sensitivity}.

Đối với tiêu chí rare, ngưỡng 1\% không xác định lớp nào là rare; tại
ngưỡng chính 5\% có 2/14 lớp, tương ứng 14,29\%; và tại ngưỡng 10\% có
5/14 lớp, tương ứng 35,71\%. Năm lớp có trạng thái rare/non-rare thay đổi
ít nhất một lần trong khoảng ngưỡng khảo sát gồm
\textit{Atelectasis}, \textit{Calcification},
\textit{Consolidation}, \textit{ILD} và
\textit{Pneumothorax}.

Khi ranh giới Small thay đổi từ 0,5\% sang 1\% và 2\% diện tích ảnh,
trong khi ranh giới Large được giữ ở 10\%, tỷ lệ bounding box thuộc nhóm
Small lần lượt là 23,52\%, 37,10\% và 58,90\%.

Ngược lại, khi ranh giới Small được giữ ở 1\% và ranh giới Large thay
đổi từ 5\% sang 10\% và 20\%, tỷ lệ bounding box thuộc nhóm Large lần
lượt là 21,78\%, 5,05\% và 0,95\%.

\begin{figure*}[htbp]
    \centering
    \includegraphics[width=\textwidth]
    {z_threshold_sensitivity.png}
    \caption{Kết quả phân tích độ nhạy của các ngưỡng chẩn đoán.
    Panel (a) biểu diễn tỷ lệ lớp được xác định là rare khi thay đổi
    ngưỡng prevalence; panel (b) biểu diễn tỷ lệ bounding box thuộc nhóm
    Small khi thay đổi ranh giới Small; panel (c) biểu diễn tỷ lệ
    bounding box thuộc nhóm Large khi thay đổi ranh giới Large. Các điểm
    được đánh dấu \textit{Primary} tương ứng với các ngưỡng chính đã
    được xác định trước.}
    \label{fig:result_threshold_sensitivity}
\end{figure*}

Kết quả cho thấy tỷ lệ các nhóm được tạo bởi những định nghĩa vận hành có
thể thay đổi theo ranh giới sử dụng. Theo thiết kế phân tích đã trình bày
tại Mục~\ref{subsec:pretraining_dataset_diagnostics}, các kết quả độ nhạy
này chỉ được sử dụng để đánh giá độ ổn định của cách mô tả và không được
dùng để lựa chọn lại các ngưỡng chính.

\paragraph{Tổng hợp các đặc trưng dữ liệu cần lưu ý trước huấn luyện.}

Các kết quả chẩn đoán trước huấn luyện cho thấy một số đặc trưng dữ liệu
cần được giữ làm cơ sở đối chiếu với các thí nghiệm mô hình sau này.

Thứ nhất, số ảnh hỗ trợ giữa các lớp biến thiên từ 66 đến 2.144 ảnh,
tương ứng tỷ số \(R_{\max/\min}=32{,}48\). Theo tiêu chí 5\%,
\textit{Atelectasis} và \textit{Pneumothorax} là hai lớp có mức hỗ trợ
thấp trong tập huấn luyện.

Thứ hai, theo các ranh giới diện tích chuẩn hóa chính, 37,10\% bounding
box thuộc nhóm Small, 57,85\% thuộc nhóm Medium và 5,05\% thuộc nhóm
Large. Phân tích độ nhạy đồng thời cho thấy các tỷ lệ này thay đổi khi
các ranh giới vận hành được thay đổi.

Thứ ba, tất cả bốn tập có nhãn đều bao phủ đủ 14 lớp và độ lệch phân bố
so với tập huấn luyện giảm khi ngân sách nhãn tăng. Tuy nhiên, ở ngân
sách 1\%, hai lớp rare mỗi lớp chỉ có một ảnh có nhãn, cho thấy class
coverage đầy đủ không đồng nghĩa với mức hỗ trợ tuyệt đối đồng đều.

Thứ tư, dữ liệu có cấu trúc đa nhãn ở cấp ảnh, với trung bình 3,131 lớp
bất thường duy nhất trên mỗi ảnh, trong khi một số ảnh chứa tới 10 lớp.
Phân tích 91 cặp lớp cho thấy số ảnh đồng xuất hiện khác nhau giữa các
cặp; trong đó
\textit{Aortic enlargement}--\textit{Cardiomegaly} có 1.373 ảnh đồng
xuất hiện.

Những quan sát trên mô tả các đặc trưng của dữ liệu trước huấn luyện và
được sử dụng làm thông tin nền khi phân tích class-wise AP, sai số theo
kích thước bounding box, hành vi ở các ngân sách nhãn và các kết quả
semi-supervised ở những phần tiếp theo. Chúng chưa cho phép kết luận rằng
một lớp có mức hỗ trợ thấp chắc chắn sẽ có AP thấp, một bounding box
Small chắc chắn khó phát hiện hơn, hoặc phương pháp học bán giám sát sẽ
tự động khắc phục các đặc trưng dữ liệu đã quan sát.
\section{Discussion}\label{sec:discussion}

\section{Conclusion}\label{sec:conclusion}

\section*{Declarations}

\begin{itemize}

\item \textbf{Funding:} This research received no external funding.

\item \textbf{Conflict of interest:} The authors declare no competing interests.

\item \textbf{Ethics approval and consent to participate:} Not applicable. This study used publicly available datasets and did not involve direct human subject recruitment by the authors.

\item \textbf{Consent for publication:} Not applicable.

\item \textbf{Data availability:} The datasets used in this study are publicly available. Dataset sources and access information are described in the manuscript.

\item \textbf{Code availability:} The code and experimental logs will be made available in a public repository upon acceptance.

\item \textbf{Author contributions:} All authors contributed to the study conception, experimental design, analysis, interpretation of results, and manuscript preparation.

\item \textbf{AI Declaration Statement:} The authors used AI-assisted tools, including ChatGPT (OpenAI, GPT-5.5 Thinking), Gemini, and Grammarly, to support language editing, manuscript structuring, LaTeX formatting, consistency checking across sections, and clarity improvement during manuscript preparation. AI tools were not used to generate experimental data, fabricate images, create false references, or produce unverified scientific claims. All AI-assisted content was reviewed, verified, revised, and approved by the authors. The authors take full responsibility for the accuracy, originality, integrity, and final content of the manuscript.

\end{itemize}

\FloatBarrier
\clearpage

\bibliography{sn-bibliography}
\end{document}
